%% GENERATED by battery/render_battery.py
%% from claims.json — do not edit by hand; edit the registry and re-render.
\subsection*{Census}
The registry holds 112 claims covering all 91 labels of the 2026-09-29 version of the draft (7 labels are section heads or bare definitions, listed as such in the registry).
\begin{center}\begin{tabular}{@{}l>{\raggedright\arraybackslash}p{11.5cm}@{}}\toprule
standing & \CERTIFIED{} 64, \DERIVED{} 3, \MEASURED{} 19, \OPEN{} 9; definition rows (no standing) 17 \\
kind & conjecture 24, definition 17, identity 20, implementation 26, not-testable 5, theorem 20 \\
cost tier & extensive 22, fast 63, none 27 \\
environment & none 27, web 77, web+local 8 \\
result records & 84 claims run, 1287 checks, 0 failing (latest per environment and depth; 2026-09-21, 2026-09-22, 2026-09-23, 2026-09-24, 2026-09-25, 2026-09-26, 2026-09-27, 2026-09-28, 2026-09-29) \\
\bottomrule\end{tabular}\end{center}
\noindent\emph{Kinds}: a \emph{definition} row is not a claim --- a definition has no truth value, and checking an axiom on a realisation tests the realisation, not the definition (the author's ruling) --- so it carries no standing; it records per tier where the axiom is enforced by construction and where it is emergent, and names the rows that test it there, the verifiers being the instruments; a \emph{theorem} is proved in the paper or the repository and checked for consistency; an \emph{identity} is a closed form checked to a stated order; a \emph{conjecture} gathers evidence; \emph{coverage} is a classification table; \emph{not-testable} has no counterpart in code.  \emph{Cost tiers}: \emph{fast} runs in minutes with the suite; \emph{extensive} runs for hours with windows and depths that grow from run to run.  \emph{Environments}: \emph{web} means the tracked repository is the whole input, so the claim runs in the web session and in continuous integration; \emph{local} means the run needs data that only the local machine holds --- the gitignored dictionary builds (the enumerated build beyond the shipped weight, $E\ge13$, which stays on the local machine, the flavoured staging builds, the weak dictionary with its mutation edges) --- so it runs there and its result records are committed from there.  A claim listing both has a web-sized population (the shipped dictionary tiers, whose coverage is read from their manifests) and a local-sized one, each named under \emph{Inputs}.  \emph{Standing} \PLANNED{} marks a test designed here and not yet implemented.

\subsection*{Section~1: introduction}
\begin{longtable}{@{}>{\raggedright\arraybackslash}p{2.4cm}>{\raggedright\arraybackslash}p{4.7cm}>{\raggedright\arraybackslash}p{4.95cm}>{\raggedright\arraybackslash}p{1.85cm}@{}}
\toprule claim & statement & test: procedure, population, controls & standing \\ \midrule \endfirsthead
\toprule claim & statement & test: procedure, population, controls & standing \\ \midrule \endhead
\bottomrule \endfoot
\code{intro/delta}\newline \peqref{eq:intro-delta} & $I_{a,b}=\Tr L_{\rho(a)}L_b=\Tr L_bL_{\rho^{-1}(a)}=\delta_{a,b}+O(\fq)$ on the canonical basis. & Not a claim: a definition has no truth value, and checking an axiom on a realisation tests the realisation, not the definition.  This row records where the axiom is \emph{enforced} by construction (a pass is a regression guard, left to the suite) and where it is \emph{emergent} (a pass is evidence), and names the rows that test it there.  The introduction's display is axiom K4 restated; see \code{def:kq/K4}.  Consistency of the definition is witnessed by its first model, \code{ex:qt}.  Instruments: \code{tests/test_kalgebra.py}, \code{every realisation suite} \emph{Evidence:} emergent (the trace is computed independently of the construction of the basis). \emph{Where:} \code{tests/test_kalgebra.py}, \code{every realisation suite}. & ---\newline {\footnotesize definition, not run} \\ \addlinespace
\code{intro/positivity}\newline Section~\pref{sec:kq}, introduction & The inner product $I_{a,b}(\fq)$ is positive-definite for real $-1<\fq<1$ (Schur quantization). & No test: the Gram matrix is a truncated series with no bound on its tail, so positive-definiteness of a truncated matrix at fixed $\fq$ would be numerical evidence at best. \emph{Population:} finite label windows of each realisation; Gram matrices truncated at order $K$. \emph{Note:} An open research goal of the source repository; no test exists, and the row stays open. & \OPEN\newline {\footnotesize conjecture, none\newline no test} \\ \addlinespace
\end{longtable}

\subsection*{Section~\pref{sec:kq}: $K_\fq$-algebras}
\begin{longtable}{@{}>{\raggedright\arraybackslash}p{2.4cm}>{\raggedright\arraybackslash}p{4.7cm}>{\raggedright\arraybackslash}p{4.95cm}>{\raggedright\arraybackslash}p{1.85cm}@{}}
\toprule claim & statement & test: procedure, population, controls & standing \\ \midrule \endfirsthead
\toprule claim & statement & test: procedure, population, controls & standing \\ \midrule \endhead
\bottomrule \endfoot
\code{def:kq/K1}\newline Def.~\pref{def:kq}, K1 & antimultiplicative bar involution, $\fq\mapsto\fq^{-1}$, $A_\fq\simeq A^{\mathrm{op}}_{\fq^{-1}}$ & Not a claim: a definition has no truth value, and checking an axiom on a realisation tests the realisation, not the definition.  This row records where the axiom is \emph{enforced} by construction (a pass is a regression guard, left to the suite) and where it is \emph{emergent} (a pass is evidence), and names the rows that test it there.  Enforced on the abelianized tier (W1 bar-invariance is a gate of the guarded solve) and on solver-built flows; analytic on the quantum torus and on RG-built algebras (from R1--R4); emergent on hand-built presentations.  Tested where emergent: \code{def:penta/existence}, \code{sec:uqsl2/algebra}, \code{sec:a1d3/algebra}; on the Coulomb-branch side \code{conj:abeKalgebra/consistency}.  Instruments: \code{src/core/kalgebra.py}, \code{tests/test_kalgebra.py} \emph{Evidence:} per tier (the audit: enforced where built in, emergent otherwise). \emph{Where:} \code{src/core/kalgebra.py}, \code{tests/test_kalgebra.py}. & ---\newline {\footnotesize definition, not run} \\ \addlinespace
\code{def:kq/K2}\newline Def.~\pref{def:kq}, K2 & a bar-invariant canonical basis containing $1$; the algebra is spanned by it & Not a claim: a definition has no truth value, and checking an axiom on a realisation tests the realisation, not the definition.  This row records where the axiom is \emph{enforced} by construction (a pass is a regression guard, left to the suite) and where it is \emph{emergent} (a pass is evidence), and names the rows that test it there.  Closure is the one axiom that is emergent on almost every tier: that two canonical elements multiply into a $\bZ[\fq^{\pm1}]$-combination of canonical elements is what the constructions must deliver.  Tested on \code{def:penta/existence}, \code{sec:uqsl2/algebra}, \code{sec:a1d3/algebra}, \code{conj:abeKalgebra/consistency}, \code{conj:upper} (Laurentness and closure on charts) and the skein rows of Section~\pref{sec:skein}.  Instruments: \code{src/core/kalgebra.py}, \code{src/abe/abe_kalgebra.py} \emph{Evidence:} per tier (the audit: enforced where built in, emergent otherwise). \emph{Where:} \code{src/core/kalgebra.py}, \code{src/abe/abe_kalgebra.py}. & ---\newline {\footnotesize definition, not run} \\ \addlinespace
\code{def:kq/K3}\newline Def.~\pref{def:kq}, K3 & an algebra automorphism $\rho$ permuting the basis and fixing $1$ & Not a claim: a definition has no truth value, and checking an axiom on a realisation tests the realisation, not the definition.  This row records where the axiom is \emph{enforced} by construction (a pass is a regression guard, left to the suite) and where it is \emph{emergent} (a pass is evidence), and names the rows that test it there.  Manifest on the quantum torus and on the abelianized tier (the label-level closed form of $\rho$); inherited on RG-built algebras; emergent on hand-built presentations.  Tested on \code{def:penta/existence}, \code{sec:uqsl2/algebra}, \code{sec:a1d3/algebra}.  Instruments: \code{src/core/kalgebra.py} \emph{Evidence:} per tier (the audit: enforced where built in, emergent otherwise). \emph{Where:} \code{src/core/kalgebra.py}. & ---\newline {\footnotesize definition, not run} \\ \addlinespace
\code{def:kq/K4}\newline Def.~\pref{def:kq}, K4 & a $\rho^2$-twisted trace with $I_{a,b}=\delta_{a,b}+O(\fq)$ in both faces & Not a claim: a definition has no truth value, and checking an axiom on a realisation tests the realisation, not the definition.  This row records where the axiom is \emph{enforced} by construction (a pass is a regression guard, left to the suite) and where it is \emph{emergent} (a pass is evidence), and names the rows that test it there.  Where the trace is bootstrapped from orthonormality (the cone / finite-type tier) the axiom \emph{defines} the trace and the test is existence and uniqueness of the bootstrap solution, \code{app:finite/trace-unique}; analytic on RG-built algebras (from R4 and the IR); a build gate on the abelianized tier (W2), where the informative statement is that the gate never rejects, \code{conj:abeKalgebra/consistency}; emergent on hand-built presentations, \code{def:penta/existence}, \code{sec:uqsl2/algebra}, \code{sec:a1d3/algebra}.  Instruments: \code{src/core/kalgebra.py} \emph{Evidence:} per tier (the audit: enforced where built in, emergent otherwise). \emph{Where:} \code{src/core/kalgebra.py}. & ---\newline {\footnotesize definition, not run} \\ \addlinespace
\code{def:kq/K5}\newline Def.~\pref{def:kq}, K5 & $\Tr L_{\rho(a)}=\Tr L_a$, hence $I_{b,a}=I_{a,b}$ & Not a claim: a definition has no truth value, and checking an axiom on a realisation tests the realisation, not the definition.  This row records where the axiom is \emph{enforced} by construction (a pass is a regression guard, left to the suite) and where it is \emph{emergent} (a pass is evidence), and names the rows that test it there.  Per tier: inherited on RG flows with the bilinear pairing; a tautology of the Nahm-sum route on the BPS pairing form (only the trace form has content); a theorem on the abelianized tier from the derived closed form of the sector measure; emergent on hand-built presentations, tested on \code{def:penta/existence}, \code{sec:uqsl2/algebra}, \code{sec:a1d3/algebra}.  Instruments: \code{src/core/kalgebra.py}, a test in the source repository \emph{Evidence:} per tier (the audit: enforced where built in, emergent otherwise). \emph{Where:} \code{src/core/kalgebra.py}. & ---\newline {\footnotesize definition, not run} \\ \addlinespace
\code{def:kq/contract}\newline Definition~\pref{def:kq} & The \code{KAlgebra} contract expresses K1--K5 as printed: each axiom has a verifier computing the paper's formula, and a realisation violating one axiom is rejected by that axiom's verifier. & \emph{Transcription}: K1 $\leftrightarrow$ \code{verify_bar_involution} ($C^c_{ab}(\fq^{-1})=C^c_{ba}(\fq)$); K2 $\leftrightarrow$ \code{verify_identity_in_basis}, \code{verify_unit_law}, \code{verify_associativity} and closure of products on canonical labels; K3 $\leftrightarrow$ \code{verify_rho_is_automorphism}, \code{verify_rho_inverse}, \code{verify_rho_fixes_identity}; K4 $\leftrightarrow$ \code{verify_rho_twisted_trace}, \code{verify_trace_pairing_faces}, \code{verify_orthonormality} (no negative $\fq$-powers, $I_{a,b}[\fq^0]=\delta_{a,b}$); K5 $\leftrightarrow$ \code{verify_trace_intertwines_rho}, \code{verify_pairing_rho_star_symmetric}.  \emph{Discrimination}: a non-bar-invariant basis $\fq^{a_1^2}X_a$ (K1), a unit acting by $\fq$ (K2), the anti-automorphism $(a_1,a_2)\mapsto(a_1,-a_2)$ (K3), the trace doubled (K4), and on the degenerate torus the even flavour trace $\Tr X_a=\mu^{|a|}$ (K5); the matrix of which verifiers fire is recorded. \emph{Population:} the transcription (five axioms, twelve verifiers); five toy countermodels on the rank-2 quantum torus and the degenerate rank-1 torus, one per axiom, plus the genuine torus as positive control; window $|\gamma_i|\le1$. \emph{Controls:} positive: the genuine quantum torus passes all twelve verifiers; negative: each countermodel is rejected by its axiom's verifier. \emph{Evidence:} n.a. (a test of the code against the text of the definition). \emph{Where:} \code{battery/checks_kq.py}. \emph{Note:} user, 2026-09-21: the first check is (a) 'does the KAlgebra contract express correctly the K\_q algebra axioms?'  K4 and K5 admit countermodels caught by their verifier alone; K1-K3 countermodels also break orthonormality or the bar symmetry, an entanglement of the axioms, not of the verifiers. \emph{Run (web, 2026-09-21):} 11 checks, 0 failing; toy realisations: 6; window: rank 2: |$\gamma$\_i| $\le$ 1 (9 labels); rank 1: |a| $\le$ 2; K: 6. & \CERTIFIED\newline {\footnotesize implementation, fast\newline web}\newline {\footnotesize run web 2026-09-21: 11/11} \\ \addlinespace
\code{def:kq/rescaling}\newline Remark 2.2 (after Def.~\pref{def:kq}) & The axioms fix the trace only up to a factor $1+O(\fq)$; physics fixes the normalisation. & not a claim to test; the normalisation is the $(\fq^2;\fq^2)_\infty^{2\rk}/|W|$ prefactor of \code{trace_residual}. & \DERIVED\newline {\footnotesize theorem, none\newline no test} \\ \addlinespace
\code{def:kq/webster}\newline Remark 2.3 (after Def.~\pref{def:kq}) & The canonical basis and the $\delta_{a,b}$ axiom play a similar role in Webster's work, with $\fq_{\mathrm{here}}=q^{-1}_{\mathrm{there}}$. & a comparison with the literature; no code. & \OPEN\newline {\footnotesize not-testable, none\newline no test} \\ \addlinespace
\code{ex:qt}\newline Ex.~\pref{ex:qt} & The quantum torus $Q_\fq(\Gamma)$ with $X_\gamma X_{\gamma'}=\fq^{\langle\gamma,\gamma'\rangle}X_{\gamma+\gamma'}$, $\rho(\gamma)=-\gamma$ and $\Tr X_\gamma=(\fq^2)^{\rk\Gamma}_\infty\delta_{\gamma,0}$ is a $K_\fq$-algebra, with $I_{\gamma,\gamma'}=(\fq^2)^{\rk\Gamma}_\infty\delta_{\gamma,\gamma'}$. & analytically obvious, so the check is that \code{QuantumTorusKAlg} implements it: relations exact as elements on every pair, $\rho=-\mathrm{id}$, the trace and $I_{\gamma,\gamma'}$ against an independent $(\fq^2;\fq^2)_\infty^{\rk}$ (controlled by the pentagonal-number theorem); then K1--K5 on the rank-2 torus as the consistency witness of Definition~\pref{def:kq}. \emph{Population:} pairings of rank 2 ($\langle e_1,e_2\rangle=1,2$) and rank 4; labels $|\gamma_i|\le1$; $K=10$. \emph{Controls:} positive: Euler's pentagonal-number theorem on the independent series; negative: a wrong prefactor exponent is detected. \emph{Evidence:} n.a. (a test of the code against the text of the example). \emph{Where:} \code{battery/checks_kq.py}, \code{tests/test_quantum_torus_kalgebra.py}. \emph{Note:} Formalised in ClusterLean (the local Lean 4 repository): Q\_q($\Gamma$,$\omega$) over any coefficient ring, abelian group $\Gamma$ and alternating $\omega$ is a K\_q-algebra (K1-K5) for every Tr 1 = c\_0 in 1+qR[[q]], the draft's (q\^{}2;q\^{}2)\_$\infty$\^{}(rk $\Gamma$) included, with I = c\_0 $\delta$ (qt\_isKq); the trace is unique up to 1+qR[[q]] when $\omega$ is nondegenerate (qt\_traceClassification).  The rank-2 case also matched Z2QTorusSampleKAlgebra and QuantumTorusZ2KAlg exactly against a proven executable. \emph{Run (web, 2026-09-21):} 16 checks, 0 failing; pairings: ["rank 2, <e1,e2> = 1 (the pentagon's IR torus)", 'rank 2, <e1,e2> = 2', 'rank 4, block (1, 3)']; labels: |$\gamma$\_i| $\le$ 1; pairs checked: 6723; K: 10. & \CERTIFIED\newline {\footnotesize implementation, fast\newline web}\newline {\footnotesize formalised: \code{ClusterLean/QTorus/Kq.lean:qt_isKq}}\newline {\footnotesize run web 2026-09-21: 16/16} \\ \addlinespace
\code{rem:maximal}\newline Rem.~\pref{rem:maximal} & The $\delta_{a,b}$ axiom holds only after the collection of line defects is maximally extended (direct summands included). & a statement about physical completeness; the testable proxy is K2's closure on each presentation. & \OPEN\newline {\footnotesize not-testable, none\newline no test} \\ \addlinespace
\code{def:penta/implementation}\newline Definition~\pref{def:penta} & The stand-alone class \code{PentagonKAlg} implements the definition: the relations $L_{i+1}L_i=\fq^2L_iL_{i+1}$, $L_{i+1}L_{i-1}=1+\fq L_i$, $L_{i-1}L_{i+1}=1+\fq^{-1}L_i$ with $L_{i+5}=L_i$; the canonical basis $L_{i;a,b}=\fq^{ab}L_i^aL_{i+1}^b$; $\rho(L_i)=L_{i+2}$ of order five. & every printed relation recomputed with \code{multiply} and compared exactly as elements; $L_i^aL_{i+1}^b=\fq^{-ab}L_{i;a,b}$ and pure powers; $\rho$, $\rho^{-1}$ on generators and on the window, $\rho^5=\mathrm{id}$, $\rho^k\neq\mathrm{id}$ for $k<5$. \emph{Population:} all five $i$; $1\le a,b\le4$ for the basis form; the window $a+b\le3$ for $\rho$. \emph{Controls:} positive: the relations reproduced on all $i$; negative: an altered relation ($1+\fq^2L_i$) is detected. \emph{Evidence:} n.a. (a test of the code against the text of the definition). \emph{Where:} \code{battery/checks_kq.py}, \code{src/abe/kalgebra_samples.py}, \code{src/cone/pentagon_cone_data.py}. \emph{Note:} the specification a Lean formalisation of the pentagon example (a local session, user 2026-09-21) has to meet. \emph{Run (web, 2026-09-21):} 11 checks, 0 failing; generators: 5; basis window: a, b $\le$ 4 for the basis form; a+b $\le$ 3 for rho; K: 24. & \CERTIFIED\newline {\footnotesize implementation, fast\newline web}\newline {\footnotesize run web 2026-09-21: 11/11} \\ \addlinespace
\code{def:penta/existence}\newline Def.~\pref{def:penta}, Remark 2.7 & The pentagon algebra is a $K_\fq$-algebra: K1--K5 hold on the stand-alone class, where all five are emergent; the trace is pinned by the axioms; the three presentations (stand-alone, cone data, BPS chart) are one algebra. & the twelve verifiers on every label pair of the window, associativity on generator triples, closure of all products on canonical labels; the trace seed perturbed by $\fq^3$ fails orthonormality on a pure power $L_i^n$ (the negative control that shows K4 pins the seeds); \code{KAlgebraIso} witnesses in \code{finite_kalgebra_objects.kalgebra_object(\"pentagon\")} pairwise and path-independent. \emph{Population:} the window $a+b\le3$, all $i$ (the identity, 15 pure powers and the mixed labels); $K=8$. \emph{Controls:} positive: all verifiers pass on the window; negative: the perturbed seed is caught by orthonormality. \emph{Evidence:} emergent (hand-built presentation: cone data from the relations, closed-form seeds). \emph{Where:} \code{battery/checks_kq.py}, \code{tests/test_pentagon_cone_data.py}, \code{tests/test_pentagon_kalg_notes.py}, \code{tests/test_finite_kalgebras.py}. \emph{Note:} user, 2026-09-21: 'is the stand-alone Pentagon class implementing correctly the definition, and is it actually a KAlgebra?'; formalised in ClusterLean (the local Lean 4 repository): the abstract algebra of pentagon\_kq\_proof.md Def. 2.1 satisfies K1-K5 with each axiom on its face, and the K4 functionals are exactly c$\cdot$Tr, c in 1+qZ[[q]] (pentagon\_KqAlgebra); of the three presentations, the injective algebra map into the quantum torus is formalised (pentagon\_proof\_document), the cone-data presentation is not; PentagonKAlg matched the proven executable exactly. \emph{Run (web, 2026-09-21):} 16 checks, 0 failing; window: 31 labels, a+b $\le$ 3, all i; K: 8. & \CERTIFIED\newline {\footnotesize theorem, fast\newline web}\newline {\footnotesize formalised: \code{ClusterLean/Pentagon/KqAlgebra.lean:pentagon_KqAlgebra}}\newline {\footnotesize run web 2026-09-21: 16/16} \\ \addlinespace
\code{def:penta/trace-seeds}\newline Def.~\pref{def:penta} & $\Tr 1=\sum_n\fq^{2n(n+1)}/(\fq^2;\fq^2)_n$ and $\Tr L_i=-\sum_n\fq^{2(n+1)^2-1}/(\fq^2;\fq^2)_n$; equivalently $\Tr 1=\chi_0(\fq^2)$, $\Tr L_i=\fq^{-1}(\chi_0-\chi_1)(\fq^2)$ with the Yang--Lee characters. & the class's $\Tr1$ and $\Tr L_i$ (all five $i$) against both Nahm sums and against $\chi_0(\fq^2)$, $\fq^{-1}(\chi_0-\chi_1)(\fq^2)$ with the Yang--Lee products, all recomputed independently to $\fq^{30}$; the Rogers--Ramanujan identity as the equality of the two forms; the printed expansions of $\Tr1$, $\Tr L_i$, $\chi_0$, $\chi_1$ verbatim. \emph{Population:} orders $K\le 40$. \emph{Controls:} positive: the printed expansions; negative: a seed perturbed at $\fq^3$ is detected. \emph{Evidence:} emergent. \emph{Where:} \code{battery/checks_kq.py}, \code{tests/test_pentagon_trace.py}. \emph{Note:} Formalised in ClusterLean (the local Lean 4 repository): Tr 1 = $\chi$\_0(q\^{}2) and Tr L\_i = q\^{}-1 ($\chi$\_0 - $\chi$\_1)(q\^{}2) for the document's trace $\chi$\_0(q\^{}2)$\cdot$Tr, with $\chi$\_0 = R\_1, $\chi$\_1 = R\_0, R\_s(p) = $\Sigma$\_k p\^{}(k\^{}2+sk)/(p;p)\_k and R\_s = R\_(s+1) + p\^{}(s+1) R\_(s+2). \emph{Run (web, 2026-09-21):} 10 checks, 0 failing; K: 30; seeds: 2; orbit: 5. & \CERTIFIED\newline {\footnotesize identity, fast\newline web}\newline {\footnotesize formalised: \code{ClusterLean/Pentagon/KqAlgebra.lean:pentagon_proof_document}}\newline {\footnotesize run web 2026-09-21: 10/10} \\ \addlinespace
\code{def:penta/recursion}\newline \peqref{eq:izero}, \peqref{eq:pentarec}, \peqref{eq:ABdef} & $\rho^2$-cyclicity gives $\Tr L_i^aL_{i+1}^b=\Tr L_i^{a+1}L_{i+1}^{b-1}$, hence $\Tr L_{i;a,b}=\fq^{ab}\Tr L_i^{a+b}$, and $\Tr L_i^n=\fq^{1-2n}\Tr L_i^{n-1}+\fq^{2-2n}\Tr L_i^{n-2}$; twisted cyclicity reduces to the five families $\Tr L_0L_{i;a,b}=\Tr L_{i;a,b}L_1$; $\Tr L_1=(-\fq+O(\fq^4))\Tr1$. & each identity evaluated on the class's traces of products of pure powers and compared to $\fq^{16}$; the five cyclicity families on 80 instances; the ratio by exact series division. \emph{Population:} $a\le3$, $b\le3$, $n\le8$, all $i$; $K=16$ (traces of products computed at $K+40$ to absorb $\fq^{-ab}$). \emph{Controls:} positive: the recursion reproduces $\Tr L_i^2=\fq^{-2}\Tr1+\fq^{-3}\Tr L_i$; negative: n.a. \emph{Evidence:} emergent. \emph{Where:} \code{battery/checks_kq.py}, \code{tests/test_pentagon_trace.py}, \code{tests/test_pentagon_kalg_notes.py}. \emph{Note:} rows n = 3..6 of the table reproduced exactly on 2026-09-21; formalised in ClusterLean (the local Lean 4 repository): Tr L\_i\^{}a L\_(i+1)\^{}b = T\_(a+b) with no phase (Prop. 7.2(iii)), the three-term recursion of T\_n, and T\_n in (-1)\^{}n q\^{}n (1 + q\^{}(2n+4) Z[[q]]) with nonnegative coefficients, from which Tr L\_1 = (-q + O(q\^{}4)) Tr 1 follows; twisted cyclicity itself is proved, by reduction to the generators. \emph{Run (web, 2026-09-21):} 6 checks, 0 failing; K: 16; powers: n $\le$ 8; cyclicity instances: 80. & \CERTIFIED\newline {\footnotesize theorem, fast\newline web}\newline {\footnotesize formalised: \code{ClusterLean/Pentagon/KqAlgebra.lean:pentagon_proof_document}}\newline {\footnotesize run web 2026-09-21: 6/6} \\ \addlinespace
\code{def:penta/AB-table}\newline \peqref{eq:ABdef}, Table~\pref{tab:pentagon-stabilization} & $\fq^{n^2-1}\Tr L_i^n=A_n\Tr1+B_n\Tr L_1$ with the tabulated $A_n,B_n$ ($n=3,\dots,6$); the ratio $\Tr L_1/\Tr1$ is determined to $O(\fq^{n^2})$; $1-\fq\Tr L_i/\Tr1=1+\fq^2/(1+\fq^4/(1+\cdots))$. & $A_n,B_n$ generated from \peqref{eq:pentarec} alone ($T_n=T_{n-1}+\fq^{2n-2}T_{n-2}$ for $T_n=\fq^{n^2-1}\Tr L^n$), compared verbatim with the table and, times the class's seeds, with the class's $\Tr L_i^n$; the first power at which $-A_n/B_n$ departs from the ratio; the valuation of $\Tr L^n$; the convergents' agreement order. \emph{Population:} $n\le10$ against the class trace; the ratio for $n\le8$; convergents of depth $\le8$. \emph{Controls:} positive: the table rows reproduced from the recursion; negative: n.a. \emph{Evidence:} emergent (the recursion is derived from the relations, the traces are computed by the class). \emph{Where:} \code{battery/checks_kq.py}, \code{tests/test_pentagon_kalg_notes.py}. \emph{Note:} Partly formalised in ClusterLean (the local Lean 4 repository) (pentagon\_miracle): eq:ABdef for every n and every twisted-cyclic trace, with A\_n, B\_n defined by their recursion, and the limits B\_n $\to$ $\chi$\_0, A\_n $\to$ p R\_2 with error O(p\^{}n), i.e. B\_$\infty$ = Tr 1 and A\_$\infty$ = -Tr L\_1; the tabulated entries (n = 3..6) and the continued fraction are not stated in Lean. \emph{Run (web, 2026-09-21):} 5 checks, 0 failing; n: $\le$ 10 for eq:ABdef, $\le$ 8 for the ratio, table rows 3..6; K: 40. & \CERTIFIED\newline {\footnotesize identity, fast\newline web}\newline {\footnotesize run web 2026-09-21: 5/5} \\ \addlinespace
\code{rem:miracle}\newline Rem.~\pref{rem:miracle} & $A_\infty=-\Tr L_1$ and $B_\infty=\Tr 1$: the recursion knows the traces themselves. & $A_n,B_n$ from the recursion for $n\le16$ against $-\Tr L_1$ and $\Tr1$ to $\fq^{80}$: the first differing power per $n$, required to grow with $n$; each row extends the one above; $A_{16},B_{16}$ agree through $\fq^{30}$. \emph{Population:} $n\le16$; $K=80$. \emph{Evidence:} emergent. \emph{Where:} \code{battery/checks_kq.py}. \emph{Note:} the trace-miracle puzzle; pentagon\_miracle\_reduction.md. \emph{Run (web, 2026-09-21):} 4 checks, 0 failing; n: $\le$ 16; K: 80. & \MEASURED\newline {\footnotesize conjecture, fast\newline web}\newline {\footnotesize run web 2026-09-21: 4/4} \\ \addlinespace
\code{def:kq-flavoured/F1}\newline Def.~\pref{def:kq-flavoured}, F1 & antimultiplicative bar involution on $A_\fq[G_f]$ & Not a claim: a definition has no truth value, and checking an axiom on a realisation tests the realisation, not the definition.  This row records where the axiom is \emph{enforced} by construction (a pass is a regression guard, left to the suite) and where it is \emph{emergent} (a pass is evidence), and names the rows that test it there.  As K1 with $R(G_f)$ coefficients; emergent on the hand-built flavoured presentation \code{sec:a1d3/algebra} and on the flavoured charts, \code{sec:rg/flavoured-quivers}.  Instruments: \code{tests/test_bps_kalgebra_flavour.py} \emph{Evidence:} per tier. \emph{Where:} \code{tests/test_bps_kalgebra_flavour.py}. & ---\newline {\footnotesize definition, not run} \\ \addlinespace
\code{def:kq-flavoured/F2}\newline Def.~\pref{def:kq-flavoured}, F2 & bar-invariant canonical basis over $\bZ[\fq^{\pm1}]$ from which a basis over $\bZ[\fq^{\pm1}]\otimes R_{G_f}$ is selected non-canonically, the ambiguity being 1d irreps & Not a claim: a definition has no truth value, and checking an axiom on a realisation tests the realisation, not the definition.  This row records where the axiom is \emph{enforced} by construction (a pass is a regression guard, left to the suite) and where it is \emph{emergent} (a pass is evidence), and names the rows that test it there.  As K2; the informative content is that products close over $R(G_f)$ with characters, never imposed: \code{sec:rg/flavoured-quivers}, \code{sec:coulomb/flavour-enhancement}.  Instruments: \code{src/core/kalgebra.py} \emph{Evidence:} per tier. \emph{Where:} \code{src/core/kalgebra.py}. & ---\newline {\footnotesize definition, not run} \\ \addlinespace
\code{def:kq-flavoured/F3}\newline Def.~\pref{def:kq-flavoured}, F3 & $\rho$ twisted-linear over $R_{G_f}$ with $\rho(\chi_r)=\chi_{r^\vee}$ & Not a claim: a definition has no truth value, and checking an axiom on a realisation tests the realisation, not the definition.  This row records where the axiom is \emph{enforced} by construction (a pass is a regression guard, left to the suite) and where it is \emph{emergent} (a pass is evidence), and names the rows that test it there.  As K3; tested on \code{sec:a1d3/algebra}, \code{sec:rg/flavoured-quivers}.  Instruments: \code{src/core/kalgebra.py}, \code{src/cone/cone_kalgebra.py} \emph{Evidence:} per tier. \emph{Where:} \code{src/core/kalgebra.py}, \code{src/cone/cone_kalgebra.py}. & ---\newline {\footnotesize definition, not run} \\ \addlinespace
\code{def:kq-flavoured/F4}\newline Def.~\pref{def:kq-flavoured}, F4 & $I_{a,b}\in R_{G_f}+\fq R_{G_f}[[\fq]]$ with identity summand $\delta_{a,b}+O(\fq)$ & Not a claim: a definition has no truth value, and checking an axiom on a realisation tests the realisation, not the definition.  This row records where the axiom is \emph{enforced} by construction (a pass is a regression guard, left to the suite) and where it is \emph{emergent} (a pass is evidence), and names the rows that test it there.  As K4 with the $R(G_f)$-valued pairing; tested where emergent, \code{sec:a1d3/algebra}, \code{sec:coulomb/flavour-enhancement}.  Instruments: \code{src/core/kalgebra.py} \emph{Evidence:} per tier. \emph{Where:} \code{src/core/kalgebra.py}. & ---\newline {\footnotesize definition, not run} \\ \addlinespace
\code{def:kq-flavoured/F5}\newline Def.~\pref{def:kq-flavoured}, F5 & $\Tr L_{\rho(a)}=\rho(\Tr L_a)$, the outer $\rho$ acting on $R_{G_f}$ as $\chi_r\mapsto\chi_{r^\vee}$ & Not a claim: a definition has no truth value, and checking an axiom on a realisation tests the realisation, not the definition.  This row records where the axiom is \emph{enforced} by construction (a pass is a regression guard, left to the suite) and where it is \emph{emergent} (a pass is evidence), and names the rows that test it there.  The $\star$-twisted form of K5; non-vacuous only with a non-abelian $G_f$: \code{sec:a1d3/algebra}, \code{sec:coulomb/flavour-enhancement}.  Instruments: a test in the source repository \emph{Evidence:} per tier. & ---\newline {\footnotesize definition, not run} \\ \addlinespace
\code{def:kq-flavoured/F6}\newline Def.~\pref{def:kq-flavoured}, F6 & the forgetful map $\chi_r\mapsto\dim r$ lands on a $K_\fq$-algebra, identifies bases, and commutes with $\rho$ and $\Tr$ & Not a claim: a definition has no truth value, and checking an axiom on a realisation tests the realisation, not the definition.  This row records where the axiom is \emph{enforced} by construction (a pass is a regression guard, left to the suite) and where it is \emph{emergent} (a pass is evidence), and names the rows that test it there.  The forgetful map's commutation with $\rho$ and $\Tr$ has no verifier today (the instrument is planned); once written it runs where the flavour is emergent, \code{sec:rg/flavoured-quivers}, \code{rem:flavour-change}.  Instruments: \code{src/core/kalgebra.py (forget, lower_flavour)} \emph{Evidence:} per tier. \emph{Where:} \code{src/core/kalgebra.py (forget, lower_flavour)}. & ---\newline {\footnotesize definition, not run} \\ \addlinespace
\code{def:kq-flavoured/contract}\newline Definition~\pref{def:kq-flavoured} & The \code{KAlgebra} contract expresses F1--F6: each axiom has a verifier computing the printed statement (F6 inline, no contract verifier yet), and a flavoured realisation violating one axiom is rejected by that axiom's check. & \emph{Transcription}: F1 $\leftrightarrow$ \code{verify_bar_involution}; F2 $\leftrightarrow$ unit, associativity, \code{verify_section_is_single_irrep}, \code{verify_embed_section_roundtrip} (the characters are in the basis, the $R$-basis is a choice of section); F3 $\leftrightarrow$ \code{verify_rho_is_automorphism}, \code{verify_embed_intertwines_rho} ($\rho(\chi_r)=\chi_{r^\vee}$); F4 $\leftrightarrow$ \code{verify_rho_twisted_trace}, \code{verify_trace_pairing_faces}, \code{verify_orthonormality} (no negative powers, identity summand $\delta$); F5 $\leftrightarrow$ \code{verify_trace_intertwines_rho_star}, \code{verify_pairing_rho_star_symmetric}; F6 $\leftrightarrow$ \code{forget()} is a homomorphism commuting with $\rho$ and, after the augmentation, with $\Tr$.  \emph{Discrimination}: a non-bar-invariant basis (F1), characters embedded as $\fq X_{\gamma_f}$ (F2), $\rho$ fixing the flavour (F3), the trace doubled (F4), an even flavour trace $\mu^{|f|}$ (F5), the prefactor $(\fq^2;\fq^2)_\infty^{\rk\Gamma}$ instead of $\rk(\Gamma/\Gamma_f)$ (F6). \emph{Population:} the transcription (six axioms, fifteen verifiers plus the inline F6 check); six toy countermodels on the degenerate rank-3 quantum torus ($\Gamma_f=\bZ e_3$, $R=\bZ[\mu^{\pm1}]$), one per axiom, plus the genuine torus; window $|\gamma_i|\le1$. \emph{Controls:} positive: the degenerate torus passes every verifier; negative: each countermodel rejected by its axiom's check. \emph{Evidence:} n.a. (a test of the code against the text of the definition). \emph{Where:} \code{battery/checks_flavoured.py}. \emph{Note:} F6-countermodel caught by the inline check alone: the prefactor exponent of the flavoured trace is fixed by the forgetful map, not by orthonormality; the inline check is the candidate KAlgebra.verify\_forget\_commutes of task T3. \emph{Run (web, 2026-09-21):} 13 checks, 0 failing; toy realisations: 7; window: |$\gamma$\_i| $\le$ 1 on Z\^{}3 (27 labels); K: 6. & \CERTIFIED\newline {\footnotesize implementation, fast\newline web}\newline {\footnotesize run web 2026-09-21: 13/13} \\ \addlinespace
\code{rem:flavour-change}\newline Rem.~\pref{rem:flavour-change} & Lowering along $H_f\to G_f$ through $R_{G_f}\to R_{H_f}$ gives an $H_f$-flavoured $K_\fq$-algebra with a map from $A_\fq[G_f]$. & \code{lower_flavour}$(\varphi)$ with $\varphi=$\code{restriction_hom} the pullback $R_{G_f}\to R_{H_f}$: the map $(s,\chi)\mapsto(s,\varphi(\chi))$ is an algebra homomorphism commuting with $\rho$ and with $\Tr$ ($\Tr'=\varphi\circ\Tr$); the lowered algebra passes F1--F5; the augmentation case reproduces \code{forget()}. \emph{Population:} the rank-4 torus with $G_f=U(1)^2$; two homs $H_f=U(1)\to G_f$ (diagonal, first factor); labels $|\gamma_i|\le1$; $K=6$. \emph{Controls:} positive: the augmentation case is the forgetful map; negative: two different lowerings are told apart on a flavour generator. \emph{Evidence:} emergent. \emph{Where:} \code{battery/checks_flavoured.py}, \code{tests/test_lower_flavour.py}. \emph{Run (web, 2026-09-21):} 8 checks, 0 failing; source: rank-4 torus, G\_f = U(1)\^{}2; homs: ['diagonal U(1) -> U(1)\^{}2 (f\_1, f\_2) -> f\_1 + f\_2', 'first factor U(1) -> U(1)\^{}2 (f\_1, f\_2) -> f\_1']; labels: |$\gamma$\_i| $\le$ 1 (81); K: 6. & \CERTIFIED\newline {\footnotesize theorem, fast\newline web}\newline {\footnotesize run web 2026-09-21: 8/8} \\ \addlinespace
\code{ex:qt-flavoured}\newline Ex.~\pref{ex:qt-flavoured} & The quantum torus with pairing kernel $\Gamma_f$ is flavoured by the torus $T_{\Gamma_f}$, $X_{\gamma_f}$ the 1d characters, with $\Tr X_\gamma=(\fq^2)^{\rk(\Gamma/\Gamma_f)}_\infty\chi_\gamma\delta_{\gamma\in\Gamma_f}$. & \code{QuantumTorusKAlg}: $\rk(\Gamma/\Gamma_f)$ against an independent matrix rank; the trace formula coefficient by coefficient with $\chi_\gamma=\mu^{\mathrm{flav}(\gamma)}$; \code{embed_R}$(\mu^f)=X_{\gamma_f}$; $I_{\gamma,\gamma'}=(\fq^2;\fq^2)_\infty^{\rk}\mu^{\mathrm{flav}(\gamma'-\gamma)}\delta$ (in $R+\fq R[[\fq]]$, identity summand $\delta$); F1--F6 all pass. \emph{Population:} three degenerate pairings (rank 3 with $\Gamma_f=\bZ e_3$; rank 3 with kernel $(1,-1,0)$; rank 4 with $\Gamma_f=\bZ^2$); labels $|\gamma_i|\le1$; $K=8$. \emph{Controls:} positive: the class's gauge rank equals the independent rank; negative: a prefactor with the full rank is detected. \emph{Evidence:} n.a. (a test of the code against the text of the example). \emph{Where:} \code{battery/checks_flavoured.py}, \code{tests/test_quantum_torus_kalgebra.py}. \emph{Run (web, 2026-09-21):} 16 checks, 0 failing; pairings: ['rank 3, $\Gamma$\_f = Z e\_3', 'rank 3, kernel (1,-1,0) (non-split coordinates)', 'rank 4, $\Gamma$\_f = Z\^{}2']; labels: |$\gamma$\_i| $\le$ 1; K: 8. & \CERTIFIED\newline {\footnotesize implementation, fast\newline web}\newline {\footnotesize run web 2026-09-21: 16/16} \\ \addlinespace
\code{rem:flavour-comments}\newline Rem.~\pref{rem:flavour-comments} & Line defects are assumed $G_f$-equivariant; lifts are canonical up to 1d irreps; $G_f$ may need a finite cover; disconnected $G_f$ left open. & assumptions; the finite-cover practice is the source repository's (SO(3) $\to$ SU(2)). & \OPEN\newline {\footnotesize not-testable, none\newline no test} \\ \addlinespace
\code{sec:uqsl2/algebra}\newline Def.~2.13 (Section~\pref{sec:uqsl2}) & \code{UqSL2PBW} implements the definition: $KE=\fq^{-2}EK$, $KF=\fq^2FK$, $EF=\chi_1+\fq K+\fq^{-1}K^{-1}$, $FE=\chi_1+\fq^{-1}K+\fq K^{-1}$ (so $[E,F]=(\fq-\fq^{-1})(K-K^{-1})$, Casimir $=\chi_1$); $E_{a,b}=\fq^{-ab}E^aK^b$, $F_{a,b}=\fq^{ab}F^aK^b$; $\rho(K)=K^{-1}$, $\rho(E)=\fq^{-1}FK^{-1}$, $\rho(F)=\fq KE$, $\rho$ fixing $\chi_k$ and of infinite order (Lusztig's braid). & every printed relation recomputed with \code{multiply} and compared exactly as elements; the normalisations $E^aK^b=\fq^{ab}E_{a,b}$, $F^aK^b=\fq^{-ab}F_{a,b}$; $\rho$ on $E,F,K,\chi_1$ as label permutations; the orbit of $E$ under $\rho$ ($F_{1,-1},E_{1,2},F_{1,-3},\dots$, all distinct); $\rho^2(K)=K$ with $\rho^2(E)=E_{1,2}\ne E$; K1--K3 verifiers on the window. \emph{Population:} the generators and a 12-label window; $\rho^k(E)$ for $k\le12$; $a\le3$, $|b|\le3$ for the normalisations. \emph{Controls:} positive: the relations reproduced as elements; negative: an altered relation ($\fq^{\pm2}$ in $EF$) is detected. \emph{Evidence:} n.a. (a test of the code against the text of the definition). \emph{Where:} \code{battery/checks_uq_a1d3.py}, \code{tests/test_uq_sl2_pbw.py}. \emph{Note:} the class carries the Casimir as a label coordinate over Z; K4/K5 live on the SQED\_2 realisation (next rows); the definition itself is formalised in ClusterLean (the local Lean 4 repository) (uq\_draft\_KqAlgebra): Def. 2.13's presentation with the basis K\^{}n, E\_(a,b), F\_(a,b), the bar fixing E, F, K and this $\rho$ is a flavoured K\_q-algebra over R(SU(2)) with a unique trace, through the isomorphism with SQED\_2 (E $\to$ u\_+, F $\to$ q\^{}-1 u\_- v\^{}-1, K $\to$ v); UqSL2PBW itself has not been compared with the Lean executable (UqSU2KAlgebra has, sec:uqsl2/trace). \emph{Run (web, 2026-09-22):} 13 checks, 0 failing; generators: E, F, K\^{}{$\pm$1}, $\chi$\_1; window: 12 labels; orbit: rho\^{}k(E), k $\le$ 12. & \CERTIFIED\newline {\footnotesize implementation, fast\newline web}\newline {\footnotesize run web 2026-09-22: 13/13} \\ \addlinespace
\code{sec:uqsl2/trace}\newline \peqref{eq:uqsl2trace} & The trace vanishes off the Cartan sector, and $\sum_n\Tr K^n x^n$ equals $(\fq^2;\fq^2)^2_\infty$ times $E_\fq(\mu x)\,E_\fq(\mu^{-1}x)$ $\times E_\fq(\mu x^{-1})\,E_\fq(\mu^{-1}x^{-1})$. & $G(x,\mu)$ expanded independently as a two-variable series ($E_\fq(y)=\prod_{k\ge0}(1+\fq^{2k+1}y)^{-1}$); the Wilson traces of SQED$_2$ (\code{GNAbeKAlgebra(u_n(1),(1,),nf=2)} after \code{base_change(un_to_sun_hom(2))}, characters expanded in $\mu$) against $[x^n]G$; symmetry of $G$ under $x\to x^{-1}$ and $\mu\to\mu^{-1}$; vanishing of the trace on the charged (magnetic) sector; measured $\Tr K^n=O(\fq^{|n|})$. \emph{Population:} $|n|\le6$ through $\fq^{12}$; 7 charged labels. \emph{Controls:} positive: the printed leading terms $\Tr K^0=1+\fq^2(\chi_2-1)+\dots$; negative: the opposite sign in $E_\fq$ is detected. \emph{Evidence:} emergent. \emph{Where:} \code{battery/checks_uq_a1d3.py}, \code{notes/trace_bootstrap.md (2026-06-13)}. \emph{Note:} the SQED\_2 realisation computes the trace from eq:measure at (U(1), 2 fundamentals), the draft's own derivation of eq:uqsl2trace; extends the 2026-09-21 check (n = 0, $\pm$1, 2 through q\^{}9) to |n| $\le$ 6 through q\^{}12; partly formalised in ClusterLean (the local Lean 4 repository): the trace vanishes off the Cartan sector and Tr K\^{}n = (q\^{}2;q\^{}2)\_$\infty$\^{}2 $\Sigma$ $\phi$\_a $\psi$\_b over a-b = n, with $\Psi$(q\^{}2 y) = (1 + q$\chi$\_1 y + q\^{}2 y\^{}2) $\Psi$(y) and $\phi$ = $\star$$\psi$, is the unique trace; the E\_q-product form is not stated in Lean.  The proven executable matched UqSU2KAlgebra and SQEDNfSampleKAlgebra(2) exactly for |n| $\le$ 6 through q\^{}16 (2026-09-24, Cluster f8543529). \emph{Run (web, 2026-09-22):} 6 checks, 0 failing; n: |n| $\le$ 6; K: 12; realisation: SQED\_2 as GNAbeKAlgebra(u\_n(1),(1,),nf=2) after base\_change(un\_to\_sun\_hom(2)). & \CERTIFIED\newline {\footnotesize identity, fast\newline web}\newline {\footnotesize run web 2026-09-22: 6/6} \\ \addlinespace
\code{sec:uqsl2/recursion}\newline \peqref{eq:uqsl2rec}, \peqref{eq:uqsl2con} & Cyclicity on $\Tr(K^nEF)$ gives the four-term recursion, whose iteration from $\Tr 1=1+O(\fq)$ has a unique solution order by order, and $G(x,\mu)$ is it. & both recursions on the class's traces; the paper's iteration --- from $\Tr1=[x^0]G$ and $\Tr K^n=O(\fq)$, order by order --- reproduces $[x^n]G$ for $1\le|n|\le4$ through $\fq^{10}$ (the unique solution); well-posedness by inspection of the coefficients. \emph{Population:} $0\le n\le4$ (\peqref{eq:uqsl2rec}) and $-4\le n\le-1$ (\peqref{eq:uqsl2con}) on the SQED$_2$ traces through $\fq^{10}$; the iteration on $|n|\le11$. \emph{Controls:} positive: the class's traces satisfy the recursion; negative: n.a. \emph{Evidence:} emergent. \emph{Where:} \code{battery/checks_uq_a1d3.py}, \code{notes/trace_bootstrap.md}. \emph{Note:} Formalised in ClusterLean (the local Lean 4 repository) except the last clause: every $\rho$\^{}2-twisted-cyclic trace satisfies the cyclicity recursion, T\_n = $\Sigma$ $\phi$\_a $\psi$\_b over a-b = n solves it for every n, and a twisted-cyclic orthonormal trace is unique up to 1+qR[[q]] by order raising (uq\_draft\_KqAlgebra); that the solution is G(x,$\mu$) in the E\_q-product form is checked numerically (sec:uqsl2/trace), not formalised. \emph{Run (web, 2026-09-22):} 4 checks, 0 failing; n: |n| $\le$ 4 (traces to |n| $\le$ 6); K: 10; iteration: |n| $\le$ 11. & \CERTIFIED\newline {\footnotesize theorem, fast\newline web}\newline {\footnotesize run web 2026-09-22: 4/4} \\ \addlinespace
\code{sec:uqsl2/ade}\newline Remark 2.14 & Some variant of $U_\fq(\fg)$ admits the structure of a $K_\fq$-algebra for every ADE $\fg$, and perhaps every reductive $\fg$. & the source repository's route through balanced linear quivers, whose higher-rank PBW engines exist there without traces. & \OPEN\newline {\footnotesize conjecture, none\newline no test} \\ \addlinespace
\code{sec:a1d3/algebra}\newline Def.~2.15 (Section~\pref{sec:a1d3}) & \code{A1D3KAlg} implements the punctured-triangle definition: the six relation families in $T_i,D_i$ with $\chi_1$, the $\fq$-commutations $T_iD_i=\fq^{-2}D_iT_i$, $T_iD_{i-1}=\fq^2D_{i-1}T_i$; the monomial basis $\fq^{ab}T_i^aD_i^b$, $\fq^{-ab}T_i^aD_{i-1}^b$ dressed by $\chi_k$; $\rho(T_i)=T_{i+1}$, $\rho(D_i)=D_{i+1}$ of order three fixing $\chi_k$; and it is an $SU(2)$-flavoured $K_\fq$-algebra. & every printed relation recomputed with \code{multiply} and compared exactly as elements; the monomial normalisations; $\chi_1\cdot L$ is the dressed label, $\chi_1\chi_1=\chi_0+\chi_2$; $\rho$ and its order; F1--F5 verifiers on the window; the presentations of \code{a1d3_object} one algebra (witnesses on every edge). \emph{Population:} all $i$; $a,b\le3$ for the monomials; a 55-label window ($a+b\le2$, $\chi$-dressing $0,1$) for the verifiers at $K=6$. \emph{Controls:} positive: the relations reproduced on all $i$; negative: an altered relation detected. \emph{Evidence:} n.a. (a test of the code against the text of the definition); F1-F5 emergent on the hand-built presentation. \emph{Where:} \code{battery/checks_uq_a1d3.py}, \code{tests/test_a1d3_kalg.py}, \code{src/bps/a1d3_object.py}. \emph{Run (web, 2026-09-22):} 10 checks, 0 failing; generators: 6; window: 38 labels; K: 6. & \CERTIFIED\newline {\footnotesize implementation, fast\newline web}\newline {\footnotesize run web 2026-09-22: 10/10} \\ \addlinespace
\code{sec:a1d3/traces}\newline \peqref{eq:a1d3trace}, Remark 2.16 & The three elementary traces over the Weyl--Kac denominator; the recursions fix the trace uniquely up to scale. & the three sums over the Weyl--Kac denominator recomputed independently as series in $\fq$ with $\mu$-Laurent coefficients, against the class's traces (characters expanded in $\mu$); $\bZ_3$ invariance; the sharp leading constraint $\Tr T_0=-\fq+O(\fq^3)$, $[\fq]\Tr D_0=-\chi_1$. \emph{Population:} $\Tr1$, $\Tr T_0$, $\Tr D_0$ and their $\bZ_3$ images through $\fq^{14}$. \emph{Controls:} positive: the printed expansion $1+\chi_2\fq^2+(\chi_0+\chi_2+\chi_4)\fq^4$; negative: a factor of the denominator omitted is detected. \emph{Evidence:} emergent. \emph{Where:} \code{battery/checks_uq_a1d3.py}, \code{tests/test_a1d3_kalg.py}. \emph{Note:} a first version of the independent sum for Tr T\_0 stopped one term early (the three exponents are not monotone in m); the class was right. \emph{Run (web, 2026-09-22):} 7 checks, 0 failing; K: 14; traces: 1, T\_0, D\_0 (and the Z\_3 images). & \CERTIFIED\newline {\footnotesize identity, fast\newline web}\newline {\footnotesize run web 2026-09-22: 7/7} \\ \addlinespace
\code{rem:a1d3-miracle}\newline Rem.~\pref{rem:a1d3-miracle} & The minors of the $2\times3$ constraint matrix, divided by $(\mu^2\fq^2;\fq^2)_\infty(\mu^{-2}\fq^2;\fq^2)_\infty$, reproduce $\Tr 1,\Tr T_0,\Tr D_0$. & the two constraint families $\Tr T_0^a=A_a\Tr1+B_a\Tr T_0+C_a\Tr D_0$ and $\Tr(T_0^aD_0)=A'_a\Tr1+B'_a\Tr T_0+C'_a\Tr D_0$ read off the class's layer-1 reduction (\code{trace_layer1}), each row normalised by its minimal $\fq$-power; the three $2\times2$ minors (Cramer's ray), divided by $(\mu^2\fq^2;\fq^2)_\infty(\mu^{-2}\fq^2;\fq^2)_\infty$ and normalised to start at $+1$, against $(\Tr1,\Tr T_0,\Tr D_0)$: the first differing power per component as $a$ grows. \emph{Population:} $a=1,\dots,8$; $K=24$. \emph{Controls:} positive: the $a=2$ rows are the trace-recursions note's table rows; negative: n.a. (a measured convergence). \emph{Evidence:} emergent. \emph{Where:} \code{battery/checks_uq_a1d3.py}. \emph{Note:} MEASURED 2026-09-22 for the first time (the cited a1d3\_trace\_recursions.md says it does not address the minors; the trace-miracle puzzle listed only the heptagon and nonagon cases): the minors approach (Tr 1, Tr T\_0, Tr D\_0) with first differing powers (2a, 2a+1, 2a+1) for a $\ge$ 3, up to (16, 17, 17) at a = 8 -- the same stabilisation shape as the pentagon's A\_n, B\_n. \emph{Run (web, 2026-09-22):} 3 checks, 0 failing; a: 1..8; K: 24. & \MEASURED\newline {\footnotesize conjecture, fast\newline web}\newline {\footnotesize run web 2026-09-22: 3/3} \\ \addlinespace
\end{longtable}

\subsection*{Section~\pref{sec:coulomb}: Coulomb branch algebras}
\begin{longtable}{@{}>{\raggedright\arraybackslash}p{2.4cm}>{\raggedright\arraybackslash}p{4.7cm}>{\raggedright\arraybackslash}p{4.95cm}>{\raggedright\arraybackslash}p{1.85cm}@{}}
\toprule claim & statement & test: procedure, population, controls & standing \\ \midrule \endfirsthead
\toprule claim & statement & test: procedure, population, controls & standing \\ \midrule \endhead
\bottomrule \endfoot
\code{sec:coulomb/labels}\newline Section~\pref{sec:coulomb}, the labelling & The canonical basis of $K_\fq[\fg,M]$ is labelled by $(m,e)\in(\Lambda_m\times\Lambda_e)/W$, $m$ dominant and $e$ dominant for the Levi of $m$. & every window label lies in the fundamental domain ($m$ dominant, $e$ dominant for the Levi of $m$); every Weyl image of every raw $(m,e)$ of the window folds to the same label; the window's elements are linearly independent, by an exact rank at a rational point (residuals evaluated at $\fq=3/11$, $v$ at rational points, flavour fugacities at generic rationals: a nonzero maximal minor there is a nonzero rational function, so the rank holds over the coefficient ring).  That every admitted label builds, and that products decompose in the family, are the rows conj:abeKalgebra/existence-uniqueness and conj:abeKalgebra/consistency. \emph{Population:} the nine conjecture windows (\code{conjecture_windows}): pure SU(2), SO(3) ($m\le2$, $|e|\le2$), SU(3), Sp(4), $G_2$; SU(2)+1, SU(2)+2, SO(3)+adjoint, SU(3)+1 (flavour 0); the extensive tier, recorded beside the fast one, runs the 24 theories of the Section~3 lineup on grown windows (\code{_extensive_windows}): pure SU(2) and SO(3) ($m\le4$, $|e|\le3$), the rank-one third global form, SU(3), PSU(3), Sp(4), Spin(5), SO(5), $G_2$; SU(2) with $N_f=1,\dots,4$; $N=2^*$ at SU(2), SO(3), the third global form and SU(3); SU(3) with a symmetric, alone and with a fundamental; Spin(5) and SO(5) with a vector; $G_2$ with a $\mathbf 7$; the SU(2)$\times$SU(2) quiver with a bifundamental; U(1)--U(2) with three fundamentals of U(2) (matter at the neutral flavour label). \emph{Controls:} negative: every window has raw pairs with a nontrivial Weyl orbit, so the identity labelling would fail the invariance test; the relation $x_0+2x_1$ appended to the rank test does not raise the rank. \emph{Evidence:} emergent (closure). \emph{Where:} \code{battery/checks_coulomb.py (check_labels)}, \code{tests/test_pure_g_abe_kalgebra.py}, \code{tests/test_g_matter_abe_kalgebra.py}. \emph{Note:} Own adapter since 2026-09-23; before that the row shared the consistency run, which did not test linear independence.  The rank harness was corrected on 2026-09-23: it evaluated too few points per cell where several lines share one (every Wilson line lives on the cell 0), which caps the rank; it now uses as many points as lines (a full rank is a certificate either way, a deficient one is not). \emph{Run (web, 2026-09-23):} 27 checks, 0 failing; pure SU(2): {'labels': 13, 'weyl\_images': 30, 'nontrivial\_orbits': 14, 'rank': 13, 'seconds': 0.2}; pure SO(3): {'labels': 13, 'weyl\_images': 30, 'nontrivial\_orbits': 14, 'rank': 13, 'seconds': 0.0}; pure SU(3): {'labels': 8, 'weyl\_images': 54, 'nontrivial\_orbits': 8, 'rank': 8, 'seconds': 0.3}; pure Sp(4): {'labels': 5, 'weyl\_images': 48, 'nontrivial\_orbits': 5, 'rank': 5, 'seconds': 0.2}; pure G2: {'labels': 4, 'weyl\_images': 48, 'nontrivial\_orbits': 3, 'rank': 4, 'seconds': 13.5}; SU(2) + 1 fund. (roster su2-nf1): {'labels': 5, 'weyl\_images': 12, 'nontrivial\_orbits': 5, 'rank': 5, 'seconds': 0.1}; SU(2) + 2 fund. (roster su2-nf2, U(2) flavour): {'labels': 5, 'weyl\_images': 12, 'nontrivial\_orbits': 5, 'rank': 5, 'seconds': 0.1}; SO(3) + adjoint (N=2*): {'labels': 5, 'weyl\_images': 12, 'nontrivial\_orbits': 5, 'rank': 5, 'seconds': 0.0}; SU(3) + 1 fund. (roster su3-nf1): {'labels': 4, 'weyl\_images': 24, 'nontrivial\_orbits': 3, 'rank': 4, 'seconds': 0.7}. \emph{Run (web (extensive), 2026-09-24):} 72 checks, 0 failing; pure SU(2): {'labels': 32, 'unreached': [], 'errors': {}, 'weyl\_images': 70, 'nontrivial\_orbits': 34, 'rank': 32, 'seconds': 12.1}; pure SO(3): {'labels': 32, 'unreached': [], 'errors': {}, 'weyl\_images': 70, 'nontrivial\_orbits': 34, 'rank': 32, 'seconds': 1.2}; pure rank 1, the third global form: {'labels': 12, 'unreached': [], 'errors': {}, 'weyl\_images': 50, 'nontrivial\_orbits': 24, 'rank': 12, 'seconds': 0.2}; pure SU(3): {'labels': 16, 'unreached': [], 'errors': {}, 'weyl\_images': 108, 'nontrivial\_orbits': 17, 'rank': 16, 'seconds': 34.5}; pure PSU(3): {'labels': 16, 'unreached': [], 'errors': {}, 'weyl\_images': 96, 'nontrivial\_orbits': 15, 'rank': 16, 'seconds': 0.5}; pure Sp(4): {'labels': 19, 'unreached': [], 'errors': {}, 'weyl\_images': 240, 'nontrivial\_orbits': 29, 'rank': 19, 'seconds': 800.9}; pure Spin(5): {'labels': 25, 'unreached': [], 'errors': {}, 'weyl\_images': 240, 'nontrivial\_orbits': 29, 'rank': 25, 'seconds': 322.6}; pure SO(5): {'labels': 16, 'unreached': [], 'errors': {}, 'weyl\_images': 128, 'nontrivial\_orbits': 15, 'rank': 16, 'seconds': 2.4}; pure G2: {'labels': 12, 'unreached': [], 'errors': {}, 'weyl\_images': 144, 'nontrivial\_orbits': 11, 'rank': 12, 'seconds': 263.2}; SU(2)+1: {'labels': 18, 'unreached': [], 'errors': {}, 'weyl\_images': 40, 'nontrivial\_orbits': 19, 'rank': 18, 'seconds': 7.9}; SU(2)+2: {'labels': 13, 'unreached': [], 'errors': {}, 'weyl\_images': 30, 'nontrivial\_orbits': 14, 'rank': 13, 'seconds': 2.3}; SU(2)+3: {'labels': 8, 'unreached': [], 'errors': {}, 'weyl\_images': 18, 'nontrivial\_orbits': 8, 'rank': 8, 'seconds': 4.7}; SU(2)+4: {'labels': 8, 'unreached': [], 'errors': {}, 'weyl\_images': 18, 'nontrivial\_orbits': 8, 'rank': 8, 'seconds': 15.8}; SU(2)+adjoint (N=2*): {'labels': 13, 'unreached': [], 'errors': {}, 'weyl\_images': 30, 'nontrivial\_orbits': 14, 'rank': 13, 'seconds': 1.7}; SO(3)+adjoint (N=2*): {'labels': 11, 'unreached': [], 'errors': {}, 'weyl\_images': 24, 'nontrivial\_orbits': 11, 'rank': 11, 'seconds': 0.7}; rank 1 + adjoint (N=2*), the third global form: {'labels': 6, 'unreached': [], 'errors': {}, 'weyl\_images': 24, 'nontrivial\_orbits': 11, 'rank': 6, 'seconds': 0.1}; SU(3)+adjoint (N=2*): {'labels': 6, 'unreached': [], 'errors': {}, 'weyl\_images': 36, 'nontrivial\_orbits': 5, 'rank': 6, 'seconds': 2.2}; SU(3)+symmetric: {'labels': 10, 'unreached': [], 'errors': {}, 'weyl\_images': 72, 'nontrivial\_orbits': 11, 'rank': 10, 'seconds': 3.9}; SU(3)+symmetric+1: {'labels': 6, 'unreached': [], 'errors': {}, 'weyl\_images': 36, 'nontrivial\_orbits': 5, 'rank': 6, 'seconds': 4.5}; Spin(5)+vector: {'labels': 9, 'unreached': [], 'errors': {}, 'weyl\_images': 72, 'nontrivial\_orbits': 8, 'rank': 9, 'seconds': 14.9}; SO(5)+vector: {'labels': 12, 'unreached': [], 'errors': {}, 'weyl\_images': 96, 'nontrivial\_orbits': 11, 'rank': 12, 'seconds': 16.7}; G2+7: {'labels': 6, 'unreached': [], 'errors': {}, 'weyl\_images': 72, 'nontrivial\_orbits': 5, 'rank': 6, 'seconds': 216.4}; SU(2)xSU(2)+bifundamental: {'labels': 16, 'unreached': [], 'errors': {}, 'weyl\_images': 64, 'nontrivial\_orbits': 15, 'rank': 16, 'seconds': 4.4}; U(1)-U(2)+3: {'labels': 24, 'unreached': [], 'errors': {}, 'weyl\_images': 48, 'nontrivial\_orbits': 12, 'rank': 24, 'seconds': 0.3}. & \CERTIFIED\newline {\footnotesize conjecture, extensive\newline web}\newline {\footnotesize run web 2026-09-23: 27/27}\newline {\footnotesize run web (extensive) 2026-09-24: 72/72} \\ \addlinespace
\code{sec:coulomb/global-forms}\newline Section~\pref{sec:coulomb}, global forms & A global form is a Lagrangian sublattice $\Lambda_G\subset\Lambda^w_m\times\Lambda^w_e$ for the Dirac pairing; conventional forms are $\Lambda^G_m\times\Lambda^G_e$; $K_\fq[\fg,M]\subset K_\fq[G,M]$; three forms for $\fsl_2$. & \code{LineLattice}: SU(2) ($m\in Q^\vee$, $e\in P$) and SO(3) ($m\in P^\vee$, $e\in Q$) admit their labels, are Dirac-integral and maximal; the three order-two subgroups of the centre-class pairs $(\bZ/2)^2$ give the three maximal Dirac-integral lattices (the full group fails integrality); the third form admits exactly $e-m$ even (half-integral coroot coordinates on the coweight torus) and is mutually local.  Census: the subgroups \code{LineLattice} reports Dirac-integral (hence Lagrangian) number exactly the independent count (a cyclic $p$-part $\bZ/p^a$ gives $1+p+\dots+p^a$, and $(\bZ/p)^n$ gives $\prod_{i=1}^n(p^i+1)$), and the product (conventional) forms among them are one per subgroup of the centre. \emph{Population:} every order-$|Z|$ subgroup of the centre-class pairs $Z\times\hat Z$ for every simply connected root datum of rank $\le4$ in the repository (SU(2)--SU(5), Sp(4)--Sp(8), Spin(5)--Spin(9), $G_2$, $D_4$) and the products SU(2)$^2$, SU(2)$\times$SU(3), SU(3)$^2$, SU(2)$^3$; the rank-one forms on labels $|m|,|e|\le3$. \emph{Evidence:} emergent. \emph{Where:} \code{battery/checks_coulomb.py}, \code{tests/test_global_form.py}. \emph{Note:} The inclusion of K\_q[g,M] in K\_q[G,M] is not checked by this row beyond the lattice level (every form contains the zero class pair, hence the coroot lattice times the root lattice). \emph{Run (web, 2026-09-23):} 8 checks, 0 failing; forms: SU(2), SO(3), the third form; window: |m|, |e| $\le$ 3; census: {'SU(2)': {'centre': [2], 'order-|Z| subgroups': 3, 'Lagrangian': 3, 'formula': 3, 'product forms': 2, 'subgroups of Z': 2}, 'SU(3)': {'centre': [3], 'order-|Z| subgroups': 4, 'Lagrangian': 4, 'formula': 4, 'product forms': 2, 'subgroups of Z': 2}, 'SU(4)': {'centre': [4], 'order-|Z| subgroups': 7, 'Lagrangian': 7, 'formula': 7, 'product forms': 3, 'subgroups of Z': 3}, 'SU(5)': {'centre': [5], 'order-|Z| subgroups': 6, 'Lagrangian': 6, 'formula': 6, 'product forms': 2, 'subgroups of Z': 2}, 'Sp(4)': {'centre': [2], 'order-|Z| subgroups': 3, 'Lagrangian': 3, 'formula': 3, 'product forms': 2, 'subgroups of Z': 2}, 'Sp(6)': {'centre': [2], 'order-|Z| subgroups': 3, 'Lagrangian': 3, 'formula': 3, 'product forms': 2, 'subgroups of Z': 2}, 'Sp(8)': {'centre': [2], 'order-|Z| subgroups': 3, 'Lagrangian': 3, 'formula': 3, 'product forms': 2, 'subgroups of Z': 2}, 'Spin(5)': {'centre': [2], 'order-|Z| subgroups': 3, 'Lagrangian': 3, 'formula': 3, 'product forms': 2, 'subgroups of Z': 2}, 'Spin(7)': {'centre': [2], 'order-|Z| subgroups': 3, 'Lagrangian': 3, 'formula': 3, 'product forms': 2, 'subgroups of Z': 2}, 'Spin(9)': {'centre': [2], 'order-|Z| subgroups': 3, 'Lagrangian': 3, 'formula': 3, 'product forms': 2, 'subgroups of Z': 2}, 'G2': {'centre': [], 'order-|Z| subgroups': 1, 'Lagrangian': 1, 'formula': 1, 'product forms': 1, 'subgroups of Z': 1}, 'D4': {'centre': [2, 2], 'order-|Z| subgroups': 35, 'Lagrangian': 15, 'formula': 15, 'product forms': 5, 'subgroups of Z': 5}, 'SU(2)xSU(2)': {'centre': [2, 2], 'order-|Z| subgroups': 35, 'Lagrangian': 15, 'formula': 15, 'product forms': 5, 'subgroups of Z': 5}, 'SU(2)xSU(3)': {'centre': [6], 'order-|Z| subgroups': 12, 'Lagrangian': 12, 'formula': 12, 'product forms': 4, 'subgroups of Z': 4}, 'SU(3)xSU(3)': {'centre': [3, 3], 'order-|Z| subgroups': 130, 'Lagrangian': 40, 'formula': 40, 'product forms': 6, 'subgroups of Z': 6}, 'SU(2)\^{}3': {'centre': [2, 2, 2], 'order-|Z| subgroups': 1395, 'Lagrangian': 135, 'formula': 135, 'product forms': 16, 'subgroups of Z': 16}}. & \CERTIFIED\newline {\footnotesize theorem, fast\newline web}\newline {\footnotesize run web 2026-09-23: 8/8} \\ \addlinespace
\code{sec:coulomb/flavour-enhancement}\newline Section~\pref{sec:coulomb}, flavour enhancement & $n$ copies of an irrep enhance $U(1)^n$ to $U(n)$; real / pseudo-real irreps to $USp(2n)$ / $SO(2n)$; $R\oplus R^\vee$ similarly. & the algebra built with its manifest $R(U(N_f))$ flavour (\code{GNAbeKAlgebra(su_2(), (1,), nf)}); the traces and pairings of the window sections through $\fq^4$ and the structure constants of their products (\code{to_R_form}), each coefficient taken to its weight diagram (\code{UNZPlusRing.to_abelian}) and twisted by the power of $\det$ that $\rho$ fixes: $\rho$ shifts the flavour charge by $\kappa_f(m)=\det^{-m}$, so the section at magnetic charge $m$ is $\det^{m/2}$ times a Spin($2N_f$) section, and a trace is twisted by $\det^{-m/2}$, a pairing $I_{a,b}$ by $\det^{(m_a-m_b)/2}$, a structure constant $c_{ab}^c$ by $\det^{(m_c-m_a-m_b)/2}$ (half-integral at odd $m$, where the lifted weights are spinor weights).  Each twisted coefficient is $W(D_{N_f})$-invariant, i.e. lifts to $R(\mathrm{Spin}(2N_f))$ (\code{so2nf_characters}), with an exact round trip.  The real case (\code{_usp_one}): at $N=2^*$ the manifest flavour is $R(U(1))$, and $\rho$ shifts the flavour charge by $\kappa_f(m)=\mu^{-2m}$ (SU(2) units), so the section at $m$ is $\mu^{m}$ times a $USp(2)$ section: a trace is twisted by $\mu^{-m}$, a pairing by $\mu^{m_a-m_b}$, a structure constant by $\mu^{m_c-m_a-m_b}$, and each twisted coefficient must be integral and invariant under $\mu\mapsto\mu^{-1}$, i.e. a $USp(2)$ character.  Every twist on nonzero data is integral, the SO(3) form's half-integral $m$ included: a trace is read at the zero cell, and a pairing or a structure constant is nonzero only when the magnetic charges agree modulo the coroot lattice. \emph{Population:} SU(2) with $N_f$ doublets (the pseudo-real case, $U(N_f)\to SO(2N_f)$): $N_f=2,3$ at the fast depth, $N_f=2,3,4$ at the extensive depth; sections with $m\le2$, $e\in\{0,1\}$ at $N_f=2$, fewer at $N_f\ge3$ (\code{_flavour_windows}); $N_f=1$ has nothing to test ($SO(2)=U(1)$).  The real case, $U(1)\to USp(2)$: $N=2^*$ (one adjoint) at the SU(2) form ($m\le2$, $e\in\{0,1\}$, products with $m_a+m_b\le2$; at the extensive depth $m\le3$, $m_a+m_b\le3$) and the SO(3) form ($m\le1$ in SU(2) units, $m=\tfrac12$ included, $e\in\{0,2\}$, $m_a+m_b\le1$; at the extensive depth $m\le\tfrac32$, $m_a+m_b\le2$).  Not yet run: the $U(n)$ enhancement for $n$ copies at other groups, the $USp(2n)$ case at $n\ge2$, Sp(4) with fundamentals, $R\oplus R^\vee$. \emph{Controls:} positive: $\mathrm{Tr}(1)$ at $\fq^2$, the flavour currents, re-expands as the SO($2N_f$) adjoint; negative: the odd-$m$ traces without the twist are not $W(D_{N_f})$-invariant; and the test rejects a theory that must not enhance: U(1) with three charge-1 hypers (a complex representation, the same $R(U(3))$) fails $W(D_3)$-invariance of its Schur index. \emph{Evidence:} emergent. \emph{Where:} \code{battery/checks_coulomb.py (check_flavour_enhancement, _usp_one)}, \code{src/abe/so2nf_characters.py}, \code{tests/test_so2nf_characters.py}, \code{tests/test_n2star_usp2_enhancement.py}. \emph{Note:} The recorded follow-up (in the source repository), re-derived on the live tier 2026-09-24 (the author approved taking it up).  The twist is the one the legacy route used (spinor shift m/2), derived here from $\kappa$\_f.  At N\_f = 2, W(D\_2)-invariance of a U(2) character is charge-conjugation symmetry about the twisted centre, so the discriminating cases are N\_f $\ge$ 3.  The extra reflection of O(2N\_f) (a single sign flip) is recorded in the population, not claimed.  The real case was added 2026-09-26; positive control: Tr(1) at $\fq$\^{}2 is the USp(2) adjoint $\mu$\^{}-2 + 1 + $\mu$\^{}2; negative control: the traces with nonzero m, without the twist. \emph{Run (web, 2026-09-26):} 21 checks, 0 failing; SU(2) + 2 doublets: {'sections': ['((0,), (0,))', '((0,), (1,))', '((1,), (0,))', '((1,), (1,))', '((2,), (0,))', '((2,), (1,))'], 'K': 4, 'max\_product\_m': 3, 'not\_reached': [], 'traces': {'objects': 6, 'orders': 12, 'not\_invariant': [], 'round\_trip\_failures': [], 'also\_single\_flip\_invariant': 6}, 'pairings': {'objects': 36, 'orders': 62, 'not\_invariant': [], 'round\_trip\_failures': [], 'also\_single\_flip\_invariant': 22}, 'products': {'objects': 18, 'orders': 29, 'not\_invariant': [], 'round\_trip\_failures': [], 'also\_single\_flip\_invariant': 24}, 'unshifted\_odd\_trace\_orders\_failing': 4, 'seconds': 288.6}; SU(2) + 3 doublets: {'sections': ['((0,), (0,))', '((1,), (0,))'], 'K': 4, 'max\_product\_m': 2, 'not\_reached': [], 'traces': {'objects': 2, 'orders': 5, 'not\_invariant': [], 'round\_trip\_failures': [], 'also\_single\_flip\_invariant': 3}, 'pairings': {'objects': 4, 'orders': 10, 'not\_invariant': [], 'round\_trip\_failures': [], 'also\_single\_flip\_invariant': 6}, 'products': {'objects': 3, 'orders': 3, 'not\_invariant': [], 'round\_trip\_failures': [], 'also\_single\_flip\_invariant': 3}, 'unshifted\_odd\_trace\_orders\_failing': 2, 'seconds': 149.5}; SU(2) + 1 doublet: SO(2) = U(1): no Weyl reflection beyond the manifest U(1), so nothing to test; SU(2) + adjoint (N=2*), USp(2): {'sections': ['((0,), (0,))', '((0,), (1,))', '((1,), (0,))', '((1,), (1,))', '((2,), (0,))', '((2,), (1,))'], 'K': 4, 'max\_product\_m': '2', 'not\_reached': [], 'traces': {'objects': 6, 'orders': 11, 'not\_invariant': []}, 'pairings': {'objects': 36, 'orders': 63, 'not\_invariant': []}, 'products': {'objects': 14, 'orders': 33, 'not\_invariant': []}, 'untwisted\_trace\_orders\_failing': '7 of 7', 'seconds': 22.2}; SO(3) + adjoint (N=2*), USp(2): {'sections': ['((0,), (0,))', '((0,), (2,))', '((Fraction(1, 2),), (0,))', '((Fraction(1, 2),), (2,))', '((1,), (0,))', '((1,), (2,))'], 'K': 4, 'max\_product\_m': '1', 'not\_reached': [], 'traces': {'objects': 6, 'orders': 16, 'not\_invariant': []}, 'pairings': {'objects': 36, 'orders': 75, 'not\_invariant': []}, 'products': {'objects': 14, 'orders': 36, 'not\_invariant': []}, 'untwisted\_trace\_orders\_failing': '8 of 8', 'seconds': 8.3}; U(1) + 3 (negative control): {'orders\_not\_invariant': [2, 4]}. \emph{Run (web (extensive), 2026-09-26):} 26 checks, 0 failing; SU(2) + 2 doublets: {'sections': ['((0,), (0,))', '((0,), (1,))', '((1,), (0,))', '((1,), (1,))', '((2,), (0,))', '((2,), (1,))'], 'K': 4, 'max\_product\_m': 3, 'not\_reached': [], 'traces': {'objects': 6, 'orders': 12, 'not\_invariant': [], 'round\_trip\_failures': [], 'also\_single\_flip\_invariant': 6}, 'pairings': {'objects': 36, 'orders': 62, 'not\_invariant': [], 'round\_trip\_failures': [], 'also\_single\_flip\_invariant': 22}, 'products': {'objects': 18, 'orders': 29, 'not\_invariant': [], 'round\_trip\_failures': [], 'also\_single\_flip\_invariant': 24}, 'unshifted\_odd\_trace\_orders\_failing': 4, 'seconds': 285.4}; SU(2) + 3 doublets: {'sections': ['((0,), (0,))', '((0,), (1,))', '((1,), (0,))', '((1,), (1,))', '((2,), (0,))', '((2,), (1,))'], 'K': 4, 'max\_product\_m': 2, 'not\_reached': ['I(((2,), (0,)), ((2,), (0,)))', 'I(((2,), (0,)), ((2,), (1,)))', 'I(((2,), (1,)), ((2,), (0,)))', 'I(((2,), (1,)), ((2,), (1,)))'], 'traces': {'objects': 6, 'orders': 12, 'not\_invariant': [], 'round\_trip\_failures': [], 'also\_single\_flip\_invariant': 6}, 'pairings': {'objects': 32, 'orders': 56, 'not\_invariant': [], 'round\_trip\_failures': [], 'also\_single\_flip\_invariant': 20}, 'products': {'objects': 14, 'orders': 28, 'not\_invariant': [], 'round\_trip\_failures': [], 'also\_single\_flip\_invariant': 23}, 'unshifted\_odd\_trace\_orders\_failing': 4, 'seconds': 6247.6}; SU(2) + 4 doublets: {'sections': ['((0,), (0,))', '((0,), (1,))', '((1,), (0,))', '((1,), (1,))'], 'K': 4, 'max\_product\_m': 2, 'not\_reached': ['I(((1,), (0,)), ((1,), (0,)))', 'I(((1,), (0,)), ((1,), (1,)))', 'I(((1,), (1,)), ((1,), (0,)))', 'I(((1,), (1,)), ((1,), (1,)))'], 'traces': {'objects': 4, 'orders': 9, 'not\_invariant': [], 'round\_trip\_failures': [], 'also\_single\_flip\_invariant': 5}, 'pairings': {'objects': 12, 'orders': 26, 'not\_invariant': [], 'round\_trip\_failures': [], 'also\_single\_flip\_invariant': 10}, 'products': {'objects': 10, 'orders': 19, 'not\_invariant': [], 'round\_trip\_failures': [], 'also\_single\_flip\_invariant': 17}, 'unshifted\_odd\_trace\_orders\_failing': 4, 'seconds': 4519.2}; SU(2) + 1 doublet: SO(2) = U(1): no Weyl reflection beyond the manifest U(1), so nothing to test; SU(2) + adjoint (N=2*), USp(2): {'sections': ['((0,), (0,))', '((0,), (1,))', '((1,), (0,))', '((1,), (1,))', '((2,), (0,))', '((2,), (1,))', '((3,), (0,))', '((3,), (1,))'], 'K': 4, 'max\_product\_m': '3', 'not\_reached': [], 'traces': {'objects': 8, 'orders': 13, 'not\_invariant': []}, 'pairings': {'objects': 64, 'orders': 99, 'not\_invariant': []}, 'products': {'objects': 22, 'orders': 65, 'not\_invariant': []}, 'untwisted\_trace\_orders\_failing': '9 of 9', 'seconds': 139.3}; SO(3) + adjoint (N=2*), USp(2): {'sections': ['((0,), (0,))', '((0,), (2,))', '((Fraction(1, 2),), (0,))', '((Fraction(1, 2),), (2,))', '((1,), (0,))', '((1,), (2,))', '((Fraction(3, 2),), (0,))', '((Fraction(3, 2),), (2,))'], 'K': 4, 'max\_product\_m': '2', 'not\_reached': [], 'traces': {'objects': 8, 'orders': 16, 'not\_invariant': []}, 'pairings': {'objects': 64, 'orders': 109, 'not\_invariant': []}, 'products': {'objects': 29, 'orders': 108, 'not\_invariant': []}, 'untwisted\_trace\_orders\_failing': '8 of 8', 'seconds': 92.9}; U(1) + 3 (negative control): {'orders\_not\_invariant': [2, 4]}. & \MEASURED\newline {\footnotesize conjecture, extensive\newline web}\newline {\footnotesize run web 2026-09-26: 21/21}\newline {\footnotesize run web (extensive) 2026-09-26: 26/26} \\ \addlinespace
\code{eq:rho-witten}\newline \peqref{eq:rho-witten} & \code{rho} on labels is $(m,e)\mapsto(-m,\kappa(m)-e)$ with $\kappa(m)=\sum_{\langle\alpha,m\rangle>0}\langle\alpha,m\rangle\alpha-\sum_{w\in N:\langle w,m\rangle>0}\langle w,m\rangle w$, the flavour shifted by $\kappa_f(m)$ in the diagonal $U(1)$ of each $U(n_f)$ (as $f\mapsto f^\vee\otimes\det^{-D}$); on the torus $(\rho f)_{-m}(v,\mu)=v^{\kappa(m)}\mu^{\kappa_f(m)}f_m(v^{-1},\mu^{-1})$. & $\kappa$ recomputed from the datum's roots and the representation's weights with the datum's pairing (the matter sum over every copy of $N$; the $\det^{-D}$ twist per copy), compared with \code{rho} after Weyl canonicalisation; the torus form on the residual vectors of the charts. \emph{Population:} pure: SU(2), SO(3), SU(3), Sp(4), $G_2$ on windows of $(m,e)$ (155 labels); with matter: SO(3)+adjoint, SU(2)+1, SU(2)+2 (U(2) flavour), SU(3)+1 (261 labels); \peqref{eq:rhotorus} on five SU(2) lines. \emph{Controls:} positive: agreement on every label; negative: dropping the matter term of $\kappa$ is detected. \emph{Evidence:} n.a. (a test of the code against the text). \emph{Where:} \code{battery/checks_coulomb.py}, \code{tests/test_abe_rho_label.py}. \emph{Run (web, 2026-09-23):} 8 checks, 0 failing; pure labels: 155; matter labels: 261; torus lines: 5; torus lines with $\mu$: 4; torus products: 41. & \CERTIFIED\newline {\footnotesize implementation, fast\newline web}\newline {\footnotesize run web 2026-09-23: 8/8} \\ \addlinespace
\code{sec:coulomb/presentation}\newline Section~\pref{sec:coulomb}, the abelianized presentation & The rational quantum torus classes implement the presentation as printed: the vector cocycle \peqref{eq:ccclosed} (and its bar image for $A<0<B$), the product over positive roots \peqref{eq:ccprod} with the matter factor $(-\fq^{r+1}x;\fq^2)_p$, the 2-cocycle identity behind \peqref{eq:cocydef}, and the product law \peqref{eq:fgprod}. & \code{CC(a,b)} and \code{CC_N} against the printed closed forms as exact rational functions (cross-multiplied numerators against the denominator factors); the printed closed form checked to satisfy $\mathrm{cocy}_{a,b}(\fq^{a+b}v)\,\mathrm{cocy}_{a+b,c}(\fq^{a+b+c}v)=\mathrm{cocy}_{b,c}(\fq^{2a+b+c}v)\,\mathrm{cocy}_{a,b+c}(\fq^{a+b+c}v)$; \peqref{eq:fgprod} evaluated with the printed cocycle on the residuals of $L_{2,0},L_{2,1},L_{0,1},L_{0,2}$ against the class's products.  Widened: the same comparisons at the non-simply-laced groups; the matter cocycle compared level by level in $\mu$ (\code{CC_N} against \code{CC} times the printed $\cocy^{\rm hyp}$ factors over the weights of $N$); \peqref{eq:fgprod} with the printed cocycle against class products at rank two and with matter; and the bar property $\overline{\cocy_{a,b}}=\cocy_{b,a}$, which makes the cellwise bar an anti-automorphism, with $\overline{xy}=\bar y\,\bar x$ checked on torus products. \emph{Population:} SU(2): 49 pairs $(a,b)$, $|a|,|b|\le3$; SU(3): 81 pairs; the 2-cocycle identity on 125 SU(2) and 729 SU(3) triples; the matter cocycle on 50 pairs (SO(3)+adjoint, SU(2)+fundamental); five products of SU(2) lines in every magnetic sector.  Widened 2026-09-23: the vector cocycle at Sp(4), Spin(5), $G_2$, SU(2)$\times$SU(2) (324 pairs) and its 2-cocycle identity at Sp(4), Spin(5), $G_2$ (2187 triples); the matter cocycle at rank two on 486 pairs (SU(3) with a fundamental, a symmetric, both; Spin(5)+vector; $G_2$+7; SU(2)$\times$SU(2) with a bifundamental); products at SU(3), Spin(5) and SU(2)+1 ($G_2$ on the local tier); the bar property on 243 pairs and on the torus products of SU(3) and SU(3)+1 lines. \emph{Controls:} positive: the class's products reproduced from the printed formula; negative: an altered cocycle detected; the weights of $N^\vee$ in place of $N$ give a different matter cocycle on the complex representations (for the real ones, Spin(5)+vector and $G_2$+7, the two coincide); the multiplicative version $\overline{xy}=\bar x\,\bar y$ fails. \emph{Evidence:} n.a. (a test of the code against the text); the 2-cocycle identity of the printed form is a theorem checked on windows. \emph{Where:} \code{battery/checks_coulomb.py}, \code{src/abe/wrq_torus.py}, \code{src/abe/matter_wrq_torus.py}. \emph{Note:} the planned 'explicit 2-cocycle identity' verifier is realised here on the printed closed form; a first version of the exact arithmetic mis-built (1 - v\^{}alpha) at k = 0 and was caught by eq:fgprod at m = $\pm$1 of L\_20\^{}2, where the (1 - v\^{}2) poles of the two terms cancel. \emph{Run (web, 2026-09-23):} 13 checks, 0 failing; SU(2) pairs: 49; SU(3) pairs: 81; cocycle triples: 854; matter pairs: 50; products: 5; widened (2026-09-23): {'vector cocycle pairs': 324, 'cocycle triples (non-simply-laced)': 2187, 'matter cocycle pairs': 486, 'products': 17, 'bar pairs': 243, 'seconds': 64.6}. & \CERTIFIED\newline {\footnotesize implementation, fast\newline web}\newline {\footnotesize run web 2026-09-23: 13/13} \\ \addlinespace
\code{eq:rhotorus}\newline \peqref{eq:rhotorus} & $(\rho f)_{-m}(v,\mu)$ $=v^{\kappa(m)}\mu^{\kappa_f(m)}\,f_m(v^{-1},\mu^{-1})$ is the algebra automorphism inducing \peqref{eq:rho-witten}. & checked in the run of \code{eq:rho-witten} (same adapter): the torus map \code{rho} of the chart is the printed formula on every magnetic sector (with matter, $\kappa$ with its matter term and the $\mu$-exponent $k\mapsto\kappa_f(m)-k$); it induces \peqref{eq:rho-witten} (the chart of the label-level image is the torus image of the chart); and it is multiplicative on the torus products of the lines. \emph{Population:} five SU(2) lines and four SU(2)+1 lines (every magnetic sector and $\mu$-level); their 25 + 16 torus products. \emph{Controls:} negative: the formula without $v\mapsto v^{-1}$ is detected. \emph{Evidence:} emergent. \emph{Where:} \code{battery/checks_coulomb.py}, \code{tests/test_abe_rho_label.py}. \emph{Note:} Until 2026-09-23 the run checked the pure formula on the five SU(2) lines only; that the map induces the label-level rho and is multiplicative, and the mu part, were added then. \emph{Run (web, 2026-09-23):} 8 checks, 0 failing; pure labels: 155; matter labels: 261; torus lines: 5; torus lines with $\mu$: 4; torus products: 41. & \CERTIFIED\newline {\footnotesize theorem, fast\newline web}\newline {\footnotesize run web 2026-09-23: 8/8} \\ \addlinespace
\code{eq:measure}\newline \peqref{eq:measure} & \code{trace} and \code{inner_product} implement \peqref{eq:measure} and \peqref{eq:Iexplicit}: the contour integral of $f_0$ (resp. of $f_m(v^{-1},\mu^{-1})g_m(v,\mu)$ summed over $m$) against the printed infinite products with the $|\langle m,\alpha\rangle|$- and $|\langle m,w\rangle|$-shifted vector and matter factors, prefactor $(\fq^2;\fq^2)_\infty^{2\rk}/|W|$. & the printed formulas evaluated directly as $\fq$-series with Laurent-monomial coefficients in $(v,\mu)$ (the contour integral is the constant term in $v$; the class's residuals expanded from their numerator and denominator factors), compared coefficient by coefficient with the classes' traces and pairings — an independent check of the repository's rational sector-measure route, which is proved equal to the printed products in \code{notes/kalgebra.md}.  At SU(3) the trace formula is checked with its rank-two prefactor $(\fq^2;\fq^2)_\infty^4/6$, which rank one cannot distinguish from other rank dependences ($2\,\rk=|W|=2$ there). \emph{Population:} pure SU(2): six lines ($1,L_{0,1},L_{0,2},L_{2,0},L_{2,1},L_{4,0}$), 36 pairs through $\fq^{10}$; pure SU(3): five lines (traces) through $\fq^6$; SO(3)+adjoint with $\mu$: five lines ($1,L_{0,1},L_{1,0},L_{2,0},\mu$), 25 pairs through $\fq^8$; widened 2026-09-23: SU(3), four lines and 16 pairs through $\fq^4$ (pairing); Spin(5), four lines (traces, $|W|=8$; $G_2$, $|W|=12$, on the local tier: 11 min through $\fq^4$); SU(3)+1 with $\mu$, three lines and 9 pairs through $\fq^3$ (a complex representation). \emph{Controls:} positive: orthonormality read off the formula ($I_{a,b}[\fq^0]=\delta_{a,b}$); negative: the unshifted measure fails at $m\ne0$; the rank-one prefactor fails at SU(3); the rank-one prefactor fails at SU(3). \emph{Evidence:} n.a. (a test of the code against the text). \emph{Where:} \code{battery/checks_coulomb.py}, \code{tests/test_wrq_sector_pairing.py}, \code{tests/test_matter_sector_pairing.py}, \code{tests/test_sector_measure_closed_form.py}. \emph{Run (web, 2026-09-23):} 13 checks, 0 failing; SU(2) lines: 6; pairs: 36; SO(3)+adjoint lines: 5; K: 10; SU(3) lines (trace): 5; trace-face pairs: 61; widened (2026-09-23): {'SU(3) pairs': 16, 'non-simply-laced traces': 4, 'SU(3)+1 pairs': 9, 'seconds': 230.1}. & \CERTIFIED\newline {\footnotesize implementation, fast\newline web}\newline {\footnotesize run web 2026-09-23: 13/13} \\ \addlinespace
\code{eq:Iexplicit}\newline \peqref{eq:Iexplicit} and the remark after it & $I_{f,g}$ as a sum over $m$ of contour integrals against the $|\langle m,\alpha\rangle|$-, $|\langle m,w\rangle|$-shifted measure; equal to $\Tr\rho(f)g$ at integrand level up to contour shifts, the pole compatibility being conjectural. & The printed sum over $m$ --- each sector's contour integral at $|v|=1$ against its own $|\langle m,\alpha\rangle|$-, $|\langle m,w\rangle|$-shifted measure --- is evaluated independently from the charts and compared with the class's pairing.  The class's pairing is $\Tr\rho(f)g$ itself: the Schur residue \peqref{eq:measure} of the magnetic-0 residual of $\rho(f)\cdot g$, assembled sector by sector (\code{inner_by_sector} is \code{trace_residual} of \code{pairing_residual}).  That comparison is therefore the conjectural half: bringing each printed sector onto the common contour is the shift $v\to\fq^{\pm m}v$ whose pole compatibility the paper conjectures.  It runs in the eq:measure adapter on all 61 rank-one pairs (SU(2), SO(3)+adjoint) and at rank two on 16 SU(3) and 9 SU(3)+1 pairs.  A second computation of $\Tr\rho(f)g$ is compared with the class's as a consistency check, not as the conjecture (the two agree by the total-charge-0 part of the cocycle convolution, a theorem): multiply-then-trace at the fast depth ($\rho(f)g$ decomposed into canonical lines, each traced), and at the extensive depth the Schur residue of the magnetic-0 residual of the torus product $\rho(f)\cdot g$ itself (pure gauge the raw exact series $|W|\Tr$, with matter the $\mu$-refined residue at the class's Nahm window, $K+2$ raised to $K$ plus the depth of the residual's $\fq$-expansion below $\fq^0$ (the source repository's audit record); equal to multiply-then-trace pair by pair at SU(2) and SU(2)+1, and it avoids the decomposition, which outgrows memory at rank two).  At the extensive depth each theory runs in its own process, 1800 s and 5 GB, the pairs ordered so that the costliest line comes last; pairs not compared within the budget are reported, never counted. \emph{Population:} the SU(2) and SO(3)+adjoint lines of \code{eq:measure} (36 and 25 pairs); widened 2026-09-23: SU(3), 16 pairs; SU(3)+1 with $\mu$, 9 pairs; the extensive tier, recorded beside the fast one, runs the 24 theories of the Section~3 lineup (listed under sec:coulomb/labels) with five lines each --- the identity, the smallest Wilson line, the two smallest magnetic charges at the smallest electric charge the form admits with them, and one dyonic line; four where the window has a single nonzero magnetic charge (SU(3) with an adjoint, with a symmetric, and with a symmetric and a fundamental) --- on all ordered pairs (25, or 16), through $\fq^4$ ($\fq^3$ with matter above rank one). \emph{Controls:} positive: at the extensive depth the printed pairing is $|W|(\delta+O(\fq))$ on every compared pair (no negative $\fq$-power, $\delta$ at $\fq^0$); negative: the unshifted measure (every sector at the $m=0$ measure) changes the printed pairing of the first magnetic line with itself (extensive tier); the fast run's controls are those of eq:measure (same adapter). \emph{Evidence:} emergent (the two routes share nothing but the chart). \emph{Where:} \code{battery/checks_coulomb.py}, \code{tests/test_matter_sector_pairing.py}, \code{tests/test_sector_measure_closed_form.py}, \code{tests/test_wrq_sector_pairing.py}. \emph{Note:} the measure is DERIVED ; the pole compatibility is the open part.  The class's pairing is Tr rho(f)g assembled sector by sector, so the printed-vs-class comparison has always been the test of the conjecture; until a review of the battery on 2026-09-23 the row named instead the comparison of the class's pairing with multiply-then-trace (two computations of Tr rho(f)g, equal by a theorem) as the conjectural half.  tests/test\_wrq\_sector\_pairing.py (72 SU(2), 32 SU(3) pairs) checks that theorem-level equality of the two computations. \emph{Run (web, 2026-09-23):} 13 checks, 0 failing; SU(2) lines: 6; pairs: 36; SO(3)+adjoint lines: 5; K: 10; SU(3) lines (trace): 5; trace-face pairs: 61; widened (2026-09-23): {'SU(3) pairs': 16, 'non-simply-laced traces': 4, 'SU(3)+1 pairs': 9, 'seconds': 226.5}. \emph{Run (web (extensive), 2026-09-24):} 93 checks, 0 failing; pure SU(2): {'lines': ['((0,), (0,))', '((0,), (1,))', '((1,), (0,))', '((1,), (1,))', '((2,), (0,))'], 'K': 4, 'weyl': 2, 'slots': 0, 'pairs': 25, 'compared': 25, 'printed\_eq\_trace': 25, 'class\_eq\_torus': 25, 'printed\_delta\_plus\_O\_q': 25, 'unshifted\_detected': True, 'stopped\_by': None, 'error': None, 'seconds': 0.1}; pure SO(3): {'lines': ['((0,), (0,))', '((0,), (1,))', '((1,), (0,))', '((1,), (1,))', '((2,), (0,))'], 'K': 4, 'weyl': 2, 'slots': 0, 'pairs': 25, 'compared': 25, 'printed\_eq\_trace': 25, 'class\_eq\_torus': 25, 'printed\_delta\_plus\_O\_q': 25, 'unshifted\_detected': True, 'stopped\_by': None, 'error': None, 'seconds': 0.1}; pure rank 1, the third global form: {'lines': ['((0,), (0,))', '((0,), (2,))', '((Fraction(1, 2),), (1,))', '((Fraction(1, 2),), (-1,))', '((1,), (0,))'], 'K': 4, 'weyl': 2, 'slots': 0, 'pairs': 25, 'compared': 25, 'printed\_eq\_trace': 25, 'class\_eq\_torus': 25, 'printed\_delta\_plus\_O\_q': 25, 'unshifted\_detected': True, 'stopped\_by': None, 'error': None, 'seconds': 0.1}; pure SU(3): {'lines': ['((0, 0), (0, 0))', '((0, 0), (1, 0))', '((1, 1), (0, 0))', '((1, 1), (1, 0))', '((2, 2), (0, 0))'], 'K': 4, 'weyl': 6, 'slots': 0, 'pairs': 25, 'compared': 25, 'printed\_eq\_trace': 25, 'class\_eq\_torus': 25, 'printed\_delta\_plus\_O\_q': 25, 'unshifted\_detected': True, 'stopped\_by': None, 'error': None, 'seconds': 25.8}; pure PSU(3): {'lines': ['((0, 0), (0, 0))', '((0, 0), (1, 1))', '((Fraction(2, 3), Fraction(1, 3)), (0, 0))', '((Fraction(2, 3), Fraction(1, 3)), (1, 1))', '((Fraction(1, 3), Fraction(2, 3)), (0, 0))'], 'K': 4, 'weyl': 6, 'slots': 0, 'pairs': 25, 'compared': 25, 'printed\_eq\_trace': 25, 'class\_eq\_torus': 25, 'printed\_delta\_plus\_O\_q': 25, 'unshifted\_detected': True, 'stopped\_by': None, 'error': None, 'seconds': 0.5}; pure Sp(4): {'lines': ['((0, 0), (0, 0))', '((0, 0), (1, 0))', '((1, 1), (0, 0))', '((1, 1), (1, 0))', '((2, 1), (0, 0))'], 'K': 4, 'weyl': 8, 'slots': 0, 'pairs': 25, 'compared': 25, 'printed\_eq\_trace': 25, 'class\_eq\_torus': 25, 'printed\_delta\_plus\_O\_q': 25, 'unshifted\_detected': True, 'stopped\_by': None, 'error': None, 'seconds': 84.7}; pure Spin(5): {'lines': ['((0, 0), (0, 0))', '((0, 0), (1, 0))', '((1, 1), (0, 0))', '((1, 1), (1, 0))', '((2, 1), (0, 0))'], 'K': 4, 'weyl': 8, 'slots': 0, 'pairs': 25, 'compared': 25, 'printed\_eq\_trace': 25, 'class\_eq\_torus': 25, 'printed\_delta\_plus\_O\_q': 25, 'unshifted\_detected': True, 'stopped\_by': None, 'error': None, 'seconds': 5.1}; pure SO(5): {'lines': ['((0, 0), (0, 0))', '((0, 0), (1, 0))', '((Fraction(1, 1), Fraction(1, 2)), (0, 0))', '((Fraction(1, 1), Fraction(1, 2)), (1, 0))', '((1, 1), (0, 0))'], 'K': 4, 'weyl': 8, 'slots': 0, 'pairs': 25, 'compared': 25, 'printed\_eq\_trace': 25, 'class\_eq\_torus': 25, 'printed\_delta\_plus\_O\_q': 25, 'unshifted\_detected': True, 'stopped\_by': None, 'error': None, 'seconds': 1.7}; pure G2: {'lines': ['((0, 0), (0, 0))', '((0, 0), (1, 0))', '((1, 2), (0, 0))', '((1, 2), (1, 0))', '((2, 3), (0, 0))'], 'K': 4, 'weyl': 12, 'slots': 0, 'pairs': 25, 'compared': 25, 'printed\_eq\_trace': 25, 'class\_eq\_torus': 25, 'printed\_delta\_plus\_O\_q': 25, 'unshifted\_detected': True, 'stopped\_by': None, 'error': None, 'seconds': 91.7}; SU(2)+1: {'lines': ['(((0,), (0,)), (0,))', '(((0,), (1,)), (0,))', '(((1,), (0,)), (0,))', '(((1,), (1,)), (0,))', '(((2,), (0,)), (0,))'], 'K': 4, 'weyl': 2, 'slots': 1, 'pairs': 25, 'compared': 25, 'printed\_eq\_trace': 25, 'class\_eq\_torus': 25, 'printed\_delta\_plus\_O\_q': 25, 'unshifted\_detected': True, 'stopped\_by': None, 'error': None, 'seconds': 5.7}; SU(2)+2: {'lines': ['(((0,), (0,)), (0, 0))', '(((0,), (1,)), (0, 0))', '(((1,), (0,)), (0, 0))', '(((1,), (1,)), (0, 0))', '(((2,), (0,)), (0, 0))'], 'K': 4, 'weyl': 2, 'slots': 2, 'pairs': 25, 'compared': 25, 'printed\_eq\_trace': 25, 'class\_eq\_torus': 25, 'printed\_delta\_plus\_O\_q': 25, 'unshifted\_detected': True, 'stopped\_by': None, 'error': None, 'seconds': 314.8}; SU(2)+3: {'lines': ['(((0,), (0,)), (0, 0, 0))', '(((0,), (1,)), (0, 0, 0))', '(((1,), (0,)), (0, 0, 0))', '(((1,), (1,)), (0, 0, 0))', '(((2,), (0,)), (0, 0, 0))'], 'K': 4, 'weyl': 2, 'slots': 3, 'pairs': 25, 'compared': 16, 'printed\_eq\_trace': 16, 'class\_eq\_torus': 16, 'printed\_delta\_plus\_O\_q': 16, 'unshifted\_detected': True, 'stopped\_by': '1800 s', 'error': None, 'seconds': 1800.2}; SU(2)+4: {'lines': ['(((0,), (0,)), (0, 0, 0, 0))', '(((0,), (1,)), (0, 0, 0, 0))', '(((1,), (0,)), (0, 0, 0, 0))', '(((1,), (1,)), (0, 0, 0, 0))', '(((2,), (0,)), (0, 0, 0, 0))'], 'K': 4, 'weyl': 2, 'slots': 4, 'pairs': 25, 'compared': 6, 'printed\_eq\_trace': 6, 'class\_eq\_torus': 6, 'printed\_delta\_plus\_O\_q': 6, 'unshifted\_detected': None, 'stopped\_by': '1800 s', 'error': None, 'seconds': 1800.2}; SU(2)+adjoint (N=2*): {'lines': ['(((0,), (0,)), (0,))', '(((0,), (1,)), (0,))', '(((1,), (0,)), (0,))', '(((1,), (1,)), (0,))', '(((2,), (0,)), (0,))'], 'K': 4, 'weyl': 2, 'slots': 1, 'pairs': 25, 'compared': 25, 'printed\_eq\_trace': 25, 'class\_eq\_torus': 25, 'printed\_delta\_plus\_O\_q': 25, 'unshifted\_detected': True, 'stopped\_by': None, 'error': None, 'seconds': 27.1}; SO(3)+adjoint (N=2*): {'lines': ['(((0,), (0,)), (0,))', '(((0,), (2,)), (0,))', '(((Fraction(1, 2),), (0,)), (0,))', '(((Fraction(1, 2),), (2,)), (0,))', '(((1,), (0,)), (0,))'], 'K': 4, 'weyl': 2, 'slots': 1, 'pairs': 25, 'compared': 25, 'printed\_eq\_trace': 25, 'class\_eq\_torus': 25, 'printed\_delta\_plus\_O\_q': 25, 'unshifted\_detected': True, 'stopped\_by': None, 'error': None, 'seconds': 15.7}; rank 1 + adjoint (N=2*), the third global form: {'lines': ['(((0,), (0,)), (0,))', '(((0,), (2,)), (0,))', '(((Fraction(1, 2),), (1,)), (0,))', '(((Fraction(1, 2),), (-1,)), (0,))', '(((1,), (0,)), (0,))'], 'K': 4, 'weyl': 2, 'slots': 1, 'pairs': 25, 'compared': 25, 'printed\_eq\_trace': 25, 'class\_eq\_torus': 25, 'printed\_delta\_plus\_O\_q': 25, 'unshifted\_detected': True, 'stopped\_by': None, 'error': None, 'seconds': 15.9}; SU(3)+adjoint (N=2*): {'lines': ['(((0, 0), (0, 0)), (0,))', '(((0, 0), (1, 0)), (0,))', '(((1, 1), (0, 0)), (0,))', '(((1, 1), (1, 0)), (0,))'], 'K': 3, 'weyl': 6, 'slots': 1, 'pairs': 16, 'compared': 8, 'printed\_eq\_trace': 8, 'class\_eq\_torus': 8, 'printed\_delta\_plus\_O\_q': 8, 'unshifted\_detected': True, 'stopped\_by': '1800 s', 'error': None, 'seconds': 1800.1}; SU(3)+symmetric: {'lines': ['(((0, 0), (0, 0)), (0,))', '(((0, 0), (1, 0)), (0,))', '(((1, 1), (0, 0)), (0,))', '(((1, 1), (1, 0)), (0,))'], 'K': 3, 'weyl': 6, 'slots': 1, 'pairs': 16, 'compared': 12, 'printed\_eq\_trace': 12, 'class\_eq\_torus': 12, 'printed\_delta\_plus\_O\_q': 12, 'unshifted\_detected': True, 'stopped\_by': '1800 s', 'error': None, 'seconds': 1800.0}; SU(3)+symmetric+1: {'lines': ['(((0, 0), (0, 0)), ((0,), (0,)))', '(((0, 0), (1, 0)), ((0,), (0,)))', '(((1, 1), (0, 0)), ((0,), (0,)))', '(((1, 1), (1, 0)), ((0,), (0,)))'], 'K': 3, 'weyl': 6, 'slots': 2, 'pairs': 16, 'compared': 4, 'printed\_eq\_trace': 4, 'class\_eq\_torus': 4, 'printed\_delta\_plus\_O\_q': 4, 'unshifted\_detected': None, 'stopped\_by': '1800 s', 'error': None, 'seconds': 1800.4}; Spin(5)+vector: {'lines': ['(((0, 0), (0, 0)), (0,))', '(((0, 0), (1, 0)), (0,))', '(((1, 1), (0, 0)), (0,))', '(((1, 1), (1, 0)), (0,))', '(((2, 1), (0, 0)), (0,))'], 'K': 3, 'weyl': 8, 'slots': 1, 'pairs': 25, 'compared': 18, 'printed\_eq\_trace': 18, 'class\_eq\_torus': 18, 'printed\_delta\_plus\_O\_q': 18, 'unshifted\_detected': True, 'stopped\_by': '1800 s', 'error': None, 'seconds': 1800.0}; SO(5)+vector: {'lines': ['(((0, 0), (0, 0)), (0,))', '(((0, 0), (1, 0)), (0,))', '(((Fraction(1, 1), Fraction(1, 2)), (0, 0)), (0,))', '(((Fraction(1, 1), Fraction(1, 2)), (1, 0)), (0,))', '(((1, 1), (0, 0)), (0,))'], 'K': 3, 'weyl': 8, 'slots': 1, 'pairs': 25, 'compared': 20, 'printed\_eq\_trace': 20, 'class\_eq\_torus': 20, 'printed\_delta\_plus\_O\_q': 20, 'unshifted\_detected': True, 'stopped\_by': '1800 s', 'error': None, 'seconds': 1800.0}; G2+7: {'lines': ['(((0, 0), (0, 0)), (0,))', '(((0, 0), (1, 0)), (0,))', '(((1, 2), (0, 0)), (0,))', '(((1, 2), (1, 0)), (0,))', '(((2, 3), (0, 0)), (0,))'], 'K': 3, 'weyl': 12, 'slots': 1, 'pairs': 25, 'compared': 6, 'printed\_eq\_trace': 6, 'class\_eq\_torus': 6, 'printed\_delta\_plus\_O\_q': 6, 'unshifted\_detected': None, 'stopped\_by': '1800 s', 'error': None, 'seconds': 1800.1}; SU(2)xSU(2)+bifundamental: {'lines': ['(((0, 0), (0, 0)), (0,))', '(((0, 0), (1, 0)), (0,))', '(((1, 0), (0, 0)), (0,))', '(((1, 0), (1, 0)), (0,))', '(((0, 1), (0, 0)), (0,))'], 'K': 3, 'weyl': 4, 'slots': 1, 'pairs': 25, 'compared': 25, 'printed\_eq\_trace': 25, 'class\_eq\_torus': 25, 'printed\_delta\_plus\_O\_q': 25, 'unshifted\_detected': True, 'stopped\_by': None, 'error': None, 'seconds': 85.3}; U(1)-U(2)+3: {'lines': ['(((0, 0, 0), (0, 0, 0)), ((0,), (0, 0, 0)))', '(((0, 0, 0), (1, 0, 0)), ((0,), (0, 0, 0)))', '(((1, 0, 0), (0, 0, 0)), ((0,), (0, 0, 0)))', '(((1, 0, 0), (1, 0, 0)), ((0,), (0, 0, 0)))', '(((0, 1, 0), (0, 0, 0)), ((0,), (0, 0, 0)))'], 'K': 3, 'weyl': 2, 'slots': 4, 'pairs': 25, 'compared': 25, 'printed\_eq\_trace': 25, 'class\_eq\_torus': 25, 'printed\_delta\_plus\_O\_q': 25, 'unshifted\_detected': True, 'stopped\_by': None, 'error': None, 'seconds': 586.0}. & \CERTIFIED\newline {\footnotesize conjecture, extensive\newline web}\newline {\footnotesize run web 2026-09-23: 13/13}\newline {\footnotesize run web (extensive) 2026-09-24: 93/93} \\ \addlinespace
\code{eq:cocha}\newline \peqref{eq:cocha} & \code{dressing_psi} is the printed $\mathrm{cocha}_m(v)$ wherever the prefactor $(-\fq)^{\frac12\sum\langle\alpha,m\rangle}$ is integral (SU(2) all $m$, SO(3) even $m$); at odd $m$ of SO(3) the prefactor is half-integral and the repository ships the floored $\psi$ with the phase restored through the cocycle . & exact rational identity (cross-multiplication) between \code{dressing_psi} and the printed product of $k$-even denominators with its $(-\fq)$ prefactor, at rank one and two, dominant and non-dominant $m$; the half-integral case of SO(3) recorded, not compared.  The draft's formulas among themselves: eq:cocydef with $f=\sum_m f_m(\fq^mv)U_m=\sum_m d_mu^m$ and $d_m=f_m(\fq^mv)\cocha_m(v)$ gives $U_m=\cocha_m(v)u^m$, hence $\cocy_{a,b}(\fq^{a+b}v)\cocha_{a+b}(v)=(-\fq)^{s(a)+s(b)-s(a+b)}\cocha_a(v)\cocha_b(\fq^{2a}v)$, $s$ the prefactor exponent; checked with the printed eq:ccclosed / eq:ccprod, matter factors included, the net exponent required to be an integer. \emph{Population:} SU(2) $m=1,2,3$; SO(3) $m=2,4$; SU(3), Sp(4), Spin(5), $G_2$ at every $m$ of the box $[-2,2]^2$; the trivialization identity on every pair of a box for 18 theories: pure SU(2), SO(3), SU(3), Sp(4), Spin(5), $G_2$, SU(2)$\times$SU(2); SU(2)+1, SU(2)+2, SO(3)+adjoint; SU(3)+1, SU(3)+symmetric, SU(3)+symmetric+1, SU(3)+adjoint, Spin(5)+vector, $G_2$+7; the SU(2)--SU(2) quiver with a bifundamental, and with two fundamentals at each node (the latter an extra case, outside the lineup). \emph{Controls:} positive: exact identities; negative: the prefactor $(+\fq)^m$ is detected; the older all-positive-roots form of eq:cocha, commented out in the draft source, fails the trivialization identity at every pure group. \emph{Evidence:} n.a. (a test of the code against the text). \emph{Where:} \code{battery/checks_coulomb.py}, \code{tests/test_pure_g_abe_kalgebra.py}. \emph{Note:} The matter factors of eq:cocha are tested twice: here, as the trivialization of the printed matter cocycle, and in conj:abeKalgebra/existence-uniqueness, where the code's lines satisfy A3 with them (2026-09-23).  Before that they were exercised only through the SO(3)+adjoint d-forms of the adjoint row. \emph{Run (web, 2026-09-23):} 7 checks, 0 failing; SU(2) m: 1..3; SO(3) m: 2, 4; rank-two cocharacters: 100; trivialization pairs: 1178; theories: {'pure SU(2)': '25/25', 'pure SO(3)': '25/25', 'pure SU(3)': '81/81', 'pure Sp(4)': '81/81', 'pure Spin(5)': '81/81', 'pure G2': '81/81', 'pure SU(2)xSU(2)': '81/81', 'SU(2)+1': '25/25', 'SU(2)+2': '25/25', 'SO(3)+adjoint': '25/25', 'SU(3)+1': '81/81', 'SU(3)+symmetric': '81/81', 'SU(3)+symmetric+1': '81/81', 'SU(3)+adjoint': '81/81', 'Spin(5)+vector': '81/81', 'G2+7': '81/81', 'SU(2)xSU(2)+bifundamental': '81/81', '[2]-SU(2)-SU(2)-[2]': '81/81'}. & \CERTIFIED\newline {\footnotesize implementation, fast\newline web}\newline {\footnotesize run web 2026-09-23: 7/7} \\ \addlinespace
\code{conj:abeKalgebra/existence-uniqueness}\newline Conj.~\pref{conj:abeKalgebra} & Conditions A1 (Weyl-covariant, bar-invariant bubbling), A2 ($O(\fq)$) and A3 (the residue rule \peqref{eq:star}) define $L_{m,e}$ uniquely. & the guarded solvers called directly on every window label.  Pure gauge: \code{joint_solve} returns `unique' and \code{solve_canonical} applies W1, the residue rule and the box certificate (re-solving one pad wider returns the same element).  With matter: \code{solve_canonical_matter_vec}, every $\mu$-sector unique and divisible by $Z$ (strict).  In both cases the result equals the production chart (with matter after dividing by $Z$: the chart stores $F/Z$).  Every window line satisfies the draft's A3 exactly as printed: $d_a=f_a(\fq^av)\cocha_a(v,\mu)$, with the printed matter factors of eq:cocha, has simple poles only and its pair residues cancel, power by power in $\mu$ (on solver-built lines the simple-pole half holds by construction: the solver's ansatz admits only simple poles, one at each wall where a residue partner lies in the support; the independent content of the re-check is the residue cancellation with the draft's printed partner $b=s_\alpha(a)-l\alpha^\vee$ and printed matter factors); also on SU(3)+1, SU(3)+symmetric and Spin(5)+vector lines.  At the extensive depth each solve has a budget of 600 s; a solve not finished within it is reported in the population, never counted as a pass.  What `unique' is relative to: the solver's ansatz assumes the support $\mathrm{conv}(W\cdot m)$, Weyl covariance (unknowns on orbit representatives), bar invariance ($\fq$-palindromic coefficients), simple poles at the predicted walls, and numerators in a finite $v$-box of bounded $\fq$-degree (the degree window derived from the residue equations' constants plus a margin of two); the box certificate widens the $v$-box by one step, not the degree window.  The support assumption is measured separately (conj:abeKalgebra/support, the enlarged-support solve). \emph{Population:} the nine conjecture windows (\code{conjecture_windows}): pure SU(2), SO(3) ($m\le2$, $|e|\le2$), SU(3), Sp(4), $G_2$; SU(2)+1, SU(2)+2, SO(3)+adjoint, SU(3)+1 (flavour 0); the extensive tier, recorded beside the fast one, runs the 24 theories of the Section~3 lineup on grown windows (\code{_extensive_windows}): pure SU(2) and SO(3) ($m\le4$, $|e|\le3$), the rank-one third global form, SU(3), PSU(3), Sp(4), Spin(5), SO(5), $G_2$; SU(2) with $N_f=1,\dots,4$; $N=2^*$ at SU(2), SO(3), the third global form and SU(3); SU(3) with a symmetric, alone and with a fundamental; Spin(5) and SO(5) with a vector; $G_2$ with a $\mathbf 7$; the SU(2)$\times$SU(2) quiver with a bifundamental; U(1)--U(2) with three fundamentals of U(2) (matter at the neutral flavour label). \emph{Controls:} positive: the Wilson lines, a single support cell and no unknowns; negative: every pure window line with bubbling, perturbed two ways (bubbling removed; bubbling doubled), keeps bar invariance, the seed, $O(\fq)$ and the support bound, and the residue rule rejects every one; for A3 as printed, the matter factor over the weights of $N^\vee$ fails on the complex representations (SU(3) fundamental and symmetric; for self-conjugate matter the two coincide), and dropping the matter factor fails. \emph{Evidence:} enforced (the solver imposes the three conditions; the finding is existence and uniqueness). \emph{Where:} \code{battery/checks_coulomb.py (check_existence_uniqueness)}, \code{tests/test_pure_g_abe_kalgebra.py}, \code{tests/test_matter_multislot.py}. \emph{Note:} Own adapter since 2026-09-23; the declared negative control was not run before.  The negative controls are pure gauge only: with matter the residue rule sits inside the solver's divisibility guard.  The code imposes the residue rule with matter on the dressed image F = Z$\cdot$f through its own trivialization Z; that the stored lines also satisfy the draft's A3 with the printed matter factors is the check that the two formulations agree (added 2026-09-23).  Corrected 2026-09-24: the check described the simple-pole clause as one `no solver imposes'; the solver's ansatz builds it in (star\_bubbling.joint\_solve fixes the denominators in advance), found in the Section 3 walkthrough. \emph{Run (web, 2026-09-24):} 34 checks, 0 failing; pure SU(2): {'labels': 13, 'unreached': [], 'solves\_unreached': [], 'errors': {}, 'pinned': 13, 'equal\_to\_chart': 13, 'a3\_as\_printed': 13, 'unknown\_cells': {'((0,), (0,))': 0, '((0,), (1,))': 0, '((0,), (2,))': 0, '((1,), (-1,))': 1, '((1,), (-2,))': 1, '((1,), (0,))': 1, '((1,), (1,))': 1, '((1,), (2,))': 1, '((2,), (-1,))': 3, '((2,), (-2,))': 3, '((2,), (0,))': 3, '((2,), (1,))': 3, '((2,), (2,))': 3}, 'wilson': 3, 'negative': {'perturbed': 20, 'other\_conditions\_hold': 20, 'star\_rejects': 20}, 'seconds': 0.4}; pure SO(3): {'labels': 13, 'unreached': [], 'solves\_unreached': [], 'errors': {}, 'pinned': 13, 'equal\_to\_chart': 13, 'a3\_as\_printed': 13, 'unknown\_cells': {'((0,), (0,))': 0, '((0,), (1,))': 0, '((0,), (2,))': 0, '((1,), (-1,))': 0, '((1,), (-2,))': 0, '((1,), (0,))': 0, '((1,), (1,))': 0, '((1,), (2,))': 0, '((2,), (-1,))': 1, '((2,), (-2,))': 1, '((2,), (0,))': 1, '((2,), (1,))': 1, '((2,), (2,))': 1}, 'wilson': 3, 'negative': {'perturbed': 10, 'other\_conditions\_hold': 10, 'star\_rejects': 10}, 'seconds': 0.1}; pure SU(3): {'labels': 8, 'unreached': [], 'solves\_unreached': [], 'errors': {}, 'pinned': 8, 'equal\_to\_chart': 8, 'a3\_as\_printed': 8, 'unknown\_cells': {'((0, 0), (0, 0))': 0, '((0, 0), (0, 1))': 0, '((0, 0), (1, 0))': 0, '((1, 1), (-1, 1))': 1, '((1, 1), (0, 0))': 1, '((1, 1), (0, 1))': 1, '((1, 1), (1, -1))': 1, '((1, 1), (1, 0))': 1}, 'wilson': 3, 'negative': {'perturbed': 10, 'other\_conditions\_hold': 10, 'star\_rejects': 10}, 'seconds': 1.1}; pure Sp(4): {'labels': 5, 'unreached': [], 'solves\_unreached': [], 'errors': {}, 'pinned': 5, 'equal\_to\_chart': 5, 'a3\_as\_printed': 5, 'unknown\_cells': {'((0, 0), (0, 0))': 0, '((0, 0), (1, 0))': 0, '((1, 0), (0, 0))': 1, '((1, 0), (0, 1))': 1, '((1, 0), (1, 0))': 1}, 'wilson': 2, 'negative': {'perturbed': 6, 'other\_conditions\_hold': 6, 'star\_rejects': 6}, 'seconds': 0.6}; pure G2: {'labels': 4, 'unreached': [], 'solves\_unreached': [], 'errors': {}, 'pinned': 4, 'equal\_to\_chart': 4, 'a3\_as\_printed': 4, 'unknown\_cells': {'((0, 0), (0, 0))': 0, '((0, 0), (1, 0))': 0, '((2, 3), (-1, 1))': 7, '((2, 3), (0, 0))': 7}, 'wilson': 2, 'negative': {'perturbed': 4, 'other\_conditions\_hold': 4, 'star\_rejects': 4}, 'seconds': 41.7}; SU(2) + 1 fund. (roster su2-nf1): {'labels': 5, 'unreached': [], 'solves\_unreached': [], 'errors': {}, 'pinned': 5, 'equal\_to\_chart': 5, 'a3\_as\_printed': 5, 'unknown\_cells': None, 'wilson': 2, 'negative': None, 'seconds': 0.1}; SU(2) + 2 fund. (roster su2-nf2, U(2) flavour): {'labels': 5, 'unreached': [], 'solves\_unreached': [], 'errors': {}, 'pinned': 5, 'equal\_to\_chart': 5, 'a3\_as\_printed': 5, 'unknown\_cells': None, 'wilson': 2, 'negative': None, 'seconds': 0.1}; SO(3) + adjoint (N=2*): {'labels': 5, 'unreached': [], 'solves\_unreached': [], 'errors': {}, 'pinned': 5, 'equal\_to\_chart': 5, 'a3\_as\_printed': 5, 'unknown\_cells': None, 'wilson': 2, 'negative': None, 'seconds': 0.0}; SU(3) + 1 fund. (roster su3-nf1): {'labels': 4, 'unreached': [], 'solves\_unreached': [], 'errors': {}, 'pinned': 4, 'equal\_to\_chart': 4, 'a3\_as\_printed': 4, 'unknown\_cells': None, 'wilson': 2, 'negative': None, 'seconds': 0.9}; A3 with printed matter factors: {'SU(3)+1': {'printed': 2, 'dual': 0, 'none': 0}, 'SU(3)+symmetric': {'printed': 2, 'dual': 0, 'none': 0}, 'Spin(5)+vector': {'printed': 2, 'dual': 2, 'none': 0}}. \emph{Run (web (extensive), 2026-09-24):} 83 checks, 0 failing; pure SU(2): {'labels': 32, 'unreached': [], 'solves\_unreached': [], 'errors': {}, 'pinned': 32, 'equal\_to\_chart': 32, 'a3\_as\_printed': 32, 'unknown\_cells': {'((0,), (0,))': 0, '((0,), (1,))': 0, '((0,), (2,))': 0, '((0,), (3,))': 0, '((1,), (-1,))': 1, '((1,), (-2,))': 1, '((1,), (-3,))': 1, '((1,), (0,))': 1, '((1,), (1,))': 1, '((1,), (2,))': 1, '((1,), (3,))': 1, '((2,), (-1,))': 3, '((2,), (-2,))': 3, '((2,), (-3,))': 3, '((2,), (0,))': 3, '((2,), (1,))': 3, '((2,), (2,))': 3, '((2,), (3,))': 3, '((3,), (-1,))': 5, '((3,), (-2,))': 5, '((3,), (-3,))': 5, '((3,), (0,))': 5, '((3,), (1,))': 5, '((3,), (2,))': 5, '((3,), (3,))': 5, '((4,), (-1,))': 7, '((4,), (-2,))': 7, '((4,), (-3,))': 7, '((4,), (0,))': 7, '((4,), (1,))': 7, '((4,), (2,))': 7, '((4,), (3,))': 7}, 'wilson': 4, 'negative': {'perturbed': 56, 'other\_conditions\_hold': 56, 'star\_rejects': 56}, 'seconds': 22.2}; pure SO(3): {'labels': 32, 'unreached': [], 'solves\_unreached': [], 'errors': {}, 'pinned': 32, 'equal\_to\_chart': 32, 'a3\_as\_printed': 32, 'unknown\_cells': {'((0,), (0,))': 0, '((0,), (1,))': 0, '((0,), (2,))': 0, '((0,), (3,))': 0, '((1,), (-1,))': 0, '((1,), (-2,))': 0, '((1,), (-3,))': 0, '((1,), (0,))': 0, '((1,), (1,))': 0, '((1,), (2,))': 0, '((1,), (3,))': 0, '((2,), (-1,))': 1, '((2,), (-2,))': 1, '((2,), (-3,))': 1, '((2,), (0,))': 1, '((2,), (1,))': 1, '((2,), (2,))': 1, '((2,), (3,))': 1, '((3,), (-1,))': 2, '((3,), (-2,))': 2, '((3,), (-3,))': 2, '((3,), (0,))': 2, '((3,), (1,))': 2, '((3,), (2,))': 2, '((3,), (3,))': 2, '((4,), (-1,))': 3, '((4,), (-2,))': 3, '((4,), (-3,))': 3, '((4,), (0,))': 3, '((4,), (1,))': 3, '((4,), (2,))': 3, '((4,), (3,))': 3}, 'wilson': 4, 'negative': {'perturbed': 42, 'other\_conditions\_hold': 42, 'star\_rejects': 42}, 'seconds': 0.8}; pure rank 1, the third global form: {'labels': 12, 'unreached': [], 'solves\_unreached': [], 'errors': {}, 'pinned': 12, 'equal\_to\_chart': 12, 'a3\_as\_printed': 12, 'unknown\_cells': {'((0,), (0,))': 0, '((0,), (2,))': 0, '((1,), (-2,))': 1, '((1,), (0,))': 1, '((1,), (2,))': 1, '((2,), (-2,))': 3, '((2,), (0,))': 3, '((2,), (2,))': 3, '((Fraction(1, 2),), (-1,))': 0, '((Fraction(1, 2),), (1,))': 0, '((Fraction(3, 2),), (-1,))': 2, '((Fraction(3, 2),), (1,))': 2}, 'wilson': 2, 'negative': {'perturbed': 16, 'other\_conditions\_hold': 16, 'star\_rejects': 16}, 'seconds': 0.4}; pure SU(3): {'labels': 16, 'unreached': [], 'solves\_unreached': [], 'errors': {}, 'pinned': 16, 'equal\_to\_chart': 16, 'a3\_as\_printed': 16, 'unknown\_cells': {'((0, 0), (0, 0))': 0, '((0, 0), (0, 1))': 0, '((0, 0), (1, 0))': 0, '((0, 0), (2, 0))': 0, '((1, 1), (-1, 1))': 1, '((1, 1), (0, 0))': 1, '((1, 1), (0, 1))': 1, '((1, 1), (1, -1))': 1, '((1, 1), (1, 0))': 1, '((1, 1), (2, 0))': 1, '((2, 2), (-1, 1))': 13, '((2, 2), (0, 0))': 13, '((2, 2), (0, 1))': 13, '((2, 2), (1, -1))': 13, '((2, 2), (1, 0))': 13, '((2, 2), (2, 0))': 13}, 'wilson': 4, 'negative': {'perturbed': 24, 'other\_conditions\_hold': 24, 'star\_rejects': 24}, 'seconds': 144.6}; pure PSU(3): {'labels': 16, 'unreached': [], 'solves\_unreached': [], 'errors': {}, 'pinned': 16, 'equal\_to\_chart': 16, 'a3\_as\_printed': 16, 'unknown\_cells': {'((0, 0), (0, 0))': 0, '((0, 0), (0, 3))': 0, '((0, 0), (1, 1))': 0, '((0, 0), (3, 0))': 0, '((1, 1), (0, 0))': 1, '((1, 1), (0, 3))': 1, '((1, 1), (1, 1))': 1, '((1, 1), (3, 0))': 1, '((Fraction(1, 3), Fraction(2, 3)), (0, 0))': 0, '((Fraction(1, 3), Fraction(2, 3)), (0, 3))': 0, '((Fraction(1, 3), Fraction(2, 3)), (1, 1))': 0, '((Fraction(1, 3), Fraction(2, 3)), (3, 0))': 0, '((Fraction(2, 3), Fraction(1, 3)), (0, 0))': 0, '((Fraction(2, 3), Fraction(1, 3)), (0, 3))': 0, '((Fraction(2, 3), Fraction(1, 3)), (1, 1))': 0, '((Fraction(2, 3), Fraction(1, 3)), (3, 0))': 0}, 'wilson': 4, 'negative': {'perturbed': 8, 'other\_conditions\_hold': 8, 'star\_rejects': 8}, 'seconds': 1.2}; pure Sp(4): {'labels': 19, 'unreached': [], 'solves\_unreached': ['((2, 2), (2, 0))'], 'errors': {}, 'pinned': 18, 'equal\_to\_chart': 18, 'a3\_as\_printed': 19, 'unknown\_cells': {'((0, 0), (0, 0))': 0, '((0, 0), (1, 0))': 0, '((0, 0), (1, 1))': 0, '((0, 0), (2, 0))': 0, '((1, 1), (0, 0))': 5, '((1, 1), (1, -1))': 5, '((1, 1), (1, 0))': 5, '((1, 1), (2, 0))': 5, '((2, 1), (-1, 1))': 13, '((2, 1), (0, 0))': 13, '((2, 1), (0, 1))': 13, '((2, 1), (0, 2))': 13, '((2, 1), (1, -1))': 13, '((2, 1), (1, 0))': 13, '((2, 1), (2, 0))': 13, '((2, 2), (0, 0))': 21, '((2, 2), (1, -1))': 21, '((2, 2), (1, 0))': 21}, 'wilson': 4, 'negative': {'perturbed': 30, 'other\_conditions\_hold': 30, 'star\_rejects': 30}, 'seconds': 3247.8}; pure Spin(5): {'labels': 25, 'unreached': [], 'solves\_unreached': [], 'errors': {}, 'pinned': 25, 'equal\_to\_chart': 25, 'a3\_as\_printed': 25, 'unknown\_cells': {'((0, 0), (0, 0))': 0, '((0, 0), (0, 1))': 0, '((0, 0), (1, 0))': 0, '((0, 0), (2, 0))': 0, '((1, 1), (0, 0))': 1, '((1, 1), (0, 1))': 1, '((1, 1), (1, -1))': 1, '((1, 1), (1, 0))': 1, '((1, 1), (2, 0))': 1, '((2, 1), (-1, 1))': 5, '((2, 1), (0, 0))': 5, '((2, 1), (0, 1))': 5, '((2, 1), (1, 0))': 5, '((2, 1), (2, 0))': 5, '((2, 2), (0, 0))': 9, '((2, 2), (0, 1))': 9, '((2, 2), (1, -1))': 9, '((2, 2), (1, 0))': 9, '((2, 2), (2, 0))': 9, '((3, 2), (-1, 1))': 13, '((3, 2), (-1, 2))': 13, '((3, 2), (-2, 4))': 13, '((3, 2), (0, 0))': 13, '((3, 2), (0, 1))': 13, '((3, 2), (1, -1))': 13}, 'wilson': 4, 'negative': {'perturbed': 42, 'other\_conditions\_hold': 42, 'star\_rejects': 42}, 'seconds': 1256.9}; pure SO(5): {'labels': 16, 'unreached': [], 'solves\_unreached': [], 'errors': {}, 'pinned': 16, 'equal\_to\_chart': 16, 'a3\_as\_printed': 16, 'unknown\_cells': {'((0, 0), (0, 0))': 0, '((0, 0), (0, 2))': 0, '((0, 0), (1, 0))': 0, '((0, 0), (2, 0))': 0, '((1, 1), (0, 0))': 1, '((1, 1), (0, 2))': 1, '((1, 1), (1, 0))': 1, '((1, 1), (2, 0))': 1, '((2, 1), (0, 0))': 5, '((2, 1), (0, 2))': 5, '((2, 1), (1, 0))': 5, '((2, 1), (2, 0))': 5, '((Fraction(1, 1), Fraction(1, 2)), (0, 0))': 0, '((Fraction(1, 1), Fraction(1, 2)), (0, 2))': 0, '((Fraction(1, 1), Fraction(1, 2)), (1, 0))': 0, '((Fraction(1, 1), Fraction(1, 2)), (2, 0))': 0}, 'wilson': 4, 'negative': {'perturbed': 16, 'other\_conditions\_hold': 16, 'star\_rejects': 16}, 'seconds': 10.6}; pure G2: {'labels': 12, 'unreached': [], 'solves\_unreached': [], 'errors': {}, 'pinned': 12, 'equal\_to\_chart': 12, 'a3\_as\_printed': 12, 'unknown\_cells': {'((0, 0), (0, 0))': 0, '((0, 0), (0, 1))': 0, '((0, 0), (1, 0))': 0, '((1, 2), (0, 0))': 1, '((1, 2), (0, 1))': 1, '((1, 2), (1, 0))': 1, '((2, 3), (0, 0))': 7, '((2, 3), (0, 1))': 7, '((2, 3), (1, 0))': 7, '((2, 4), (0, 0))': 13, '((2, 4), (0, 1))': 13, '((2, 4), (1, 0))': 13}, 'wilson': 3, 'negative': {'perturbed': 18, 'other\_conditions\_hold': 18, 'star\_rejects': 18}, 'seconds': 1298.6}; SU(2)+1: {'labels': 18, 'unreached': [], 'solves\_unreached': [], 'errors': {}, 'pinned': 18, 'equal\_to\_chart': 18, 'a3\_as\_printed': 18, 'unknown\_cells': None, 'wilson': 3, 'negative': None, 'seconds': 7.9}; SU(2)+2: {'labels': 13, 'unreached': [], 'solves\_unreached': [], 'errors': {}, 'pinned': 13, 'equal\_to\_chart': 13, 'a3\_as\_printed': 13, 'unknown\_cells': None, 'wilson': 3, 'negative': None, 'seconds': 3.0}; SU(2)+3: {'labels': 8, 'unreached': [], 'solves\_unreached': [], 'errors': {}, 'pinned': 8, 'equal\_to\_chart': 8, 'a3\_as\_printed': 8, 'unknown\_cells': None, 'wilson': 2, 'negative': None, 'seconds': 5.0}; SU(2)+4: {'labels': 8, 'unreached': [], 'solves\_unreached': [], 'errors': {}, 'pinned': 8, 'equal\_to\_chart': 8, 'a3\_as\_printed': 8, 'unknown\_cells': None, 'wilson': 2, 'negative': None, 'seconds': 16.8}; SU(2)+adjoint (N=2*): {'labels': 13, 'unreached': [], 'solves\_unreached': [], 'errors': {}, 'pinned': 13, 'equal\_to\_chart': 13, 'a3\_as\_printed': 13, 'unknown\_cells': None, 'wilson': 3, 'negative': None, 'seconds': 2.1}; SO(3)+adjoint (N=2*): {'labels': 11, 'unreached': [], 'solves\_unreached': [], 'errors': {}, 'pinned': 11, 'equal\_to\_chart': 11, 'a3\_as\_printed': 11, 'unknown\_cells': None, 'wilson': 2, 'negative': None, 'seconds': 0.8}; rank 1 + adjoint (N=2*), the third global form: {'labels': 6, 'unreached': [], 'solves\_unreached': [], 'errors': {}, 'pinned': 6, 'equal\_to\_chart': 6, 'a3\_as\_printed': 6, 'unknown\_cells': None, 'wilson': 2, 'negative': None, 'seconds': 0.1}; SU(3)+adjoint (N=2*): {'labels': 6, 'unreached': [], 'solves\_unreached': [], 'errors': {}, 'pinned': 6, 'equal\_to\_chart': 6, 'a3\_as\_printed': 6, 'unknown\_cells': None, 'wilson': 3, 'negative': None, 'seconds': 3.4}; SU(3)+symmetric: {'labels': 10, 'unreached': [], 'solves\_unreached': [], 'errors': {}, 'pinned': 10, 'equal\_to\_chart': 10, 'a3\_as\_printed': 10, 'unknown\_cells': None, 'wilson': 4, 'negative': None, 'seconds': 5.0}; SU(3)+symmetric+1: {'labels': 6, 'unreached': [], 'solves\_unreached': [], 'errors': {}, 'pinned': 6, 'equal\_to\_chart': 6, 'a3\_as\_printed': 6, 'unknown\_cells': None, 'wilson': 3, 'negative': None, 'seconds': 4.7}; Spin(5)+vector: {'labels': 9, 'unreached': [], 'solves\_unreached': [], 'errors': {}, 'pinned': 9, 'equal\_to\_chart': 9, 'a3\_as\_printed': 9, 'unknown\_cells': None, 'wilson': 3, 'negative': None, 'seconds': 15.6}; SO(5)+vector: {'labels': 12, 'unreached': [], 'solves\_unreached': [], 'errors': {}, 'pinned': 12, 'equal\_to\_chart': 12, 'a3\_as\_printed': 12, 'unknown\_cells': None, 'wilson': 3, 'negative': None, 'seconds': 18.1}; G2+7: {'labels': 6, 'unreached': [], 'solves\_unreached': [], 'errors': {}, 'pinned': 6, 'equal\_to\_chart': 6, 'a3\_as\_printed': 6, 'unknown\_cells': None, 'wilson': 2, 'negative': None, 'seconds': 220.4}; SU(2)xSU(2)+bifundamental: {'labels': 16, 'unreached': [], 'solves\_unreached': [], 'errors': {}, 'pinned': 16, 'equal\_to\_chart': 16, 'a3\_as\_printed': 16, 'unknown\_cells': None, 'wilson': 4, 'negative': None, 'seconds': 4.0}; U(1)-U(2)+3: {'labels': 24, 'unreached': [], 'solves\_unreached': [], 'errors': {}, 'pinned': 24, 'equal\_to\_chart': 24, 'a3\_as\_printed': 24, 'unknown\_cells': None, 'wilson': 4, 'negative': None, 'seconds': 0.1}; A3 with printed matter factors: {'SU(3)+1': {'printed': 2, 'dual': 0, 'none': 0}, 'SU(3)+symmetric': {'printed': 2, 'dual': 0, 'none': 0}, 'Spin(5)+vector': {'printed': 2, 'dual': 2, 'none': 0}}. & \CERTIFIED\newline {\footnotesize conjecture, extensive\newline web}\newline {\footnotesize run web 2026-09-24: 34/34}\newline {\footnotesize run web (extensive) 2026-09-24: 83/83} \\ \addlinespace
\code{conj:abeKalgebra/consistency}\newline Conj.~\pref{conj:abeKalgebra}, ``consistent $K_\fq$-algebra'' & The family so defined is a $K_\fq$-algebra: closed under multiplication with $\bZ[\fq^{\pm1}]$ coefficients and satisfying K1--K5. & every label of every window builds by the axiom route (\code{route} `star'/`solve') or a declared optimization asserted equal to it (`twist', `monoid'), passes \code{certify_canonical} and, in pure gauge, the residue rule \code{criterion}.  Fast tier: at rank one the twelve K1--K5 verifiers on all pairs ($K=4$) and closure of the products with $\bZ[\fq^{\pm1}]$ coefficients (with matter, products of $m_a+m_b\le2$); at rank two the pair checks are the extensive tier's (products reach $m$ up to $2m$, e.g. SU(3) $m=(2,2)$ at 68 s per window).  Extensive tier, at every rank: the K1--K5 verifiers on all pairs of each window's first eight labels ($K=4$; the three-argument verifier on triples of the first four) and closure of their products with $\bZ[\fq^{\pm1}]$ coefficients (64 for an eight-label window, 36 for a six-label one), under a budget of 2400 s and 5 GB per theory, each theory in its own process; a label or battery not finished within the time or memory budget is reported in the population, never counted as a pass or a failure. \emph{Population:} web tier: pure SU(2), SO(3) ($m\le2$, $|e|\le2$), SU(3), Sp(4), $G_2$ (small $m$, $e$ windows), SU(2)+1, SU(2)+2, SO(3)+adjoint, SU(3)+1 (flavour 0); the extensive tier, recorded beside the fast one, runs the 24 theories of the Section~3 lineup on grown windows (\code{_extensive_windows}): pure SU(2) and SO(3) ($m\le4$, $|e|\le3$), the rank-one third global form, SU(3), PSU(3), Sp(4), Spin(5), SO(5), $G_2$; SU(2) with $N_f=1,\dots,4$; $N=2^*$ at SU(2), SO(3), the third global form and SU(3); SU(3) with a symmetric, alone and with a fundamental; Spin(5) and SO(5) with a vector; $G_2$ with a $\mathbf 7$; the SU(2)$\times$SU(2) quiver with a bifundamental; U(1)--U(2) with three fundamentals of U(2) (matter at the neutral flavour label). \emph{Evidence:} emergent (nothing in the solver imposes closure or orthonormality). \emph{Where:} \code{battery/checks_coulomb.py}, \code{tests/test_pure_g_abe_kalgebra.py}, \code{tests/test_g_matter_abe_kalgebra.py}. \emph{Note:} Since 2026-09-23 the rows sec:coulomb/labels, conj:abeKalgebra/existence-uniqueness and conj:abeKalgebra/support have their own adapters on the same windows (conjecture\_windows in checks\_coulomb.py); before that they shared this run.  2026-09-22 web run (checks\_coulomb.check\_conjecture\_windows): nine windows (pure SU(2), SO(3), SU(3), Sp(4), G2; SU(2)+1, SU(2)+2, SO(3)+adjoint, SU(3)+1) build, certify and satisfy ($\star$) on every label; K1-K5 on all pairs at K = 4 and closure at rank one.  Routes: the axiom route ('star' / 'solve') plus the declared optimizations theta\_twist and monoid, asserted equal to it -- a first version of the check demanded 'star' literally and flagged them.  Cost: the SU(2)+2 window's pair battery (rank-2 flavour ring) took 2364 s at K = 4 and is the local environment's; the rest of the adapter is $\sim$2.5 min.  2026-09-23 extensive run: axiom 5 in pairing form (verify\_pairing\_rho\_star\_symmetric) failed at SU(2)+1 on (L\_(1),(-2), L\_(1),(2)) at K = 4 -- not the algebra but the matter pairing's truncation: the Nahm window W = K + 2 is exact only for residuals whose q-expansion starts at q\^{}($\ge$ -2), and this pair's starts at q\^{}-4, so I(a,b) carried a spurious (mu\^{}6 + mu\^{}-6) q\^{}4 (the printed sum over m, which has no Nahm window, agrees with W $\ge$ 7).  Fixed in MatterWRQTorus.\_mu\_trace\_rows (window raised to K + pad; the audit) before this row's extensive record; the run from the unfixed code was stopped, not recorded. \emph{Run (web, 2026-09-23):} 13 checks, 0 failing; pure SU(2): {'labels': 13, 'unreached': [], 'errors': {}, 'built': 13, 'routes': {'star': 8, 'twist[k=-1 from (0,)]': 1, 'twist[k=1 from (-1,)]': 1, 'twist[k=1 from (0,)]': 1, 'monoid[(1,)+(1,)]': 1, 'twist[k=1 from (-2,)]': 1}, 'certified': 13, 'star': 13, 'verifier\_failures': 0, 'closure': True, 'seconds': 19.3}; pure SO(3): {'labels': 13, 'unreached': [], 'errors': {}, 'built': 13, 'routes': {'star': 8, 'twist[k=-2 from (0,)]': 1, 'twist[k=-1 from (0,)]': 1, 'monoid[(1,)+(1,)]': 1, 'twist[k=1 from (0,)]': 1, 'twist[k=1 from (1,)]': 1}, 'certified': 13, 'star': 13, 'verifier\_failures': 0, 'closure': True, 'seconds': 3.4}; pure SU(3): {'labels': 8, 'unreached': [], 'errors': {}, 'built': 8, 'routes': {'star': 8}, 'certified': 8, 'star': 8, 'verifier\_failures': None, 'closure': None, 'seconds': 0.3}; pure Sp(4): {'labels': 5, 'unreached': [], 'errors': {}, 'built': 5, 'routes': {'star': 5}, 'certified': 5, 'star': 5, 'verifier\_failures': None, 'closure': None, 'seconds': 0.1}; pure G2: {'labels': 4, 'unreached': [], 'errors': {}, 'built': 4, 'routes': {'star': 4}, 'certified': 4, 'star': 4, 'verifier\_failures': None, 'closure': None, 'seconds': 9.1}; SU(2) + 1 fund. (roster su2-nf1): {'labels': 5, 'unreached': [], 'errors': {}, 'built': 5, 'routes': {'solve': 4, 'twist[k=1 from (-1,)]': 1}, 'certified': 5, 'star': 5, 'verifier\_failures': 0, 'closure': True, 'seconds': 35.7}; SU(2) + 2 fund. (roster su2-nf2, U(2) flavour): {'labels': 5, 'unreached': [], 'errors': {}, 'built': 5, 'routes': {'solve': 4, 'twist[k=1 from (-1,)]': 1}, 'certified': 5, 'star': 5, 'verifier\_failures': None, 'closure': None, 'seconds': 0.1}; SO(3) + adjoint (N=2*): {'labels': 5, 'unreached': [], 'errors': {}, 'built': 5, 'routes': {'solve': 5}, 'certified': 5, 'star': 5, 'verifier\_failures': 0, 'closure': True, 'seconds': 53.4}; SU(3) + 1 fund. (roster su3-nf1): {'labels': 4, 'unreached': [], 'errors': {}, 'built': 4, 'routes': {'solve': 4}, 'certified': 4, 'star': 4, 'verifier\_failures': None, 'closure': None, 'seconds': 0.4}. \emph{Run (web (extensive), 2026-09-25):} 35 checks, 0 failing; pure SU(2): {'labels': 32, 'unreached': [], 'errors': {}, 'built': 32, 'routes': {'star': 20, 'twist[k=-1 from (0,)]': 1, 'twist[k=1 from (-3,)]': 3, 'twist[k=1 from (-1,)]': 2, 'twist[k=1 from (0,)]': 1, 'twist[k=1 from (1,)]': 1, 'monoid[(1,)+(1,)]': 1, 'twist[k=1 from (-2,)]': 1, 'monoid[(1,)+(2,)]': 1, 'monoid[(1,)+(3,)]': 1}, 'certified': 32, 'star': 32, 'verifier\_failures': 0, 'closure': True, 'seconds': 7.9}; pure SO(3): {'labels': 32, 'unreached': [], 'errors': {}, 'built': 32, 'routes': {'star': 18, 'twist[k=-3 from (0,)]': 1, 'twist[k=1 from (-3,)]': 2, 'twist[k=-1 from (0,)]': 2, 'monoid[(1,)+(1,)]': 1, 'twist[k=1 from (0,)]': 2, 'twist[k=1 from (1,)]': 2, 'twist[k=1 from (2,)]': 1, 'monoid[(1,)+(2,)]': 1, 'monoid[(1,)+(3,)]': 1, 'twist[k=1 from (-1,)]': 1}, 'certified': 32, 'star': 32, 'verifier\_failures': 0, 'closure': True, 'seconds': 1.7}; pure rank 1, the third global form: {'labels': 12, 'unreached': [], 'errors': {}, 'built': 12, 'routes': {'star': 6, 'twist[k=-1 from (0,)]': 2, 'twist[k=1 from (0,)]': 2, 'monoid[(1,)+(1,)]': 1, 'twist[k=1 from (-2,)]': 1}, 'certified': 12, 'star': 12, 'verifier\_failures': 0, 'closure': True, 'seconds': 4.6}; pure SU(3): {'labels': 16, 'unreached': [], 'errors': {}, 'built': 16, 'routes': {'star': 14, 'twist[k=1 from (1, -1)]': 1, 'monoid[(1, 1)+(1, 1)]': 1}, 'certified': 16, 'star': 16, 'verifier\_failures': 0, 'closure': True, 'seconds': 449.5}; pure PSU(3): {'labels': 16, 'unreached': [], 'errors': {}, 'built': 16, 'routes': {'star': 13, 'twist[k=3 from (0, 0)]': 2, 'twist[k=1 from (0, 0)]': 1}, 'certified': 16, 'star': 16, 'verifier\_failures': 0, 'closure': True, 'seconds': 21.5}; pure Sp(4): {'labels': 19, 'unreached': [], 'errors': {}, 'built': 19, 'routes': {'star': 17, 'monoid[(1, 1)+(1, 0)]': 1, 'monoid[(1, 1)+(1, 1)]': 1}, 'certified': 19, 'star': 19, 'verifier\_failures': 'not reached within 2400 s', 'closure': 'not reached within 2400 s', 'seconds': 3204.5}; pure Spin(5): {'labels': 25, 'unreached': [], 'errors': {}, 'built': 25, 'routes': {'star': 22, 'twist[k=1 from (0, 0)]': 1, 'monoid[(1, 1)+(1, 1)]': 1, 'monoid[(2, 1)+(1, 1)]': 1}, 'certified': 25, 'star': 25, 'verifier\_failures': 0, 'closure': True, 'seconds': 523.4}; pure SO(5): {'labels': 16, 'unreached': [], 'errors': {}, 'built': 16, 'routes': {'star': 12, 'twist[k=1 from (0, 0)]': 3, 'twist[k=1 from (1, 0)]': 1}, 'certified': 16, 'star': 16, 'verifier\_failures': 0, 'closure': True, 'seconds': 40.9}; pure G2: {'labels': 12, 'unreached': [], 'errors': {}, 'built': 12, 'routes': {'star': 10, 'twist[k=1 from (0, 0)]': 1, 'monoid[(1, 2)+(1, 2)]': 1}, 'certified': 12, 'star': 12, 'verifier\_failures': 'not reached within 5 GB', 'closure': 'not reached within 5 GB', 'seconds': 629.6}; SU(2)+1: {'labels': 18, 'unreached': [], 'errors': {}, 'built': 18, 'routes': {'solve': 14, 'twist[k=-1 from (0,)]': 1, 'twist[k=1 from (-1,)]': 1, 'twist[k=1 from (0,)]': 1, 'twist[k=1 from (-2,)]': 1}, 'certified': 18, 'star': 18, 'verifier\_failures': 0, 'closure': True, 'seconds': 195.9}; SU(2)+2: {'labels': 13, 'unreached': [], 'errors': {}, 'built': 13, 'routes': {'solve': 9, 'twist[k=-1 from (0,)]': 1, 'twist[k=1 from (-1,)]': 1, 'twist[k=1 from (0,)]': 1, 'twist[k=1 from (-2,)]': 1}, 'certified': 13, 'star': 13, 'verifier\_failures': 'not reached within 2400 s', 'closure': 'not reached within 2400 s', 'seconds': 2402.1}; SU(2)+3: {'labels': 8, 'unreached': [], 'errors': {}, 'built': 8, 'routes': {'solve': 7, 'twist[k=1 from (-1,)]': 1}, 'certified': 8, 'star': 8, 'verifier\_failures': 'not reached within 2400 s', 'closure': 'not reached within 2400 s', 'seconds': 2404.5}; SU(2)+4: {'labels': 8, 'unreached': [], 'errors': {}, 'built': 8, 'routes': {'solve': 7, 'twist[k=1 from (-1,)]': 1}, 'certified': 8, 'star': 8, 'verifier\_failures': 'not reached within 2400 s', 'closure': 'not reached within 2400 s', 'seconds': 2414.9}; SU(2)+adjoint (N=2*): {'labels': 13, 'unreached': [], 'errors': {}, 'built': 13, 'routes': {'solve': 9, 'twist[k=-1 from (0,)]': 1, 'twist[k=1 from (-1,)]': 1, 'twist[k=1 from (0,)]': 1, 'twist[k=1 from (-2,)]': 1}, 'certified': 13, 'star': 13, 'verifier\_failures': 0, 'closure': True, 'seconds': 1013.3}; SO(3)+adjoint (N=2*): {'labels': 11, 'unreached': [], 'errors': {}, 'built': 11, 'routes': {'solve': 7, 'twist[k=-2 from (0,)]': 1, 'twist[k=2 from (0,)]': 1, 'twist[k=-1 from (0,)]': 1, 'twist[k=1 from (0,)]': 1}, 'certified': 11, 'star': 11, 'verifier\_failures': 0, 'closure': True, 'seconds': 636.1}; rank 1 + adjoint (N=2*), the third global form: {'labels': 6, 'unreached': [], 'errors': {}, 'built': 6, 'routes': {'solve': 4, 'twist[k=2 from (-1,)]': 1, 'twist[k=1 from (0,)]': 1}, 'certified': 6, 'star': 6, 'verifier\_failures': 0, 'closure': True, 'seconds': 186.8}; SU(3)+adjoint (N=2*): {'labels': 6, 'unreached': [], 'errors': {}, 'built': 6, 'routes': {'solve': 5, 'twist[k=1 from (0, 0)]': 1}, 'certified': 6, 'star': 6, 'verifier\_failures': 'not reached within 2400 s', 'closure': 'not reached within 2400 s', 'seconds': 2402.3}; SU(3)+symmetric: {'labels': 10, 'unreached': [], 'errors': {}, 'built': 10, 'routes': {'solve': 9, 'twist[k=1 from (1, -1)]': 1}, 'certified': 10, 'star': 10, 'verifier\_failures': 'not reached within 2400 s', 'closure': 'not reached within 2400 s', 'seconds': 2402.8}; SU(3)+symmetric+1: {'labels': 6, 'unreached': [], 'errors': {}, 'built': 6, 'routes': {'solve': 6}, 'certified': 6, 'star': 6, 'verifier\_failures': 'not reached within 2400 s', 'closure': 'not reached within 2400 s', 'seconds': 2403.5}; Spin(5)+vector: {'labels': 9, 'unreached': [], 'errors': {}, 'built': 9, 'routes': {'solve': 9}, 'certified': 9, 'star': 9, 'verifier\_failures': 'not reached within 2400 s', 'closure': 'not reached within 2400 s', 'seconds': 2413.5}; SO(5)+vector: {'labels': 12, 'unreached': [], 'errors': {}, 'built': 12, 'routes': {'solve': 10, 'twist[k=1 from (0, 0)]': 2}, 'certified': 12, 'star': 12, 'verifier\_failures': 'not reached within 2400 s', 'closure': 'not reached within 2400 s', 'seconds': 2412.0}; G2+7: {'labels': 6, 'unreached': [], 'errors': {}, 'built': 6, 'routes': {'solve': 6}, 'certified': 6, 'star': 6, 'verifier\_failures': 'not reached within 5 GB', 'closure': 'not reached within 5 GB', 'seconds': 2537.8}; SU(2)xSU(2)+bifundamental: {'labels': 16, 'unreached': [], 'errors': {}, 'built': 16, 'routes': {'solve': 16}, 'certified': 16, 'star': 16, 'verifier\_failures': 'not reached within 2400 s', 'closure': 'not reached within 2400 s', 'seconds': 2402.8}; U(1)-U(2)+3: {'labels': 24, 'unreached': [], 'errors': {}, 'built': 24, 'routes': {'solve': 24}, 'certified': 24, 'star': 24, 'verifier\_failures': 'not reached within 2400 s', 'closure': 'not reached within 2400 s', 'seconds': 2400.2}. & \CERTIFIED\newline {\footnotesize conjecture, extensive\newline web}\newline {\footnotesize run web 2026-09-23: 13/13}\newline {\footnotesize run web (extensive) 2026-09-25: 35/35} \\ \addlinespace
\code{conj:abeKalgebra/support}\newline note after Conj.~\pref{conj:abeKalgebra} & The bubbling is supported on Weyl orbits subdominant to $m$, and this appears to follow from the axioms. & the bubbling of every window line lies in $\mathrm{conv}(W\cdot m)$ (\code{tropical_support}): structural, since the solver's unknowns live there and the declared optimizations preserve support.  The measurement of the remark, in pure gauge: the residue rule, bar invariance and $O(\fq)$, with Weyl covariance, solved on the cells of $\mathrm{conv}(W\cdot(m+k\beta))$ (the hull grown by $k$ steps of the smallest dominant coroot $\beta$; $k=2$ when one step adds only the Weyl orbit of $m+\beta$, the extremal orbit of the larger hull, always so at rank one, else $k=1$) have a unique solution, it vanishes outside $\mathrm{conv}(W\cdot m)$, and it equals the production line (\code{enlarged_support_one}; \code{star_bubbling.joint_solve} with \code{support=}).  With matter the same at every $\mu$-sector (\code{enlarged_support_one_matter}; \code{matter_multislot.solve_canonical_matter_vec} with \code{support=}): each sector is solved on the larger hull, the $O(\fq)$ condition on the quotient by the dressing $Z$ as always, and none is left to the degree law, which is a statement about the hull's cells and is not let decide the cells it is tested against; the tower vanishes outside $\mathrm{conv}(W\cdot m)$ at every sector, and on the hull its quotient by $Z$ equals the production line.  Control: without the residue rows the $O(\fq)$ conditions are cell-local, and the added Weyl-orbit representatives strictly inside the larger hull admit nonzero bar-invariant $O(\fq)$ fillings (non-pivot columns, \code{stats=}), so it is the residue rule that empties them; the outermost added layer, the extremal orbit of the larger hull, has no poles of its own and is emptied by bar invariance and $O(\fq)$ alone, which is why a step that adds only that orbit is doubled.  Grafts: a bubbling residual of a donor line (palindromic, $\mathrm{val}_\fq\ge1$) grafted onto a genuine line at a cell outside $\mathrm{conv}(W\cdot m)$ keeps bar invariance, the seed and $O(\fq)$ and, symmetrised over the Weyl orbit of the cell, Weyl covariance; the residue rule rejects every graft.  Each genuine line passes all five conditions of \code{star_bubbling.verify_axioms} and is Weyl-covariant. \emph{Population:} the nine conjecture windows (\code{conjecture_windows}): pure SU(2), SO(3) ($m\le2$, $|e|\le2$), SU(3), Sp(4), $G_2$; SU(2)+1, SU(2)+2, SO(3)+adjoint, SU(3)+1 (flavour 0); the extensive tier, recorded beside the fast one, runs the 24 theories of the Section~3 lineup on grown windows (\code{_extensive_windows}): pure SU(2) and SO(3) ($m\le4$, $|e|\le3$), the rank-one third global form, SU(3), PSU(3), Sp(4), Spin(5), SO(5), $G_2$; SU(2) with $N_f=1,\dots,4$; $N=2^*$ at SU(2), SO(3), the third global form and SU(3); SU(3) with a symmetric, alone and with a fundamental; Spin(5) and SO(5) with a vector; $G_2$ with a $\mathbf 7$; the SU(2)$\times$SU(2) quiver with a bifundamental; U(1)--U(2) with three fundamentals of U(2) (matter at the neutral flavour label) (containment); grafts: the five cases of a probe in the source repository (SU(2) $L_{2,0}$, $L_{2,1}$, $L_{1,0}$; SU(3) $L_{(1,1),0}$; Spin(5) $L_{(1,1),0}$), six grafts each, single-cell and Weyl-symmetrised; enlarged support: every window line of pure SU(2), SO(3), SU(3) and Sp(4) and of SU(2)+1, SU(2)+2, SO(3)+adjoint and SU(3)+1 at the fast depth, of all 24 theories of the extensive lineup at the extensive depth, every $\mu$-sector with matter ($G_2$ is extensive-only: one step of its smallest dominant coroot adds only the outermost orbit, and two steps at $L_{(1,0),0}$ outgrew a 3.5 GB memory cap). \emph{Controls:} positive: the genuine lines; the enlarged-support solve returns the production line at SU(2) $L_{2,0}$, and at SU(2)+adjoint $L_{2,0}$ every $\mu$-sector, before any theory is scanned; negative: the grafts (the support condition fails for them by construction; the finding is that the residue rule alone rejects them), and for the enlarged support the same system without the residue rows, which leaves the added cells strictly inside the larger hull free (with matter, sector by sector, each with its own offset). \emph{Evidence:} emergent. \emph{Where:} \code{battery/checks_coulomb.py (check_support, enlarged_support_one, enlarged_support_one_matter)}, \code{src/gn/star_bubbling.py (joint_solve: support=, stats=)}, \code{src/gn/matter_multislot.py (solve_canonical_matter_vec: support=, stats=)}, \code{src/gn/matter_star_bubbling.py (solve_level: support=, stats=, admits_e=)}, \code{kalgebra.md, the L_{m,e} construction}. \emph{Note:} Own adapter since 2026-09-23.  The repository's measurement (kalgebra.md, 30/30) grafts onto a single cell, which also breaks Weyl covariance, part of the paper's A1; so it does not isolate the residue rule against the paper's conditions.  The Weyl-symmetrised grafts, added here, do.  The enlarged-support measurement was added 2026-09-24 (the author approved the proposal made in the Section 3 walkthrough); with it the containment check is recorded as structural, not as evidence.  Extended to matter 2026-09-26: every $\mu$-sector is solved on the larger hull, including the sectors the degree law forces on the hull; measured the same day, those solve at the default window to the forced answer (all 47 at U(2)+2, m = (0,-2), (0,-3), (0,-4)). \emph{Run (web, 2026-09-26):} 37 checks, 0 failing; pure SU(2): {'labels': 13, 'unreached': [], 'contained': True, 'enlarged\_support': {'labels': 13, 'unreached': [], 'errors': {}, 'steps': ['2×(1,)'], 'unique': 13, 'zero': 13, 'equal': 13, 'added': 52, 'reps': 26, 'free': 13, 'free\_lines': 13, 'seconds': 2.4}}; pure SO(3): {'labels': 13, 'unreached': [], 'contained': True, 'enlarged\_support': {'labels': 13, 'unreached': [], 'errors': {}, 'steps': ['2×(2,)'], 'unique': 13, 'zero': 13, 'equal': 13, 'added': 52, 'reps': 26, 'free': 13, 'free\_lines': 13, 'seconds': 0.4}}; pure SU(3): {'labels': 8, 'unreached': [], 'contained': True, 'enlarged\_support': {'labels': 8, 'unreached': [], 'errors': {}, 'steps': ['1×(1, 1)', '2×(1, 1)'], 'unique': 8, 'zero': 8, 'equal': 8, 'added': 114, 'reps': 27, 'free': 19, 'free\_lines': 8, 'seconds': 14.7}}; pure Sp(4): {'labels': 5, 'unreached': [], 'contained': True, 'enlarged\_support': {'labels': 5, 'unreached': [], 'errors': {}, 'steps': ['1×(1, 0)', '2×(1, 0)'], 'unique': 5, 'zero': 5, 'equal': 5, 'added': 48, 'reps': 12, 'free': 7, 'free\_lines': 5, 'seconds': 8.5}}; pure G2: {'labels': 4, 'unreached': [], 'contained': True}; SU(2) + 1 fund. (roster su2-nf1): {'labels': 5, 'unreached': [], 'contained': True, 'enlarged\_support': {'labels': 5, 'unreached': [], 'errors': {}, 'steps': ['2×(1,)'], 'unique': 5, 'zero': 5, 'equal': 5, 'added': 20, 'reps': 16, 'free': 8, 'free\_lines': 5, 'sectors': 8, 'forced': 0, 'seconds': 0.5}}; SU(2) + 2 fund. (roster su2-nf2, U(2) flavour): {'labels': 5, 'unreached': [], 'contained': True, 'enlarged\_support': {'labels': 5, 'unreached': [], 'errors': {}, 'steps': ['2×(1,)'], 'unique': 5, 'zero': 5, 'equal': 5, 'added': 20, 'reps': 28, 'free': 14, 'free\_lines': 5, 'sectors': 14, 'forced': 0, 'seconds': 1.1}}; SO(3) + adjoint (N=2*): {'labels': 5, 'unreached': [], 'contained': True, 'enlarged\_support': {'labels': 5, 'unreached': [], 'errors': {}, 'steps': ['2×(2,)'], 'unique': 5, 'zero': 5, 'equal': 5, 'added': 20, 'reps': 16, 'free': 8, 'free\_lines': 5, 'sectors': 8, 'forced': 3, 'seconds': 0.2}}; SU(3) + 1 fund. (roster su3-nf1): {'labels': 4, 'unreached': [], 'contained': True, 'enlarged\_support': {'labels': 4, 'unreached': [], 'errors': {}, 'steps': ['1×(1, 1)', '2×(1, 1)'], 'unique': 4, 'zero': 4, 'equal': 4, 'added': 60, 'reps': 20, 'free': 14, 'free\_lines': 4, 'sectors': 6, 'forced': 0, 'seconds': 11.6}}; grafts: {'SU(2) L\_((2,), (0,)) (donor L\_((4,), (0,)))': {'positive\_control': True, 'single-cell': {'tried': 6, 'other\_conditions\_hold': 6, 'star\_rejects': 6}, 'Weyl-symmetrised': {'tried': 6, 'other\_conditions\_hold': 6, 'star\_rejects': 6}}, 'SU(2) L\_((2,), (1,)) (donor L\_((4,), (0,)))': {'positive\_control': True, 'single-cell': {'tried': 6, 'other\_conditions\_hold': 6, 'star\_rejects': 6}, 'Weyl-symmetrised': {'tried': 6, 'other\_conditions\_hold': 6, 'star\_rejects': 6}}, 'SU(2) L\_((1,), (0,)) (donor L\_((3,), (0,)))': {'positive\_control': True, 'single-cell': {'tried': 6, 'other\_conditions\_hold': 6, 'star\_rejects': 6}, 'Weyl-symmetrised': {'tried': 6, 'other\_conditions\_hold': 6, 'star\_rejects': 6}}, 'SU(3) L\_((1, 1), (0, 0)) (donor L\_((2, 2), (0, 0)))': {'positive\_control': True, 'single-cell': {'tried': 6, 'other\_conditions\_hold': 6, 'star\_rejects': 6}, 'Weyl-symmetrised': {'tried': 6, 'other\_conditions\_hold': 6, 'star\_rejects': 6}}, 'Spin(5) L\_((1, 1), (0, 0)) (donor L\_((2, 2), (0, 0)))': {'positive\_control': True, 'single-cell': {'tried': 6, 'other\_conditions\_hold': 6, 'star\_rejects': 6}, 'Weyl-symmetrised': {'tried': 6, 'other\_conditions\_hold': 6, 'star\_rejects': 6}}}. \emph{Run (web (extensive), 2026-09-26):} 84 checks, 0 failing; pure SU(2): {'labels': 32, 'unreached': [], 'contained': True, 'enlarged\_support': {'labels': 32, 'unreached': [], 'errors': {}, 'steps': ['2×(1,)'], 'unique': 32, 'zero': 32, 'equal': 32, 'added': 128, 'reps': 64, 'free': 32, 'free\_lines': 32, 'seconds': 76.6}}; pure SO(3): {'labels': 32, 'unreached': [], 'contained': True, 'enlarged\_support': {'labels': 32, 'unreached': [], 'errors': {}, 'steps': ['2×(2,)'], 'unique': 32, 'zero': 32, 'equal': 32, 'added': 128, 'reps': 64, 'free': 32, 'free\_lines': 32, 'seconds': 4.3}}; pure rank 1, the third global form: {'labels': 12, 'unreached': [], 'contained': True, 'enlarged\_support': {'labels': 12, 'unreached': [], 'errors': {}, 'steps': ['2×(1,)'], 'unique': 12, 'zero': 12, 'equal': 12, 'added': 48, 'reps': 24, 'free': 12, 'free\_lines': 12, 'seconds': 2.5}}; pure SU(3): {'labels': 16, 'unreached': [], 'contained': True, 'enlarged\_support': {'labels': 16, 'unreached': [], 'errors': {}, 'steps': ['1×(1, 1)', '2×(1, 1)'], 'unique': 16, 'zero': 16, 'equal': 16, 'added': 252, 'reps': 52, 'free': 36, 'free\_lines': 16, 'seconds': 1436.7}}; pure PSU(3): {'labels': 16, 'unreached': [], 'contained': True, 'enlarged\_support': {'labels': 16, 'unreached': [], 'errors': {}, 'steps': ['1×(1, 1)', '2×(1, 1)'], 'unique': 16, 'zero': 16, 'equal': 16, 'added': 192, 'reps': 44, 'free': 28, 'free\_lines': 16, 'seconds': 35.0}}; pure Sp(4): {'labels': 19, 'unreached': [], 'contained': True, 'enlarged\_support': {'labels': 19, 'unreached': ['((2, 1), (-1, 1))', '((2, 1), (0, 0))', '((2, 1), (0, 1))', '((2, 1), (0, 2))', '((2, 1), (1, -1))', '((2, 1), (1, 0))', '((2, 1), (2, 0))', '((2, 2), (0, 0))', '((2, 2), (1, -1))', '((2, 2), (1, 0))', '((2, 2), (2, 0))'], 'errors': {}, 'steps': ['1×(1, 0)', '2×(1, 0)'], 'unique': 8, 'zero': 8, 'equal': 8, 'added': 96, 'reps': 20, 'free': 12, 'free\_lines': 8, 'seconds': 2512.6}}; pure Spin(5): {'labels': 25, 'unreached': [], 'contained': True, 'enlarged\_support': {'labels': 25, 'unreached': ['((3, 2), (-2, 4))', '((3, 2), (-1, 1))', '((3, 2), (-1, 2))', '((3, 2), (0, 0))', '((3, 2), (0, 1))', '((3, 2), (1, -1))'], 'errors': {}, 'steps': ['1×(1, 1)', '2×(1, 1)'], 'unique': 19, 'zero': 19, 'equal': 19, 'added': 208, 'reps': 42, 'free': 23, 'free\_lines': 19, 'seconds': 1970.9}}; pure SO(5): {'labels': 16, 'unreached': [], 'contained': True, 'enlarged\_support': {'labels': 16, 'unreached': [], 'errors': {}, 'steps': ['1×(1, 1)', '2×(1, 1)'], 'unique': 16, 'zero': 16, 'equal': 16, 'added': 208, 'reps': 40, 'free': 24, 'free\_lines': 16, 'seconds': 829.5}}; pure G2: {'labels': 12, 'unreached': [], 'contained': True, 'enlarged\_support': {'labels': 12, 'unreached': ['((2, 3), (0, 0))', '((2, 3), (0, 1))', '((2, 3), (1, 0))', '((2, 4), (0, 0))', '((2, 4), (0, 1))', '((2, 4), (1, 0))'], 'errors': {}, 'steps': ['1×(1, 2)', '2×(1, 2)'], 'unique': 6, 'zero': 6, 'equal': 6, 'added': 90, 'reps': 15, 'free': 9, 'free\_lines': 6, 'seconds': 1999.5}}; SU(2)+1: {'labels': 18, 'unreached': [], 'contained': True, 'enlarged\_support': {'labels': 18, 'unreached': [], 'errors': {}, 'steps': ['2×(1,)'], 'unique': 18, 'zero': 18, 'equal': 18, 'added': 72, 'reps': 96, 'free': 48, 'free\_lines': 18, 'sectors': 48, 'forced': 0, 'seconds': 74.7}}; SU(2)+2: {'labels': 13, 'unreached': [], 'contained': True, 'enlarged\_support': {'labels': 13, 'unreached': [], 'errors': {}, 'steps': ['2×(1,)'], 'unique': 13, 'zero': 13, 'equal': 13, 'added': 52, 'reps': 136, 'free': 68, 'free\_lines': 13, 'sectors': 68, 'forced': 0, 'seconds': 44.6}}; SU(2)+3: {'labels': 8, 'unreached': [], 'contained': True, 'enlarged\_support': {'labels': 8, 'unreached': [], 'errors': {}, 'steps': ['2×(1,)'], 'unique': 8, 'zero': 8, 'equal': 8, 'added': 32, 'reps': 214, 'free': 107, 'free\_lines': 8, 'sectors': 107, 'forced': 0, 'seconds': 97.2}}; SU(2)+4: {'labels': 8, 'unreached': [], 'contained': True, 'enlarged\_support': {'labels': 8, 'unreached': [], 'errors': {}, 'steps': ['2×(1,)'], 'unique': 8, 'zero': 8, 'equal': 8, 'added': 32, 'reps': 586, 'free': 293, 'free\_lines': 8, 'sectors': 293, 'forced': 0, 'seconds': 402.3}}; SU(2)+adjoint (N=2*): {'labels': 13, 'unreached': [], 'contained': True, 'enlarged\_support': {'labels': 13, 'unreached': [], 'errors': {}, 'steps': ['2×(1,)'], 'unique': 13, 'zero': 13, 'equal': 13, 'added': 52, 'reps': 86, 'free': 43, 'free\_lines': 13, 'sectors': 43, 'forced': 0, 'seconds': 54.8}}; SO(3)+adjoint (N=2*): {'labels': 11, 'unreached': [], 'contained': True, 'enlarged\_support': {'labels': 11, 'unreached': [], 'errors': {}, 'steps': ['2×(1,)'], 'unique': 11, 'zero': 11, 'equal': 11, 'added': 44, 'reps': 58, 'free': 29, 'free\_lines': 11, 'sectors': 29, 'forced': 6, 'seconds': 18.8}}; rank 1 + adjoint (N=2*), the third global form: {'labels': 6, 'unreached': [], 'contained': True, 'enlarged\_support': {'labels': 6, 'unreached': [], 'errors': {}, 'steps': ['2×(1,)'], 'unique': 6, 'zero': 6, 'equal': 6, 'added': 24, 'reps': 24, 'free': 12, 'free\_lines': 6, 'sectors': 12, 'forced': 2, 'seconds': 1.4}}; SU(3)+adjoint (N=2*): {'labels': 6, 'unreached': [], 'contained': True, 'enlarged\_support': {'labels': 6, 'unreached': [], 'errors': {}, 'steps': ['1×(1, 1)', '2×(1, 1)'], 'unique': 6, 'zero': 6, 'equal': 6, 'added': 90, 'reps': 57, 'free': 39, 'free\_lines': 6, 'sectors': 18, 'forced': 0, 'seconds': 165.5}}; SU(3)+symmetric: {'labels': 10, 'unreached': [], 'contained': True, 'enlarged\_support': {'labels': 10, 'unreached': [], 'errors': {}, 'steps': ['1×(1, 1)', '2×(1, 1)'], 'unique': 10, 'zero': 10, 'equal': 10, 'added': 144, 'reps': 88, 'free': 60, 'free\_lines': 10, 'sectors': 28, 'forced': 0, 'seconds': 157.1}}; SU(3)+symmetric+1: {'labels': 6, 'unreached': [], 'contained': True, 'enlarged\_support': {'labels': 6, 'unreached': [], 'errors': {}, 'steps': ['1×(1, 1)', '2×(1, 1)'], 'unique': 6, 'zero': 6, 'equal': 6, 'added': 90, 'reps': 84, 'free': 57, 'free\_lines': 6, 'sectors': 27, 'forced': 0, 'seconds': 167.5}}; Spin(5)+vector: {'labels': 9, 'unreached': [], 'contained': True, 'enlarged\_support': {'labels': 9, 'unreached': [], 'errors': {}, 'steps': ['1×(1, 1)', '2×(1, 1)'], 'unique': 9, 'zero': 9, 'equal': 9, 'added': 96, 'reps': 45, 'free': 24, 'free\_lines': 9, 'sectors': 21, 'forced': 0, 'seconds': 555.8}}; SO(5)+vector: {'labels': 12, 'unreached': [], 'contained': True, 'enlarged\_support': {'labels': 12, 'unreached': [], 'errors': {}, 'steps': ['1×(1, 1)', '2×(1, 1)'], 'unique': 12, 'zero': 12, 'equal': 12, 'added': 156, 'reps': 63, 'free': 36, 'free\_lines': 12, 'sectors': 27, 'forced': 3, 'seconds': 1771.0}}; G2+7: {'labels': 6, 'unreached': [], 'contained': True, 'enlarged\_support': {'labels': 6, 'unreached': ['(((2, 3), (0, 0)), (0,))', '(((2, 3), (1, 0)), (0,))'], 'errors': {}, 'steps': ['1×(1, 2)', '2×(1, 2)'], 'unique': 4, 'zero': 4, 'equal': 4, 'added': 60, 'reps': 18, 'free': 10, 'free\_lines': 4, 'sectors': 8, 'forced': 0, 'seconds': 761.1}}; SU(2)xSU(2)+bifundamental: {'labels': 16, 'unreached': [], 'contained': True, 'enlarged\_support': {'labels': 16, 'unreached': [], 'errors': {}, 'steps': ['1×(1, 0)', '2×(1, 0)'], 'unique': 16, 'zero': 16, 'equal': 16, 'added': 80, 'reps': 80, 'free': 40, 'free\_lines': 16, 'sectors': 40, 'forced': 0, 'seconds': 23.4}}; U(1)-U(2)+3: {'labels': 24, 'unreached': [], 'contained': True, 'enlarged\_support': {'labels': 24, 'unreached': [], 'errors': {}, 'steps': ['2×(0, 1, -1)'], 'unique': 24, 'zero': 24, 'equal': 24, 'added': 96, 'reps': 88, 'free': 44, 'free\_lines': 24, 'sectors': 44, 'forced': 20, 'seconds': 0.8}}; grafts: {'SU(2) L\_((2,), (0,)) (donor L\_((4,), (0,)))': {'positive\_control': True, 'single-cell': {'tried': 6, 'other\_conditions\_hold': 6, 'star\_rejects': 6}, 'Weyl-symmetrised': {'tried': 6, 'other\_conditions\_hold': 6, 'star\_rejects': 6}}, 'SU(2) L\_((2,), (1,)) (donor L\_((4,), (0,)))': {'positive\_control': True, 'single-cell': {'tried': 6, 'other\_conditions\_hold': 6, 'star\_rejects': 6}, 'Weyl-symmetrised': {'tried': 6, 'other\_conditions\_hold': 6, 'star\_rejects': 6}}, 'SU(2) L\_((1,), (0,)) (donor L\_((3,), (0,)))': {'positive\_control': True, 'single-cell': {'tried': 6, 'other\_conditions\_hold': 6, 'star\_rejects': 6}, 'Weyl-symmetrised': {'tried': 6, 'other\_conditions\_hold': 6, 'star\_rejects': 6}}, 'SU(3) L\_((1, 1), (0, 0)) (donor L\_((2, 2), (0, 0)))': {'positive\_control': True, 'single-cell': {'tried': 6, 'other\_conditions\_hold': 6, 'star\_rejects': 6}, 'Weyl-symmetrised': {'tried': 6, 'other\_conditions\_hold': 6, 'star\_rejects': 6}}, 'Spin(5) L\_((1, 1), (0, 0)) (donor L\_((2, 2), (0, 0)))': {'positive\_control': True, 'single-cell': {'tried': 6, 'other\_conditions\_hold': 6, 'star\_rejects': 6}, 'Weyl-symmetrised': {'tried': 6, 'other\_conditions\_hold': 6, 'star\_rejects': 6}}}. & \MEASURED\newline {\footnotesize conjecture, extensive\newline web}\newline {\footnotesize run web 2026-09-26: 37/37}\newline {\footnotesize run web (extensive) 2026-09-26: 84/84} \\ \addlinespace
\code{sec:coulomb/axioms-subsection}\newline Section~\pref{sec:coulomb}, ``The $K_\fq$ algebra axioms'' & The residue rule is equivalent to a well-defined action of the rational torus on the characters $\chi_e(v)$; $O(\fq)$ bubbling with \peqref{eq:Iexplicit} implies $\delta_{a,b}+O(\fq)$; the pole cancellations behind $I=\Tr\rho(f)g$ occur. & on 13 SU(2) lines: \code{criterion}$(x)$ holds iff \code{neumann_row}$(x)$ has no denominator (the action on characters is well defined), both true on every built line and both false on the bare $U_2+U_{-2}$; the bubbling residuals are $O(\fq)$ as series (A2) and $I_{a,b}=\delta_{a,b}+O(\fq)$ on all pairs; the two faces $\Tr\rho(f)g=\Tr g\rho^{-1}(f)$ agree on all pairs.  Widened 2026-09-23: the draft's statement proper, the residue rule as a well-defined action on the space of $G$ characters, $u^m\chi_e(v)=\chi_e(\fq^{2m}v)$: $\sum_a d_a(v)\chi_e(\fq^{2a}v)$ is a Laurent polynomial for every character of a window (the Neumann row above tests $\chi_0=1$ only, which the repository calls strictly weaker), at SU(2), SU(3), Spin(5), SU(2)+1 and SU(3)+1 (with matter, the $d$-form with the printed matter factors, power by power in $\mu$). \emph{Population:} random elements and every built chart. \emph{Controls:} negative: lines with their bubbling removed or doubled act on no character. \emph{Evidence:} emergent. \emph{Where:} \code{battery/checks_coulomb.py}, \code{src/gn/star_bubbling.py}, \code{notes/residue_cancellation_recognition.md}. \emph{Note:} the equivalence of the two readings is stated in the repo notes; the battery makes it a running check  2026-09-22 web run: on 13 SU(2) lines the residue rule is equivalent to a denominator-free action on characters (criterion $\le$> neumann\_row), the bubbling corrections are O(q) with the pairing $\delta$ + O(q) on all pairs, and the two faces I\_fg = Tr $\rho$(f) g = Tr g $\rho$\^{}-1(f) agree on all pairs. \emph{Run (web, 2026-09-23):} 5 checks, 0 failing; window: 13 SU(2) lines; action on characters: {'pure SU(2)': '5/5 lines x 4 characters', 'pure SU(3)': '3/3 lines x 5 characters', 'pure Spin(5)': '2/2 lines x 3 characters', 'SU(2)+1': '3/3 lines x 3 characters', 'SU(3)+1': '2/2 lines x 3 characters'}. & \CERTIFIED\newline {\footnotesize theorem, fast\newline web}\newline {\footnotesize run web 2026-09-23: 5/5} \\ \addlinespace
\code{sec:coulomb/kp-category}\newline Section~\pref{sec:coulomb}, the categorical remarks & The category $\mathrm{KP}(G,N)$ lifts the product, trace and $\rho$; the bar involution is known only at $N=0$. & no categorical code. & \OPEN\newline {\footnotesize not-testable, none\newline no test} \\ \addlinespace
\code{sec:coulomb/su2-examples}\newline \peqref{eq:su2thooft}--\peqref{eq:su2dyonic} & The printed $L_{0,2}$, $L_{2,0}$ (with its bubbling $f_0$), the $d$-form and its three residue cancellations, and $L_{2,1}$; $f_0$ is rigid; no Laurent polynomial is both bar-invariant and $O(\fq)$. & residual by residual against the printed $L_{0,2}$, $L_{2,0}$, $L_{2,1}$ (exact rational identities); $d_a=f_a(\fq^av)\mathrm{cocha}_a(v)$ recomputed from the printed $f$'s and $\mathrm{cocha}$ against the printed $d$'s, and against \code{star_bubbling.conventional}; the three printed sums of \peqref{eq:su2res}; the residue rule by exact division (\code{criterion}) for $L_{2,0}$, failing for the bare $U_2+U_{-2}$; $f_0$ bar-invariant, $O(\fq)$ and not a Laurent polynomial; $L_{2,0}=L_{1,0}^2$ in the SO(3) form. \emph{Population:} pure $SU(2)$ and $SO(3)$. \emph{Controls:} positive: the printed residuals reproduced exactly; negative: the bare 't Hooft line fails the residue rule. \emph{Evidence:} emergent. \emph{Where:} \code{battery/checks_coulomb.py}, \code{tests/test_pure_g_abe_kalgebra.py}. \emph{Run (web, 2026-09-22):} 10 checks, 0 failing; lines: L\_02, L\_20, L\_21 (SU(2)); L\_10 (SO(3)). & \CERTIFIED\newline {\footnotesize identity, fast\newline web}\newline {\footnotesize run web 2026-09-22: 10/10} \\ \addlinespace
\code{sec:coulomb/adjoint}\newline \peqref{eq:adjsquare}, \peqref{eq:adjbubble} & For $N=\mathrm{Adj}$: $\kappa\equiv0$, $\rho$ is the antipode, $\kappa_f(m)=-m$; $L_{1,0}^2=L_{2,0}+\mu$ and the printed $f_0(L_{2,0})$. & \code{GNAbeKAlgebra(so_n(3),(1,),nf=1)}: $L_{1,0}\cdot L_{1,0}=L_{2,0}+\mu$ as elements; $f_0(L_{2,0})$ level by level in $\mu$ against the printed formula (SO(3) coordinates $v^2\to v$); $\kappa\equiv0$ and $\kappa_f(m)=-m$ from the weights; $\rho$ the antipode with the flavour shift on 175 labels; the class's certification of $L_{2,0}$.  The $SU(2)$ form (added 2026-09-24; before, only the $SO(3)$ form was run although the population named both): \peqref{eq:adjbubble} verbatim in $SU(2)$ coordinates, where the draft's $L_{2,0}$ is the coroot cell and its $v^2$ is $v^\alpha$; the $U$-coefficients 1; $\rho$ the antipode with the flavour shifted by $-2m$ in $SU(2)$ units (the draft's $\kappa_f(m)=-m$ in its units, where the coroot is $m=2$), on 140 labels. \emph{Population:} $SO(3)$ and $SU(2)$ forms of the $N=2^*$ theory: \code{GNAbeKAlgebra(so_n(3), (1,), nf=1)} and \code{GNAbeKAlgebra(su_2(), highest_root, nf=1)}; \peqref{eq:adjsquare} in the $SO(3)$ form only, where $L_{1,0}$ exists. \emph{Evidence:} emergent. \emph{Where:} \code{battery/checks_coulomb.py}, \code{src/gn/g_matter_roster.py}. \emph{Run (web, 2026-09-24):} 8 checks, 0 failing; lines: SO(3) form: L\_10, L\_20, $\mu$, 175 labels for rho; SU(2) form: L\_20, 140 labels for rho. & \CERTIFIED\newline {\footnotesize identity, fast\newline web}\newline {\footnotesize run web 2026-09-24: 8/8} \\ \addlinespace
\end{longtable}

\subsection*{Section~\pref{sec:rg}: RG flows}
\begin{longtable}{@{}>{\raggedright\arraybackslash}p{2.4cm}>{\raggedright\arraybackslash}p{4.7cm}>{\raggedright\arraybackslash}p{4.95cm}>{\raggedright\arraybackslash}p{1.85cm}@{}}
\toprule claim & statement & test: procedure, population, controls & standing \\ \midrule \endfirsthead
\toprule claim & statement & test: procedure, population, controls & standing \\ \midrule \endhead
\bottomrule \endfoot
\code{def:rg/R1}\newline Def.~\pref{def:rg}, R1 & a lattice $\Gamma$ and a $\Gamma$-graded IR $K_\fq$-algebra & Not a claim: a definition has no truth value, and checking an axiom on a realisation tests the realisation, not the definition.  This row records where the axiom is \emph{enforced} by construction (a pass is a regression guard, left to the suite) and where it is \emph{emergent} (a pass is evidence), and names the rows that test it there.  On solver-built flows bar-invariance of $\RG(a)$ is enforced; multiplicativity is emergent and tested on the flow constructions, \code{conj:upper}, \code{conj:factored}, \code{conj:cluster}.  Instruments: \code{src/rg/rgkalgebra.py} \emph{Evidence:} per flow (emergent for multiplicativity, discovery, twist, pairing; enforced for bar-invariance on solver-built flows). \emph{Where:} \code{src/rg/rgkalgebra.py}. & ---\newline {\footnotesize definition, not run} \\ \addlinespace
\code{def:rg/R2}\newline Def.~\pref{def:rg}, R2 & an injective algebra map $\RG$ with bar-invariant images that are integral combinations of $[n]_\fq L^\IR_b$ & Not a claim: a definition has no truth value, and checking an axiom on a realisation tests the realisation, not the definition.  This row records where the axiom is \emph{enforced} by construction (a pass is a regression guard, left to the suite) and where it is \emph{emergent} (a pass is evidence), and names the rows that test it there.  Injectivity is analytic on BPS charts (distinct labels have distinct leading monomials by R3); the planned instrument is for flows built otherwise.  The map's multiplicativity is the content, tested on \code{conj:upper}, \code{conj:factored}.  Instruments: \code{src/rg/rgkalgebra.py} \emph{Evidence:} per flow (emergent for multiplicativity, discovery, twist, pairing; enforced for bar-invariance on solver-built flows). \emph{Where:} \code{src/rg/rgkalgebra.py}. & ---\newline {\footnotesize definition, not run} \\ \addlinespace
\code{def:rg/R3}\newline Def.~\pref{def:rg}, R3 & $S\in\bZ[\fq,(1-\fq^{2n})^{-1}]\otimes K_\fq[\IR]_{\Gamma_+}$, $\RG_aS=S\rho^{-1}_\IR(\RG_{\rho_\UV(a)})=L^\IR_{\ell(a)}+O(\fq)$, $\ell$ bijective, $S=1+O(\fq)$ & Not a claim: a definition has no truth value, and checking an axiom on a realisation tests the realisation, not the definition.  This row records where the axiom is \emph{enforced} by construction (a pass is a regression guard, left to the suite) and where it is \emph{emergent} (a pass is evidence), and names the rows that test it there.  On BPS charts $\RG(a)S=X_a+O(\fq)$ is the F-solver's defining relation (enforced); what is emergent is that the solution exists and is unique, \code{sec:rg/F-unique}, and that $S$ is a spectrum generator at all, \code{conj:cluster}, \code{conj:S-unique}.  Instruments: \code{src/rg/rgkalgebra.py} \emph{Evidence:} per flow (emergent for multiplicativity, discovery, twist, pairing; enforced for bar-invariance on solver-built flows). \emph{Where:} \code{src/rg/rgkalgebra.py}. & ---\newline {\footnotesize definition, not run} \\ \addlinespace
\code{def:rg/R4}\newline Def.~\pref{def:rg}, R4 & $(L^\UV_a|L^\UV_b)_\UV=(\RG_aS|\RG_bS)_\IR$ & Not a claim: a definition has no truth value, and checking an axiom on a realisation tests the realisation, not the definition.  This row records where the axiom is \emph{enforced} by construction (a pass is a regression guard, left to the suite) and where it is \emph{emergent} (a pass is evidence), and names the rows that test it there.  Where the UV trace is \emph{defined} through the IR the axiom is enforced; emergent on flows with an independent UV trace, and on matter removal where the IR is the Coulomb-branch tier: \code{conj:factored}, \code{sec:coulomb/labels}.  Instruments: \code{src/rg/rgkalgebra.py} \emph{Evidence:} per flow (emergent for multiplicativity, discovery, twist, pairing; enforced for bar-invariance on solver-built flows). \emph{Where:} \code{src/rg/rgkalgebra.py}. & ---\newline {\footnotesize definition, not run} \\ \addlinespace
\code{def:rg/no-exotic}\newline footnote to R2 & (Not imposed.)  The coefficients of $\RG_a$ are non-negative combinations of $[n]_\fq$. & each coefficient of $\RG_a$ written in $\fq$-numbers $[n]_\fq$ --- the BPS solver's native form, and for matter removal the decomposition of the palindromic Laurent polynomial; a negative coefficient is a recorded counterexample.  The census also counts the coefficients carrying a $\fq$-number beyond $[1]_\fq$, so that a census of $1$'s alone is visible as such. \emph{Population:} the Section~4 lineup's BPS charts, the Markov quiver included (node charges, pairwise sums, minus the first node; the Markov quiver's nodes); the $11$ dictionary-sample charts of the \pref{conj:upper} row; three matter-removal flows ($SU(2)$ with $N_f=1,2$ and $N=2^*$ $SU(2)$: a Wilson line, the minimal monopole and every label in the expansions of their products). \emph{Controls:} positive: $\fq^{-2}+\fq^{2}=[3]_\fq-[1]_\fq$ is flagged, $\fq^{-1}+\fq=[2]_\fq$ is not. \emph{Evidence:} emergent. \emph{Where:} \code{battery/checks_rg.py (check_no_exotic)}. \emph{Inputs:} \emph{web} --- \code{dictionaries/enumerated/} (tracked; E, accept\_depth, verify\_depth read from \code{manifest.json}), one chart per entry, a label window each; \emph{local} --- the enumerated build beyond the shipped weight (E $\ge$ 13: the weight-13 tier, which stays on the local machine by the author's ruling, at a package of the source repository) and the flavoured builds beyond the shipped weights (gitignored while in progress: \code{.flavoured_*stage/}); \code{dictionaries_weak/} (gitignored, local machine only: 579,293 entries with a spec each, ranks up to 15, several known specs per quiver; a valid source by the author's ruling of 2026-09-21, being integrated into the enumerated build as harvested seeds, so its weight $\le$ E quivers reach the shipped tiers as they ship while its tail beyond E stays a local population). \emph{Note:} 2026-09-27 web run: 0 negative among 1,333 coefficients, 774 of them carrying a q-number beyond [1]\_q -- the lineup 277 (the Markov quiver 18), the dictionary sample 989, matter removal 67.  On the pentagon and [A\_1,D\_4] every coefficient is 1, so there the statement holds trivially. \emph{Run (web, 2026-09-27):} 2 checks, 0 failing; coefficients: {'pentagon': 8, 'A\_3': 17, '[A\_1,D\_4]': 39, 'Kronecker (pure SU(2))': 16, 'SU(2)+1': 42, 'SU(2)+2': 54, 'SU(2)+3': 83, 'Markov': 18, 'dictionary sample': 989, 'matter removal SU(2)+1': 13, 'matter removal SU(2)+2': 26, 'matter removal SU(2)+adjoint (N=2*)': 28}; matter-removal labels: {'SU(2)+1': 8, 'SU(2)+2': 9, 'SU(2)+adjoint (N=2*)': 10}; seconds: 35.6. & \MEASURED\newline {\footnotesize conjecture, fast\newline web, local}\newline {\footnotesize run web 2026-09-27: 2/2} \\ \addlinespace
\code{def:rg/reconstruction}\newline Remark 4.2 & $\RG_aS=L^\IR_{\ell(a)}+O(\fq)$ determines $\RG$; the UV algebra is reconstructed from $(K_\fq[\IR],S)$. & the \code{KAlgebraIso} battery (unit, round trip, products, $\rho$, traces through $\fq^8$, $\fq^{10}$ for $SU(2)$) between the UV algebra rebuilt from $(K_\fq[\IR],S)$ by the F-solver ($F(a)S=X_a+O(\fq)$) and the independent presentation. \emph{Population:} three flows whose UV algebra has a presentation that does not use $S$: the pentagon chart (against its closed-form, cone, $[A_1,A_{2k}]$, SQED$_1$-flow and skein presentations), the Kronecker chart of pure $SU(2)$ (against the abelianized presentation), the SQED$_2$ chart (against $U_\fq(\fsl_2)$ built from its definition). \emph{Controls:} positive: the three isomorphisms; negative: the pentagon rebuilt from the reversed product $E_\fq(X_{\gamma_2})E_\fq(X_{\gamma_1})$, whose $S$ violates \peqref{eq:quiver} and whose pairing does not return within the budget, while its node products equal the correct ones. \emph{Evidence:} emergent. \emph{Where:} \code{battery/checks_rg.py}, \code{tests/test_kalgebra_object.py}, \code{tests/test_pure_su2_bps_iso.py}, \code{tests/test_uq_su2_bps_iso.py}. \emph{Note:} 2026-09-26: adapter checks\_rg.check\_reconstruction.  A wrong S is invisible to products at low labels (measured on the pentagon: the reversed factor order leaves the node lines and their products unchanged), so the pairings carry the test.  That the relation determines RG uniquely is sec:rg/F-unique, not this row. \emph{Run (web, 2026-09-26):} 4 checks, 0 failing; presentations: ['pentagon: BPS chart vs closed form, cone, a1a2k, sqed1 flow, skein (traces q\^{}8)', 'pure SU(2): Kronecker chart vs abelianized (traces q\^{}10)', 'U\_q(sl\_2) vs the SQED\_2 chart (traces q\^{}8)']; seconds: 13.4. & \CERTIFIED\newline {\footnotesize theorem, fast\newline web}\newline {\footnotesize run web 2026-09-26: 4/4} \\ \addlinespace
\code{def:rg/composition}\newline Remark 4.3 & Flows compose: $\RG^{1\to3}=\RG^{2\to3}\circ\RG^{1\to2}$ and $S^{1\to3}=\RG^{2\to3}(S^{1\to2})S^{2\to3}$. & \code{RGKAlgebra.then} against the direct flow: $\RG$ images on a label window, $S^{1\to3}$ through grading height 3, products of the composite against the UV chart, unital and multiplicative. \emph{Population:} node deletions: $A_4$ deleting $\gamma_1$ then $\gamma_2$ against both at once (directional flows); $A_3$ likewise (BPS node-deletion flows). \emph{Controls:} positive: the direct two-node deletion; negative: $S$ composed in the other order, $S^{2\to3}\RG^{2\to3}(S^{1\to2})$, differs from it. \emph{Evidence:} emergent. \emph{Where:} \code{battery/checks_rg.py}, \code{tests/test_directional_subquiver_rg.py}. \emph{Note:} 2026-09-26: adapter checks\_rg.check\_composition.  The law follows from R3 and R4 applied term by term, so this row tests the code (RGKAlgebra.then, ComposedRG), not the statement.  The matter-removal tower SU(2)+3 $\to$ +2 $\to$ +1 $\to$ pure does not compose through then (each rung's IR carries the removed flavour as a spectator that the next rung's UV lacks); its composition is checked by hand in def:rg/matter-removal. \emph{Run (web, 2026-09-26):} 5 checks, 0 failing; flows: ['A\_4: delete $\gamma$1 then $\gamma$2 vs both (directional)', 'A\_3: delete $\gamma$1 then $\gamma$2 vs both (BPS node-deletion flows)']; S window: grading height $\le$ 3; seconds: 0.1. & \CERTIFIED\newline {\footnotesize theorem, fast\newline web}\newline {\footnotesize run web 2026-09-26: 5/5} \\ \addlinespace
\code{def:rg/rho-family}\newline Remark 4.4 & $(\RG,S)\mapsto(\rho^{-1}_\IR(\RG_{\rho_\UV(a)}),\rho^{-1}_\IR(S))$ is again a flow; $I_{a,b}=I_{\rho^{-1}(b),\rho^{-1}(a)}$. & A regression check of the realisation, the statement being derived from the axioms.  On a BPS chart $\rho_\IR(X_\gamma)=X_{-\gamma}$ (\pref{ex:qt}), so the transformed flow is the flow of the same quiver with node charges $-\gamma_i$, relabelled by $a\mapsto-\rho_\UV(a)$.  Checked: the F-solver's symmetry under negating every charge; $I_{a,b}=I_{\rho^{-1}(b),\rho^{-1}(a)}$; R4 of the transformed flow, computed on the negated chart. \emph{Population:} the pentagon chart and the Kronecker chart of pure $SU(2)$, $8$ labels, $25$ pairs, through $\fq^6$. \emph{Controls:} positive: the pentagon chart; negative: a label permutation other than $\rho$ breaks $I_{a,b}=I_{\rho^{-1}(b),\rho^{-1}(a)}$. \emph{Evidence:} enforced (derived from the axioms; the check guards the realisation). \emph{Where:} \code{battery/checks_rg.py}. \emph{Note:} Derived.  Write RG' = $\rho$\_IR$^{-1}$$\circ$RG$\circ$$\rho$\_UV and S' = $\rho$\_IR$^{-1}$(S).  R2: a composite of algebra maps, and $\rho$ preserves bar invariance.  R3: apply $\rho$\_IR$^{-1}$ to both equalities of R3 at the label $\rho$\_UV(a); the new $\ell$ is $\rho$\_IR$^{-1}$$\circ$$\ell$$\circ$$\rho$\_UV.  R4: (RG'\_a S'|RG'\_b S') = I(RG\_$\rho$(b) S, RG\_$\rho$(a) S) = I\_UV($\rho$(b), $\rho$(a)) = I\_UV(a, b), using I(x,y) = I($\rho$$^{-1}$y, $\rho$$^{-1}$x) twice; that identity is the two forms of I in K4.  The one input beyond the axioms is that the identity holds for the formal elements RG\_a S, which is the point of the remark. \emph{Run (web, 2026-09-26):} 3 checks, 0 failing; charts: ['pentagon', 'Kronecker (pure SU(2))']; labels: [[0, 0], [1, 0], [0, 1], [1, 1], [2, 1], [1, 2], [-1, 1], [1, -1]]; K: 6; seconds: 39.9. & \DERIVED\newline {\footnotesize theorem, fast\newline web}\newline {\footnotesize run web 2026-09-26: 3/3} \\ \addlinespace
\code{def:wallcrossing}\newline Def.~\pref{def:wallcrossing} & Mutation-related flows satisfy $S^{1\bar1}=S^{12}S^{2\bar1}$, $S^{2\bar2}=S^{2\bar1}\rho^{-1}_\IR(S^{12})$ with the two intertwining relations; the partial factors need not be $O(\fq)$-good. & A mutation at the spec head $\gamma_k$ factors $S^{1\bar1}=S^{12}S^{2\bar1}$ with $S^{12}=E_\fq(X_{\gamma_k})$.  Checked: the mutated chart's spec is the old one with its head moved to the end and negated, i.e. $S^{2\bar2}=S^{2\bar1}\rho^{-1}_\IR(S^{12})$; and $\RG^{(1)}_aS^{12}=S^{12}\RG^{(2)}_a$: $F$ conjugated by $E_\fq(X_{\gamma_k})$ (\code{lattice_mutation.solve}, which never expands $E_\fq$) equals the mutated chart's $F$, that chart rebuilt from its node charges and spec.  The second intertwining relation and R3 of the new flow follow from these and R3 of the old flow.  The note: pure $SU(2)$ at weak coupling carries the W boson, whose factor $E^{(1/2)}_\fq(X_{\gamma_1+\gamma_2})$ has $X$-coefficient $-(1+\fq^2)/(1-\fq^2)$. \emph{Population:} the Section 4 lineup's BPS charts with a spec: the pentagon (both specs), $A_3$, $[A_1,D_4]$, the Kronecker quiver (pure $SU(2)$), $SU(2)$ with $1$, $2$, $3$ fundamentals; one mutation at the head of each spec, on windows of $8$--$41$ labels. \emph{Controls:} positive: the pentagon (both specs); negative: conjugating by another node's factor changes the image on most labels of every chart. \emph{Evidence:} emergent. \emph{Where:} \code{battery/checks_rg.py}, \code{src/bps/bps_atlas.py}, \code{src/bps/lattice_mutation.py}. \emph{Note:} 2026-09-26 (the design record lineup).  The draft's note that partial factors need not behave well as q $\to$ 0 is witnessed by the W boson's factor: order 1, not 1 + O(q), while the whole S is.  Further populations (every adjacent chart pair of an atlas, the weak dictionary's mutation edges on the local machine) are not run. \emph{Run (web, 2026-09-27):} 10 checks, 0 failing; charts: {'pentagon': 'head (1, 0); spec rotated True; 8/8 labels', 'pentagon, 3-factor spec': 'head (0, 1); spec rotated True; 8/8 labels', 'A\_3': 'head (1, 0, 0); spec rotated True; 16/16 labels', '[A\_1,D\_4]': 'head (0, 1, 0, 0); spec rotated True; 27/27 labels', 'Kronecker (pure SU(2))': 'head (1, 0); spec rotated True; 8/8 labels', 'SU(2)+1': 'head (0, -1, 1); spec rotated True; 16/16 labels', 'SU(2)+2': 'head (1, 0, 1, 0); spec rotated True; 27/27 labels', 'SU(2)+3': 'head (1, 0, 0, 0, 0); spec rotated True; 41/41 labels'}; seconds: 0.3. & \CERTIFIED\newline {\footnotesize theorem, fast\newline web}\newline {\footnotesize run web 2026-09-27: 10/10} \\ \addlinespace
\code{def:rg/central-charge}\newline Remark 4.6 & Rotating the central charge gives a locally constant family of flows; $\rho^{\pm1}$ is the rotation by $\pm\pi$. & \code{BPSAtlas.monodromy}, the full rotation of the central charge, equals $\rho^2$ on the node labels; the factor engine ordered by the phases of a generic linear central charge reproduces the chart's $S$ in each of two chambers (the node phases in one order and in the reverse).  The rotation by $\pi$ is \code{def:rg/rho-family}. \emph{Population:} monodromy: the pentagon, $A_3$ and $[A_1,D_4]$ atlases (finite type, so the rotation loop closes); chambers: the pentagon, $A_3$, the Kronecker quiver and $SU(2)+1$, two chambers each, through cone depth $4$. \emph{Controls:} positive: the pentagon (2 and 3 factor charges in its two chambers). \emph{Evidence:} emergent. \emph{Where:} \code{battery/checks_rg.py}, \code{tests/test_bps_atlas*.py}. \emph{Note:} 2026-09-26: [A\_1,D\_4] added to the monodromy and the chambers check added (the design record lineup).  The Kronecker quiver's rotation loop does not close (infinite type), so its monodromy is not computed by the chart graph; its weak-coupling chamber is checked through depth 4 (the two nodes, the W boson, two dyons). \emph{Run (web, 2026-09-27):} 3 checks, 0 failing; atlases (monodromy): ['pentagon (A\_2)', 'A\_3', '[A\_1,D\_4]']; chambers: {'pentagon': ['2 factor charges, S equal True', '3 factor charges, S equal True'], 'A\_3': ['3 factor charges, S equal True', '6 factor charges, S equal True'], 'Kronecker (pure SU(2))': ['2 factor charges, S equal True', '5 factor charges, S equal True'], 'SU(2)+1': ['5 factor charges, S equal True', '8 factor charges, S equal True']}; depth: 4; seconds: 0.1. & \CERTIFIED\newline {\footnotesize theorem, fast\newline web}\newline {\footnotesize run web 2026-09-27: 3/3} \\ \addlinespace
\code{def:rg/matter-removal}\newline Remark 4.7 & $K_\fq[G,N]\to K_\fq[G,\emptyset]\otimes Q_\fq[\Gamma]$ with $S=\prod_{w\in N}E_\fq(\mu v^w)$; one flavour charge per summand factors the flow. & \code{GMatterOverPure}, the flow over pure $G$ with the removed flavours adjoined, against the remark: its $S_\RG$, characters expanded into monomials, equals $\prod_i\prod_{w\in N_i}E_\fq(\mu_iv^w)$ multiplied out over the weights, through $\fq^6$; the RG verifiers (unital, bar-invariant, multiplicative, $\RG_aS=L+O(\fq)$, the twist) on the identity, a Wilson line and the minimal monopole, and R4 against the $(G,N)$ tier (\code{GNAbeKAlgebra}) on their $9$ pairs at $K=4$: at $SU(2)$ with one or two fundamentals or an adjoint in the fast tier, on every theory at the extensive depth; the factorisation over summands: the tower removing one hypermultiplet per rung, composed by hand through the flavour dictionary, equals removing all at once. \emph{Population:} thirteen matter theories: the twelve of the Section 3 lineup that \code{GMatterOverPure} reaches ($SU(2)$ with $1$--$4$ fundamentals, $SU(2)$ and $SU(3)$ with an adjoint, $SU(3)$ with a symmetric, alone or with a fundamental, $Spin(5)$ with a vector, $G_2$ with a $\mathbf 7$, the $SU(2)\times SU(2)$ bifundamental, $U(1)$--$U(2)$ with three fundamentals) and $SU(3)$ with one fundamental; the removal towers from $SU(2)+2$ and $SU(2)+3$ to pure $SU(2)$.  Not reachable: the non-simply-connected forms of the lineup ($SO(3)$ and the third rank-one form with an adjoint, $SO(5)$ with a vector), since the class takes no line lattice. \emph{Controls:} positive: the hand-coded two-variable series at $SU(2)+1$; the Weyl symmetry of the product; negative: the tower with each rung's removed level filed under the wrong slot disagrees with the direct removal. \emph{Evidence:} emergent. \emph{Where:} \code{battery/checks_rg.py}, \code{tests/test_g_matter_over_pure.py}, \code{tests/test_g_matter_over_matter.py}. \emph{Note:} 2026-09-26: extended to the lineup and the towers.  The tower does not compose through RGKAlgebra.then (each rung's IR carries the removed flavour as a spectator that the next rung's UV lacks), so the check composes the rungs by hand: rung i's IR presents its identical slots natively with U(r) flavour, rung i+1's UV splits them as U(r$-$1)×U(1), which for r $\le$ 2 is the Cartan expansion; hence towers up to SU(2)+3.  R4's dictionary is GNAbeKAlgebra.flow\_iso's: the same label tuples for distinct irreps, the Cartan expansion for n copies of one irrep; the mixed U(1)–U(2)+3 has no shipped dictionary and is not compared.  Measured costs: the S check is seconds everywhere; the verifiers take minutes per step at SU(2) with 3 or 4 fundamentals and at rank two.  Extensive record 2026-09-27 (15 minutes per theory, memory-capped workers): every RG check and R4 reached at SU(2)+1, SU(2)+2, N=2* SU(2) and SU(3)+1; the verifiers but not R4 at SU(2)+3 and the SU(2)xSU(2) bifundamental; the verifiers at U(1)-U(2)+3 (R4 not compared); only the unital and bar-invariance steps (and multiplicativity at Spin(5)+vector) at SU(2)+4, SU(3) with an adjoint, a symmetric, or both, Spin(5)+vector, G2+7.  No failure anywhere; unreached steps are reported as such. \emph{Run (web, 2026-09-27):} 23 checks, 0 failing; SU(2)+1, the 2026-09-22 checks: {'flow': 'GMatterOverPure(su\_2(), fundamental, nf=1)', 'K': 8, 'labels': 5, 'seconds': 9.3}; lineup theories: {'SU(2)+1': {'S monomials': 47, 'RG checks': {'unital': True, 'bar-invariant': True, 'multiplicative': True, 'RG\_a S = L + O(q)': True, 'twist': True, 'R4': True}, 'RG checks seconds': 3.5, 'seconds': 3.6}, 'SU(2)+2': {'S monomials': 207, 'RG checks': {'unital': True, 'bar-invariant': True, 'multiplicative': True, 'RG\_a S = L + O(q)': True, 'twist': True, 'R4': True}, 'RG checks seconds': 90.4, 'seconds': 90.4}, 'SU(2)+3': {'S monomials': 625, 'seconds': 0.1}, 'SU(2)+4': {'S monomials': 1549, 'seconds': 0.2}, 'SU(2)+adjoint (N=2*)': {'S monomials': 81, 'RG checks': {'unital': True, 'bar-invariant': True, 'multiplicative': True, 'RG\_a S = L + O(q)': True, 'twist': True, 'R4': True}, 'RG checks seconds': 59.3, 'seconds': 59.3}, 'SU(3)+adjoint (N=2*)': {'S monomials': 493, 'seconds': 0.4}, 'SU(3)+1': {'S monomials': 127, 'seconds': 0.1}, 'SU(3)+symmetric': {'S monomials': 367, 'seconds': 0.2}, 'SU(3)+symmetric+1': {'S monomials': 1209, 'seconds': 0.5}, 'Spin(5)+vector': {'S monomials': 333, 'seconds': 0.2}, 'G2+7': {'S monomials': 493, 'seconds': 0.4}, 'SU(2)xSU(2)+bifundamental': {'S monomials': 207, 'seconds': 0.1}, 'U(1)-U(2)+3': {'S monomials': 1549, 'seconds': 0.3}}; not reachable (GMatterOverPure takes no line lattice): ['SO(3)+adjoint (N=2*)', 'rank 1 + adjoint (N=2*), the third global form', 'SO(5)+vector']; towers: {'SU(2)+2': '24/24 agree, 5.9 s', 'SU(2)+3': '36/36 agree, 99.0 s'}; depth: fast; seconds: 276.0. \emph{Run (web (extensive), 2026-09-27):} 33 checks, 0 failing; SU(2)+1, the 2026-09-22 checks: {'flow': 'GMatterOverPure(su\_2(), fundamental, nf=1)', 'K': 8, 'labels': 5, 'seconds': 9.5}; lineup theories: {'SU(2)+1': {'S monomials': 47, 'RG checks': {'unital': True, 'bar-invariant': True, 'multiplicative': True, 'RG\_a S = L + O(q)': True, 'twist': True, 'R4': True}, 'RG checks seconds': 3.6, 'seconds': 3.6}, 'SU(2)+2': {'S monomials': 207, 'RG checks': {'unital': True, 'bar-invariant': True, 'multiplicative': True, 'RG\_a S = L + O(q)': True, 'twist': True, 'R4': True}, 'RG checks seconds': 89.0, 'seconds': 89.1}, 'SU(2)+3': {'S monomials': 625, 'RG checks': {'unital': True, 'bar-invariant': True, 'multiplicative': True, 'RG\_a S = L + O(q)': True, 'twist': True, 'R4': 'not reached in 900 s'}, 'RG checks seconds': 900.0, 'seconds': 900.1}, 'SU(2)+4': {'S monomials': 1549, 'RG checks': {'unital': True, 'bar-invariant': True, 'multiplicative': True, 'RG\_a S = L + O(q)': 'not reached in 900 s', 'twist': 'not reached in 900 s', 'R4': 'not reached in 900 s'}, 'RG checks seconds': 900.0, 'seconds': 900.3}, 'SU(2)+adjoint (N=2*)': {'S monomials': 81, 'RG checks': {'unital': True, 'bar-invariant': True, 'multiplicative': True, 'RG\_a S = L + O(q)': True, 'twist': True, 'R4': True}, 'RG checks seconds': 58.9, 'seconds': 58.9}, 'SU(3)+adjoint (N=2*)': {'S monomials': 493, 'RG checks': {'unital': True, 'bar-invariant': True, 'multiplicative': 'not reached in 900 s', 'RG\_a S = L + O(q)': 'not reached in 900 s', 'twist': 'not reached in 900 s', 'R4': 'not reached in 900 s'}, 'RG checks seconds': 900.1, 'seconds': 900.5}, 'SU(3)+1': {'S monomials': 127, 'RG checks': {'unital': True, 'bar-invariant': True, 'multiplicative': True, 'RG\_a S = L + O(q)': True, 'twist': True, 'R4': True}, 'RG checks seconds': 883.6, 'seconds': 883.6}, 'SU(3)+symmetric': {'S monomials': 367, 'RG checks': {'unital': True, 'bar-invariant': True, 'multiplicative': 'not reached in 900 s', 'RG\_a S = L + O(q)': 'not reached in 900 s', 'twist': 'not reached in 900 s', 'R4': 'not reached in 900 s'}, 'RG checks seconds': 900.1, 'seconds': 900.3}, 'SU(3)+symmetric+1': {'S monomials': 1209, 'RG checks': {'unital': True, 'bar-invariant': True, 'multiplicative': 'not reached in 900 s', 'RG\_a S = L + O(q)': 'not reached in 900 s', 'twist': 'not reached in 900 s', 'R4': 'not reached in 900 s'}, 'RG checks seconds': 900.1, 'seconds': 900.5}, 'Spin(5)+vector': {'S monomials': 333, 'RG checks': {'unital': True, 'bar-invariant': True, 'multiplicative': True, 'RG\_a S = L + O(q)': 'not reached in 900 s', 'twist': 'not reached in 900 s', 'R4': 'not reached in 900 s'}, 'RG checks seconds': 900.0, 'seconds': 900.2}, 'G2+7': {'S monomials': 493, 'RG checks': {'unital': True, 'bar-invariant': True, 'multiplicative': 'not reached in 900 s', 'RG\_a S = L + O(q)': 'not reached in 900 s', 'twist': 'not reached in 900 s', 'R4': 'not reached in 900 s'}, 'RG checks seconds': 900.0, 'seconds': 900.4}, 'SU(2)xSU(2)+bifundamental': {'S monomials': 207, 'RG checks': {'unital': True, 'bar-invariant': True, 'multiplicative': True, 'RG\_a S = L + O(q)': True, 'twist': True, 'R4': 'not reached in 900 s'}, 'RG checks seconds': 900.0, 'seconds': 900.1}, 'U(1)-U(2)+3': {'S monomials': 1549, 'RG checks': {'unital': True, 'bar-invariant': True, 'multiplicative': True, 'RG\_a S = L + O(q)': True, 'twist': True, 'R4': 'not compared (mixed flavour: no shipped dictionary)'}, 'RG checks seconds': 164.0, 'seconds': 164.3}}; not reachable (GMatterOverPure takes no line lattice): ['SO(3)+adjoint (N=2*)', 'rank 1 + adjoint (N=2*), the third global form', 'SO(5)+vector']; towers: {'SU(2)+2': '24/24 agree, 6.1 s', 'SU(2)+3': '36/36 agree, 98.1 s'}; depth: extensive; seconds: 8522.3. & \CERTIFIED\newline {\footnotesize theorem, extensive\newline web}\newline {\footnotesize run web 2026-09-27: 23/23}\newline {\footnotesize run web (extensive) 2026-09-27: 33/33} \\ \addlinespace
\code{sec:rg/E-group}\newline Def.~4.8, Remark 4.9 & Every $g\in\cE$ satisfies $g^{-1}(\fq)=g(\fq^{-1})$; $\cE$ is smaller than the group of such elements ($E_{\fq^2}(x)\notin\cE$). & the identity to depth $D$; PBW-factorise $E_{\fq^2}(x)$ and record non-integral or unbounded exponents. \emph{Population:} random products of $E^{(s)}_\fq(X_\gamma)^{\pm1}$ to cone depth $D$; the single element $E_{\fq^2}(x)$. \emph{Evidence:} emergent. \emph{Where:} \code{battery/checks_rg.py}, \code{notes/pronilpotent_group_conjecture.md}. \emph{Note:} the S$^{-1}$($\fq$) = S($\fq$$^{-1}$) identity  2026-09-22: g\^{}-1(q) = g(q\^{}-1) verified EXACTLY (Laurent polynomials over (q\^{}2;q\^{}2)\_n\^{}(2s+1), x-degree $\le$ 6) for the generators E\^{}(s), s = 0, 1/2, 1, 3/2, with two negative controls (asymmetric shifts fail; mixed exponents are bar-invariant instead); the statement for every group element follows because the bar involution is antimultiplicative on the quantum torus.  E\_{q\^{}2} not in the group: the degree-one argument at an indecomposable charge, both computational inputs verified. \emph{Run (web, 2026-09-22):} 5 checks, 0 failing; K: 24; n: $\le$ 6; s: 0, 1/2. & \CERTIFIED\newline {\footnotesize theorem, fast\newline web}\newline {\footnotesize run web 2026-09-22: 5/5} \\ \addlinespace
\code{def:sw}\newline Def.~\pref{def:sw}, \peqref{eq:quiver} & A BPS-quiver flow has $K_\fq[\IR]=Q_\fq(\Gamma)$, $S\in\cE$ and $S=1-\fq\sum_iX_{\gamma_i}+O(\fq^2)$. & Not a claim: a definition has no truth value, and checking an axiom on a realisation tests the realisation, not the definition.  This row records where the axiom is \emph{enforced} by construction (a pass is a regression guard, left to the suite) and where it is \emph{emergent} (a pass is evidence), and names the rows that test it there.  The constraint defines a spectrum generator.  Its acceptance census over the dictionaries is a test of the \emph{entries}: for a green spec it is evidence for \code{conj:cluster} (moved there); for a crystalline-only entry the factor engine imposes it, so a pass is a regression guard.  Instruments: \code{src/bps/spec_acceptance.py}, \code{dictionaries/enumerated/manifest.json} \emph{Evidence:} emergent (to depth). \emph{Where:} \code{battery/checks_rg.py (run recorded under conj:cluster)}, \code{src/bps/spec_acceptance.py}, \code{dictionaries/enumerated/manifest.json}. \emph{Note:} The acceptance census of the shipped tier's (quiver, spec) pairs is recorded under conj:cluster (checks\_rg.check\_sw\_sample). & ---\newline {\footnotesize definition, not run} \\ \addlinespace
\code{conj:S-unique}\newline Conj.~\pref{conj:S-unique} & There is a unique $S\in\cE$ satisfying \peqref{eq:quiver}. & On the independent arithmetic of \code{experiments/pbw_e_group.py}, which shares no code with the repo engine: $S$ built by the forced recursion in five total orders (degree-lex; anti-degree, every top-degree factor before any node factor, which no central charge produces; lex-reversed; spin-first; shuffled) must be one element.  Order independence is the computable half of uniqueness. \emph{Population:} the Section 4 lineup (the pentagon, $A_3$, $[A_1,D_4]$, the Kronecker quiver, $SU(2)$ with $1$--$3$ fundamentals, the Markov quiver) and the first $120$ entries with $n\le4$ of the shipped enumerated tier (web); the entries of the local builds and the weak dictionary (local). \emph{Controls:} positive: the pentagon's $S$; agreement with the repo engine on four quivers; negative: the vacuity guard, the number of factors moves with the order. \emph{Evidence:} emergent. \emph{Where:} \code{battery/checks_rg.py}, \code{experiments/sw_flow_conjectures.py}, \code{experiments/pbw_e_group.py}, \code{src/bps/bps_factor_spectrum.py}. \emph{Inputs:} \emph{web} --- the quivers of \code{dictionaries/enumerated/} (tracked; E, accept\_depth, verify\_depth read from \code{manifest.json}) and of a package of the source repository (tracked; orbit\_size, max\_weight, coverage read from the ten manifests); the computation needs the quiver alone; \emph{local} --- the quivers of the enumerated build beyond the shipped weight (E $\ge$ 13: the weight-13 tier, which stays on the local machine by the author's ruling, at a package of the source repository) and the flavoured builds beyond the shipped weights (gitignored while in progress: \code{.flavoured_*stage/}); \code{dictionaries_weak/} (gitignored, local machine only: 579,293 entries with a spec each, ranks up to 15, several known specs per quiver; a valid source by the author's ruling of 2026-09-21, being integrated into the enumerated build as harvested seeds, so its weight $\le$ E quivers reach the shipped tiers as they ship while its tail beyond E stays a local population). \emph{Note:} 2026-09-27: own adapter, on the independent arithmetic (until then a check inside conj:cluster's adapter, on the repo engine).  Prior sweep: 1108/1108 quivers of the seed-closure dictionary n\_001..n\_008 (2026-09-04, pronilpotent\_group\_conjecture.md 6c). \emph{Run (web, 2026-09-27):} 4 checks, 0 failing; lineup: ['pentagon', 'A\_3', '[A\_1,D\_4]', 'Kronecker (pure SU(2))', 'SU(2)+1', 'SU(2)+2', 'SU(2)+3', 'Markov']; sample: first 120 entries with n $\le$ 4 of dictionaries/enumerated (shard format 3, E = 11: 111987 strongly connected quivers, 104086 with a green spec, 7901 crystalline-only, 0 uncertified); seconds: 58.0. & \MEASURED\newline {\footnotesize conjecture, extensive\newline web, local}\newline {\footnotesize run web 2026-09-27: 4/4} \\ \addlinespace
\code{conj:pbw}\newline Conj.~\pref{conj:pbw}, Remarks 4.11--4.12 & Every $g\in\cE$ factorises as $\prod E^{(s)}_\fq(X_\gamma)^{\Omega(\gamma,s)}$ in any total order on $\Gamma_+\times\tfrac12\mathbb N$ with integral exponents; the factorisation is unique along the filtration. & On the independent arithmetic: the read-off \code{pbw_factorize} factorises a given element in a given total order and raises when a $(1-\fq^{2n})$ denominator survives or a residual is not palindromic; the exponents must be integral and the ordered product must rebuild the element exactly.  The repo engine's integral $\Omega$ is imposed by its construction (the source repository's audit record) and is not used as evidence. \emph{Population:} random elements of $\mathcal E$ (three words per quiver) and $S$ itself, on the Section 4 lineup, in five total orders (web); every $S$ of the local builds (local). \emph{Controls:} positive: generators refactorise to themselves, round trips are exact, the engine's $S$ is reproduced; negative: $E_{\fq^2}(X)$, bar-unitary but outside $\mathcal E$, does not factorise; a $\fq^0$ bump fails ($96$ controls, \code{experiments/pbw_controls.py}). \emph{Evidence:} emergent. \emph{Where:} \code{battery/checks_rg.py}, \code{experiments/pbw_e_group.py}, \code{experiments/pbw_controls.py}. \emph{Inputs:} \emph{web} --- the (quiver, spec) pairs of \code{dictionaries/enumerated/} (tracked; E, accept\_depth, verify\_depth read from \code{manifest.json}); \emph{local} --- the (quiver, spec) pairs of the enumerated build beyond the shipped weight (E $\ge$ 13: the weight-13 tier, which stays on the local machine by the author's ruling, at a package of the source repository) and the flavoured builds beyond the shipped weights (gitignored while in progress: \code{.flavoured_*stage/}); \code{dictionaries_weak/} (gitignored, local machine only: 579,293 entries with a spec each, ranks up to 15, several known specs per quiver; a valid source by the author's ruling of 2026-09-21, being integrated into the enumerated build as harvested seeds, so its weight $\le$ E quivers reach the shipped tiers as they ship while its tail beyond E stays a local population). \emph{Note:} 2026-09-27: rebuilt on the independent arithmetic; the 2026-09-22 check read the engine's Omega, whose integrality A44 shows is imposed.  Prior campaign on the same arithmetic: ranks 1-8, brackets -8..8, orders no central charge produces, spin content to 45.5, no counterexample (pronilpotent\_group\_conjecture.md 6b). \emph{Run (web, 2026-09-27):} 4 checks, 0 failing; lineup: ['pentagon', 'A\_3', '[A\_1,D\_4]', 'Kronecker (pure SU(2))', 'SU(2)+1', 'SU(2)+2', 'SU(2)+3', 'Markov']; words per quiver: 3; orders: 5; seconds: 52.9. & \MEASURED\newline {\footnotesize conjecture, extensive\newline web, local}\newline {\footnotesize run web 2026-09-27: 4/4} \\ \addlinespace
\code{conj:cluster}\newline Conj.~\pref{conj:cluster} & $S_{\mathrm{cluster}}=S$ whenever a maximal green sequence exists; different green sequences give the same $S$. & the repo's acceptance check (\code{spec_acceptance.accept_spec}: exact green replay, the finite-depth constraint \peqref{eq:quiver}, agreement with the crystalline $S$) and, independently, the arithmetic of \code{experiments/pbw_e_group.py}: $S_{\mathrm{cluster}}$ satisfies the constraint and equals $S$; different green sequences of one quiver give one $S$.  A green spec passing is evidence; a crystalline-only entry passing is a regression guard on the factor engine. \emph{Population:} the first $120$ entries with $n\le4$ of the shipped enumerated tier, its size and certification read from \code{dictionaries/enumerated/manifest.json}; the Section 4 lineup's specs, the pentagon's two chambers included (web); the rich build's spec families, several green sequences per quiver (local). \emph{Controls:} positive: the pentagon's two chambers; negative: five wrong pentagon specs (reversed, a node repeated, a node missing, a spurious charge, the 3-state chamber reversed) and a doubled node factor are caught. \emph{Evidence:} emergent (to depth). \emph{Where:} \code{battery/checks_rg.py}, \code{src/bps/spec_acceptance.py}, \code{experiments/sw_flow_conjectures.py}, \code{dictionaries/enumerated/manifest.json}. \emph{Inputs:} \emph{web} --- \code{dictionaries/enumerated/} (tracked; n\_green, family\_mismatches, axiom\_mismatches read from \code{manifest.json}); \emph{local} --- a package of the source repository (the family pairs need every spec found per quiver, which only the rich build keeps); \code{dictionaries_weak/} (gitignored, local machine only: 579,293 entries with a spec each, ranks up to 15, several known specs per quiver; a valid source by the author's ruling of 2026-09-21, being integrated into the enumerated build as harvested seeds, so its weight $\le$ E quivers reach the shipped tiers as they ship while its tail beyond E stays a local population); the enumerated build beyond the shipped weight (E $\ge$ 13: the weight-13 tier, which stays on the local machine by the author's ruling, at a package of the source repository) and the flavoured builds beyond the shipped weights (gitignored while in progress: \code{.flavoured_*stage/}). \emph{Note:} 2026-09-27: the population is read from the manifest (shard format 3 holds the strongly connected quivers only), and the lineup's specs are added.  The manifest's own census at the time of writing: 104,086 green, 7,901 crystalline-only, 0 uncertified, no axiom mismatch, at E = 11.  Since 2026-09-27 the shipped tier is E = 12: 674,440 green, 54,481 crystalline-only, 0 uncertified (manifest).  Prior sweep: 2381/2381 specs of n\_001..n\_008 satisfy the constraint and equal S (2026-09-04). \emph{Run (web, 2026-09-27):} 4 checks, 0 failing; sample: first 120 entries with n $\le$ 4 of dictionaries/enumerated (shard format 3, E = 11: 111987 strongly connected quivers, 104086 with a green spec, 7901 crystalline-only, 0 uncertified), cell order; lineup: ['pentagon', 'A\_3', '[A\_1,D\_4]', 'Kronecker (pure SU(2))', 'SU(2)+1', 'SU(2)+2', 'SU(2)+3']; depth: 4; seconds: 41.6. & \CERTIFIED\newline {\footnotesize conjecture, extensive\newline web, local}\newline {\footnotesize run web 2026-09-27: 4/4} \\ \addlinespace
\code{conj:dt}\newline Conj.~\pref{conj:dt} & $S_{\mathrm{DT}}=S$ for quivers of character varieties. & No test: reference DT invariants from the literature are not encoded . \emph{Evidence:} emergent. \emph{Note:} 2026-09-26: the conjecture stays open, with no test planned. & \OPEN\newline {\footnotesize conjecture, none\newline no test} \\ \addlinespace
\code{conj:coha}\newline Conj.~\pref{conj:coha} & $S_{\mathrm{CoHA}}=S$ for a suitable potential. & no CoHA characters encoded; candidate first cases are quivers without potential where the character is known. & \OPEN\newline {\footnotesize conjecture, none\newline no test} \\ \addlinespace
\code{sec:rg/F-unique}\newline paragraph before Conj.~\pref{conj:upper} & Given $S$, $\RG_\gamma S=X_\gamma+O(\fq)$ has a unique bar-invariant solution supported on $\gamma+\Gamma_+$. & Derived: the statement follows from a recursion: $\RG_\gamma$ has $1$ at $\gamma$, and at cone degree $d$ its coefficient at $\gamma+\delta$ is minus the palindromic completion of the non-positive part of $(\RG^{<d}_\gamma S)_{\gamma+\delta}$ --- forced degree by degree, hence existence and uniqueness of the formal solution (two solutions differ first at some degree, where the difference is bar-invariant and $O(\fq)$, hence zero).  Checked: the repo's F-solver equals this recursion, run on the independent arithmetic of \code{experiments/pbw_e_group.py}, through its depth; \code{verify_F_S_leading} as the regression guard. \emph{Population:} the Section~4 lineup's BPS charts with a spec (the pentagon, $A_3$, $[A_1,D_4]$, the Kronecker quiver, $SU(2)$ with $N_f=1,2,3$) and the $11$ dictionary-sample charts of the \pref{conj:upper} row: node charges, their pairwise sums and minus the first node ($137$ labels), through cone degree $6$ (rank two), $5$ (rank three), $4$ (ranks four and five). \emph{Controls:} positive: the pentagon; negative: the pentagon's $S$ in the reversed order changes the recursion's $\RG_\gamma$ on $4$ of $4$ labels. \emph{Evidence:} enforced (derived from the recursion; the check guards the F-solver, which it shares no code with). \emph{Where:} \code{battery/checks_rg.py (check_f_unique)}, \code{experiments/pbw_e_group.py}. \emph{Inputs:} \emph{web} --- \code{dictionaries/enumerated/} (tracked; E, accept\_depth, verify\_depth read from \code{manifest.json}), one chart per entry, a label window each; \emph{local} --- the enumerated build beyond the shipped weight (E $\ge$ 13: the weight-13 tier, which stays on the local machine by the author's ruling, at a package of the source repository) and the flavoured builds beyond the shipped weights (gitignored while in progress: \code{.flavoured_*stage/}); \code{dictionaries_weak/} (gitignored, local machine only: 579,293 entries with a spec each, ranks up to 15, several known specs per quiver; a valid source by the author's ruling of 2026-09-21, being integrated into the enumerated build as harvested seeds, so its weight $\le$ E quivers reach the shipped tiers as they ship while its tail beyond E stays a local population). \emph{Note:} Derived: the recursion is forced degree by degree, so the formal solution exists and is unique; the check guards the repo's F-solver.  2026-09-27 web run: the repo's F equals the independent recursion on all 137 labels through its depth, 791 terms compared.  That the solution is finite is conj:upper's. \emph{Run (web, 2026-09-27):} 3 checks, 0 failing; charts: {'pentagon': '4 labels, depth 6', 'A\_3': '7 labels, depth 5', '[A\_1,D\_4]': '11 labels, depth 4', 'Kronecker (pure SU(2))': '4 labels, depth 6', 'SU(2)+1': '7 labels, depth 5', 'SU(2)+2': '11 labels, depth 4', 'SU(2)+3': '16 labels, depth 4', 'dictionary sample (11 charts)': '77 labels, depth 5-6'}; terms compared: 791; seconds: 34.3. & \DERIVED\newline {\footnotesize theorem, fast\newline web, local}\newline {\footnotesize run web 2026-09-27: 3/3} \\ \addlinespace
\code{conj:upper}\newline Conj.~\pref{conj:upper} & The $\RG_\gamma$ are Laurent polynomials forming a $\bZ[\fq^{\pm1}]$-subalgebra; a label permutation $\rho_\UV$ makes it a $K_\fq$-algebra; as an algebra it is the upper cluster algebra of $(\Gamma,Q)$. & (i) Finiteness: the F-solver searches a window, so its $\RG_\gamma$ being finite is not evidence (the source repository's audit record); the independent recursion of the uniqueness row (the paragraph before \pref{conj:upper}) must stop where the solver's $\RG_\gamma$ does, i.e.\ have no term in the degrees above it through its depth.  (ii) The first half of the conjecture's operational statement: along every mutation sequence of bounded length (no node twice in a row --- that is a monomial change of chart), conjugate $\RG_\gamma$ by $E_\fq(X_{\gamma_k})$ at the current chart's node (\code{lattice_mutation.solve}, exact, raising on a non-Laurent result) and mutate the chart; every conjugate must be a Laurent polynomial.  (iii) The $K_\fq$-algebra axioms on the labels $1$ and three nodes: \pref{ax:bar}, \pref{ax:rho}, $\RG$-multiplicativity (closure over $\bZ[\fq^{\pm1}]$), \pref{ax:rhotr} (a regression guard), and the orthonormality and $\rho^2$-twisted cyclicity of \pref{ax:trace}, within a per-chart time budget whose unreached verifiers are named.  (iv) The second half of that statement, on a box of charges: every $\RG_\gamma$ supported in the box stays Laurent after one mutation at each node of the root chart, and the conditions for a combination of the other box points' monomials to stay Laurent there (the $\fq$-binomial-packet criterion of \code{lattice_mutation}, at the exactly found roots) have full rank at $\fq=2$, which bounds the rank over $\mathbb{Q}(\fq)$ from below.  Since $\RG_\gamma=X_\gamma+$ higher terms, the box elements Laurent in the root chart and its immediate neighbours are then exactly the span of the $\RG_\gamma$ inside the box: the upper cluster algebra meets the box inside that span, and equals it there given (ii). \emph{Population:} the Section~4 lineup's charts --- the pentagon, $A_3$, $[A_1,D_4]$, the Kronecker quiver, $SU(2)$ with $N_f=1,2,3$, and the Markov quiver, built spec-free --- on node charges, their pairwise sums and minus the first node (the Markov quiver: its nodes); every mutation sequence of length $\le5$ (rank two), $4$ (rank three), $3$ (rank four), $2$ (rank five), one more at the extensive depth; boxes $[-1,1]^n$ on every lineup chart, and $[-2,2]^n$ for $n\le3$ at the extensive depth (the Markov quiver's chart built to depth $10$ there); the $11$ dictionary-sample charts (web); the weak dictionary's mutation edges (local). \emph{Controls:} positive: $\RG_\gamma S=X_\gamma+O(\fq)$ (the uniqueness row); negative: a bare node monomial $X_{\gamma_j}$ fails after one mutation on every chart, and every $\RG_\gamma$ with its top-degree terms dropped fails on some sequence.  For the box test: $X_0=1$, Laurent in every chart, added to the other points drops the rank; $\RG_\gamma$ inside the box with the top-degree terms dropped fail one mutation. \emph{Evidence:} emergent. \emph{Where:} \code{experiments/pbw_e_group.py}, \code{experiments/upper_cluster_box_certificate.py}, \code{src/bps/lattice_mutation.py}, \code{battery/checks_rg.py (check_upper_sample)}, \code{tests/test_bps_atlas*.py}, \code{tests/test_bps_kalgebra.py}. \emph{Inputs:} \emph{web} --- \code{dictionaries/enumerated/} (tracked; one BPS chart per sampled entry with a stored spec, node charges and pairwise sums as the label window); the atlases of the zoo (tracked code); \emph{local} --- \code{dictionaries/n_*.json} and \code{graph_n*.json}, the seed-closure tree with its mutation edges (untracked; built per machine in $\sim$75 s by PYTHONPATH=. python \code{dictionaries/build.py}) for Laurentness ACROSS mutated charts; \code{dictionaries_weak/mutation_edges_shard_*.json} (gitignored, local machine only: 1,115,973 mutation edges) and \code{dictionaries_final_data/} (26 GB). \emph{Note:} 2026-09-27 web runs (fast and extensive).  Finiteness: wherever the solver's F ends below the recursion's depth the independent recursion ends there too, on every lineup chart (the Markov quiver's nodes included).  D7, first half: every RG\_gamma on the label window stays Laurent along every mutation sequence -- at the extensive depth 48/48 conjugations on the pentagon (length $\le$ 6), 651/651 on A\_3, 1760/1760 on [A\_1,D\_4], 48/48 on the Kronecker quiver, 651/651, 1760/1760 and 1680/1680 on SU(2)+1, +2, +3, 279/279 on the Markov quiver; a set of node charges is not a chart (the pentagon's 1,2,1,2 returns its node charges through the monodromy), so the counts of distinct node-charge sets only bound the charts from below.  D7, second half: the box test certifies 13 (chart, box) cases -- [-1,1]\^{}n on all eight charts, [-2,2]\^{}n on the pentagon, A\_3, the Kronecker quiver, SU(2)+1 and the Markov quiver (its chart built to depth 10, the F's inside the box unchanged at depth 12).  The K\_q axioms with a 600 s budget per chart: all hold on all seven charts with a spec, except rho\^{}2-twisted cyclicity on SU(2)+3, not reached.  The Markov quiver has no finite rho on the BPS route, so its axioms are not run. \emph{Run (web, 2026-09-27):} 12 checks, 0 failing; lineup: ['pentagon', 'A\_3', '[A\_1,D\_4]', 'Kronecker (pure SU(2))', 'SU(2)+1', 'SU(2)+2', 'SU(2)+3', 'Markov']; labels: node charges, pairwise sums and minus the first node (Markov: the nodes); boxes: ['pentagon [-1,1]\^{}n', 'A\_3 [-1,1]\^{}n', '[A\_1,D\_4] [-1,1]\^{}n', 'Kronecker (pure SU(2)) [-1,1]\^{}n', 'SU(2)+1 [-1,1]\^{}n', 'SU(2)+2 [-1,1]\^{}n', 'SU(2)+3 [-1,1]\^{}n', 'Markov [-1,1]\^{}n']; mutation sequences: {'pentagon': 'length $\le$ 5, 8 distinct sets of node charges reached', 'A\_3': 'length $\le$ 4, 30 distinct sets of node charges reached', '[A\_1,D\_4]': 'length $\le$ 3, 36 distinct sets of node charges reached', 'Kronecker (pure SU(2))': 'length $\le$ 5, 9 distinct sets of node charges reached', 'SU(2)+1': 'length $\le$ 4, 36 distinct sets of node charges reached', 'SU(2)+2': 'length $\le$ 3, 39 distinct sets of node charges reached', 'SU(2)+3': 'length $\le$ 2, 20 distinct sets of node charges reached', 'Markov': 'length $\le$ 4, 43 distinct sets of node charges reached'}; K\_q axioms: {'pentagon': {'ax:bar': True, 'ax:rho': True, 'RG multiplicative': True, 'ax:rhotr': True, 'ax:trace orthonormality': True, 'ax:trace rho\^{}2-twisted cyclicity': True}, 'A\_3': {'ax:bar': True, 'ax:rho': True, 'RG multiplicative': True, 'ax:rhotr': True, 'ax:trace orthonormality': True, 'ax:trace rho\^{}2-twisted cyclicity': True}, '[A\_1,D\_4]': {'ax:bar': True, 'ax:rho': True, 'RG multiplicative': True, 'ax:rhotr': True, 'ax:trace orthonormality': 'not reached in 30 s', 'ax:trace rho\^{}2-twisted cyclicity': 'not reached in 30 s'}, 'Kronecker (pure SU(2))': {'ax:bar': True, 'ax:rho': True, 'RG multiplicative': True, 'ax:rhotr': True, 'ax:trace orthonormality': True, 'ax:trace rho\^{}2-twisted cyclicity': True}, 'SU(2)+1': {'ax:bar': True, 'ax:rho': True, 'RG multiplicative': True, 'ax:rhotr': True, 'ax:trace orthonormality': True, 'ax:trace rho\^{}2-twisted cyclicity': 'not reached in 30 s'}, 'SU(2)+2': {'ax:bar': True, 'ax:rho': True, 'RG multiplicative': True, 'ax:rhotr': True, 'ax:trace orthonormality': True, 'ax:trace rho\^{}2-twisted cyclicity': 'not reached in 30 s'}, 'SU(2)+3': {'ax:bar': True, 'ax:rho': True, 'RG multiplicative': True, 'ax:rhotr': True, 'ax:trace orthonormality': True, 'ax:trace rho\^{}2-twisted cyclicity': 'not reached in 30 s'}}; dictionary sample: 11 charts; skipped (no stored spec, so no finite rho on the BPS route): [[[0, 2, -2], [-2, 0, 2], [2, -2, 0]]]; seconds: 178.0. \emph{Run (web (extensive), 2026-09-27):} 12 checks, 0 failing; lineup: ['pentagon', 'A\_3', '[A\_1,D\_4]', 'Kronecker (pure SU(2))', 'SU(2)+1', 'SU(2)+2', 'SU(2)+3', 'Markov']; labels: node charges, pairwise sums and minus the first node (Markov: the nodes); boxes: ['pentagon [-1,1]\^{}n', 'pentagon [-2,2]\^{}n', 'A\_3 [-1,1]\^{}n', 'A\_3 [-2,2]\^{}n', '[A\_1,D\_4] [-1,1]\^{}n', 'Kronecker (pure SU(2)) [-1,1]\^{}n', 'Kronecker (pure SU(2)) [-2,2]\^{}n', 'SU(2)+1 [-1,1]\^{}n', 'SU(2)+1 [-2,2]\^{}n', 'SU(2)+2 [-1,1]\^{}n', 'SU(2)+3 [-1,1]\^{}n', 'Markov [-1,1]\^{}n', 'Markov [-2,2]\^{}n']; mutation sequences: {'pentagon': 'length $\le$ 6, 8 distinct sets of node charges reached', 'A\_3': 'length $\le$ 5, 54 distinct sets of node charges reached', '[A\_1,D\_4]': 'length $\le$ 4, 89 distinct sets of node charges reached', 'Kronecker (pure SU(2))': 'length $\le$ 6, 10 distinct sets of node charges reached', 'SU(2)+1': 'length $\le$ 5, 64 distinct sets of node charges reached', 'SU(2)+2': 'length $\le$ 4, 105 distinct sets of node charges reached', 'SU(2)+3': 'length $\le$ 3, 64 distinct sets of node charges reached', 'Markov': 'length $\le$ 5, 79 distinct sets of node charges reached'}; K\_q axioms: {'pentagon': {'ax:bar': True, 'ax:rho': True, 'RG multiplicative': True, 'ax:rhotr': True, 'ax:trace orthonormality': True, 'ax:trace rho\^{}2-twisted cyclicity': True}, 'A\_3': {'ax:bar': True, 'ax:rho': True, 'RG multiplicative': True, 'ax:rhotr': True, 'ax:trace orthonormality': True, 'ax:trace rho\^{}2-twisted cyclicity': True}, '[A\_1,D\_4]': {'ax:bar': True, 'ax:rho': True, 'RG multiplicative': True, 'ax:rhotr': True, 'ax:trace orthonormality': True, 'ax:trace rho\^{}2-twisted cyclicity': True}, 'Kronecker (pure SU(2))': {'ax:bar': True, 'ax:rho': True, 'RG multiplicative': True, 'ax:rhotr': True, 'ax:trace orthonormality': True, 'ax:trace rho\^{}2-twisted cyclicity': True}, 'SU(2)+1': {'ax:bar': True, 'ax:rho': True, 'RG multiplicative': True, 'ax:rhotr': True, 'ax:trace orthonormality': True, 'ax:trace rho\^{}2-twisted cyclicity': True}, 'SU(2)+2': {'ax:bar': True, 'ax:rho': True, 'RG multiplicative': True, 'ax:rhotr': True, 'ax:trace orthonormality': True, 'ax:trace rho\^{}2-twisted cyclicity': True}, 'SU(2)+3': {'ax:bar': True, 'ax:rho': True, 'RG multiplicative': True, 'ax:rhotr': True, 'ax:trace orthonormality': True, 'ax:trace rho\^{}2-twisted cyclicity': 'not reached in 600 s'}}; dictionary sample: 11 charts; skipped (no stored spec, so no finite rho on the BPS route): [[[0, 2, -2], [-2, 0, 2], [2, -2, 0]]]; seconds: 1842.4. & \MEASURED\newline {\footnotesize conjecture, extensive\newline web, local}\newline {\footnotesize run web 2026-09-27: 12/12}\newline {\footnotesize run web (extensive) 2026-09-27: 12/12} \\ \addlinespace
\code{conj:mut}\newline Conj.~\pref{conj:mut} & For a BPS quiver $Q$ and its mutation $Q'=\mu_j(Q)$ at a node $j$, with the node charges \peqref{eq:mut-charges}: $S=E_\fq(X_{\gamma_j})S''$ and $S'=S''E_\fq(X_{-\gamma_j})$ for some $S''\in\cE\cap\cE'$. & Eliminating $S''$ gives $S'=E_\fq(X_{\gamma_j})^{-1}S\,E_\fq(X_{-\gamma_j})$.  \code{mutation_s_relation.compare_edge} compares it exactly on every charge the two cones determine, and records a support violation wherever $E_\fq(X_{\gamma_j})^{-1}S$ leaves the cone of $Q'$, the half of the statement that says $S''\in\cE'$.  Each $S$ is built from its own quiver's constraint, in its own node frame.  On the lineup, an independent construction of $S$ (the forced recursion of \file{experiments/pbw_e_group.py}) gives the same verdict. \emph{Population:} the Section~4 lineup at every node --- the pentagon, $A_3$, $[A_1,D_4]$, the Kronecker quiver, $SU(2)$ with $N_f=1,2,3$, and the Markov quiver, which has no maximal green sequence --- to cone depth $6$, $5$, $4$, $4$ at ranks $2$, $3$, $4$, $5$; every node of the first 120 entries with $n\le4$ of the shipped enumerated tier. \emph{Controls:} positive: the pentagon, where the relation is the pentagon identity; negative: the two mixed variants, each pairing one side with the wrong lattice identification, fail on every lineup edge. \emph{Evidence:} emergent: each S is fixed by its own quiver's constraint; that the two are related by the mutation is not imposed. \emph{Where:} \code{src/bps/mutation_s_relation.py}, \code{tests/test_mutation_s_relation.py}, \code{battery/checks_rg.py (check_conj_mut)}, \code{experiments/pbw_e_group.py}. \emph{Inputs:} \emph{web} --- \code{dictionaries/enumerated/} (tracked; its first entries with n $\le$ 4). \emph{Note:} Added with the draft of 2026-09-29.  What a comparison tests: a compared charge tests the relation only on the subquiver where it is nonzero, and it agrees automatically when the mutated node is a source or a sink there -- so every comparison on the pentagon and the Kronecker quivers, and at a source or sink of an acyclic quiver, is the reflection identity, and a check reaches the quiver itself only through charges nonzero on every node.  Informative, reaching every node, and agreeing, in a sweep in the source repository (2026-09-29): A3 at its middle node, the 3-cycle, Markov, and six crystalline tier quivers with 5-7 nodes, each at the depth it needs; the mixed variants fail everywhere.  A run on Perimeter Institute's Symmetry cluster checks every edge of the crystalline tier entries with up to 7 nodes on the whole quiver (depth up to 10), and those with 8-12 nodes on smaller quivers only.  The records stay on the local machine. \emph{Run (web, 2026-09-29):} 4 checks, 0 failing; lineup: {'pentagon': {'depth': 6, 'edges': '2/2', 'charges compared beyond the leading data': [25, 19], 'support violations': 0}, 'A\_3': {'depth': 5, 'edges': '3/3', 'charges compared beyond the leading data': [52, 38, 38], 'support violations': 0}, '[A\_1,D\_4]': {'depth': 4, 'edges': '4/4', 'charges compared beyond the leading data': [30, 65, 65, 65], 'support violations': 0}, 'Kronecker (pure SU(2))': {'depth': 6, 'edges': '2/2', 'charges compared beyond the leading data': [25, 15], 'support violations': 0}, 'SU(2)+1': {'depth': 5, 'edges': '3/3', 'charges compared beyond the leading data': [38, 29, 38], 'support violations': 0}, 'SU(2)+2': {'depth': 4, 'edges': '4/4', 'charges compared beyond the leading data': [65, 65, 37, 37], 'support violations': 0}, 'SU(2)+3': {'depth': 4, 'edges': '5/5', 'charges compared beyond the leading data': [120, 92, 92, 92, 92], 'support violations': 0}, 'Markov': {'depth': 5, 'edges': '3/3', 'charges compared beyond the leading data': [29, 29, 29], 'support violations': 0}}; sample: first 120 entries with n $\le$ 4 of dictionaries/enumerated (shard format 3, E = 12: 728921 strongly connected quivers, 674597 with a green spec, 54324 crystalline-only, 0 uncertified), every node; seconds: 30.1. & \MEASURED\newline {\footnotesize conjecture, fast\newline web}\newline {\footnotesize run web 2026-09-29: 4/4} \\ \addlinespace
\code{conj:factored}\newline Conj.~\pref{conj:factored} & $S_Q=\RG^{Q',\Gamma}(S_Q^{Q'})\,S_{Q'}$ for every full subquiver, defining a flow $K_\fq(Q,\Gamma)\to K_\fq(Q',\Gamma)$. & (A) On the independent arithmetic (\code{experiments/sw_factored_flows.py}): $S_Q$ built with the charges of $Q'$ last equals $R\,S_{Q'}$, $R$ the product of the other factors; $R=\RG^{Q'}(S_Q^{Q'})$ with $S_Q^{Q'}\in1+\fq\,\bZ[\fq,(1-\fq^{2n})^{-1}]$, the ring of \pref{ax:S}; $\RG^Q$ at every node decomposes in the $\RG^{Q'}$ basis over $\bZ[\fq^{\pm1}]$; the pairs where negative $\fq$-powers had to cancel for the ring condition are counted.  (B) \code{DirectionalSubquiverRG}, with a spec whose factors carrying the dropped node come first (found by $S$-preserving local moves): $\RG^Q=\RG^{Q'}\circ\RG^{Q,Q'}$ exactly on node charges and pairwise sums; its $S_\RG$ (a peel of vacuum kets) against (A) as exact $\fq$-series through $\fq^8$; the flow's axioms; its UV algebra, derived from the IR chart and $S_\RG$ alone, against the independently built $Q$-chart (\code{certify_directional_vs_bps}). \emph{Population:} (A) every full subquiver $Q'$ of the Section~4 lineup's quivers (the Markov quiver included) and of the $11$ dictionary-sample charts of the \pref{conj:upper} row, through cone degree $6$, $5$, $4$ by rank; (B) every node deletion of the lineup's charts with a spec that \code{DirectionalSubquiverRG} can build (web); the local builds (local). \emph{Controls:} positive: $Q'=Q$ and $Q'=\emptyset$; the pentagon by hand.  negative: a perturbed $S_{Q'}$ breaks the factorisation; the wrong sub-quiver's $\RG$ basis leaves non-Laurent coefficients; the pentagon's $S$ in the wrong order leaves the ring of \pref{ax:S}; the pentagon's flow mapped through the other deletion's IR chart fails to reproduce $\RG^Q$. \emph{Evidence:} emergent. \emph{Where:} \code{battery/checks_rg.py (check_factored)}, \code{experiments/sw_factored_flows.py}, \code{src/bps/directional_subquiver_rg.py}, \code{tests/test_directional_subquiver_rg.py}, \code{scripts/certify_directional_on_dictionary.py}. \emph{Inputs:} \emph{web} --- \code{dictionaries/enumerated/} (tracked; E, accept\_depth, verify\_depth read from \code{manifest.json}): every (Q, Q') with Q' a full subquiver; \emph{local} --- the enumerated build beyond the shipped weight (E $\ge$ 13: the weight-13 tier, which stays on the local machine by the author's ruling, at a package of the source repository) and the flavoured builds beyond the shipped weights (gitignored while in progress: \code{.flavoured_*stage/}). \emph{Note:} the pro-nilpotent group hosting S\_Q\^{}{Q'} is open (pronilpotent\_group\_conjecture.md); on the 80 lineup pairs every coefficient of S\_Q\^{}{Q'} but the identity's needs a (1-q\^{}2n) denominator.  2026-09-27 web run: (A) 146 pairs, every full subquiver of the eight lineup quivers (80, the Markov quiver included) and of 11 sample charts (66) -- all pass; on 115 of them the RG\^{}{Q'} elements used carry negative powers of q, so the ring condition had terms to cancel.  (B) 21 node deletions build (the class's arranger called with a longer cap, which finds the 11-factor spec [A\_1,D\_4] needs to drop its centre); all 21 factorise RG exactly, match (A)'s S\_RG and pass the flow axioms; the UV certificate holds on 16, the five SU(2)+3 deletions needing more than 30 s.  2026-09-27 extensive run (600 s per deletion, the untraced checks first): the flow axioms and the UV certificate hold on all 21, the five SU(2)+3 deletions included; with traces to the third power of q both hold on the 6 deletions of the pentagon, A\_3 and the Kronecker quiver, and are not reached within 600 s on the other 15.  Not buildable: the Kronecker quiver and SU(2)+1 dropping node 1, where only the weak-coupling chamber -- infinitely many factors, the W boson among them -- puts the dropped node's factors first; (A) covers them. \emph{Run (web, 2026-09-27):} 8 checks, 0 failing; (A) pairs: 146; (A) quivers: ['pentagon', 'A\_3', '[A\_1,D\_4]', 'Kronecker (pure SU(2))', 'SU(2)+1', 'SU(2)+2', 'SU(2)+3', 'Markov', 'sample [[0, 1, -1], [-1, 0, 1], [1, -1, 0]]', 'sample [[0, 1, -2], [-1, 0, 1], [2, -1, 0]]', 'sample [[0, -2, 1], [2, 0, -2], [-1, 2, 0]]', 'sample [[0, 1, -3], [-1, 0, 1], [3, -1, 0]]', 'sample [[0, -3, 1], [3, 0, -2], [-1, 2, 0]]', 'sample [[0, 1, -4], [-1, 0, 1], [4, -1, 0]]', 'sample [[0, 2, -3], [-2, 0, 1], [3, -1, 0]]', 'sample [[0, 2, -2], [-2, 0, 2], [2, -2, 0]]', 'sample [[0, -4, 1], [4, 0, -2], [-1, 2, 0]]', 'sample [[0, -3, 1], [3, 0, -3], [-1, 3, 0]]', 'sample [[0, 1, -5], [-1, 0, 1], [5, -1, 0]]']; (B) node deletions: {'pentagon dropping node 0': {'RG factorises': [3, 3], 'S\_RG = the independent decomposition': [7, 7], 'flow axioms': True, 'UV = the Q-chart': True}, 'pentagon dropping node 1': {'RG factorises': [3, 3], 'S\_RG = the independent decomposition': [7, 7], 'flow axioms': True, 'UV = the Q-chart': True}, 'A\_3 dropping node 0': {'RG factorises': [6, 6], 'S\_RG = the independent decomposition': [6, 6], 'flow axioms': True, 'UV = the Q-chart': True}, 'A\_3 dropping node 1': {'RG factorises': [6, 6], 'S\_RG = the independent decomposition': [6, 6], 'flow axioms': True, 'UV = the Q-chart': True}, 'A\_3 dropping node 2': {'RG factorises': [6, 6], 'S\_RG = the independent decomposition': [6, 6], 'flow axioms': True, 'UV = the Q-chart': True}, '[A\_1,D\_4] dropping node 0': {'RG factorises': [10, 10], 'S\_RG = the independent decomposition': [5, 5], 'flow axioms': True, 'UV = the Q-chart': True}, '[A\_1,D\_4] dropping node 1': {'RG factorises': [10, 10], 'S\_RG = the independent decomposition': [5, 5], 'flow axioms': True, 'UV = the Q-chart': True}, '[A\_1,D\_4] dropping node 2': {'RG factorises': [10, 10], 'S\_RG = the independent decomposition': [5, 5], 'flow axioms': True, 'UV = the Q-chart': True}, '[A\_1,D\_4] dropping node 3': {'RG factorises': [10, 10], 'S\_RG = the independent decomposition': [5, 5], 'flow axioms': True, 'UV = the Q-chart': True}, 'Kronecker (pure SU(2)) dropping node 0': {'RG factorises': [3, 3], 'S\_RG = the independent decomposition': [7, 7], 'flow axioms': True, 'UV = the Q-chart': True}, 'SU(2)+1 dropping node 0': {'RG factorises': [6, 6], 'S\_RG = the independent decomposition': [6, 6], 'flow axioms': True, 'UV = the Q-chart': True}, 'SU(2)+1 dropping node 2': {'RG factorises': [6, 6], 'S\_RG = the independent decomposition': [8, 8], 'flow axioms': True, 'UV = the Q-chart': True}, 'SU(2)+2 dropping node 0': {'RG factorises': [10, 10], 'S\_RG = the independent decomposition': [5, 5], 'flow axioms': True, 'UV = the Q-chart': True}, 'SU(2)+2 dropping node 1': {'RG factorises': [10, 10], 'S\_RG = the independent decomposition': [5, 5], 'flow axioms': True, 'UV = the Q-chart': True}, 'SU(2)+2 dropping node 2': {'RG factorises': [10, 10], 'S\_RG = the independent decomposition': [5, 5], 'flow axioms': True, 'UV = the Q-chart': True}, 'SU(2)+2 dropping node 3': {'RG factorises': [10, 10], 'S\_RG = the independent decomposition': [5, 5], 'flow axioms': True, 'UV = the Q-chart': True}, 'SU(2)+3 dropping node 0': {'RG factorises': [15, 15], 'S\_RG = the independent decomposition': [5, 5], 'flow axioms': True, 'UV = the Q-chart': 'not reached in 30 s'}, 'SU(2)+3 dropping node 1': {'RG factorises': [15, 15], 'S\_RG = the independent decomposition': [5, 5], 'flow axioms': True, 'UV = the Q-chart': 'not reached in 30 s'}, 'SU(2)+3 dropping node 2': {'RG factorises': [15, 15], 'S\_RG = the independent decomposition': [5, 5], 'flow axioms': True, 'UV = the Q-chart': 'not reached in 30 s'}, 'SU(2)+3 dropping node 3': {'RG factorises': [15, 15], 'S\_RG = the independent decomposition': [5, 5], 'flow axioms': True, 'UV = the Q-chart': 'not reached in 30 s'}, 'SU(2)+3 dropping node 4': {'RG factorises': [15, 15], 'S\_RG = the independent decomposition': [5, 5], 'flow axioms': True, 'UV = the Q-chart': 'not reached in 30 s'}}; (B) not buildable: ['Kronecker (pure SU(2)) dropping node 1', 'SU(2)+1 dropping node 1']; seconds: 217.0. \emph{Run (web (extensive), 2026-09-27):} 10 checks, 0 failing; (A) pairs: 146; (A) quivers: ['pentagon', 'A\_3', '[A\_1,D\_4]', 'Kronecker (pure SU(2))', 'SU(2)+1', 'SU(2)+2', 'SU(2)+3', 'Markov', 'sample [[0, 1, -1], [-1, 0, 1], [1, -1, 0]]', 'sample [[0, 1, -2], [-1, 0, 1], [2, -1, 0]]', 'sample [[0, -2, 1], [2, 0, -2], [-1, 2, 0]]', 'sample [[0, 1, -3], [-1, 0, 1], [3, -1, 0]]', 'sample [[0, -3, 1], [3, 0, -2], [-1, 2, 0]]', 'sample [[0, 1, -4], [-1, 0, 1], [4, -1, 0]]', 'sample [[0, 2, -3], [-2, 0, 1], [3, -1, 0]]', 'sample [[0, 2, -2], [-2, 0, 2], [2, -2, 0]]', 'sample [[0, -4, 1], [4, 0, -2], [-1, 2, 0]]', 'sample [[0, -3, 1], [3, 0, -3], [-1, 3, 0]]', 'sample [[0, 1, -5], [-1, 0, 1], [5, -1, 0]]']; (B) node deletions: {'pentagon dropping node 0': {'RG factorises': [3, 3], 'S\_RG = the independent decomposition': [7, 7], 'flow axioms': True, 'UV = the Q-chart': True, 'flow axioms with traces to q\^{}3': True, 'UV = the Q-chart with traces to q\^{}3': True}, 'pentagon dropping node 1': {'RG factorises': [3, 3], 'S\_RG = the independent decomposition': [7, 7], 'flow axioms': True, 'UV = the Q-chart': True, 'flow axioms with traces to q\^{}3': True, 'UV = the Q-chart with traces to q\^{}3': True}, 'A\_3 dropping node 0': {'RG factorises': [6, 6], 'S\_RG = the independent decomposition': [6, 6], 'flow axioms': True, 'UV = the Q-chart': True, 'flow axioms with traces to q\^{}3': True, 'UV = the Q-chart with traces to q\^{}3': True}, 'A\_3 dropping node 1': {'RG factorises': [6, 6], 'S\_RG = the independent decomposition': [6, 6], 'flow axioms': True, 'UV = the Q-chart': True, 'flow axioms with traces to q\^{}3': True, 'UV = the Q-chart with traces to q\^{}3': True}, 'A\_3 dropping node 2': {'RG factorises': [6, 6], 'S\_RG = the independent decomposition': [6, 6], 'flow axioms': True, 'UV = the Q-chart': True, 'flow axioms with traces to q\^{}3': True, 'UV = the Q-chart with traces to q\^{}3': True}, '[A\_1,D\_4] dropping node 0': {'RG factorises': [10, 10], 'S\_RG = the independent decomposition': [5, 5], 'flow axioms': True, 'UV = the Q-chart': True, 'flow axioms with traces to q\^{}3': 'not reached in 600 s', 'UV = the Q-chart with traces to q\^{}3': 'not reached in 600 s'}, '[A\_1,D\_4] dropping node 1': {'RG factorises': [10, 10], 'S\_RG = the independent decomposition': [5, 5], 'flow axioms': True, 'UV = the Q-chart': True, 'flow axioms with traces to q\^{}3': 'not reached in 600 s', 'UV = the Q-chart with traces to q\^{}3': 'not reached in 600 s'}, '[A\_1,D\_4] dropping node 2': {'RG factorises': [10, 10], 'S\_RG = the independent decomposition': [5, 5], 'flow axioms': True, 'UV = the Q-chart': True, 'flow axioms with traces to q\^{}3': 'not reached in 600 s', 'UV = the Q-chart with traces to q\^{}3': 'not reached in 600 s'}, '[A\_1,D\_4] dropping node 3': {'RG factorises': [10, 10], 'S\_RG = the independent decomposition': [5, 5], 'flow axioms': True, 'UV = the Q-chart': True, 'flow axioms with traces to q\^{}3': 'not reached in 600 s', 'UV = the Q-chart with traces to q\^{}3': 'not reached in 600 s'}, 'Kronecker (pure SU(2)) dropping node 0': {'RG factorises': [3, 3], 'S\_RG = the independent decomposition': [7, 7], 'flow axioms': True, 'UV = the Q-chart': True, 'flow axioms with traces to q\^{}3': True, 'UV = the Q-chart with traces to q\^{}3': True}, 'SU(2)+1 dropping node 0': {'RG factorises': [6, 6], 'S\_RG = the independent decomposition': [6, 6], 'flow axioms': True, 'UV = the Q-chart': True, 'flow axioms with traces to q\^{}3': 'not reached in 600 s', 'UV = the Q-chart with traces to q\^{}3': 'not reached in 600 s'}, 'SU(2)+1 dropping node 2': {'RG factorises': [6, 6], 'S\_RG = the independent decomposition': [8, 8], 'flow axioms': True, 'UV = the Q-chart': True, 'flow axioms with traces to q\^{}3': 'not reached in 600 s', 'UV = the Q-chart with traces to q\^{}3': 'not reached in 600 s'}, 'SU(2)+2 dropping node 0': {'RG factorises': [10, 10], 'S\_RG = the independent decomposition': [5, 5], 'flow axioms': True, 'UV = the Q-chart': True, 'flow axioms with traces to q\^{}3': 'not reached in 600 s', 'UV = the Q-chart with traces to q\^{}3': 'not reached in 600 s'}, 'SU(2)+2 dropping node 1': {'RG factorises': [10, 10], 'S\_RG = the independent decomposition': [5, 5], 'flow axioms': True, 'UV = the Q-chart': True, 'flow axioms with traces to q\^{}3': 'not reached in 600 s', 'UV = the Q-chart with traces to q\^{}3': 'not reached in 600 s'}, 'SU(2)+2 dropping node 2': {'RG factorises': [10, 10], 'S\_RG = the independent decomposition': [5, 5], 'flow axioms': True, 'UV = the Q-chart': True, 'flow axioms with traces to q\^{}3': 'not reached in 600 s', 'UV = the Q-chart with traces to q\^{}3': 'not reached in 600 s'}, 'SU(2)+2 dropping node 3': {'RG factorises': [10, 10], 'S\_RG = the independent decomposition': [5, 5], 'flow axioms': True, 'UV = the Q-chart': True, 'flow axioms with traces to q\^{}3': 'not reached in 600 s', 'UV = the Q-chart with traces to q\^{}3': 'not reached in 600 s'}, 'SU(2)+3 dropping node 0': {'RG factorises': [15, 15], 'S\_RG = the independent decomposition': [5, 5], 'flow axioms': True, 'UV = the Q-chart': True, 'flow axioms with traces to q\^{}3': 'not reached in 600 s', 'UV = the Q-chart with traces to q\^{}3': 'not reached in 600 s'}, 'SU(2)+3 dropping node 1': {'RG factorises': [15, 15], 'S\_RG = the independent decomposition': [5, 5], 'flow axioms': True, 'UV = the Q-chart': True, 'flow axioms with traces to q\^{}3': 'not reached in 600 s', 'UV = the Q-chart with traces to q\^{}3': 'not reached in 600 s'}, 'SU(2)+3 dropping node 2': {'RG factorises': [15, 15], 'S\_RG = the independent decomposition': [5, 5], 'flow axioms': True, 'UV = the Q-chart': True, 'flow axioms with traces to q\^{}3': 'not reached in 600 s', 'UV = the Q-chart with traces to q\^{}3': 'not reached in 600 s'}, 'SU(2)+3 dropping node 3': {'RG factorises': [15, 15], 'S\_RG = the independent decomposition': [5, 5], 'flow axioms': True, 'UV = the Q-chart': True, 'flow axioms with traces to q\^{}3': 'not reached in 600 s', 'UV = the Q-chart with traces to q\^{}3': 'not reached in 600 s'}, 'SU(2)+3 dropping node 4': {'RG factorises': [15, 15], 'S\_RG = the independent decomposition': [5, 5], 'flow axioms': True, 'UV = the Q-chart': True, 'flow axioms with traces to q\^{}3': 'not reached in 600 s', 'UV = the Q-chart with traces to q\^{}3': 'not reached in 600 s'}}; (B) not buildable: ['Kronecker (pure SU(2)) dropping node 1', 'SU(2)+1 dropping node 1']; seconds: 9168.3. & \CERTIFIED\newline {\footnotesize conjecture, extensive\newline web, local}\newline {\footnotesize run web 2026-09-27: 8/8}\newline {\footnotesize run web (extensive) 2026-09-27: 10/10} \\ \addlinespace
\code{sec:rg/flavoured-quivers}\newline \peqref{eq:fleq}, \peqref{eq:flavouredquiver}, Def.~4.13 & $N$ identical nodes make the algebra $SU(N)$-flavoured; the flavoured constraint \peqref{eq:flavouredquiver}; the PBW factorisation extends to $\cE_{G_f}$; the flavoured flow lands on $R(G_f)\otimes Q_\fq(\Gamma)$. & $\Omega(\gamma,r)$ comes out a $G_f$-character at every charge (the engine forces it weight by weight and never imposes Weyl invariance); the covariant spec replays at depth $4$, or, where the tier stores only $\Omega$, a fresh run reproduces it; $S$ equals the $S$ of the unfolded quiver, built by the unflavoured engine; the PBW factorisation extended to the group generated by the $E^{(s;R)}_\fq(X_\gamma)^{\pm1}$: $S$ is unchanged under random orders on (charge, spin, irrep), with $\Omega$ a character in each, while the order moves the factor content; \peqref{eq:flavouredquiver} as a regression guard (the engine imposes it).  The flavoured $K_\fq$-algebra: \code{GBPSKAlgebra} (products on the unfolded chart, irreps recovered by character decomposition) against \code{GfBPSKAlgebra} (Weyl-character diagrams in the F-basis), which share no code, product for product; the bar involution, $\rho$ an automorphism and $\rho\circ\rho^{-1}$ on \code{GBPSKAlgebra}. \emph{Population:} samples of the ten shipped flavoured tiers, which since 2026-09-28 store only their strongly connected quivers --- three entries each (up to $12$ at the extensive depth): stored entries from the smallest stored reduced rank up, topped up where a tier stores fewer with up to three composed answers, a $3$-cycle with the tier's orbits attached as identical sinks (cyclic, not strongly connected).  $SU(6)$, $SU(3)\times SU(3)$, $SU(4)\times SU(2)$ and $SU(2)^3$ store nothing at their weight and are sampled by composed answers alone, $SU(5)$ stores one entry; at the extensive depth that gives $12$ entries on $SU(2)$, $SU(3)$, $SU(4)$ and $SU(2)\times SU(2)$, $11$ on $SU(3)\times SU(2)$, $4$ on $SU(5)$ and $3$ on the other four.  Cutoff $3$; $[A_1,D_3]$ with $SU(2)$ and $[A_1,D_4]$ with $SU(3)$, each presented twice (web); the staging builds (local). \emph{Controls:} positive: the covariant spec replays and $S$ equals the unfolded $S$; negative: ordering the individual weights $E^{(s)}_\fq(\mu^w X_\gamma)$ leaves $S$ unchanged but breaks the character property. \emph{Evidence:} emergent. \emph{Where:} \code{battery/checks_rg.py (check_flavoured_quivers)}, \code{src/bps/flavoured_factor_spectrum.py}, \code{src/bps/flavoured_spec.py}, \code{src/bps/gbps_kalgebra.py}, \code{src/bps/gf_bps_kalgebra.py}, \code{tests/test_gbps_kalgebra.py}, \code{tests/test_flavoured_factor_spectrum.py}, \code{tests/test_gf_bps_kalgebra.py}. \emph{Inputs:} \emph{web} --- a package of the source repository (tracked; orbit\_size, max\_weight, coverage read from the ten manifests); \emph{local} --- the flavoured staging builds \code{.flavoured_*stage/} (gitignored while in progress; every shipped tier is complete through its weight, 10, 11 or 12). \emph{Note:} 0 crystalline\_failed across the ten shipped tiers (the manifests).  2026-09-27 web run: 30 entries, three per tier -- Omega a character 30/30, spec replay (or, for the SU(4) entries that store only Omega, a fresh run reproducing it) 30/30, S = the unfolded S 30/30, S order-independent 30/30 with the order moving the factor content on all 30; the weight-level order breaks the character property in 4 of 6 orders; GBPSKAlgebra = GfBPSKAlgebra on 100/100 products at [A\_1,D\_3] with SU(2) and 100/100 at [A\_1,D\_4] with SU(3).  the design record (2026-09-28) moved the tiers to shard format 3, and the local session topped up \_flav\_sample with composed answers.  On the first staging this check caught carried crystalline Omega that did not match a fresh run: the content is that of the engine's order, which depends on the node labelling, so the tiers now recompute Omega in the stored labelling.  The 2026-09-29 records are on the format-3 tiers. \emph{Run (web, 2026-09-29):} 6 checks, 0 failing; tiers: {'SU(2)': 3, 'SU(3)': 3, 'SU(4)': 3, 'SU(5)': 3, 'SU(6)': 3, 'SU(2)xSU(2)': 3, 'SU(3)xSU(2)': 3, 'SU(3)xSU(3)': 3, 'SU(4)xSU(2)': 3, 'SU(2)\^{}3': 3}; cutoff: 3; per tier: 3 stored entries at the tier's smallest stored reduced rank; presentations: ['[A\_1,D\_3] with SU(2)', '[A\_1,D\_4] with SU(3)']; seconds: 18.1. \emph{Run (web (extensive), 2026-09-29):} 6 checks, 0 failing; tiers: {'SU(2)': 12, 'SU(3)': 12, 'SU(4)': 12, 'SU(5)': 4, 'SU(6)': 3, 'SU(2)xSU(2)': 12, 'SU(3)xSU(2)': 11, 'SU(3)xSU(3)': 3, 'SU(4)xSU(2)': 3, 'SU(2)\^{}3': 3}; cutoff: 3; per tier: 12 stored entries at the tier's smallest stored reduced rank; presentations: ['[A\_1,D\_3] with SU(2)', '[A\_1,D\_4] with SU(3)']; seconds: 25.6. & \MEASURED\newline {\footnotesize conjecture, extensive\newline web, local}\newline {\footnotesize run web 2026-09-29: 6/6}\newline {\footnotesize run web (extensive) 2026-09-29: 6/6} \\ \addlinespace
\end{longtable}

\subsection*{Section~\pref{sec:skein}: skein algebras}
\begin{longtable}{@{}>{\raggedright\arraybackslash}p{2.4cm}>{\raggedright\arraybackslash}p{4.7cm}>{\raggedright\arraybackslash}p{4.95cm}>{\raggedright\arraybackslash}p{1.85cm}@{}}
\toprule claim & statement & test: procedure, population, controls & standing \\ \midrule \endfirsthead
\toprule claim & statement & test: procedure, population, controls & standing \\ \midrule \endhead
\bottomrule \endfoot
\code{conj:junctions}\newline Conj.~\pref{conj:junctions}, Remarks 5.1--5.2 & The invariant spaces $\mathrm{In}_{(\lambda_a)}$ have bar-invariant canonical bases with cyclic rotation a permutation and palindromic collapse coefficients; non-negativity fails. & palindromicity of the skein structure constants (\code{verify_bar_involution}); no witness of a negative coefficient (Remark 5.1) is pinned yet; higher-rank webs not encoded . \emph{Population:} $\fsl_2$ skeins on punctured spheres and the annulus (curves only). \emph{Evidence:} emergent. \emph{Note:} The scope is sl\_2 curves; canonical junctions, the collapse of webs and higher rank belong to a separate skein algebra toolset project, and no witness for Remark 5.1 is pinned here. & \MEASURED\newline {\footnotesize conjecture, fast\newline web} \\ \addlinespace
\code{sec:skein/kq-structure}\newline Section~\pref{sec:skein}, the paragraphs after Conj.~\pref{conj:junctions} & $\mathrm{Sk}_\fq(\fg,C)$ is a $K_\fq$-algebra: the bar involution is the reflection of $[-1,1]$, the canonical basis is the irreducible skeins drawn at $C\times\{0\}$ (other skeins reducing to them with $\fq$-number coefficients), $\rho$ is a dualization, and $\rho^2=1$ in the absence of irregular punctures. & (a) Every product on each window: on the nine bordered instances the stated-skein computation reproduces the returned product's coefficients inside the multiplication (up to a normalisation per label, fixed at its first occurrence and asserted at every later one); on the sphere the product is the Kauffman bracket, and there --- product and trace both skein-side --- the axioms emergent on it: \pref{ax:bar} (the mirror) and the orthonormality and cyclicity of the Askey--Wilson trace. (b) The canonical basis: on the four-punctured sphere's chart the curves carry the adjoint colouring, and $L_\gamma L_{k\gamma}=L_{(k+1)\gamma}+L_{k\gamma}+L_{(k-1)\gamma}$ with unit coefficients ($V_2\otimes V_{2k}$ with its three summands), channels $a,b,c$, $k\le 2$ ($k\le 3$ extensive). (c) $\rho$ on both sphere charts: the spec-derived $\sigma$ fixes every label of the box up to a flavour charge, and $\rho^2=1$ there. (d) With an irregular puncture $\rho^2\neq 1$: the pentagon's box, and a window label of $\rho$-order above $2$ on every skein-computing roster instance but the sphere. \emph{Population:} $\fg=\fsl_2$ (curves; no webs are encoded): the $10$ instances of the package's roster that carry a skein computation --- discs with $5$ to $11$ marked points and their $U(1)$-gaugings, $SU(2)$ on the annulus, the four-punctured sphere --- on a window of $5$ to $6$ labels each (the other $7$ return an external presentation's product with no skein computation); the charts of the four- and five-punctured spheres, on the boxes $[-2,2]^6$ and $[-1,1]^9$ of their charge lattices. \emph{Controls:} negative: the package's fundamental-character sample class obeys the two-term law $L^2=L_2+1$ on all three channels and fails (b); the pentagon's chart has $\rho^2\neq 1$ on $48$ of the $49$ labels of $[-3,3]^2$. \emph{Evidence:} per clause: (a) the skein computation checks an independent presentation's product (on the sphere it is the product); (b) and (c) emergent (the chart's product and $\sigma$ come from its quiver and spec); the flavour inversion in $\rho$ is imposed by the code (enforced). \emph{Where:} \code{battery/checks_skein.py (check_skein_structure)}, \code{src/skein/skein_kalgebra.py}, \code{src/skein/skein_atlas.py}. \emph{Note:} Added 2026-09-27 for the Section 5 walkthrough and kept with these light checks (a skein algebra toolset is a separate project).  The draft's axioms are checked only on the sphere, where product and trace are both skein-side; on the other instances they are the underlying presentation's, tested in Sections 2 to 4.  Dualization is trivial on sl\_2 curves (every V\_k is self-dual), so (c) measures that $\rho$ acts as the identity on curves; the flavour inversion is imposed by BPSKAlgebra.\_sec\_rectified\_map.  On the roster's sphere instance $\rho$ is wired as the identity, so (c) reads the sphere through its atlas charts.  Not tested: the global forms (root-lattice labels for the small skein algebra, a Lagrangian in the first cohomology with Z(g) coefficients for the others). \emph{Run (web, 2026-09-27):} 10 checks, 0 failing; skein-computing roster windows (labels): {'u1square': 5, 'pentagon': 5, 'u1hexagon': 5, 'heptagon': 5, 'u1octagon': 5, 'nonagon': 5, 'u1decagon': 5, 'hendecagon': 5, 'annulus-pure-su2': 5, 'sphere-su2-nf4': 6}; rho-orders on the windows: {'u1square': ['>12', '>12', 2, '>12'], 'pentagon': [5, 5, 5, 5], 'u1hexagon': ['>12', '>12', '>12', '>12'], 'heptagon': [7, 7, 7, 7], 'u1octagon': ['>12', '>12', '>12', '>12'], 'nonagon': [9, 9, 9, 9], 'u1decagon': ['>12', '>12', '>12', '>12'], 'hendecagon': [11, 11, 11, 11], 'annulus-pure-su2': ['>12', 1, '>12', 1], 'sphere-su2-nf4': [1, 1, 1, 1, 1]}; roster instances with no skein computation: ['u1gauged-su2-nf1 (rgflow-anchored)', 'a1d3 (cone-anchored)', 'u1a1d4 (cone-anchored)', 'annulus-su2-nf1 (bordered-chart-anchored)', 'disk-su2-nf2 (bordered-chart-anchored)', 'disk-su2-nf3 (bordered-chart-anchored)', 'su2-2punct-a1d3 (bordered-chart-anchored)']; tower: channels a, b, c, k $\le$ 2; rho boxes: [-2,2]\^{}6 (four punctures), [-1,1]\^{}9 (five), [-3,3]\^{}2 (pentagon); seconds: 207.8. \emph{Run (web (extensive), 2026-09-27):} 10 checks, 0 failing; skein-computing roster windows (labels): {'u1square': 6, 'pentagon': 6, 'u1hexagon': 6, 'heptagon': 6, 'u1octagon': 6, 'nonagon': 6, 'u1decagon': 6, 'hendecagon': 6, 'annulus-pure-su2': 5, 'sphere-su2-nf4': 6}; rho-orders on the windows: {'u1square': ['>12', '>12', 2, 2, '>12'], 'pentagon': [5, 5, 5, 5, 5], 'u1hexagon': ['>12', '>12', '>12', '>12', '>12'], 'heptagon': [7, 7, 7, 7, 7], 'u1octagon': ['>12', '>12', '>12', '>12', '>12'], 'nonagon': [9, 9, 9, 9, 9], 'u1decagon': ['>12', '>12', '>12', '>12', '>12'], 'hendecagon': [11, 11, 11, 11, 11], 'annulus-pure-su2': ['>12', 1, '>12', 1], 'sphere-su2-nf4': [1, 1, 1, 1, 1]}; roster instances with no skein computation: ['u1gauged-su2-nf1 (rgflow-anchored)', 'a1d3 (cone-anchored)', 'u1a1d4 (cone-anchored)', 'annulus-su2-nf1 (bordered-chart-anchored)', 'disk-su2-nf2 (bordered-chart-anchored)', 'disk-su2-nf3 (bordered-chart-anchored)', 'su2-2punct-a1d3 (bordered-chart-anchored)']; tower: channels a, b, c, k $\le$ 3; rho boxes: [-2,2]\^{}6 (four punctures), [-1,1]\^{}9 (five), [-3,3]\^{}2 (pentagon); seconds: 762.5. & \MEASURED\newline {\footnotesize conjecture, extensive\newline web}\newline {\footnotesize run web 2026-09-27: 10/10}\newline {\footnotesize run web (extensive) 2026-09-27: 10/10} \\ \addlinespace
\code{sec:skein/trace}\newline Section~\pref{sec:skein}, the trace & The $K_\fq$ trace on $\mathrm{Sk}_\fq(\fg,C)$ is the Kapustin--Witten path integral on $D^2\times C$; in particular it is independent of the triangulation used to compute it. & $\Tr$ of each sample in the root chart equals $\Tr$ of its image in the flipped chart, each chart's trace computed from its own quiver and transported spec; and each flip commutes with $\rho$. \emph{Population:} the bordered flip atlases, every flip from the root chart: the pentagon ($2$ flips), the heptagon ($4$), $[A_1,D_3]$ ($3$), the $U(1)$-gauged square ($1$) and hexagon ($3$), and $[A_1,D_5]$ ($5$) at the extensive depth; samples $1$, four rays and their sixteen products, to $\fq^4$. \emph{Controls:} negative: the pentagon's flipped chart rebuilt with its spec reversed (a wrong chamber) disagrees with the root chart's trace on $15$ of $21$ samples. \emph{Evidence:} emergent (each chart's trace is computed from its own quiver and spec). \emph{Where:} \code{battery/checks_skein.py (check_skein_trace)}, \code{src/skein/skein_atlas.py}, \code{notes/KQ_CONJECTURE.md}. \emph{Note:} The Kapustin-Witten path integral on D2 x C has no independent computation in the repository; the row tests its consequence, independence of the triangulation.  The closed surfaces are not in the sweep: one flip of the four-punctured sphere's chart did not finish its trace comparison through the second power of q within 900 s; through the first power it finished (1066 s) and agreed, but there the four sample traces are 1, 0, 0 and the flavour monomial on any chart (orthonormality at the leading order, and the odd-power selection rule KQ\_CONJECTURE.md records at four punctures), so that agreement carries almost no information (the recorded chamber comparison is a test in the source repository).  Reusing labels verbatim across a flip is not caught on the pentagon (that map preserves the trace there), so the control is a wrong trace.  KQ\_CONJECTURE.md records, at four punctures, the trace equal to the SU(2) N\_f = 4 Schur contour functional through the third power of q, and a trace space of dimension 5.  Comparisons on closed surfaces belong to the separate skein algebra toolset project. \emph{Run (web, 2026-09-27):} 3 checks, 0 failing; atlases (flips, trace-equivariant, rho-equivariant, samples): {'pentagon': [2, 2, 2, 21], 'heptagon': [4, 4, 4, 21], '[A\_1,D\_3]': [3, 3, 3, 21], 'U(1)-gauged square': [1, 1, 1, 21], 'U(1)-gauged hexagon': [3, 3, 3, 21]}; K: 4; seconds: 82.0. \emph{Run (web (extensive), 2026-09-27):} 3 checks, 0 failing; atlases (flips, trace-equivariant, rho-equivariant, samples): {'pentagon': [2, 2, 2, 21], 'heptagon': [4, 4, 4, 21], '[A\_1,D\_3]': [3, 3, 3, 21], 'U(1)-gauged square': [1, 1, 1, 21], 'U(1)-gauged hexagon': [3, 3, 3, 21], '[A\_1,D\_5]': [5, 5, 5, 21]}; K: 4; seconds: 2337.9. & \MEASURED\newline {\footnotesize conjecture, extensive\newline web}\newline {\footnotesize run web 2026-09-27: 3/3}\newline {\footnotesize run web (extensive) 2026-09-27: 3/3} \\ \addlinespace
\end{longtable}

\subsection*{Appendix~\pref{app:physics}: physical motivations}
\begin{longtable}{@{}>{\raggedright\arraybackslash}p{2.4cm}>{\raggedright\arraybackslash}p{4.7cm}>{\raggedright\arraybackslash}p{4.95cm}>{\raggedright\arraybackslash}p{1.85cm}@{}}
\toprule claim & statement & test: procedure, population, controls & standing \\ \midrule \endfirsthead
\toprule claim & statement & test: procedure, population, controls & standing \\ \midrule \endhead
\bottomrule \endfoot
\code{app:physics}\newline App.~\pref{app:physics} & Physical origin of the axioms; the RG-interface pairing $(L^\IR_a,\RG_bS)_\IR$; the interface index \peqref{eq:int}. & physics; the two indices have repository counterparts only as candidates (\code{RGInterface}; the one-dressed-slot pairing) and \peqref{eq:int} is unverified. & \OPEN\newline {\footnotesize not-testable, none\newline no test} \\ \addlinespace
\end{longtable}

\subsection*{Appendix~\pref{app:finite}: finite type}
\begin{longtable}{@{}>{\raggedright\arraybackslash}p{2.4cm}>{\raggedright\arraybackslash}p{4.7cm}>{\raggedright\arraybackslash}p{4.95cm}>{\raggedright\arraybackslash}p{1.85cm}@{}}
\toprule claim & statement & test: procedure, population, controls & standing \\ \midrule \endfirsthead
\toprule claim & statement & test: procedure, population, controls & standing \\ \midrule \endhead
\bottomrule \endfoot
\code{tab:finite_type}\newline Table~\pref{tab:finite_type} & Flavour symmetry and order of $\rho$ for every $[A_1,ADE]$ family. & The order of $\rho$ on a window of labels (the generators, the terms of products of generator pairs, generators dressed by flavour characters), three ways: on sections (the unflavoured algebra), on labels, and as an automorphism of the flavoured algebra ($\rho$ acting on the flavour ring by the duality, the contract's axiom). Over the table's flavour group the flavoured order must be the table's; over an equal-rank subgroup $H$ it must be the least common multiple of the table's order and the order of the duality on $R(H)$, derived from $A_H=A_G\otimes R(H)$. The flavour group of each presentation is the table's or an equal-rank subgroup; the rank is the corank of the Dynkin quiver's exchange matrix; the table's two entries at $D_4$ (and at $A_3=D_3$) are checked against each other. \emph{Population:} every ungauged class serving an $[A_1,ADE]$ entry (the cone classes, the ungauged polygons, \code{SU3ADKAlg}, the zoo with its aliases once) and the BPS chart of every Dynkin quiver, $n\le 3$ ($n\le 5$ extensive). \emph{Controls:} positive: $\rho^{\mathrm{order}}$ fixes a second window on every presentation. negative: every single-entry perturbation of the table --- each order by a proper divisor or its double, each group by another --- is rejected, except a larger group of the same rank, which no measurement here can refute. \emph{Evidence:} emergent. \emph{Where:} \code{battery/checks_finite.py (check_finite_table)}, \code{battery/test_ade_serving_guard.py}. \emph{Note:} The order of $\rho$ depends on the flavour group a presentation is taken over: the duality V $\to$ V* is trivial on R(SU(2)) and inverts U(1) charges, so over the Cartan U(1) of the enhanced SU(2) the table's odd orders double -- the U(1)-flavoured presentations of [A1,A3] give 6, the BPS charts of [A1,D5] and [A1,D7] (flavour U(1) from ker B) give 10 and 14. On the unflavoured algebra the order is the table's on every presentation. The zoo's section-labelled classes carry the flavour in coefficients, where the duality enters through the contract's axiom (verify\_embed\_intertwines\_rho), not through their labels. \emph{Run (web, 2026-09-27):} 9 checks, 0 failing; presentations: {'A1A2kKAlg(1)': {'window': 16, 'dressed': 0, 'labels': 5, 'sections': 5, 'flavoured': 5, 'group': '---', 'star': 1, 'control\_moved': 0, 'control': 13}, 'A1A2kKAlg(2)': {'window': 56, 'dressed': 0, 'labels': 7, 'sections': 7, 'flavoured': 7, 'group': '---', 'star': 1, 'control\_moved': 0, 'control': 27}, 'A1A2kKAlg(3)': {'window': 72, 'dressed': 0, 'labels': 9, 'sections': 9, 'flavoured': 9, 'group': '---', 'star': 1, 'control\_moved': 0, 'control': 40}, 'PentagonKAlg': {'window': 16, 'dressed': 0, 'labels': 5, 'sections': 5, 'flavoured': 5, 'group': '---', 'star': 1, 'control\_moved': 0, 'control': 13}, 'HeptagonKAlg': {'window': 56, 'dressed': 0, 'labels': 7, 'sections': 7, 'flavoured': 7, 'group': '---', 'star': 1, 'control\_moved': 0, 'control': 26}, 'ungauge\_u1a1aodd(1)': {'window': 39, 'dressed': 7, 'labels': 6, 'sections': 3, 'flavoured': 6, 'group': 'U(1)', 'star': 2, 'control\_moved': 0, 'control': 22}, 'ungauge\_u1a1aodd(2)': {'window': 96, 'dressed': 17, 'labels': 8, 'sections': 8, 'flavoured': 8, 'group': 'U(1)', 'star': 2, 'control\_moved': 0, 'control': 53}, 'ungauge\_u1a1aodd(3)': {'window': 134, 'dressed': 23, 'labels': 10, 'sections': 10, 'flavoured': 10, 'group': 'U(1)', 'star': 2, 'control\_moved': 0, 'control': 47}, 'HexagonKAlg': {'window': 39, 'dressed': 7, 'labels': 6, 'sections': 3, 'flavoured': 6, 'group': 'U(1)', 'star': 2, 'control\_moved': 0, 'control': 22}, 'OctagonKAlg': {'window': 96, 'dressed': 17, 'labels': 8, 'sections': 8, 'flavoured': 8, 'group': 'U(1)', 'star': 2, 'control\_moved': 0, 'control': 53}, 'DecagonKAlg': {'window': 134, 'dressed': 23, 'labels': 10, 'sections': 10, 'flavoured': 10, 'group': 'U(1)', 'star': 2, 'control\_moved': 0, 'control': 47}, 'DodecagonKAlg': {'window': 183, 'dressed': 0, 'labels': 12, 'sections': 12, 'flavoured': 12, 'group': 'U(1)', 'star': 2, 'control\_moved': 0, 'control': 47}, 'A1DnKAlg(3)': {'window': 23, 'dressed': 0, 'labels': 3, 'sections': 3, 'flavoured': 3, 'group': 'SU(2)', 'star': 1, 'control\_moved': 0, 'control': 17}, 'A1DnKAlg(5)': {'window': 65, 'dressed': 0, 'labels': 5, 'sections': 5, 'flavoured': 5, 'group': 'SU(2)', 'star': 1, 'control\_moved': 0, 'control': 55}, 'A1DnKAlg(7)': {'window': 88, 'dressed': 0, 'labels': 7, 'sections': 7, 'flavoured': 7, 'group': 'SU(2)', 'star': 1, 'control\_moved': 0, 'control': 52}, 'A1DoddConeKAlg(0)': {'window': 23, 'dressed': 0, 'labels': 3, 'sections': 3, 'flavoured': 3, 'group': 'SU(2)', 'star': 1, 'control\_moved': 0, 'control': 20}, 'A1DoddConeKAlg(1)': {'window': 72, 'dressed': 0, 'labels': 5, 'sections': 5, 'flavoured': 5, 'group': 'SU(2)', 'star': 1, 'control\_moved': 0, 'control': 52}, 'A1DoddConeKAlg(2)': {'window': 97, 'dressed': 0, 'labels': 7, 'sections': 7, 'flavoured': 7, 'group': 'SU(2)', 'star': 1, 'control\_moved': 0, 'control': 50}, 'A1DevenKAlg(1)': {'window': 64, 'dressed': 6, 'labels': 4, 'sections': 4, 'flavoured': 4, 'group': 'SU(2)xU(1)', 'star': 2, 'control\_moved': 0, 'control': 26}, 'A1DevenKAlg(2)': {'window': 98, 'dressed': 19, 'labels': 6, 'sections': 6, 'flavoured': 6, 'group': 'SU(2)xU(1)', 'star': 2, 'control\_moved': 0, 'control': 78}, 'A1DevenKAlg(3)': {'window': 175, 'dressed': 11, 'labels': 8, 'sections': 8, 'flavoured': 8, 'group': 'SU(2)xU(1)', 'star': 2, 'control\_moved': 0, 'control': 67}, 'SU3ADKAlg': {'window': 52, 'dressed': 8, 'labels': 4, 'sections': 4, 'flavoured': 4, 'group': 'SU(3)', 'star': 2, 'control\_moved': 0, 'control': 33}, 'zoo pentagon': {'window': 16, 'dressed': 0, 'labels': 5, 'sections': 5, 'flavoured': 5, 'group': '---', 'star': 1, 'control\_moved': 0, 'control': 13}, 'zoo heptagon': {'window': 56, 'dressed': 0, 'labels': 7, 'sections': 7, 'flavoured': 7, 'group': '---', 'star': 1, 'control\_moved': 0, 'control': 28}, 'zoo a3': {'window': 19, 'dressed': 0, 'labels': 3, 'sections': 3, 'flavoured': 6, 'group': 'U(1)', 'star': 2, 'control\_moved': 0, 'control': 14}, 'zoo a5': {'window': 81, 'dressed': 0, 'labels': 8, 'sections': 8, 'flavoured': 8, 'group': 'U(1)', 'star': 2, 'control\_moved': 0, 'control': 48}, 'zoo a7': {'window': 134, 'dressed': 0, 'labels': 10, 'sections': 10, 'flavoured': 10, 'group': 'U(1)', 'star': 2, 'control\_moved': 0, 'control': 48}, 'zoo a1d3': {'window': 19, 'dressed': 0, 'labels': 3, 'sections': 3, 'flavoured': 3, 'group': 'SU(2)', 'star': 1, 'control\_moved': 0, 'control': 15}, 'zoo a1d4': {'window': 29, 'dressed': 0, 'labels': 4, 'sections': 4, 'flavoured': 4, 'group': 'SU(2)xU(1)', 'star': 2, 'control\_moved': 0, 'control': 18}, 'zoo a1d5': {'window': 69, 'dressed': 0, 'labels': 5, 'sections': 5, 'flavoured': 5, 'group': 'SU(2)', 'star': 1, 'control\_moved': 0, 'control': 43}, 'zoo a1d7': {'window': 99, 'dressed': 0, 'labels': 7, 'sections': 7, 'flavoured': 7, 'group': 'SU(2)', 'star': 1, 'control\_moved': 0, 'control': 55}, 'FiniteE6KAlgebra (zoo e6)': {'window': 101, 'dressed': 0, 'labels': 14, 'sections': 14, 'flavoured': 14, 'group': '---', 'star': 1, 'control\_moved': 0, 'control': 56}, 'FiniteE8KAlgebra (zoo e8)': {'window': 269, 'dressed': 0, 'labels': 16, 'sections': 16, 'flavoured': 16, 'group': '---', 'star': 1, 'control\_moved': 0, 'control': 176}, 'zoo a1d6': {'window': 105, 'dressed': 0, 'labels': 6, 'sections': 6, 'flavoured': 6, 'group': 'SU(2)xU(1)', 'star': 2, 'control\_moved': 0, 'control': 52}, 'zoo a1d8': {'window': 260, 'dressed': 0, 'labels': 8, 'sections': 8, 'flavoured': 8, 'group': 'SU(2)xU(1)', 'star': 2, 'control\_moved': 0, 'control': 244}, 'FiniteE7KAlgebra (zoo e7)': {'window': 169, 'dressed': 0, 'labels': 10, 'sections': 10, 'flavoured': 10, 'group': 'U(1)', 'star': 2, 'control\_moved': 0, 'control': 72}, 'BPS chart A2': {'window': 5, 'dressed': 0, 'labels': 5, 'sections': 5, 'flavoured': 5, 'group': '---', 'star': 1, 'control\_moved': 0, 'control': 0}, 'BPS chart A3': {'window': 9, 'dressed': 0, 'labels': 6, 'sections': 3, 'flavoured': 6, 'group': 'U(1)', 'star': 2, 'control\_moved': 0, 'control': 1}, 'BPS chart A4': {'window': 14, 'dressed': 0, 'labels': 7, 'sections': 7, 'flavoured': 7, 'group': '---', 'star': 1, 'control\_moved': 0, 'control': 4}, 'BPS chart A5': {'window': 20, 'dressed': 0, 'labels': 8, 'sections': 8, 'flavoured': 8, 'group': 'U(1)', 'star': 2, 'control\_moved': 0, 'control': 10}, 'BPS chart A6': {'window': 27, 'dressed': 0, 'labels': 9, 'sections': 9, 'flavoured': 9, 'group': '---', 'star': 1, 'control\_moved': 0, 'control': 20}, 'BPS chart A7': {'window': 35, 'dressed': 0, 'labels': 10, 'sections': 10, 'flavoured': 10, 'group': 'U(1)', 'star': 2, 'control\_moved': 0, 'control': 35}, 'BPS chart A9': {'window': 54, 'dressed': 0, 'labels': 12, 'sections': 12, 'flavoured': 12, 'group': 'U(1)', 'star': 2, 'control\_moved': 0, 'control': 40}, 'BPS chart D4': {'window': 14, 'dressed': 0, 'labels': 4, 'sections': 4, 'flavoured': 4, 'group': 'U(1)\^{}2', 'star': 2, 'control\_moved': 0, 'control': 4}, 'BPS chart D5': {'window': 20, 'dressed': 0, 'labels': 10, 'sections': 5, 'flavoured': 10, 'group': 'U(1)', 'star': 2, 'control\_moved': 0, 'control': 10}, 'BPS chart D6': {'window': 27, 'dressed': 0, 'labels': 6, 'sections': 6, 'flavoured': 6, 'group': 'U(1)\^{}2', 'star': 2, 'control\_moved': 0, 'control': 20}, 'BPS chart D7': {'window': 35, 'dressed': 0, 'labels': 14, 'sections': 7, 'flavoured': 14, 'group': 'U(1)', 'star': 2, 'control\_moved': 0, 'control': 35}, 'BPS chart D8': {'window': 44, 'dressed': 0, 'labels': 8, 'sections': 8, 'flavoured': 8, 'group': 'U(1)\^{}2', 'star': 2, 'control\_moved': 0, 'control': 40}, 'BPS chart E6': {'window': 27, 'dressed': 0, 'labels': 14, 'sections': 14, 'flavoured': 14, 'group': '---', 'star': 1, 'control\_moved': 0, 'control': 20}, 'BPS chart E7': {'window': 35, 'dressed': 0, 'labels': 10, 'sections': 10, 'flavoured': 10, 'group': 'U(1)', 'star': 2, 'control\_moved': 0, 'control': 35}, 'BPS chart E8': {'window': 44, 'dressed': 0, 'labels': 16, 'sections': 16, 'flavoured': 16, 'group': '---', 'star': 1, 'control\_moved': 0, 'control': 40}}; table entries: {'A2': [['[A1,A\_2n]', '---', 5]], 'A3': [['[A1,D\_2n+1]', 'SU(2)', 3], ['[A1,A\_3]', 'SU(2)', 3]], 'A4': [['[A1,A\_2n]', '---', 7]], 'A5': [['[A1,A\_2n+3]', 'U(1)', 8]], 'A6': [['[A1,A\_2n]', '---', 9]], 'A7': [['[A1,A\_2n+3]', 'U(1)', 10]], 'A9': [['[A1,A\_2n+3]', 'U(1)', 12]], 'D4': [['[A1,D\_2n+2]', 'SU(2)xU(1)', 4], ['[A1,D\_4]', 'SU(3)', 4]], 'D5': [['[A1,D\_2n+1]', 'SU(2)', 5]], 'D6': [['[A1,D\_2n+2]', 'SU(2)xU(1)', 6]], 'D7': [['[A1,D\_2n+1]', 'SU(2)', 7]], 'D8': [['[A1,D\_2n+2]', 'SU(2)xU(1)', 8]], 'E6': [['[A1,E\_6]', '---', 14]], 'E7': [['[A1,E\_7]', 'U(1)', 10]], 'E8': [['[A1,E\_8]', '---', 16]]}; n: 1..3; seconds: 9.3. \emph{Run (web (extensive), 2026-09-27):} 9 checks, 0 failing; presentations: {'A1A2kKAlg(1)': {'window': 16, 'dressed': 0, 'labels': 5, 'sections': 5, 'flavoured': 5, 'group': '---', 'star': 1, 'control\_moved': 0, 'control': 13}, 'A1A2kKAlg(2)': {'window': 56, 'dressed': 0, 'labels': 7, 'sections': 7, 'flavoured': 7, 'group': '---', 'star': 1, 'control\_moved': 0, 'control': 27}, 'A1A2kKAlg(3)': {'window': 72, 'dressed': 0, 'labels': 9, 'sections': 9, 'flavoured': 9, 'group': '---', 'star': 1, 'control\_moved': 0, 'control': 40}, 'PentagonKAlg': {'window': 16, 'dressed': 0, 'labels': 5, 'sections': 5, 'flavoured': 5, 'group': '---', 'star': 1, 'control\_moved': 0, 'control': 13}, 'HeptagonKAlg': {'window': 56, 'dressed': 0, 'labels': 7, 'sections': 7, 'flavoured': 7, 'group': '---', 'star': 1, 'control\_moved': 0, 'control': 26}, 'ungauge\_u1a1aodd(1)': {'window': 39, 'dressed': 7, 'labels': 6, 'sections': 3, 'flavoured': 6, 'group': 'U(1)', 'star': 2, 'control\_moved': 0, 'control': 22}, 'ungauge\_u1a1aodd(2)': {'window': 96, 'dressed': 17, 'labels': 8, 'sections': 8, 'flavoured': 8, 'group': 'U(1)', 'star': 2, 'control\_moved': 0, 'control': 53}, 'ungauge\_u1a1aodd(3)': {'window': 134, 'dressed': 23, 'labels': 10, 'sections': 10, 'flavoured': 10, 'group': 'U(1)', 'star': 2, 'control\_moved': 0, 'control': 47}, 'HexagonKAlg': {'window': 39, 'dressed': 7, 'labels': 6, 'sections': 3, 'flavoured': 6, 'group': 'U(1)', 'star': 2, 'control\_moved': 0, 'control': 22}, 'OctagonKAlg': {'window': 96, 'dressed': 17, 'labels': 8, 'sections': 8, 'flavoured': 8, 'group': 'U(1)', 'star': 2, 'control\_moved': 0, 'control': 53}, 'DecagonKAlg': {'window': 134, 'dressed': 23, 'labels': 10, 'sections': 10, 'flavoured': 10, 'group': 'U(1)', 'star': 2, 'control\_moved': 0, 'control': 47}, 'DodecagonKAlg': {'window': 183, 'dressed': 0, 'labels': 12, 'sections': 12, 'flavoured': 12, 'group': 'U(1)', 'star': 2, 'control\_moved': 0, 'control': 47}, 'A1DnKAlg(3)': {'window': 23, 'dressed': 0, 'labels': 3, 'sections': 3, 'flavoured': 3, 'group': 'SU(2)', 'star': 1, 'control\_moved': 0, 'control': 17}, 'A1DnKAlg(5)': {'window': 65, 'dressed': 0, 'labels': 5, 'sections': 5, 'flavoured': 5, 'group': 'SU(2)', 'star': 1, 'control\_moved': 0, 'control': 55}, 'A1DnKAlg(7)': {'window': 88, 'dressed': 0, 'labels': 7, 'sections': 7, 'flavoured': 7, 'group': 'SU(2)', 'star': 1, 'control\_moved': 0, 'control': 52}, 'A1DoddConeKAlg(0)': {'window': 23, 'dressed': 0, 'labels': 3, 'sections': 3, 'flavoured': 3, 'group': 'SU(2)', 'star': 1, 'control\_moved': 0, 'control': 20}, 'A1DoddConeKAlg(1)': {'window': 72, 'dressed': 0, 'labels': 5, 'sections': 5, 'flavoured': 5, 'group': 'SU(2)', 'star': 1, 'control\_moved': 0, 'control': 52}, 'A1DoddConeKAlg(2)': {'window': 97, 'dressed': 0, 'labels': 7, 'sections': 7, 'flavoured': 7, 'group': 'SU(2)', 'star': 1, 'control\_moved': 0, 'control': 50}, 'A1DevenKAlg(1)': {'window': 64, 'dressed': 6, 'labels': 4, 'sections': 4, 'flavoured': 4, 'group': 'SU(2)xU(1)', 'star': 2, 'control\_moved': 0, 'control': 26}, 'A1DevenKAlg(2)': {'window': 98, 'dressed': 19, 'labels': 6, 'sections': 6, 'flavoured': 6, 'group': 'SU(2)xU(1)', 'star': 2, 'control\_moved': 0, 'control': 78}, 'A1DevenKAlg(3)': {'window': 175, 'dressed': 11, 'labels': 8, 'sections': 8, 'flavoured': 8, 'group': 'SU(2)xU(1)', 'star': 2, 'control\_moved': 0, 'control': 67}, 'SU3ADKAlg': {'window': 52, 'dressed': 8, 'labels': 4, 'sections': 4, 'flavoured': 4, 'group': 'SU(3)', 'star': 2, 'control\_moved': 0, 'control': 33}, 'zoo pentagon': {'window': 16, 'dressed': 0, 'labels': 5, 'sections': 5, 'flavoured': 5, 'group': '---', 'star': 1, 'control\_moved': 0, 'control': 13}, 'zoo heptagon': {'window': 56, 'dressed': 0, 'labels': 7, 'sections': 7, 'flavoured': 7, 'group': '---', 'star': 1, 'control\_moved': 0, 'control': 28}, 'zoo a3': {'window': 19, 'dressed': 0, 'labels': 3, 'sections': 3, 'flavoured': 6, 'group': 'U(1)', 'star': 2, 'control\_moved': 0, 'control': 14}, 'zoo a5': {'window': 81, 'dressed': 0, 'labels': 8, 'sections': 8, 'flavoured': 8, 'group': 'U(1)', 'star': 2, 'control\_moved': 0, 'control': 48}, 'zoo a7': {'window': 134, 'dressed': 0, 'labels': 10, 'sections': 10, 'flavoured': 10, 'group': 'U(1)', 'star': 2, 'control\_moved': 0, 'control': 48}, 'zoo a1d3': {'window': 19, 'dressed': 0, 'labels': 3, 'sections': 3, 'flavoured': 3, 'group': 'SU(2)', 'star': 1, 'control\_moved': 0, 'control': 15}, 'zoo a1d4': {'window': 29, 'dressed': 0, 'labels': 4, 'sections': 4, 'flavoured': 4, 'group': 'SU(2)xU(1)', 'star': 2, 'control\_moved': 0, 'control': 18}, 'zoo a1d5': {'window': 69, 'dressed': 0, 'labels': 5, 'sections': 5, 'flavoured': 5, 'group': 'SU(2)', 'star': 1, 'control\_moved': 0, 'control': 43}, 'zoo a1d7': {'window': 99, 'dressed': 0, 'labels': 7, 'sections': 7, 'flavoured': 7, 'group': 'SU(2)', 'star': 1, 'control\_moved': 0, 'control': 55}, 'FiniteE6KAlgebra (zoo e6)': {'window': 101, 'dressed': 0, 'labels': 14, 'sections': 14, 'flavoured': 14, 'group': '---', 'star': 1, 'control\_moved': 0, 'control': 56}, 'FiniteE8KAlgebra (zoo e8)': {'window': 269, 'dressed': 0, 'labels': 16, 'sections': 16, 'flavoured': 16, 'group': '---', 'star': 1, 'control\_moved': 0, 'control': 176}, 'zoo a1d6': {'window': 105, 'dressed': 0, 'labels': 6, 'sections': 6, 'flavoured': 6, 'group': 'SU(2)xU(1)', 'star': 2, 'control\_moved': 0, 'control': 52}, 'zoo a1d8': {'window': 260, 'dressed': 0, 'labels': 8, 'sections': 8, 'flavoured': 8, 'group': 'SU(2)xU(1)', 'star': 2, 'control\_moved': 0, 'control': 244}, 'FiniteE7KAlgebra (zoo e7)': {'window': 169, 'dressed': 0, 'labels': 10, 'sections': 10, 'flavoured': 10, 'group': 'U(1)', 'star': 2, 'control\_moved': 0, 'control': 72}, 'A1A2kKAlg(4)': {'window': 91, 'dressed': 0, 'labels': 11, 'sections': 11, 'flavoured': 11, 'group': '---', 'star': 1, 'control\_moved': 0, 'control': 39}, 'A1A2kKAlg(5)': {'window': 112, 'dressed': 0, 'labels': 13, 'sections': 13, 'flavoured': 13, 'group': '---', 'star': 1, 'control\_moved': 0, 'control': 40}, 'ungauge\_u1a1aodd(5)': {'window': 319, 'dressed': 0, 'labels': 14, 'sections': 14, 'flavoured': 14, 'group': 'U(1)', 'star': 2, 'control\_moved': 0, 'control': 44}, 'A1DnKAlg(9)': {'window': 125, 'dressed': 0, 'labels': 9, 'sections': 9, 'flavoured': 9, 'group': 'SU(2)', 'star': 1, 'control\_moved': 0, 'control': 66}, 'A1DoddConeKAlg(3)': {'window': 128, 'dressed': 0, 'labels': 9, 'sections': 9, 'flavoured': 9, 'group': 'SU(2)', 'star': 1, 'control\_moved': 0, 'control': 44}, 'A1DoddConeKAlg(4)': {'window': 166, 'dressed': 0, 'labels': 11, 'sections': 11, 'flavoured': 11, 'group': 'SU(2)', 'star': 1, 'control\_moved': 0, 'control': 46}, 'A1DevenKAlg(4)': {'window': 344, 'dressed': 12, 'labels': 10, 'sections': 10, 'flavoured': 10, 'group': 'SU(2)xU(1)', 'star': 2, 'control\_moved': 0, 'control': 79}, 'A1DevenKAlg(5)': {'window': 641, 'dressed': 0, 'labels': 12, 'sections': 12, 'flavoured': 12, 'group': 'SU(2)xU(1)', 'star': 2, 'control\_moved': 0, 'control': 61}, 'BPS chart A2': {'window': 5, 'dressed': 0, 'labels': 5, 'sections': 5, 'flavoured': 5, 'group': '---', 'star': 1, 'control\_moved': 0, 'control': 0}, 'BPS chart A3': {'window': 9, 'dressed': 0, 'labels': 6, 'sections': 3, 'flavoured': 6, 'group': 'U(1)', 'star': 2, 'control\_moved': 0, 'control': 1}, 'BPS chart A4': {'window': 14, 'dressed': 0, 'labels': 7, 'sections': 7, 'flavoured': 7, 'group': '---', 'star': 1, 'control\_moved': 0, 'control': 4}, 'BPS chart A5': {'window': 20, 'dressed': 0, 'labels': 8, 'sections': 8, 'flavoured': 8, 'group': 'U(1)', 'star': 2, 'control\_moved': 0, 'control': 10}, 'BPS chart A6': {'window': 27, 'dressed': 0, 'labels': 9, 'sections': 9, 'flavoured': 9, 'group': '---', 'star': 1, 'control\_moved': 0, 'control': 20}, 'BPS chart A7': {'window': 35, 'dressed': 0, 'labels': 10, 'sections': 10, 'flavoured': 10, 'group': 'U(1)', 'star': 2, 'control\_moved': 0, 'control': 35}, 'BPS chart A8': {'window': 44, 'dressed': 0, 'labels': 11, 'sections': 11, 'flavoured': 11, 'group': '---', 'star': 1, 'control\_moved': 0, 'control': 40}, 'BPS chart A9': {'window': 54, 'dressed': 0, 'labels': 12, 'sections': 12, 'flavoured': 12, 'group': 'U(1)', 'star': 2, 'control\_moved': 0, 'control': 40}, 'BPS chart A10': {'window': 65, 'dressed': 0, 'labels': 13, 'sections': 13, 'flavoured': 13, 'group': '---', 'star': 1, 'control\_moved': 0, 'control': 40}, 'BPS chart A11': {'window': 77, 'dressed': 0, 'labels': 14, 'sections': 14, 'flavoured': 14, 'group': 'U(1)', 'star': 2, 'control\_moved': 0, 'control': 40}, 'BPS chart A13': {'window': 104, 'dressed': 0, 'labels': 16, 'sections': 16, 'flavoured': 16, 'group': 'U(1)', 'star': 2, 'control\_moved': 0, 'control': 40}, 'BPS chart D4': {'window': 14, 'dressed': 0, 'labels': 4, 'sections': 4, 'flavoured': 4, 'group': 'U(1)\^{}2', 'star': 2, 'control\_moved': 0, 'control': 4}, 'BPS chart D5': {'window': 20, 'dressed': 0, 'labels': 10, 'sections': 5, 'flavoured': 10, 'group': 'U(1)', 'star': 2, 'control\_moved': 0, 'control': 10}, 'BPS chart D6': {'window': 27, 'dressed': 0, 'labels': 6, 'sections': 6, 'flavoured': 6, 'group': 'U(1)\^{}2', 'star': 2, 'control\_moved': 0, 'control': 20}, 'BPS chart D7': {'window': 35, 'dressed': 0, 'labels': 14, 'sections': 7, 'flavoured': 14, 'group': 'U(1)', 'star': 2, 'control\_moved': 0, 'control': 35}, 'BPS chart D8': {'window': 44, 'dressed': 0, 'labels': 8, 'sections': 8, 'flavoured': 8, 'group': 'U(1)\^{}2', 'star': 2, 'control\_moved': 0, 'control': 40}, 'BPS chart D9': {'window': 54, 'dressed': 0, 'labels': 18, 'sections': 9, 'flavoured': 18, 'group': 'U(1)', 'star': 2, 'control\_moved': 0, 'control': 40}, 'BPS chart D10': {'window': 65, 'dressed': 0, 'labels': 10, 'sections': 10, 'flavoured': 10, 'group': 'U(1)\^{}2', 'star': 2, 'control\_moved': 0, 'control': 40}, 'BPS chart D11': {'window': 77, 'dressed': 0, 'labels': 22, 'sections': 11, 'flavoured': 22, 'group': 'U(1)', 'star': 2, 'control\_moved': 0, 'control': 40}, 'BPS chart D12': {'window': 90, 'dressed': 0, 'labels': 12, 'sections': 12, 'flavoured': 12, 'group': 'U(1)\^{}2', 'star': 2, 'control\_moved': 0, 'control': 40}, 'BPS chart E6': {'window': 27, 'dressed': 0, 'labels': 14, 'sections': 14, 'flavoured': 14, 'group': '---', 'star': 1, 'control\_moved': 0, 'control': 20}, 'BPS chart E7': {'window': 35, 'dressed': 0, 'labels': 10, 'sections': 10, 'flavoured': 10, 'group': 'U(1)', 'star': 2, 'control\_moved': 0, 'control': 35}, 'BPS chart E8': {'window': 44, 'dressed': 0, 'labels': 16, 'sections': 16, 'flavoured': 16, 'group': '---', 'star': 1, 'control\_moved': 0, 'control': 40}}; table entries: {'A2': [['[A1,A\_2n]', '---', 5]], 'A3': [['[A1,D\_2n+1]', 'SU(2)', 3], ['[A1,A\_3]', 'SU(2)', 3]], 'A4': [['[A1,A\_2n]', '---', 7]], 'A5': [['[A1,A\_2n+3]', 'U(1)', 8]], 'A6': [['[A1,A\_2n]', '---', 9]], 'A7': [['[A1,A\_2n+3]', 'U(1)', 10]], 'A8': [['[A1,A\_2n]', '---', 11]], 'A9': [['[A1,A\_2n+3]', 'U(1)', 12]], 'A10': [['[A1,A\_2n]', '---', 13]], 'A11': [['[A1,A\_2n+3]', 'U(1)', 14]], 'A13': [['[A1,A\_2n+3]', 'U(1)', 16]], 'D4': [['[A1,D\_2n+2]', 'SU(2)xU(1)', 4], ['[A1,D\_4]', 'SU(3)', 4]], 'D5': [['[A1,D\_2n+1]', 'SU(2)', 5]], 'D6': [['[A1,D\_2n+2]', 'SU(2)xU(1)', 6]], 'D7': [['[A1,D\_2n+1]', 'SU(2)', 7]], 'D8': [['[A1,D\_2n+2]', 'SU(2)xU(1)', 8]], 'D9': [['[A1,D\_2n+1]', 'SU(2)', 9]], 'D10': [['[A1,D\_2n+2]', 'SU(2)xU(1)', 10]], 'D11': [['[A1,D\_2n+1]', 'SU(2)', 11]], 'D12': [['[A1,D\_2n+2]', 'SU(2)xU(1)', 12]], 'E6': [['[A1,E\_6]', '---', 14]], 'E7': [['[A1,E\_7]', 'U(1)', 10]], 'E8': [['[A1,E\_8]', '---', 16]]}; n: 1..5; seconds: 17.9. & \MEASURED\newline {\footnotesize identity, extensive\newline web}\newline {\footnotesize run web 2026-09-27: 9/9}\newline {\footnotesize run web (extensive) 2026-09-27: 9/9} \\ \addlinespace
\code{app:finite/canonical-basis}\newline App.~\pref{app:finite} & The canonical basis consists of the normalised monomials in $\fq$-commuting generators, all distinct; the algebra is presented by the remaining products. & (i) On the entries whose $\fq$-commuting sets are at most the rank of $B$, those sets in $\Gamma/\ker B$ (rays equal there identified): every maximal cone unimodular, no two small monomials sharing a charge, and every point of a box exactly one monomial (rank $\le 6$; at extensive depth also every wall in exactly two maximal cones, on opposite sides). On the entries whose $\fq$-commuting sets are larger, the charges shared by distinct monomials, a sample confirmed as one element in the zoo and in the twin. (ii) In the twin, products of rays of one cone --- pairs, squares, a triple --- are the single canonical element at the sum with the $\fq$-power the zoo's normalisation predicts, and the zoo's exponents are antisymmetric. (iii) The zoo's products of generators in no common cone equal the twin's: exactly on the unflavoured entries, with the flavour forgotten (augmentation, labels modulo $\ker B$) on the flavoured ones. \emph{Population:} every distinct zoo entry ($14$ classes; hexagon, octagon, decagon are $a_3$, $a_5$, $a_7$), each against its BPS twin (\code{BPSKAlgebra} on the entry's embedded quiver, default spec). \emph{Controls:} negative: the heptagon's fan with one maximal cone removed leaves box points without a monomial; the pentagon's product against the twin's in the other order differs. \emph{Evidence:} emergent: the twin's canonical basis comes from its quiver and spec, not from the zoo's tables. \emph{Where:} \code{battery/checks_finite.py (check_finite_canonical_basis)}, \code{tests/test_cone_kalgebra.py}, \code{tests/test_finite_kalgebras.py}. \emph{Note:} On decagon ([A1,A7] over U(1)) and a1d8 ([A1,D8] over SU(2)xU(1)) the zoo's q-commuting sets are larger than the rank of B, and distinct normalised monomials are one canonical element: L\_46 L\_57 = L\_47 L\_56 = q\^{}-1 L\_$\gamma$ in decagon, the same in its independent BPS twin. Every canonical element is still a normalised monomial, but not uniquely; the author expects that some monomials may coincide: two normalised monomials are one element exactly when their charges agree, which needs a q-commuting set larger than the rank of B, and the two entries are recorded apart from the twelve. The twin checks (ii) and (iii) run under separate time budgets per entry (60 s fast, 900 s extensive), each with its own fresh BPS twin; on the slowest twins a clause can stop before the end of its sample (a1d7 and a1d8, and at extensive depth decagon's (ii), by memory), and the record gives what it compared. \emph{Run (web, 2026-09-28):} 5 checks, 0 failing; fan: {'pentagon': {'rays': 5, 'rank mod ker B': 2, 'largest q-commuting set': 2, 'maximal cones': 5, 'distinct rays mod ker B': 5, 'lattice ok': True, 'unimodular': 5, 'box points': 25, 'exactly one monomial': 25, 'charges shared by two monomials of degree $\le$ 2': 0, 'seconds (ii)+(iii)': 0.2}, 'hexagon': {'rays': 6, 'rank mod ker B': 2, 'largest q-commuting set': 2, 'maximal cones': 6, 'distinct rays mod ker B': 6, 'lattice ok': True, 'unimodular': 6, 'box points': 25, 'exactly one monomial': 25, 'charges shared by two monomials of degree $\le$ 2': 0, 'seconds (ii)+(iii)': 0.2}, 'heptagon': {'rays': 14, 'rank mod ker B': 4, 'largest q-commuting set': 4, 'maximal cones': 42, 'distinct rays mod ker B': 14, 'lattice ok': True, 'unimodular': 42, 'box points': 625, 'exactly one monomial': 625, 'charges shared by two monomials of degree $\le$ 2': 0, 'seconds (ii)+(iii)': 0.3}, 'octagon': {'rays': 24, 'rank mod ker B': 4, 'largest q-commuting set': 4, 'maximal cones': 88, 'distinct rays mod ker B': 24, 'lattice ok': True, 'unimodular': 88, 'box points': 625, 'exactly one monomial': 625, 'charges shared by two monomials of degree $\le$ 2': 0, 'seconds (ii)+(iii)': 0.4}, 'decagon': {'rays': 65, 'rank mod ker B': 6, 'largest q-commuting set': 10, 'maximal cones': 700, 'distinct rays mod ker B': 65, 'lattice ok': True, 'unimodular': 700, 'charges shared by two monomials of degree $\le$ 2': 80, 'seconds (ii)+(iii)': 10.5}, 'a1d3': {'rays': 6, 'rank mod ker B': 2, 'largest q-commuting set': 2, 'maximal cones': 6, 'distinct rays mod ker B': 6, 'lattice ok': True, 'unimodular': 6, 'box points': 25, 'exactly one monomial': 25, 'charges shared by two monomials of degree $\le$ 2': 0, 'seconds (ii)+(iii)': 0.2}, 'a1d4': {'rays': 8, 'rank mod ker B': 2, 'largest q-commuting set': 2, 'maximal cones': 8, 'distinct rays mod ker B': 8, 'lattice ok': True, 'unimodular': 8, 'box points': 25, 'exactly one monomial': 25, 'charges shared by two monomials of degree $\le$ 2': 0, 'seconds (ii)+(iii)': 0.8}, 'a1d5': {'rays': 20, 'rank mod ker B': 4, 'largest q-commuting set': 4, 'maximal cones': 70, 'distinct rays mod ker B': 20, 'lattice ok': True, 'unimodular': 70, 'box points': 625, 'exactly one monomial': 625, 'charges shared by two monomials of degree $\le$ 2': 0, 'seconds (ii)+(iii)': 3.6}, 'a1d6': {'rays': 39, 'rank mod ker B': 4, 'largest q-commuting set': 4, 'maximal cones': 162, 'distinct rays mod ker B': 39, 'lattice ok': True, 'unimodular': 162, 'box points': 625, 'exactly one monomial': 625, 'charges shared by two monomials of degree $\le$ 2': 0, 'seconds (ii)+(iii)': 1.7}, 'a1d7': {'rays': 42, 'rank mod ker B': 6, 'largest q-commuting set': 6, 'maximal cones': 924, 'distinct rays mod ker B': 42, 'lattice ok': True, 'unimodular': 924, 'box points': 729, 'exactly one monomial': 729, 'charges shared by two monomials of degree $\le$ 2': 0, 'seconds (ii)+(iii)': 120.2}, 'a1d8': {'rays': 120, 'rank mod ker B': 6, 'largest q-commuting set': 10, 'maximal cones': 1736, 'distinct rays mod ker B': 120, 'lattice ok': True, 'unimodular': 1664, 'charges shared by two monomials of degree $\le$ 2': 196, 'seconds (ii)+(iii)': 120.2}, 'e6': {'rays': 42, 'rank mod ker B': 6, 'largest q-commuting set': 6, 'maximal cones': 833, 'distinct rays mod ker B': 42, 'lattice ok': True, 'unimodular': 833, 'box points': 729, 'exactly one monomial': 729, 'charges shared by two monomials of degree $\le$ 2': 0, 'seconds (ii)+(iii)': 0.4}, 'e7': {'rays': 90, 'rank mod ker B': 6, 'largest q-commuting set': 6, 'maximal cones': 2600, 'distinct rays mod ker B': 90, 'lattice ok': True, 'unimodular': 2600, 'charges shared by two monomials of degree $\le$ 2': 0, 'seconds (ii)+(iii)': 0.4}, 'e8': {'rays': 128, 'rank mod ker B': 8, 'largest q-commuting set': 8, 'maximal cones': 25080, 'distinct rays mod ker B': 128, 'lattice ok': True, 'unimodular': 25080, 'charges shared by two monomials of degree $\le$ 2': 0, 'seconds (ii)+(iii)': 0.8}}; normalisation (ii): {'pentagon': '12/12', 'hexagon': '12/12', 'heptagon': '44/44', 'octagon': '44/44', 'decagon': '52/52', 'a1d3': '12/12', 'a1d4': '12/12', 'a1d5': '44/44', 'a1d6': '44/44', 'a1d7': '7/7 before the stop (not reached in 60 s)', 'a1d8': '10/10 before the stop (not reached in 60 s)', 'e6': '52/52', 'e7': '52/52', 'e8': '52/52'}; presentation (iii): {'pentagon': '10/10 (exact)', 'hexagon': '18/18 (flavour forgotten)', 'heptagon': '24/24 (exact)', 'octagon': '24/24 (flavour forgotten)', 'decagon': '24/24 (flavour forgotten)', 'a1d3': '18/18 (flavour forgotten)', 'a1d4': '24/24 (flavour forgotten)', 'a1d5': '24/24 (flavour forgotten)', 'a1d6': '24/24 (flavour forgotten)', 'a1d7': '4/4 before the stop (not reached in 60 s) (flavour forgotten)', 'a1d8': '19/19 before the stop (not reached in 60 s) (flavour forgotten)', 'e6': '24/24 (exact)', 'e7': '24/24 (flavour forgotten)', 'e8': '8/8 (exact)'}; coincident monomials: {'decagon': {'charges shared': 80, 'sample': [[[[4, 1], [44, 1]], [[5, 1], [42, 1]]], [[[2, 1], [57, 1]], [[19, 1], [42, 1]]], [[[5, 1], [62, 1]], [[26, 1], [44, 1]]]], 'one element in the zoo': '3/3', 'one element in the twin': '3/3'}, 'a1d8': {'charges shared': 196, 'sample': [[[[4, 1], [45, 1]], [[5, 1], [42, 1]]], [[[2, 1], [84, 1]], [[19, 1], [42, 1]]], [[[5, 1], [96, 1]], [[26, 1], [45, 1]]]], 'one element in the zoo': '3/3', 'one element in the twin': '3/3'}}; seconds: 301.0. \emph{Run (web (extensive), 2026-09-28):} 5 checks, 0 failing; fan: {'pentagon': {'rays': 5, 'rank mod ker B': 2, 'largest q-commuting set': 2, 'maximal cones': 5, 'distinct rays mod ker B': 5, 'lattice ok': True, 'unimodular': 5, 'box points': 25, 'exactly one monomial': 25, 'walls': [5, 5, 5], 'charges shared by two monomials of degree $\le$ 2': 0, 'seconds (ii)+(iii)': 0.1}, 'hexagon': {'rays': 6, 'rank mod ker B': 2, 'largest q-commuting set': 2, 'maximal cones': 6, 'distinct rays mod ker B': 6, 'lattice ok': True, 'unimodular': 6, 'box points': 25, 'exactly one monomial': 25, 'walls': [6, 6, 6], 'charges shared by two monomials of degree $\le$ 2': 0, 'seconds (ii)+(iii)': 0.2}, 'heptagon': {'rays': 14, 'rank mod ker B': 4, 'largest q-commuting set': 4, 'maximal cones': 42, 'distinct rays mod ker B': 14, 'lattice ok': True, 'unimodular': 42, 'box points': 625, 'exactly one monomial': 625, 'walls': [84, 84, 84], 'charges shared by two monomials of degree $\le$ 2': 0, 'seconds (ii)+(iii)': 0.3}, 'octagon': {'rays': 24, 'rank mod ker B': 4, 'largest q-commuting set': 4, 'maximal cones': 88, 'distinct rays mod ker B': 24, 'lattice ok': True, 'unimodular': 88, 'box points': 625, 'exactly one monomial': 625, 'walls': [176, 176, 176], 'charges shared by two monomials of degree $\le$ 2': 0, 'seconds (ii)+(iii)': 0.5}, 'decagon': {'rays': 65, 'rank mod ker B': 6, 'largest q-commuting set': 10, 'maximal cones': 700, 'distinct rays mod ker B': 65, 'lattice ok': True, 'unimodular': 700, 'charges shared by two monomials of degree $\le$ 2': 80, 'seconds (ii)+(iii)': 516.2}, 'a1d3': {'rays': 6, 'rank mod ker B': 2, 'largest q-commuting set': 2, 'maximal cones': 6, 'distinct rays mod ker B': 6, 'lattice ok': True, 'unimodular': 6, 'box points': 25, 'exactly one monomial': 25, 'walls': [6, 6, 6], 'charges shared by two monomials of degree $\le$ 2': 0, 'seconds (ii)+(iii)': 0.3}, 'a1d4': {'rays': 8, 'rank mod ker B': 2, 'largest q-commuting set': 2, 'maximal cones': 8, 'distinct rays mod ker B': 8, 'lattice ok': True, 'unimodular': 8, 'box points': 25, 'exactly one monomial': 25, 'walls': [8, 8, 8], 'charges shared by two monomials of degree $\le$ 2': 0, 'seconds (ii)+(iii)': 2.7}, 'a1d5': {'rays': 20, 'rank mod ker B': 4, 'largest q-commuting set': 4, 'maximal cones': 70, 'distinct rays mod ker B': 20, 'lattice ok': True, 'unimodular': 70, 'box points': 625, 'exactly one monomial': 625, 'walls': [140, 140, 140], 'charges shared by two monomials of degree $\le$ 2': 0, 'seconds (ii)+(iii)': 4.9}, 'a1d6': {'rays': 39, 'rank mod ker B': 4, 'largest q-commuting set': 4, 'maximal cones': 162, 'distinct rays mod ker B': 39, 'lattice ok': True, 'unimodular': 162, 'box points': 625, 'exactly one monomial': 625, 'walls': [324, 324, 324], 'charges shared by two monomials of degree $\le$ 2': 0, 'seconds (ii)+(iii)': 2.2}, 'a1d7': {'rays': 42, 'rank mod ker B': 6, 'largest q-commuting set': 6, 'maximal cones': 924, 'distinct rays mod ker B': 42, 'lattice ok': True, 'unimodular': 924, 'box points': 729, 'exactly one monomial': 729, 'walls': [2772, 2772, 2772], 'charges shared by two monomials of degree $\le$ 2': 0, 'seconds (ii)+(iii)': 1439.4}, 'a1d8': {'rays': 120, 'rank mod ker B': 6, 'largest q-commuting set': 10, 'maximal cones': 1736, 'distinct rays mod ker B': 120, 'lattice ok': True, 'unimodular': 1664, 'charges shared by two monomials of degree $\le$ 2': 196, 'seconds (ii)+(iii)': 1094.6}, 'e6': {'rays': 42, 'rank mod ker B': 6, 'largest q-commuting set': 6, 'maximal cones': 833, 'distinct rays mod ker B': 42, 'lattice ok': True, 'unimodular': 833, 'box points': 729, 'exactly one monomial': 729, 'walls': [2499, 2499, 2499], 'charges shared by two monomials of degree $\le$ 2': 0, 'seconds (ii)+(iii)': 0.8}, 'e7': {'rays': 90, 'rank mod ker B': 6, 'largest q-commuting set': 6, 'maximal cones': 2600, 'distinct rays mod ker B': 90, 'lattice ok': True, 'unimodular': 2600, 'box points': 729, 'exactly one monomial': 729, 'walls': [7800, 7800, 7800], 'charges shared by two monomials of degree $\le$ 2': 0, 'seconds (ii)+(iii)': 0.7}, 'e8': {'rays': 128, 'rank mod ker B': 8, 'largest q-commuting set': 8, 'maximal cones': 25080, 'distinct rays mod ker B': 128, 'lattice ok': True, 'unimodular': 25080, 'walls': [100320, 100320, 100320], 'charges shared by two monomials of degree $\le$ 2': 0, 'seconds (ii)+(iii)': 1.4}}; normalisation (ii): {'pentagon': '15/15', 'hexagon': '18/18', 'heptagon': '88/88', 'octagon': '88/88', 'decagon': '279/279 before the stop (not reached: out of memory)', 'a1d3': '18/18', 'a1d4': '24/24', 'a1d5': '88/88', 'a1d6': '88/88', 'a1d7': '8/8 before the stop (not reached in 900 s)', 'a1d8': '339/339 before the stop (not reached in 900 s)', 'e6': '176/176', 'e7': '176/176', 'e8': '200/200'}; presentation (iii): {'pentagon': '10/10 (exact)', 'hexagon': '18/18 (flavour forgotten)', 'heptagon': '60/60 (exact)', 'octagon': '60/60 (flavour forgotten)', 'decagon': '60/60 (flavour forgotten)', 'a1d3': '18/18 (flavour forgotten)', 'a1d4': '40/40 (flavour forgotten)', 'a1d5': '60/60 (flavour forgotten)', 'a1d6': '60/60 (flavour forgotten)', 'a1d7': '60/60 (flavour forgotten)', 'a1d8': '60/60 (flavour forgotten)', 'e6': '60/60 (exact)', 'e7': '60/60 (flavour forgotten)', 'e8': '24/24 (exact)'}; coincident monomials: {'decagon': {'charges shared': 80, 'sample': [[[[4, 1], [44, 1]], [[5, 1], [42, 1]]], [[[2, 1], [57, 1]], [[19, 1], [42, 1]]], [[[5, 1], [62, 1]], [[26, 1], [44, 1]]]], 'one element in the zoo': '3/3', 'one element in the twin': '3/3'}, 'a1d8': {'charges shared': 196, 'sample': [[[[4, 1], [45, 1]], [[5, 1], [42, 1]]], [[[2, 1], [84, 1]], [[19, 1], [42, 1]]], [[[5, 1], [96, 1]], [[26, 1], [45, 1]]]], 'one element in the zoo': '3/3', 'one element in the twin': '3/3'}}; seconds: 3119.7. & \CERTIFIED\newline {\footnotesize theorem, extensive\newline web}\newline {\footnotesize run web 2026-09-28: 5/5}\newline {\footnotesize run web (extensive) 2026-09-28: 5/5} \\ \addlinespace
\code{app:finite/trace-unique}\newline App.~\pref{app:finite} & The $I_{a,b}$ axioms determine $I_{a,b}$ uniquely up to $1+O(\fq)$ and are self-consistent on these examples. & \code{trace_uniqueness.verify_unique_up_to_rescale} on windows of the zoo's labels (a cone's degree-$\le2$ or $\le3$ monomials; or the generators, in-cone pairs and squares of the first cones): the canonical trace solves every equation (self-consistency), and the admissible traces, projected to the window and to orders $\le K$, are exactly its $1+O(\fq)$ rescales. \emph{Population:} the unflavoured zoo entries: pentagon (to $\fq^2$ and $\fq^4$), heptagon ($\fq^2$, $\fq^3$), $E_6$ ($\fq$); at extensive depth also $E_6$ to $\fq^2$ on two larger windows. \emph{Controls:} negative: on the pentagon, a window without squares leaves a third mode, so the window matters; dropping the $\fq^0$-orthonormality enlarges the solution space. \emph{Evidence:} emergent. \emph{Where:} \code{battery/checks_finite.py (check_finite_trace_unique)}, \code{notes/finite_type_kalgebras.md sec. 5}. \emph{Note:} Formalised in ClusterLean (the local Lean 4 repository) on the examples formalised so far: the quantum torus with nondegenerate pairing, the pentagon, SQED\_N with SU(N) flavour for every N, hence U\_q(sl\_2): the functionals satisfying K4 are exactly c$\cdot$Tr with c in 1+qR[[q]]  The measured population is the unflavoured entries: trace\_uniqueness takes trivial-R coefficients only (a TypeError on an R-valued coefficient, and the zoo's Z-form wrappers still carry a nontrivial ring), so the flavoured zoo entries are not covered.  E8 is out of reach of this solver: its dense kernel exceeded 12 GB at order 1 on a window of 8 cones and squares. \emph{Run (web, 2026-09-27):} 2 checks, 0 failing; windows: {'pentagon, degree $\le$ 2': {'window': 16, 'K\_eq': 2, 'canonical trace solves': True, 'projected dimension': 2, 'rescale line': True, 'seconds': 0.1}, 'pentagon, degree $\le$ 3': {'window': 31, 'K\_eq': 4, 'canonical trace solves': True, 'projected dimension': 4, 'rescale line': True, 'seconds': 0.7}, 'heptagon, 3 cones + squares': {'window': 25, 'K\_eq': 2, 'canonical trace solves': True, 'projected dimension': 2, 'rescale line': True, 'seconds': 0.7}, 'heptagon, 5 cones + squares': {'window': 31, 'K\_eq': 3, 'canonical trace solves': True, 'projected dimension': 3, 'rescale line': True, 'seconds': 1.7}, 'E6, 8 cones + squares': {'window': 59, 'K\_eq': 1, 'canonical trace solves': True, 'projected dimension': 1, 'rescale line': True, 'seconds': 19.5}}; seconds: 22.7. \emph{Run (web (extensive), 2026-09-28):} 2 checks, 0 failing; windows: {'pentagon, degree $\le$ 2': {'window': 16, 'K\_eq': 2, 'canonical trace solves': True, 'projected dimension': 2, 'rescale line': True, 'seconds': 0.1}, 'pentagon, degree $\le$ 3': {'window': 31, 'K\_eq': 4, 'canonical trace solves': True, 'projected dimension': 4, 'rescale line': True, 'seconds': 0.8}, 'heptagon, 3 cones + squares': {'window': 25, 'K\_eq': 2, 'canonical trace solves': True, 'projected dimension': 2, 'rescale line': True, 'seconds': 1.0}, 'heptagon, 5 cones + squares': {'window': 31, 'K\_eq': 3, 'canonical trace solves': True, 'projected dimension': 3, 'rescale line': True, 'seconds': 2.2}, 'E6, 8 cones + squares': {'window': 59, 'K\_eq': 1, 'canonical trace solves': True, 'projected dimension': 1, 'rescale line': True, 'seconds': 20.5}, 'E6, 16 cones + squares': {'window': 71, 'K\_eq': 2, 'canonical trace solves': True, 'projected dimension': 2, 'rescale line': True, 'seconds': 34.6}, 'E6, 24 cones + squares': {'window': 87, 'K\_eq': 2, 'canonical trace solves': True, 'projected dimension': 2, 'rescale line': True, 'seconds': 46.8}}; seconds: 106.3. & \CERTIFIED\newline {\footnotesize theorem, extensive\newline web}\newline {\footnotesize run web 2026-09-27: 2/2}\newline {\footnotesize run web (extensive) 2026-09-28: 2/2} \\ \addlinespace
\code{app:finite/class-s-coincidences}\newline App.~\pref{app:finite} & $[A_1,E_6]\simeq[A_2,A_3]$ and $[A_1,E_8]\simeq[A_2,A_4]$. & A \code{KAlgebraIso} from the $A_2\times A_k$ chart (\code{a2ak_induction.member}) to the zoo's BPS twin of $[A_1,E_6]$ or $[A_1,E_8]$: a mutation path to the twin's quiver (a seeded random walk to an acyclic quiver, then sink and source mutations to its orientation), labels transported by the forward-necklace rule $\mu_g$ along it and then by the lattice map onto the twin's node charges; certified on products, $\rho$ and traces over a window. \emph{Population:} $[A_2,A_3]\to[A_1,E_6]$; $[A_2,A_4]\to[A_1,E_8]$ at extensive depth. \emph{Controls:} negative: the same transport with the path's last mutation step omitted fails the product check. \emph{Evidence:} emergent. \emph{Where:} \code{battery/checks_finite.py (check_finite_class_s)}, \code{src/bps/a2ak_induction.py}, \code{tests/test_a2ak_induction.py}. \emph{Note:} The quiver-level mutation equivalence is in a2ak\_induction.py, which also corrects the catalogue's attribution of [A2,A4] to E7. BPSAtlas cannot take these paths: in its spec mode a forward mutation needs the node at the spec's head, and it refuses the first step even with 16 local moves; the transport is composed from the chart graph's own rule instead, and the battery is the certificate. At extensive depth the E8 battery runs under an 1800 s budget, rho and traces first and then the products from the smallest transported charges; the record gives what it compared. \emph{Run (web, 2026-09-28):} 2 checks, 0 failing; isomorphisms: {'[A2,A3] -> E6 (window 15x4, traces to q\^{}3)': {'mutations': 27, 'permutation': [1, 3, 4, 0, 5, 2], 'products': '60/60', 'rho': '15/15', 'traces': '4/4', 'seconds': 49.4}}; seconds: 95.7. \emph{Run (web (extensive), 2026-09-28):} 4 checks, 0 failing; isomorphisms: {'[A2,A3] -> E6 (window 15x4, traces to q\^{}3)': {'mutations': 27, 'permutation': [1, 3, 4, 0, 5, 2], 'products': '60/60', 'rho': '15/15', 'traces': '4/4', 'seconds': 46.6}, '[A2,A3] -> E6 (window 15x8, traces to q\^{}4)': {'mutations': 27, 'permutation': [1, 3, 4, 0, 5, 2], 'products': '120/120', 'rho': '15/15', 'traces': '7/7', 'seconds': 356.5}, '[A2,A4] -> E8 (window 12x4, traces to q\^{}2)': {'mutations': 33, 'permutation': [3, 4, 2, 1, 7, 5, 0, 6], 'stopped': 'not reached in 1800 s', 'products': '37/37 of 48', 'rho': '24/24 of 24', 'traces': '3/3 of 3', 'seconds': 1804.6}}; seconds: 2256.6. & \MEASURED\newline {\footnotesize theorem, extensive\newline web}\newline {\footnotesize run web 2026-09-28: 2/2}\newline {\footnotesize run web (extensive) 2026-09-28: 4/4} \\ \addlinespace
\code{app:a1a2k/rules}\newline App.~\pref{app:a1a2k} & The chord rule for $c_{ab}\in\{0,\pm1\}$, quantum Ptolemy with $(\alpha,\beta)\in\{(1,0),(0,-1)\}$ decided by the odd arcs, and the canonical basis of non-crossing monomials. & The draft's $c_{ab}$ rule and quantum Ptolemy rule, transcribed from its text, against \code{A1A2kKAlg}$(k)$ on every ordered pair of letters, $k\le6$; every term of every letter product a non-crossing multiset; the draft's $m(a)$ against the class's; and the class's letter products against the BPS chart of the $A_{2k}$ quiver through the additive charge map (\code{a1a2k_induction.complete_block_iso}) on every ordered pair, $k\le3$ ($k\le4$ extensive). \emph{Population:} $k\le 6$, all chord pairs. \emph{Controls:} negative: the rules read with the opposite orientation fail on every crossing pair. \emph{Evidence:} emergent. \emph{Where:} \code{battery/checks_finite.py (check_a1a2k_rules)}, \code{tests/test_a1a2k_cone_data.py}, \code{tests/test_a1a2k_kalg.py}, \code{src/abe/a1a2k_induction.py}. \emph{Note:} The rules hold with the text's 'CCW' read as increasing vertex index. The draft's heptagon figure (fig:heptagon-orbits) places vertex i at angle 90 - 360 i/7 degrees, so its indices increase clockwise: in that figure the rules run clockwise. \emph{Run (web, 2026-09-27):} 5 checks, 0 failing; rules: {'k=1': {'q-commuting pairs': '15/15', 'Ptolemy pairs': '10/10', 'opposite orientation': '5/15, 0/10', 'm(a) draft': [1], 'm(a) code': [1]}, 'k=2': {'q-commuting pairs': '126/126', 'Ptolemy pairs': '70/70', 'opposite orientation': '70/126, 0/70', 'm(a) draft': [2, 1], 'm(a) code': [2, 1]}, 'k=3': {'q-commuting pairs': '477/477', 'Ptolemy pairs': '252/252', 'opposite orientation': '297/477, 0/252', 'm(a) draft': [3, 1, 2], 'm(a) code': [3, 1, 2]}, 'k=4': {'q-commuting pairs': '1276/1276', 'Ptolemy pairs': '660/660', 'opposite orientation': '836/1276, 0/660', 'm(a) draft': [4, 1, 3, 2], 'm(a) code': [4, 1, 3, 2]}, 'k=5': {'q-commuting pairs': '2795/2795', 'Ptolemy pairs': '1430/1430', 'opposite orientation': '1885/2795, 0/1430', 'm(a) draft': [5, 1, 4, 2, 3], 'm(a) code': [5, 1, 4, 2, 3]}, 'k=6': {'q-commuting pairs': '5370/5370', 'Ptolemy pairs': '2730/2730', 'opposite orientation': '3690/5370, 0/2730', 'm(a) draft': [6, 1, 5, 2, 4, 3], 'm(a) code': [6, 1, 5, 2, 4, 3]}}; BPS chart: {'k=1': '25/25', 'k=2': '196/196', 'k=3': '729/729'}; seconds: 17.3. \emph{Run (web (extensive), 2026-09-28):} 5 checks, 0 failing; rules: {'k=1': {'q-commuting pairs': '15/15', 'Ptolemy pairs': '10/10', 'opposite orientation': '5/15, 0/10', 'm(a) draft': [1], 'm(a) code': [1]}, 'k=2': {'q-commuting pairs': '126/126', 'Ptolemy pairs': '70/70', 'opposite orientation': '70/126, 0/70', 'm(a) draft': [2, 1], 'm(a) code': [2, 1]}, 'k=3': {'q-commuting pairs': '477/477', 'Ptolemy pairs': '252/252', 'opposite orientation': '297/477, 0/252', 'm(a) draft': [3, 1, 2], 'm(a) code': [3, 1, 2]}, 'k=4': {'q-commuting pairs': '1276/1276', 'Ptolemy pairs': '660/660', 'opposite orientation': '836/1276, 0/660', 'm(a) draft': [4, 1, 3, 2], 'm(a) code': [4, 1, 3, 2]}, 'k=5': {'q-commuting pairs': '2795/2795', 'Ptolemy pairs': '1430/1430', 'opposite orientation': '1885/2795, 0/1430', 'm(a) draft': [5, 1, 4, 2, 3], 'm(a) code': [5, 1, 4, 2, 3]}, 'k=6': {'q-commuting pairs': '5370/5370', 'Ptolemy pairs': '2730/2730', 'opposite orientation': '3690/5370, 0/2730', 'm(a) draft': [6, 1, 5, 2, 4, 3], 'm(a) code': [6, 1, 5, 2, 4, 3]}}; BPS chart: {'k=1': '25/25', 'k=2': '196/196', 'k=3': '729/729', 'k=4': '1936/1936'}; seconds: 292.4. & \CERTIFIED\newline {\footnotesize theorem, extensive\newline web}\newline {\footnotesize run web 2026-09-27: 5/5}\newline {\footnotesize run web (extensive) 2026-09-28: 5/5} \\ \addlinespace
\code{app:a1a2k/traces}\newline App.~\pref{app:a1a2k}, the trace & The reduction algorithm terminates on elementary traces $T_a$, $i$-independent; $T_0=\chi_1(\fq^2)$ and $T_a=(-1)^{m+1}\fq^{-m}$ $\times(\chi_m-\chi_{m+1})(\fq^2)$ with $m(a)$ as given. & The draft's $T_0,\dots,T_k$, with the $M(2,2k+3)$ characters as Andrews--Gordon products, against: the class's elementary traces (a transcription check, since those are this formula); the BPS chart of the $A_{2k}$ quiver, independently --- $\Tr(1)$ and $\Tr L_{a;i}$ for every rotation index $i$; and the reduction algorithm (the class's Layer~1) on every label of at most two letters, $\sum_s c_sT_s$ against the chart's trace of the label's image. \emph{Population:} $k\le6$ against the class's elementary traces (to $\fq^{40}$); $k\le2$ against the BPS chart, $k=3$ at extensive depth. \emph{Controls:} positive: the Andrews--Gordon products equal the Andrews--Gordon sums ($k\le3$, to $q^{30}$). negative: the formula with $m(a)=a$ disagrees with the chart. \emph{Evidence:} emergent. \emph{Where:} \code{battery/checks_finite.py (check_a1a2k_traces)}, \code{tests/test_a1a2k_kalg.py}, \code{notes/a1a2k_kq_proof.md}. \emph{Run (web, 2026-09-27):} 4 checks, 0 failing; BPS chart: {'k=1': {'Tr(1) to fq\^{}12': True, 'Tr(L\_(a;i)) = T\_a for every i': '[5] of 5 each', 'reduction on labels of $\le$ 2 letters to fq\^{}10': '16/16'}, 'k=2': {'Tr(1) to fq\^{}10': True, 'Tr(L\_(a;i)) = T\_a for every i': '[7, 7] of 7 each', 'reduction on labels of $\le$ 2 letters to fq\^{}8': '85/85'}}; seconds: 111.9. \emph{Run (web (extensive), 2026-09-28):} 4 checks, 0 failing; BPS chart: {'k=1': {'Tr(1) to fq\^{}12': True, 'Tr(L\_(a;i)) = T\_a for every i': '[5] of 5 each', 'reduction on labels of $\le$ 2 letters to fq\^{}10': '16/16'}, 'k=2': {'Tr(1) to fq\^{}10': True, 'Tr(L\_(a;i)) = T\_a for every i': '[7, 7] of 7 each', 'reduction on labels of $\le$ 2 letters to fq\^{}8': '85/85'}, 'k=3': {'Tr(1) to fq\^{}8': True, 'Tr(L\_(a;i)) = T\_a for every i': '[9, 9, 9] of 9 each', 'reduction on labels of $\le$ 2 letters, the first at rotation 0, to fq\^{}6': '80/80'}}; seconds: 805.2. & \MEASURED\newline {\footnotesize identity, extensive\newline web}\newline {\footnotesize run web 2026-09-27: 4/4}\newline {\footnotesize run web (extensive) 2026-09-28: 4/4} \\ \addlinespace
\code{app:a1a2k/minors}\newline App.~\pref{app:a1a2k}, ``The minors miracle'' & The normalised coefficient vectors stabilise to $M^{(r)}_s(\fq)$ with $\sum_sM^{(r)}_sT_s=0$, linearly independent, and $(-1)^s\det\mathcal M^{(\hat s)}=T_s$. & The rows: the reduction algorithm's coefficient vectors of $\Tr L_{2;0}^a$ and $\Tr L_{2;0}^aD_r$ ($D_1=L_{1;0}$, $D_r=L_{r+1;0}$), each divided by its most negative power of $\fq$, at two consecutive $a$. Measured: the order to which they stabilise; that each stabilised row annihilates the draft's $(T_0,\dots,T_k)$; that the maximal minors reproduce the $T_s$, and with which sign. \emph{Population:} $k=2,3,4$; at extensive depth also $k=5$ and $k=3$ at larger powers $a$. \emph{Controls:} negative: with $D_1=L_{1;0}^2$ two rows coincide and every maximal minor vanishes to the stabilised order. \emph{Evidence:} emergent. \emph{Where:} \code{battery/checks_finite.py (check_a1a2k_minors)}, \code{tests/test_a1a2k_miracle.py}. \emph{Note:} The minors reproduce T\_s with the overall sign (-1)\^{}(s+1), the same at k = 2, 3, 4 (as tests/test\_a1a2k\_miracle.py found at k = 3); the draft prints (-1)\^{}s. Every row's leading coefficient is +1, so the sign is fixed by the row order the draft states. T\_s is the draft's formula, not the class's traces. \emph{Run (web, 2026-09-27):} 3 checks, 0 failing; minors: {'k=2, a=8,9': {'rows stable to fq\^{}': 15, 'rows annihilate T to fq\^{}': 15, '(-1)\^{}s det = T\_s to fq\^{}': -1, '(-1)\^{}(s+1) det = T\_s to fq\^{}': 15}, 'k=3, a=11,12': {'rows stable to fq\^{}': 21, 'rows annihilate T to fq\^{}': 21, '(-1)\^{}s det = T\_s to fq\^{}': -1, '(-1)\^{}(s+1) det = T\_s to fq\^{}': 21}, 'k=4, a=8,9': {'rows stable to fq\^{}': 15, 'rows annihilate T to fq\^{}': 15, '(-1)\^{}s det = T\_s to fq\^{}': -1, '(-1)\^{}(s+1) det = T\_s to fq\^{}': 15}}; sign: {'2': '-', '3': '-', '4': '-'}; seconds: 4.4. \emph{Run (web (extensive), 2026-09-28):} 3 checks, 0 failing; minors: {'k=2, a=8,9': {'rows stable to fq\^{}': 15, 'rows annihilate T to fq\^{}': 15, '(-1)\^{}s det = T\_s to fq\^{}': -1, '(-1)\^{}(s+1) det = T\_s to fq\^{}': 15}, 'k=3, a=11,12': {'rows stable to fq\^{}': 21, 'rows annihilate T to fq\^{}': 21, '(-1)\^{}s det = T\_s to fq\^{}': -1, '(-1)\^{}(s+1) det = T\_s to fq\^{}': 21}, 'k=4, a=8,9': {'rows stable to fq\^{}': 15, 'rows annihilate T to fq\^{}': 15, '(-1)\^{}s det = T\_s to fq\^{}': -1, '(-1)\^{}(s+1) det = T\_s to fq\^{}': 15}, 'k=5, a=8,9': {'rows stable to fq\^{}': 15, 'rows annihilate T to fq\^{}': 15, '(-1)\^{}s det = T\_s to fq\^{}': -1, '(-1)\^{}(s+1) det = T\_s to fq\^{}': 15}, 'k=3, a=14,15': {'rows stable to fq\^{}': 27, 'rows annihilate T to fq\^{}': 27, '(-1)\^{}s det = T\_s to fq\^{}': -1, '(-1)\^{}(s+1) det = T\_s to fq\^{}': 27}}; sign: {'2': '-', '3': '-', '4': '-', '5': '-'}; seconds: 14.3. & \MEASURED\newline {\footnotesize conjecture, extensive\newline web}\newline {\footnotesize run web 2026-09-27: 3/3}\newline {\footnotesize run web (extensive) 2026-09-28: 3/3} \\ \addlinespace
\end{longtable}

\subsection*{The companion's own examples (\S\ref{sec:extra-examples} of this companion, not in the paper)}
\begin{longtable}{@{}>{\raggedright\arraybackslash}p{2.4cm}>{\raggedright\arraybackslash}p{4.7cm}>{\raggedright\arraybackslash}p{4.95cm}>{\raggedright\arraybackslash}p{1.85cm}@{}}
\toprule claim & statement & test: procedure, population, controls & standing \\ \midrule \endfirsthead
\toprule claim & statement & test: procedure, population, controls & standing \\ \midrule \endhead
\bottomrule \endfoot
\code{comp:u1a1aodd/implementation}\newline companion \S\ref{comp:u1a1aodd} & \code{U1A1AoddKAlg(k)} implements the presentation verbatim: the chord letters and $E^{\pm1}$, $\rho$'s $E$-drift \eqref{eq:comp-drift}, $\fq$-commutation iff non-crossing with the exponent \eqref{eq:comp-qcomm} from the charges \eqref{eq:comp-charge}, $E$ $\fq$-commuting with every chord by the same formula (the even types commute with it, the odd ones with $\fq^{\pm2}$), the two-resolution unit-coefficient shape of \eqref{eq:comp-ptolemy}, and the embedding $\Phi$. & Independent transcription of the printed rules (chords, rotation law's $E$-powers, crossing test, the interior and wrap charge formulas, the chain pairing, the resolutions of a crossing quadrilateral, $\Phi$) compared with the class's \code{rho}, \code{multiply}, \code{cocycle}, \code{q_commute} and \code{phi} letter by letter and pair by pair. \emph{Population:} $k = 1, 2, 3$: every letter and every ordered pair of letters. \emph{Controls:} positive: the printed $k=1$ values reproduced; negative: the drift with a wrong anchor and the charges without the $-2\mu$ exception are detected. \emph{Evidence:} transcription; the products are emergent (the class computes them by the analytic rules, not from these formulas). \emph{Where:} \code{battery/checks_extra_examples.py}, \code{src/cone/u1a1aodd_kalg.py}. \emph{Note:} The companion's own example; the write-up is the author's, in the source repository, whose status the companion repeats verbatim. \emph{Run (web, 2026-09-26):} 19 checks, 0 failing; k: [1, 2, 3]; census: {'1': {'letters': 9, 'pairs': 72, 'qcommuting': 42, 'crossing': 30, 'drift\_ok': True, 'letters\_ok': True, 'qc\_bad': 0, 'prod\_bad': 0, 'ptolemy\_bad': 0, 'phi\_bad': 0, 'E\_ok': True, 'wrong\_anchor\_detected': True}, '2': {'letters': 20, 'pairs': 380, 'qcommuting': 240, 'crossing': 140, 'drift\_ok': True, 'letters\_ok': True, 'qc\_bad': 0, 'prod\_bad': 0, 'ptolemy\_bad': 0, 'phi\_bad': 0, 'E\_ok': True, 'wrong\_anchor\_detected': True}, '3': {'letters': 35, 'pairs': 1190, 'qcommuting': 770, 'crossing': 420, 'drift\_ok': True, 'letters\_ok': True, 'qc\_bad': 0, 'prod\_bad': 0, 'ptolemy\_bad': 0, 'phi\_bad': 0, 'E\_ok': True, 'wrong\_anchor\_detected': True}}; seconds: 0.2. & \CERTIFIED\newline {\footnotesize implementation, fast\newline web}\newline {\footnotesize run web 2026-09-26: 19/19} \\ \addlinespace
\code{comp:u1a1aodd/is-kalgebra}\newline companion \S\ref{comp:u1a1aodd} & The class is a $K_\fq$-algebra: K1--K5 hold on a window of letters, $E^{\pm1}$ and products at $k=1,2$; all five axioms are emergent (the tables are analytic, not built to satisfy them). & The contract's verifier battery (\code{run_k_verifiers}) on all pairs of the window, plus closure of the products with $\bZ[\fq^{\pm1}]$ coefficients. \emph{Population:} $k = 1, 2$: the letters, $E^{\pm1}$ and a set of products; $K = 4$. \emph{Controls:} positive: the identity row; negative: n.a. (the verifiers are the instruments). \emph{Evidence:} emergent. \emph{Where:} \code{battery/checks_extra_examples.py}, \code{src/cone/u1a1aodd_kalg.py}. \emph{Run (web, 2026-09-26):} 2 checks, 0 failing; k: [1, 2]; K: 4; seconds: 1.4. & \CERTIFIED\newline {\footnotesize implementation, fast\newline web}\newline {\footnotesize run web 2026-09-26: 2/2} \\ \addlinespace
\code{comp:u1a1aodd/embedding}\newline companion \S\ref{comp:u1a1aodd} & $\Phi$ is an algebra map: $\Phi(g)\Phi(h)=\sum_cC^c_{gh}\Phi(c)$ on every ordered letter pair, the left side computed in $A^{\mathrm{even}}_\fq\otimes Q_\fq(\bZ^2)$ by the even algebra's own product. & For each pair the product of the images in the tensor algebra (\code{TensorKAlgebra} of \code{A1A2kKAlg} and the rank-2 quantum torus) against the image of the class's product; the even family is an independent algebra. \emph{Population:} $k = 1, 2$ all ordered letter pairs ($k = 3$ when within the web budget); the author's experiment covers $k = 1..6$ (21,068 identities). \emph{Controls:} positive: the identity pairs; negative: n.a. \emph{Evidence:} emergent. \emph{Where:} \code{battery/checks_extra_examples.py}. \emph{Note:} Held out at k = 6 in the author's experiment; the web tier runs the small k. \emph{Run (web, 2026-09-26):} 3 checks, 0 failing; census: {'1': {'pairs': 121, 'bad': 0}, '2': {'pairs': 484, 'bad': 0}, '3': {'pairs': 1369, 'bad': 0}}; seconds: 0.5. & \CERTIFIED\newline {\footnotesize identity, fast\newline web}\newline {\footnotesize run web 2026-09-26: 3/3} \\ \addlinespace
\code{comp:u1a1aodd/traces}\newline companion \S\ref{comp:u1a1aodd} & The trace seeds are the printed characters: $\Tr(E^n)$ is \eqref{eq:comp-singlet}, $\Tr(L_{2m;i}E^e)$ is the one rule \eqref{eq:comp-longchord} for every even type at the character index $n=-(-1)^ig$ at the positions $i=0,1$, and at the other positions after the $\rho^2$ steps of \eqref{eq:comp-drift}; the odd-type single chords have $\Tr=0$. & The printed singlet and chord rule transcribed independently of \file{src/cone/u1_pgon_layer2.py} and compared coefficient by coefficient with the class's \code{trace}; the transcription compared with the served \code{singlet_chord_trace} deep. \emph{Population:} $k = 1, 2, 3$ ($p = 3, 4, 5$): $n = -3..3$ for the singlet; every even type, positions $i = 0, 1$, $e = -4..4$ for the chord rule, and every other position, $e = -2..2$, through $\fq^{12}$; every odd-type letter; through $\fq^{40}$; the printed rule against the served code for $p = 3..7$ through $\fq^{80}$. \emph{Controls:} positive: the printed rule equals the served code through $\fq^{80}$ (and \eqref{eq:comp-singlet} at $m = 0$); negative: the retired printed long chord differs from $\fq^{39}$; the index rule read at face value at the chord $\{2,5\}$ ($k=1$) fails, as the text says. \emph{Evidence:} the seeds are closed forms served by the class; the reduction to them is cone-data cyclicity. \emph{Where:} \code{src/cone/u1_pgon_layer2.py}, \code{battery/checks_extra_examples.py}. \emph{Note:} Companion corrected 2026-09-23: the four-partial-theta long chord printed before, wrong from q\^{}39 at p = 4, is replaced by the singlet rule for every even type. The rule is measured, not derived: two exact routes agree with it (finite\_type\_trace\_machinery.md section H). Corrected 2026-09-26: the index rule was printed 'for every type and every k' at every position; it holds at i = 0, 1 (one position per rho\^{}2-orbit) and fails at face value at some positions near vertex H-1 (k=1: {2,5}; k=2: types 2 at i = 4, 5, 6; k=3: type 2 at i = 6, 7, 8 and type 4 at i = 4); the class applies it at the representative. \emph{Run (web, 2026-09-26):} 15 checks, 0 failing; k: [1, 2, 3]; K: 40; seconds: 0.1. & \CERTIFIED\newline {\footnotesize identity, fast\newline web}\newline {\footnotesize run web 2026-09-26: 15/15} \\ \addlinespace
\code{comp:a1aodd/implementation}\newline companion \S\ref{comp:a1aodd} & \code{ungauge_u1a1aodd(k)} implements the presentation: $E$ $\fq$-commutes with every chord with the magnetic charge \eqref{eq:comp-aodd-mag}; a label is in the algebra iff it is balanced; the generators are the balanced chords and the non-crossing (even--even, odd--odd) pairs, $k(k+2)+(k+2)\binom{k+2}{3}$ of them; the product is the gauged product, every term balanced; the fugacity \eqref{eq:comp-aodd-lift}; $\rho$ \eqref{eq:comp-aodd-rho}, with $\rho^H=1$ and order $3$, $8$, $10$ modulo powers of $z$; two words on one label at $k=3$. & The parity rule, the balance condition, the generator enumeration, the lift and $\rho$ with the drift transcribed from the companion and the gauged example's printed charges and drift, and compared with the $E$-commutator in \code{U1A1AoddKAlg}, \code{in_centralizer} and the multiply-based centralizer test, \code{mult_generators}, \code{geometric_label}, \code{multiply} against the gauged class, \code{r_label_decompose}, \code{trace}, \code{rho}, \code{rho_inverse}, and the zoo's $\rho$-permutations. \emph{Population:} $k = 1, 2, 3$: every letter ($9$, $20$, $35$); a window of canonical labels ($645$, $8{,}763$, $4{,}833$: up to three letters at $k=1,2$, two at $k=3$, exponents $1, 2$, $E$-powers $-1, 0, 2$); every generator; all ordered generator pairs at $k=1,2$ ($36$, $576$) and $300$ at $k=3$; for $\rho$ the generators and product terms with $E$-shifts ($79$, $229$, $474$ labels). \emph{Controls:} positive: at $k=1$ the transcribed generators are \code{HexagonKAlg}'s six hand-named ones; negative: an off-by-one parity rule, $\rho$ without the drift, and the lift with the opposite sign are each caught. \emph{Evidence:} transcription; the products are emergent (the gauged class computes them by its analytic rules, not from these formulas). \emph{Where:} \code{battery/checks_extra_examples.py}, \code{src/cone/ungauge_kalgebra.py}, \code{src/cone/u1a1aodd_kalg.py}, \code{tests/test_ungauged_aodd_geometry.py}, \code{tests/test_ungauge_kalgebra.py}. \emph{Note:} The companion's own example, added 2026-09-26 at the author's request that the ungauged [A\_1,A\_{2k+1}] join Section 2; its geometric labels are the balanced multisets of the family. \emph{Run (web, 2026-09-26):} 26 checks, 0 failing; k: [1, 2, 3]; census: {'1': {'letters': 9, 'window': 645, 'balanced': 111, 'generators': 6, 'product\_pairs': 36, 'rho\_labels': 79, 'rho\_without\_drift\_differs': 28, 'lift\_opposite\_sign\_fails': 21}, '2': {'letters': 20, 'window': 8763, 'balanced': 1239, 'generators': 24, 'product\_pairs': 576, 'rho\_labels': 229, 'rho\_without\_drift\_differs': 38, 'lift\_opposite\_sign\_fails': 21}, '3': {'letters': 35, 'window': 4833, 'balanced': 1053, 'generators': 65, 'product\_pairs': 300, 'rho\_labels': 474, 'rho\_without\_drift\_differs': 69, 'lift\_opposite\_sign\_fails': 21}}; seconds: 3.8. & \CERTIFIED\newline {\footnotesize implementation, fast\newline web}\newline {\footnotesize run web 2026-09-26: 26/26} \\ \addlinespace
\code{comp:a1aodd/hexagon}\newline companion \S\ref{comp:a1aodd} & At $k=1$ the printed presentation --- the $\fq$-commutations \eqref{eq:comp-aodd-hexqc}, the three relations \eqref{eq:comp-aodd-hex} with their bar images, and $\rho$ \eqref{eq:comp-aodd-hexrho} --- gives all $36$ ordered products of the six generators, and every balanced label is a monomial of the six cones. & The three relations transcribed, carried to $j=1,2$ by the printed $\rho$ with its powers of $z$ and reversed by bar, and compared with \code{multiply}; the $\fq$-commutations against the reversed products; the printed $\rho$ against \code{rho}; the cone statement by enumeration. \emph{Population:} $k=1$: the $36$ ordered products of the six generators; the $73$ balanced labels with up to four letters and exponents up to $3$. \emph{Controls:} positive: the printed $\rho$ on the six generators equals the class's; negative: $\rho$ without its powers of $z$ breaks $8$ of the $36$ products, the two $\fq$-commutation powers exchanged break $12$. \emph{Evidence:} transcription; the products are emergent. \emph{Where:} \code{battery/checks_extra_examples.py}, \code{src/cone/ungauge_kalgebra.py}, \code{src/cone/hexagon_kalg.py}, \code{tests/test_hexagon_kalg.py}. \emph{Note:} The k = 1 algebra is the paper's [A\_1,D\_3] with its SU(2) flavour seen through the Cartan U(1); the trace-level identification is checked in comp:a1aodd/traces. \emph{Run (web, 2026-09-26):} 5 checks, 0 failing; k: 1; products: 36; balanced labels: 73; seconds: 0.0. & \CERTIFIED\newline {\footnotesize identity, fast\newline web}\newline {\footnotesize run web 2026-09-26: 5/5} \\ \addlinespace
\code{comp:a1aodd/is-kalgebra}\newline companion \S\ref{comp:a1aodd} & The class is a $K_\fq$-algebra: K1--K5 hold on windows of the identity, $E$, generators, a generator times $E^{-1}$ and products at $k=1,2$; all five axioms are emergent. & The contract's verifier battery (\code{run_k_verifiers}) on all pairs of each window, and closure of the products over $\bZ[\fq^{\pm1}]$; the same algebra with $\rho$ stripped of its drift run through \code{verify_rho_is_automorphism} as the negative control. \emph{Population:} $k=1$: $11$ labels (all six generators); $k=2$: $12$ labels ($7$ sampled generators); all pairs; $K=4$. \emph{Controls:} positive: the identity row; negative: $\rho$ without the drift fails \code{verify_rho_is_automorphism} on $14$ of the $49$ pairs of the identity and the generators at $k=1$. \emph{Evidence:} emergent. \emph{Where:} \code{battery/checks_extra_examples.py}, \code{src/cone/ungauge_kalgebra.py}. \emph{Note:} Orthonormality on generator pairs is also pinned at k = 4, 5 by tests/test\_ungauged\_aodd\_geometry.py (40/40 and 20/20 at K = 4). \emph{Run (web, 2026-09-26):} 3 checks, 0 failing; window sizes: {'1': 11, '2': 12}; K: 4; seconds: 9.9. & \CERTIFIED\newline {\footnotesize implementation, fast\newline web}\newline {\footnotesize run web 2026-09-26: 3/3} \\ \addlinespace
\code{comp:a1aodd/traces}\newline companion \S\ref{comp:a1aodd} & The trace is \eqref{eq:comp-aodd-trace}, the gauged traces summed over the gauge charge with $z^n$ and divided by $(\fq^2;\fq^2)_\infty^2$; the vacuum and the balanced chords are \eqref{eq:comp-aodd-closed}, with $z\mapsto z^{-1}$ at $i=0$ and $\rho^2$ for the other positions; the heads are \eqref{eq:comp-aodd-heads}; at $k=1$ the traces are the paper's $[A_1,D_3]$ traces restricted to the Cartan $U(1)$; the class's window covers $|g+n|\le K+1$, where $\Tr_{\rm gauged}L_{(F,n)}=O(\fq^{|g+n|})$. & The formula transcribed with the gauged traces summed over a wide window and an independently expanded $(\fq^2;\fq^2)_\infty^{-2}$; the closed forms from the transcriptions of \eqref{eq:comp-singlet} and \eqref{eq:comp-longchord} used by the gauged row; $g$ from the printed charges; the heads as literals; at $k=1$ the traces of \code{A1D3KAlg} restricted, $\chi_n\mapsto z^{-n/2}+\dots+z^{n/2}$. \emph{Population:} $k=1,2,3$: the formula on the identity, $E^2$, every generator, $E$-shifted generators and terms of generator products (through $\fq^{10}$); the window on the powers $2..4$ of the two $k=3$ pair generators at $K=1..7$; $O(\fq^{|g+n|})$ on the generators and their powers $2..4$; the closed forms and $\rho^2$ on every balanced chord position through $\fq^{30}$; the $[A_1,D_3]$ comparison through $\fq^{30}$. \emph{Controls:} positive: at $k=1$ the traces are \code{A1D3KAlg}'s (the paper's $[A_1,D_3]$) restricted, through $\fq^{30}$; negative: the measure $(\fq^2;\fq^2)_\infty^{-1}$, the orientation of $L_{2;1}$ used for $L_{2;0}$, and the $E$-power window alone (which loses $-\fq^3z^{-6}$ of the $k=3$ pair cube at $K=4$) are each caught. \emph{Evidence:} transcription (the formula and the closed forms recomputed from the gauged seeds and an independent measure); cross-presentation at $k=1$ (\textbackslash{}code{A1D3KAlg}). \emph{Where:} \code{battery/checks_extra_examples.py}, \code{src/cone/ungauge_kalgebra.py}, \code{src/cone/u1_pgon_layer2.py}, \code{src/cone/a1d3_kalg.py}, \code{tests/test_ungauge_kalgebra.py}. \emph{Note:} The gauged seeds eq:comp-singlet / eq:comp-longchord are measured, not derived (comp:u1a1aodd/traces). The gauge-charge window was added to UngaugedKAlgebra.trace on 2026-09-26, while this example was written: the E-power window alone dropped terms of the k = 3 pair powers a = 2..4 at a $\le$ K $\le$ 2a - 2 (6 of 862 labels measured, all at k = 3). The gauged rule's index -(-1)\^{}i g is used here only at i = 0, 1, where g = 0; at chords ending on vertex H-1 or wrapping it does not hold as printed (the class applies it on the rho\^{}2-canonical representative). \emph{Run (web, 2026-09-26):} 9 checks, 0 failing; formula labels k=1: 32; formula labels k=2: 46; formula labels k=3: 88; window label-orders: 42; slack labels: 380; closed-form comparisons: 37; seconds: 96.6. & \CERTIFIED\newline {\footnotesize identity, fast\newline web}\newline {\footnotesize run web 2026-09-26: 9/9} \\ \addlinespace
\code{comp:a1aodd/presentations}\newline companion \S\ref{comp:a1aodd} & The zoo's exported cone tables \code{a3}, \code{a5}, \code{a7} present the same algebra: the \code{KAlgebraIso} of \code{aodd_seeds.kalgebra_iso} passes, and the zoo's own trace of composite labels agrees; \code{HexagonKAlg} and \code{UngaugedPolygonKAlg(k)} are this algebra with $E$ inverted, respectively in the same labels, with the fugacity read as $z^{-1}$; $\Tr1$ equals the RG flow \code{A1AoddToEvenRGKAlgebra(k)}'s; the zoo's \code{a3} characters and the $u(1)$ orthonormality bootstrap give the same traces. & \code{KAlgebraIso.verify_all} (unit, round trip, multiplicative, $\rho$-equivariant, trace-equivariant) from the zoo's Z-form wrapper \code{FiniteU1ZKAlgebra}, base-changed to the fugacity $z$, onto \code{ungauge_u1a1aodd(k)}; the zoo's own Layer-1 trace of mapped composite labels; the named classes' products and traces against the class under their label and fugacity conventions; \code{A1AoddToEvenRGKAlgebra(k).trace}, \code{ad_characters.a3_elem_entry}, \code{u1_bootstrap.generate_u1}. \emph{Population:} the isomorphism at $k=1,2,3$ on the identity, every generator and the flavour characters, $64$, $676$, $4{,}489$ ordered pairs each way, traces through $\fq^{12}$; $5$ composite labels; \code{HexagonKAlg}: all $36$ generator products and $10$ traces; \code{UngaugedPolygonKAlg(2,3,4)}: $25$ products and $6$ traces each; the RG flow at $k=1,2,3$ through $\fq^{16}$; the \code{a3} characters through $\fq^{48}$; the bootstrap for \code{a5} through $\fq^{10}$, \code{a7} through $\fq^{6}$. \emph{Controls:} positive: the certified isomorphism and $\Tr1$ against the RG flow; negative: a flipped flavour normalisation and two transposed zoo generators each fail the isomorphism, and the \code{a3} characters read back with $z\mapsto z$ differ on $4$ of $6$ generator traces. \emph{Evidence:} cross-presentation (exported BPS cone tables; an RG flow; closed-form characters; an orthonormality bootstrap). \emph{Where:} \code{battery/checks_extra_examples.py}, \code{src/cone/aodd_seeds.py}, \code{src/cone/ad_characters.py}, \code{src/cone/u1_bootstrap.py}, \code{src/rg/a1aodd_to_even_rgkalgebra.py}, \code{src/cone/hexagon_kalg.py}, \code{src/cone/ungauged_polygon_kalg.py}, \code{tests/test_aodd_seeds.py}. \emph{Note:} The zoo's a3/a5/a7 seed traces are served from ungauge\_u1a1aodd through the same map, so the isomorphism's trace check is not independent there (finite\_type\_trace\_machinery.md section N); its product and rho checks are, and so are the zoo's own reduction of composite labels, the RG flow, the a3 characters and the bootstrap. Recorded further than the battery runs: the bootstrap for a5 through q\^{}20 and a7 through q\^{}14 (src/cone/aodd\_seeds.py). \emph{Run (web, 2026-09-26):} 9 checks, 0 failing; a3: {'pairs each way': 64, 'seconds': 0.2}; a5: {'pairs each way': 676, 'seconds': 1.1}; a7: {'pairs each way': 4489, 'seconds': 9.8}; rg: {'1': [True, 2.6], '2': [True, 7.6], '3': [True, 13.6]}; bootstrap: {'a5': [True, 24, 24, 4.0], 'a7': [True, 65, 65, 5.3]}; seconds: 48.2. & \CERTIFIED\newline {\footnotesize identity, fast\newline web}\newline {\footnotesize run web 2026-09-26: 9/9} \\ \addlinespace
\code{comp:su3ad/implementation}\newline companion \S\ref{comp:su3ad} & \code{SU3ADKAlg} implements the relations of the definition verbatim for every $i$ with the $\star$-parity rule, the two $\fq$-commutations, and $\rho$ as the shift composed with $\star$ on $R(SU(3))$. & Every printed right-hand side built as an \code{Element} from the class's labels and compared with \code{multiply}; $\rho$ on generators and on $\chi_{(1,0)}$. \emph{Population:} all four $i$, all ten relation families and the two $\fq$-commutations. \emph{Controls:} positive: the printed relations; negative: the relations without the $\star$ swap at odd $i$ fail. \emph{Evidence:} transcription. \emph{Where:} \code{battery/checks_extra_examples.py}, \code{src/cone/su3_ad_kalg.py}, \code{tests/test_su3_ad_kalg.py}. \emph{Note:} The constant term of T\_i T\_{i+2} is 2 + chi\_(1,1); an older tabulation had 4 + 2 chi\_(1,1) (su3ad\_flavour\_layer1\_defect.md). \emph{Run (web, 2026-09-26):} 4 checks, 0 failing; relations: 40; seconds: 0.0. & \CERTIFIED\newline {\footnotesize implementation, fast\newline web}\newline {\footnotesize run web 2026-09-26: 4/4} \\ \addlinespace
\code{comp:su3ad/is-kalgebra}\newline companion \S\ref{comp:su3ad} & The class is a flavoured $K_\fq$-algebra: K1--K5 and F5 on a window (the identity, the eight generators, their $\chi$-dressings and the tile monomials), all emergent. & \code{run_k_verifiers} plus \code{verify_trace_intertwines_rho_star}. \emph{Population:} 16 labels, all pairs, $K = 4$. \emph{Controls:} positive: the identity row; negative: n.a. \emph{Evidence:} emergent. \emph{Where:} \code{battery/checks_extra_examples.py}. \emph{Run (web, 2026-09-26):} 2 checks, 0 failing; labels: 15; K: 4; seconds: 8.1. & \CERTIFIED\newline {\footnotesize implementation, fast\newline web}\newline {\footnotesize run web 2026-09-26: 2/2} \\ \addlinespace
\code{comp:su3ad/traces}\newline companion \S\ref{comp:su3ad} & The printed expansions \eqref{eq:comp-su3-traces}; the served $\Tr1$, $\Tr T_0$, $\Tr D_0$ are the even-$D$ $k=1$ formulas of \S\ref{comp:a1deven} (the ungauged trace \eqref{eq:comp-deven-ungauge} of the images, restricted along $3\to2_{+1}\oplus1_{-2}$); the Kac--Wakimoto character equals the served $\Tr1$ through $\fq^{100}$ and the orthonormality bootstrap the served $\Tr T_0$, $\Tr D_0$ through $\fq^{80}$; $\Tr1$ equals the BPS chart's $\Tr1$ (an independent engine); the bootstrap's defining constraint $\Tr(T_0^a),\Tr(D_0^a)=O(\fq^a)$ holds on the served traces. & Transcription of the printed series; the even-$D$ tower and seed rules transcribed in the \code{comp:a1deven} adapters, summed over the gauge charge with the measure restored, against the class's traces restricted through the $SU(3)$ weight diagrams; the two witnesses (\code{vacuum_character}, \code{SU3ElemTraces}) compared with the served seed series in Cartan fugacities, as \file{tests/test_su3ad_deven_map.py} does, those series tied to the class's traces through $\fq^{24}$; the BPS chart's trace of the identity (the exact Habiro engine); valuations of the seed powers. \emph{Population:} the printed heads through $\fq^6$; the even-$D$ formulas through $\fq^{16}$; the Kac--Wakimoto witness through $\fq^{100}$, the bootstrap through $\fq^{80}$; $\Tr1$ through $\fq^6$ against \code{SU3BPSKAlgebra}; $a = 1, 2, 3$. \emph{Controls:} positive: the printed heads, and the served $\Tr1$ against the Kac--Wakimoto character; negative: $D_0$'s offset moved by $3$, $D_1$'s image in place of $D_0$'s and $z$ read with $U(1)$ charge $-3$ each disagree with the served traces; a vacuum character without the adjoint at $\fq^2$ fails the BPS comparison. \emph{Evidence:} cross-presentation (the even-$D$ closed forms against two BPS-free witnesses and the BPS chart). \emph{Where:} \code{battery/checks_extra_examples.py}, \code{src/cone/sl3_su3_traces.py}, \code{src/cone/su3_ad_kalg.py}, \code{src/cone/u1a1deven_seed_characters.py}, \code{tests/test_su3ad_deven_map.py}. \emph{Note:} Updated 2026-09-26: since 2026-09-24 the served elementary traces are the even-D k = 1 closed forms (Creutzig's tower for Tr 1, the seeds through the curve map for Tr T\_0, Tr D\_0; sl3\_su3\_traces docstring); the Kac-Wakimoto character and the forward bootstrap SU3ElemTraces, which served them before, are the witnesses. The comparison at q\^{}100 / q\^{}80 is made on the served seed series in Cartan fugacities: the Weyl symmetrization of the class's traces at those orders would cost more than the witnesses. \emph{Run (web, 2026-09-26):} 9 checks, 0 failing; printed heads K: 6; even-D formulas K: 16; Kac-Wakimoto witness K: 100; bootstrap witness K: 80; seconds: 146.7. & \CERTIFIED\newline {\footnotesize identity, fast\newline web}\newline {\footnotesize run web 2026-09-26: 9/9} \\ \addlinespace
\code{comp:a1deven/implementation}\newline companion \S\ref{comp:a1deven} & \code{U1A1DevenConeKAlgebra(k)} implements the curves and labels, the relations \eqref{eq:comp-deven-qcomm}--\eqref{eq:comp-deven-norm} with the cocycle read off \code{A1DnKAlg(2k+1)} through the collapse of the edge $\{1,2\}$, the printed exponents of the skein rule and the products \eqref{eq:comp-deven-products}, and $\rho$ with its full turn \eqref{eq:comp-deven-rho}; its products equal the flow \code{U1A1DevenViaDoddRG(k)}'s on every ordered generator pair; \code{A1DevenKAlg(k)}'s generators are the printed ones and $\rho(z)=z^{-1}$; at $k=1$, \code{SU3ADKAlg}'s geometric labels are the printed curves of \code{A1DevenKAlg(1)} ($T_i$ the loop at $i$ with the curve $(i+1,2)$, $D_i$ the curve $(i+1,3)$). & The crossing test (on the doubled polygon), the parity rule, the collapse of the edge $\{1,2\}$ and $\rho$ transcribed from the text; the cocycle compared with \code{A1DnKAlg(2k+1)}'s products of the collapsed curves; the printed exponents applied to the model's resolution states (\code{u1a1deven_geometric_frame._resolutions}) and compared with \code{multiply}; the products against the flow through the closed-form bijection \code{u1a1deven_cone_dodd_section_iso(k)} (\code{verify_multiplicative}). \emph{Population:} $k=1,2,3$: all curves ($12$, $30$, $56$) and maximal cones ($20$, $252$, $3432$); every ordered pair of curves; the monomials of $25$ cones at $e=-1,2$ for \eqref{eq:comp-deven-norm}; every curve at $e=-1,0,2$, $\kappa=0,1$ for $\rho$ and the full turn; every ordered crossing pair ($72$, $450$, $1568$) for the exponents; every ordered generator pair ($196$, $1024$, $3364$) and $72$, $180$, $336$ pairs with $SU(2)$ weights against the flow; the five printed products at $k=1$; \code{SU3ADKAlg}'s geometric labels on $128$ tile monomials (the eight tiles, powers up to $2$, two $SU(3)$ weights). \emph{Controls:} positive: the printed $k=1$ products, and the flow's products; negative: $\rho$ with its drift at marked point $3$, and $E$ on the boundary edge $\{2,3\}$ instead of $\{1,2\}$, each caught. \emph{Evidence:} transcription and cross-presentation (the flow is an independent RG presentation, with its own solver). \emph{Where:} \code{battery/checks_extra_examples.py}, \code{src/cone/u1a1deven_cone_kalgebra.py}, \code{src/cone/u1a1deven_geometric_frame.py}, \code{src/cone/a1deven_kalg.py}, \code{src/rg/u1a1deven_via_dodd_rg.py}, \code{tests/test_u1a1deven_geometric_frame.py}, \code{tests/test_u1a1deven_cone_kalgebra.py}, \code{src/cone/su3_ad_kalg.py}. \emph{Note:} The resolution states are taken from the frame module, not transcribed; they are certified by the products' agreement with the flow. Deeper records (finite\_type\_trace\_machinery.md section J): the product rule equals the derived peel on every ordered crossing pair at k = 1..5, every generator product at k = 4 equals the flow's, and rho equals the flow's on every curve at k = 1..5. \emph{Run (web, 2026-09-26):} 36 checks, 0 failing; k: {'1': {'curves': 12, 'cones': 20, 'noncrossing pairs': 60, 'crossing pairs': 72, 'flow pairs': 268, 'ungauged generators': [4, 4]}, '2': {'curves': 30, 'cones': 252, 'noncrossing pairs': 420, 'crossing pairs': 450, 'flow pairs': 1204, 'ungauged generators': [12, 27]}, '3': {'curves': 56, 'cones': 3432, 'noncrossing pairs': 1512, 'crossing pairs': 1568, 'flow pairs': 3700, 'ungauged generators': [24, 96]}}; seconds: 5.5. & \CERTIFIED\newline {\footnotesize implementation, fast\newline web}\newline {\footnotesize run web 2026-09-26: 36/36} \\ \addlinespace
\code{comp:a1deven/is-kalgebra}\newline companion \S\ref{comp:a1deven} & Both classes are $K_\fq$-algebras: K1--K5 hold on windows of \code{U1A1DevenConeKAlgebra(k)} and \code{A1DevenKAlg(k)} at $k=1,2,3$, and $I_{a,b}=\delta_{a,b}+O(\fq)$ on every ordered pair of the identity and the generators of \code{A1DevenKAlg(k)} ($k=1,2$; every fourth generator at $k=3$); emergent. & The contract's verifier battery (\code{run_k_verifiers}) on all pairs of each window; \code{verify_orthonormality} on the generator pairs. \emph{Population:} $k=1,2,3$: windows of $9$ labels (gauged: the identity, $E^{\pm1}$, odd curves, an $SU(2)$-weighted $E$-dressed curve, charged curves, a loop) and $8$ (ungauged: the identity, five generators, $\chi_1$, $z^{-1}$); $81$, $1600$, $961$ ordered pairs for the orthonormality; $K=3$. \emph{Controls:} positive: the identity rows; negative: n.a. (the verifiers are the instruments). \emph{Evidence:} emergent. \emph{Where:} \code{battery/checks_extra_examples.py}, \code{src/cone/u1a1deven_cone_kalgebra.py}, \code{src/cone/a1deven_kalg.py}. \emph{Note:} The complete k = 3 generator set (14,641 ordered pairs, 144 s) is recorded in finite\_type\_trace\_machinery.md section J (the Layer-1 addendum) and run, in its deep mode, by a test in the source repository. \emph{Run (web, 2026-09-26):} 9 checks, 0 failing; k: {'1': {'gauged window': 9, 'ungauged window': 8, 'orthonormality pairs': 81}, '2': {'gauged window': 9, 'ungauged window': 8, 'orthonormality pairs': 1600}, '3': {'gauged window': 9, 'ungauged window': 8, 'orthonormality pairs': 961}}; K: 3; seconds: 38.4. & \CERTIFIED\newline {\footnotesize implementation, fast\newline web}\newline {\footnotesize run web 2026-09-26: 9/9} \\ \addlinespace
\code{comp:a1deven/traces}\newline companion \S\ref{comp:a1deven} & The class's traces are the printed rules: the gauge tower \eqref{eq:comp-deven-tower} and the seed rules \eqref{eq:comp-deven-seeds} at the representative of the seed's $\rho$-orbit; every other label reduces by Layer 1 onto powers of $E$ and single odd curves times powers of $E$; each pair seed so reduced equals its own closed form; and the transport route (\code{seed_closed_forms=False}) agrees with the seeds and with composite labels. & The tower and the seed rules transcribed from the text independently of \file{src/cone/u1a1deven_seed_characters.py}, with the representative, $\rho$ with its drift and the configurations (lengths and gaps), and compared with the class's \code{trace} as series in $\fq$ and the $SU(2)$ fugacity; the transport route \code{U1A1DevenConeKAlgebra(k, seed_closed_forms=False)} as the witness; the leaves of the Layer-1 reductions inspected. \emph{Population:} the tower at $n=-3..3$ through $\fq^{24}$, $\fq^{20}$, $\fq^{16}$ ($k=1,2,3$); every seed ($8$, $39$, $120$) at $E^{-2..2}$ through $\fq^{16}$, $\fq^{12}$, $\fq^{10}$; the transport route on every seed at $E^{-1,0,1}$ through $\fq^{10}$, $\fq^6$, $\fq^4$ and on $20$, $15$, $10$ composite labels through $\fq^8$, $\fq^6$, $\fq^4$; the pair seeds by Layer 1 at $E^{-1,0,1}$ through $\fq^{12}$, $\fq^{12}$, $\fq^{10}$; at $k=4$ the $50$ side-by-side seeds at $E^{-1..2}$ through $\fq^{10}$. \emph{Controls:} positive: the printed gauge tower (Creutzig's formula) against the class; negative: the nested rule in place of the curve rule, the curve's support moved to $e\ge|n|+2$, the nested gaps swapped, $\rho$ without its drift, and at $k=4$ the retired side-by-side power $\fq^{ng/2}$, each disagreeing with the served traces. \emph{Evidence:} transcription (enforced: the class serves the closed forms) and a witness route (the exact transport down the flow, whose stopping rule is a measured hypothesis, guarded). \emph{Where:} \code{battery/checks_extra_examples.py}, \code{src/cone/u1a1deven_seed_characters.py}, \code{src/cone/exact_characters.py}, \code{src/cone/u1a1deven_trace_transport.py}, \code{src/cone/u1a1deven_cone_kalgebra.py}, \code{tests/test_u1a1deven_seed_characters.py}, \code{tests/test_u1a1deven_cone_kalgebra.py}. \emph{Note:} The seed rules are measured, not derived: against the transport route 1,583 evaluations at k = 1..5 and 108 at k = 6, none disagreeing, every evaluation at k = 5, 6 held out from the fits (finite\_type\_trace\_machinery.md section J; the research record is a probe in the source repository); they are recorded as the starting point of a separate formal certification. The battery's windows are the fast tier's. \emph{Run (web, 2026-09-26):} 15 checks, 0 failing; seeds: {'1': {'seeds': 8, 'kinds': {'curve': 4, 'nested': 4}, 'K': 16}, '2': {'seeds': 39, 'kinds': {'curve': 12, 'side': 3, 'nested': 24}, 'K': 12}, '3': {'seeds': 120, 'kinds': {'curve': 24, 'side': 16, 'nested': 80}, 'K': 10}}; composites: {'1': {'labels': 20, 'K': 8, 'leaves': 23}, '2': {'labels': 15, 'K': 6, 'leaves': 41}, '3': {'labels': 10, 'K': 4, 'leaves': 48}}; witness seconds: {'1': 1.1, '2': 2.4, '3': 6.7}; k=4 side seeds: 50; seconds: 20.9. & \MEASURED\newline {\footnotesize identity, fast\newline web}\newline {\footnotesize run web 2026-09-26: 15/15} \\ \addlinespace
\code{comp:a1deven/ungauging}\newline companion \S\ref{comp:a1deven} & \code{A1DevenKAlg(k)}'s trace is the printed ungauging sum \eqref{eq:comp-deven-ungauge} with the lift \eqref{eq:comp-deven-lift}; the heads \eqref{eq:comp-deven-heads}; $\Tr1$ equals the flow \code{A1DevenRGKAlgebra(k)}'s; the zoo's \code{a1d4}, \code{a1d6}, \code{a1d8} are isomorphic to it through \code{a1deven_seeds.kalgebra_iso}. & The sum over the gauge charge, with $(\fq^2;\fq^2)^{-2}_\infty$, transcribed from the text over the transcribed gauged rules of \code{comp:a1deven/traces}; the lift's sign on a generator at $e=1$, $\kappa=1$; \code{A1DevenRGKAlgebra(k).trace}; the five checks of \code{KAlgebraIso.verify_all} on \code{a1deven_seeds.kalgebra_iso}. \emph{Population:} $k=1,2,3$: the identity and every generator through $\fq^{10}$, $\fq^8$, $\fq^6$; the heads through $\fq^4$; the flow's $\Tr1$ through $\fq^6$; the isomorphisms on the identity, every generator and the flavour characters $\chi_1$, $\mu^{\pm1}$, products on every ordered pair (\code{a1d4}: $121$, \code{a1d6}: $1764$) or those of every fourth generator (\code{a1d8}: $1089$), traces through $\fq^8$. \emph{Controls:} positive: the flow's $\Tr1$, and the zoo's certified isomorphism; negative: the opposite sign of the lift, the measure restored once, and the zoo map precomposed with $\mu\mapsto\mu^{-1}$, each caught. \emph{Evidence:} transcription and cross-presentation (an independent RG flow; the zoo's exported cone tables through a certified isomorphism). \emph{Where:} \code{battery/checks_extra_examples.py}, \code{src/cone/a1deven_kalg.py}, \code{src/cone/ungauge_kalgebra.py}, \code{src/rg/a1deven_rgkalgebra.py}, \code{src/cone/a1deven_seeds.py}, \code{src/cone/finite_su2u1_zform.py}, \code{tests/test_a1deven_seeds.py}. \emph{Note:} At a1d4 the zoo is served by SU3ADKAlg, so the isomorphism's trace check compares independent routes; at a1d6 and a1d8 the zoo's seeds are served through this map, so there only the products and rho are independent data (finite\_type\_trace\_machinery.md section N). The full a1d8 product check (15,129 pairs, 242 s) is tests/test\_a1deven\_seeds.py --slow. \emph{Run (web, 2026-09-26):} 20 checks, 0 failing; k: {'1': {'labels': 9, 'K': 10}, '2': {'labels': 40, 'K': 8}, '3': {'labels': 121, 'K': 6}, 'a1d4': {'samples': 12, 'product pairs each way': 121}, 'a1d6': {'samples': 43, 'product pairs each way': 1764}, 'a1d8': {'samples': 124, 'product pairs each way': 1089}}; seconds: 40.6. & \CERTIFIED\newline {\footnotesize identity, fast\newline web}\newline {\footnotesize run web 2026-09-26: 20/20} \\ \addlinespace
\code{comp:a1dodd/implementation}\newline companion \S\ref{comp:a1dodd} & \code{FiniteA1D5KAlgebra} and \code{FiniteA1D7KAlgebra} present $K_\fq([A_1,D_5])$ and $K_\fq([A_1,D_7])$: $20$ and $42$ generators, $\rho$ of order $5$ and $7$, the $SU(2)$ flavour in the coefficients; every generator product without an $SU(2)$ character agrees with the BPS chart, and every product the hand-written classes cover agrees with them through the flavour dictionary. & Counts and the order of the $\rho$-permutation from the literals; \code{cone_to_bps_iso} multiplicative on the flavour-free products; the products with a character compared with \code{A1D5KAlg} / \code{A1D7KAlg}, $\chi_n L_\gamma$ against the label that \code{r_label_decompose} splits as $(\gamma,n)$. \emph{Population:} D5: all 400 generator pairs (BPS chart and hand-written class); D7: every ordered product of the 8 generators whose BPS charge is a single node charge up to sign (64 pairs) against the chart, every pair inside \code{A1D7KAlg}'s atomic table. \emph{Controls:} positive: the flavour-free products against the BPS chart; negative: a doublet coefficient demoted to a singlet is caught. \emph{Evidence:} cross-presentation (BPS chart; independent hand-written classes). \emph{Where:} \code{battery/checks_extra_examples.py}, \code{src/cone/finite_a1d5_kalg.py}, \code{src/cone/finite_a1d7_kalg.py}, \code{src/cone/a1d5_kalg.py}, \code{src/cone/a1d7_kalg.py}. \emph{Note:} The zoo's cone-to-BPS witness is abelian: it compares U(1) weights and cannot read an SU(2) character, so it covers the flavour-free products only (300 of the 400 D5 generator pairs); the other products are certified against the independent hand-written classes. For D7 a generic BPS product costs minutes, so the chart witness runs on the node-charge generators only (about ten minutes in all); timings in finite\_type\_trace\_machinery.md. \emph{Run (web, 2026-09-26):} 7 checks, 0 failing; classes: {'a1d5': {'generators': 20, 'abelian witness sample': 'all generator pairs', 'abelian witness pairs': 300, 'hand-written pairs': 400, 'uncovered': 0}, 'a1d7': {'generators': 42, 'abelian witness sample': 'every ordered product of the 8 generators whose BPS charge is a single node charge up to sign', 'abelian witness pairs': 62, 'hand-written pairs': 196, 'uncovered': 1568}}; seconds: 662.0. & \CERTIFIED\newline {\footnotesize implementation, fast\newline web}\newline {\footnotesize run web 2026-09-26: 7/7} \\ \addlinespace
\code{comp:a1dodd/is-kalgebra}\newline companion \S\ref{comp:a1dodd} & K1--K5 on a window of the identity, the seeds and products, for both classes; emergent. & \code{run_k_verifiers}. \emph{Population:} the identity, the 4 (D5) and 6 (D7) seeds and a few products; all pairs; $K=4$. \emph{Controls:} positive: the identity row; negative: n.a. \emph{Evidence:} emergent. \emph{Where:} \code{battery/checks_extra_examples.py}. \emph{Run (web, 2026-09-26):} 2 checks, 0 failing; window sizes: {'a1d5': 7, 'a1d7': 9}; K: 4; seconds: 5.4. & \CERTIFIED\newline {\footnotesize implementation, fast\newline web}\newline {\footnotesize run web 2026-09-26: 2/2} \\ \addlinespace
\code{comp:a1dodd/traces}\newline companion \S\ref{comp:a1dodd} & Every elementary trace of the two classes equals its printed combination of admissible $\widehat{\fsl}(2)$ characters \eqref{eq:comp-a1dodd-kappa}--\eqref{eq:comp-a1dodd-seeds}, at levels $-8/5$ and $-12/7$; the same transcription at level $-4/3$ gives the paper's $[A_1,D_3]$ traces. & The characters transcribed independently of \file{src/cone/a1d5_layer2.py} / \file{src/cone/a1d7_layer2.py} and compared as series in $\fq$ and the $SU(2)$ fugacity; three independent witnesses: the Nahm-sum vacuum on the embedded BPS spectrum, the orthonormality bootstrap \code{su2_reliable.reliable_seeds}, and the per-seed BPS engine (D5, through $\fq^6$). \emph{Population:} Tr 1 and the 4 seeds of D5 through $\fq^{16}$, of D7 through $\fq^{12}$; the $[A_1,D_3]$ control through $\fq^{12}$. \emph{Controls:} positive: the level $-4/3$ transcription against \code{A1D3KAlg}; negative: a seed recipe with one term dropped fails. \emph{Evidence:} transcription and cross-presentation. \emph{Where:} \code{battery/checks_extra_examples.py}, \code{src/cone/a1d5_layer2.py}, \code{src/cone/a1d7_layer2.py}, \code{src/cone/su2_reliable.py}, \code{tests/test_a1d5_layer2.py}, \code{tests/test_a1d7_layer2.py}. \emph{Note:} The D7 seeds were also checked once against the per-seed BPS engine through q\^{}5 (635 s), recorded in finite\_type\_trace\_machinery.md. Deep witness (2026-09-23, not run by the fast check): the coverage-checked orthonormality bootstrap a1dodd\_trace\_bootstrap.bootstrap(k, 40) pins every seed of D5 (k = 1, amax 12, 54 s) and D7 (k = 2, amax 10, 137 s) through q\^{}40, equal to a1dodd\_layer2, which equals the printed recipes through q\^{}40 (finite\_type\_trace\_machinery.md section I). \emph{Run (web, 2026-09-26):} 11 checks, 0 failing; classes: {'a1d5': {'K': 16, 'bootstrap K': 12}, 'a1d7': {'K': 12, 'bootstrap K': 8}}; seconds: 252.9. & \CERTIFIED\newline {\footnotesize identity, fast\newline web}\newline {\footnotesize run web 2026-09-26: 11/11} \\ \addlinespace
\code{comp:a1dodd/geometry}\newline companion \S\ref{comp:a1dodd} & The geometric labels of the paragraph ``Geometric labels'': the curves $(x,\ell)$ of the once-punctured $(2n+1)$-gon, $x\in\bZ/(2n+1)$, $2\le\ell\le2n+1$, name the $20$ and $42$ generators one each; $\rho$ is the rotation $(x,\ell)\mapsto(x+1,\ell)$; a cone monomial is named by the multiset of its generators' curves; the generator map onto the family class \code{A1DoddConeKAlg} is a certified isomorphism of $K_\fq$-algebras. & The printed curve set and rotation transcribed and compared with the zoo classes' \code{geometric_label} and \code{rho}; \code{a1dodd_seeds.kalgebra_iso} run through \code{KAlgebraIso.verify_all}. \emph{Population:} $D_5$ and $D_7$: every generator, every $\fq$-commuting generator pair; the isomorphism on the identity, every generator and $\chi_1$, all ordered pairs of the non-identity samples each way, traces through $\fq^6$. \emph{Controls:} positive: the printed curve set reproduced, the isomorphism's five checks; negative: the generator map with two $\rho$-orbits swapped fails multiplicativity. \emph{Evidence:} transcription and cross-presentation (the exported table against the family class). \emph{Where:} \code{battery/checks_extra_examples.py}, \code{src/cone/a1dodd_seeds.py}, \code{src/cone/zoo_geometry.py}, \code{src/cone/a1dodd_kalg.py}, \code{tests/test_a1dodd_seeds.py}, \code{tests/test_zoo_geometry.py}. \emph{Note:} The isomorphism's full record (traces through q\^{}12, the section check, the round trip on product labels, and the D5/D7 negative controls): finite\_type\_trace\_machinery.md section N. \emph{Run (web, 2026-09-26):} 10 checks, 0 failing; classes: {'a1d5': {'generators': 20, 'q-commuting pairs': 90, 'iso samples': 22, 'trace K': 6}, 'a1d7': {'generators': 42, 'q-commuting pairs': 420, 'iso samples': 44, 'trace K': 6}}; seconds: 2.9. & \CERTIFIED\newline {\footnotesize implementation, fast\newline web}\newline {\footnotesize run web 2026-09-26: 10/10} \\ \addlinespace
\code{comp:e6/implementation}\newline companion \S\ref{comp:e6} & \code{FiniteE6KAlgebra} is the cone presentation exported from the $E_6$ BPS chart: $42$ generators, $833$ cones of size $6$, $\rho$ of order $14$, trivial flavour; the isomorphism witness with the chart passes on a sample. & Counts and the permutation order from the literals; \code{cone_to_bps_iso} unit, $\rho$-equivariant and multiplicative. \emph{Population:} the literals; 40 sampled generator pairs; all 42 generators for $\rho$. \emph{Controls:} positive: the literals' counts; negative: n.a. \emph{Evidence:} cross-presentation (BPS chart). \emph{Where:} \code{battery/checks_extra_examples.py}, \code{src/cone/finite_e6_kalg.py}, \code{src/bps/finite_kalgebra_objects.py}. \emph{Run (web, 2026-09-26):} 2 checks, 0 failing; generators: 42; cones: 833; witness pairs: 40; seconds: 6.8. & \CERTIFIED\newline {\footnotesize implementation, fast\newline web}\newline {\footnotesize run web 2026-09-26: 2/2} \\ \addlinespace
\code{comp:e6/is-kalgebra}\newline companion \S\ref{comp:e6} & K1--K5 on a window of the identity, the six seeds and a product, emergent. & \code{run_k_verifiers}. \emph{Population:} 8 labels, all pairs, $K = 4$. \emph{Controls:} positive: the identity row; negative: n.a. \emph{Evidence:} emergent. \emph{Where:} \code{battery/checks_extra_examples.py}. \emph{Run (web, 2026-09-26):} 1 checks, 0 failing; labels: 8; K: 4; seconds: 0.9. & \CERTIFIED\newline {\footnotesize implementation, fast\newline web}\newline {\footnotesize run web 2026-09-26: 1/1} \\ \addlinespace
\code{comp:e6/traces}\newline companion \S\ref{comp:e6} & The class serves $\Tr1$ and the six seeds as the printed combinations \eqref{eq:comp-e6-w3} of $\mathcal W_3(3,7)$ characters; with those combinations absent (and no frozen trace table, removed, and no BPS chart), the exact Nahm-sum vacuum and the orthonormality bootstrap give the same values \eqref{eq:comp-e6-traces}; they equal the per-seed BPS engine through $\fq^6$ and, for $\Tr1$, the $E_6$ RG flow through $\fq^{12}$. & The served values from \code{FiniteE6KAlgebra} against \code{w3_seeds.E6}; a subprocess checks that no \code{e6} table exists, removes the e6 entry of \code{elem_traces._W3_SEEDS} (the witness route) and computes the values; a profiler confirms no \file{src/bps/bps_kalgebra.py} code runs there; the per-seed BPS engine (\code{elem_traces.generate(\"e6\", 6, method=\"bps\")}) and \code{E6RGKAlgebra} run afterwards, as witnesses. \emph{Population:} Tr 1 and the six seeds through $\fq^{12}$; a window of 10 labels for the valuations and pairings. \emph{Controls:} positive: the per-seed BPS engine through $\fq^6$, and the W3(3,7) combinations against the witness route; negative: a seed perturbed by one unit is caught. \emph{Evidence:} cross-presentation (the W3(3,7) combinations against the Nahm-sum + bootstrap route; the per-seed BPS engine; the RG flow). \emph{Where:} \code{src/cone/elem_traces.py}, \code{src/cone/w3_seeds.py}, \code{src/rg/e6_rgkalgebra.py}, \code{battery/checks_extra_examples.py}, \code{tests/test_w3_seeds.py}, \code{src/cone/vacuum_nahm.py}. \emph{Note:} Switched 2026-09-23: the class serves the W3(3,7) combinations; the Nahm-sum + bootstrap route is the witness. The combinations equal the ansatz-free solve of the orthonormality equations through q\^{}37-q\^{}40 (finite\_type\_trace\_machinery.md section K; tests/test\_w3\_seeds.py). Measured further than the battery runs: the witness route to q\^{}16 in 144 s (section G); before the frozen e6 table was removed (2026-09-23) the seeds reproduced it through its whole window q\^{}12. \emph{Run (web, 2026-09-26):} 8 checks, 0 failing; K: 12; window: 10; seconds: 457.7. & \CERTIFIED\newline {\footnotesize identity, fast\newline web}\newline {\footnotesize run web 2026-09-26: 8/8} \\ \addlinespace
\code{comp:e7/implementation}\newline companion \S\ref{comp:e7} & \code{FiniteE7KAlgebra} is the cone presentation exported from the $E_7$ BPS chart: $90$ generators, $3300$ cones of which the $2600$ of size $6$ are exactly the maximal sets of pairwise $\fq$-commuting generators, the $U(1)$ flavour ($\ker B$ spanned by $e_0-e_2+e_6$) in the coefficients, the $5880$ non-$\fq$-commuting ordered products in the table with powers of $\mu$ from $-2$ to $2$, $\rho$ of order $10$ on the generators with the twist \eqref{eq:comp-e7-rho} ($41$ nonzero $\delta_i$) and $\rho^{10}$ the identity on elements; all $8100$ generator products agree with \code{U1E7ConeKAlgebra}'s stored tables under the label map $\varphi$; the object layer's BPS witness reads only the products without $\mu$. & Counts, a maximal-clique enumeration of the $\fq$-commutation graph, the rank and kernel of $B$ and the $\rho$-cycles from the literals; $\rho^{10}$ through \code{rho_element}; $\varphi^{-1}$ written from the printed statement (chord multisets add, $\mu^t\mapsto E^t$) and every generator product of the zoo class compared with the gauged class's product of the preimages; \code{kalgebra_object("e7").iso("cone-frozen", "bps")} on that deterministic sample (a generic $E_7$ BPS product can exhaust $6$\,GB). \emph{Population:} the literals; all 8100 ordered generator products (gauged tables through $\varphi$); for the BPS witness the 36 products of the six generators whose charge is a single node charge up to sign (all without $\mu$), 4 products with $\mu$, and all 90 generators for $\rho$. \emph{Controls:} positive: the literals' counts and the clique enumeration; negative: a $\delta_i$ moved by one breaks $\rho^{10}=\mathrm{id}$, and $\varphi$ with $E\mapsto\mu^{-1}$ fails the product comparison. \emph{Evidence:} cross-presentation (the gauged cone tables, built from the flow \textbackslash{}code{U1A1E7RGKAlgebra}; the BPS chart on the $\textbackslash{}mu$-free products). \emph{Where:} \code{battery/checks_extra_examples.py}, \code{src/cone/finite_e7_kalg.py}, \code{src/cone/u1e7_cone_kalgebra.py}, \code{src/bps/finite_kalgebra_objects.py}, \code{src/cone/cone_kalgebra.py}. \emph{Note:} The object layer's cone-to-BPS witness sends a label to its lattice charge and carries the coefficient across unchanged, so it cannot move a power of mu into a label (its rho check reads the bare label rho): it is multiplicative on the mu-free products (3540 of 8100) and rho-equivariant on the 49 generators with delta = 0 only. The products with mu are certified through the gauged class, whose tables were built by a different engine; its own test (tests/test\_u1e7\_cone\_kalgebra.py) adds the 2600 zoo 6-cones and four perturbed maps. \emph{Run (web, 2026-09-26):} 8 checks, 0 failing; generators: 90; cones: 3300; generator products (gauged witness): 8100; BPS witness sample: {'mu-free': 36, 'mu-carrying': 4, 'rho': 90}; seconds: 15.9. & \CERTIFIED\newline {\footnotesize implementation, fast\newline web}\newline {\footnotesize run web 2026-09-26: 8/8} \\ \addlinespace
\code{comp:e7/is-kalgebra}\newline companion \S\ref{comp:e7} & K1--K5 on a window of the identity, generators with $\delta_i=0,-1,-2$ and two cone monomials, with orthonormality, $\rho^2$-twisted cyclicity and the $\rho$-equivariance of the trace in the element form \eqref{eq:comp-e7-rho}; emergent. & the battery's K1--K3 verifiers (the iteration of \code{run_k_verifiers}) for bar, unit, associativity and $\rho$ (the contract's automorphism check uses \code{rho_element}); $\Tr(\rho(L_a)L_b)=\delta_{a,b}+O(\fq)$, $\Tr(L_aL_b)=\Tr(\rho^2(L_b)L_a)$, $\Tr\rho(L_a)=\star\Tr L_a$ and $I_{b,a}=\star I_{a,b}$ with $\rho=$ \code{rho_element}. \emph{Population:} 12 labels, all 144 pairs, $K = 4$. \emph{Controls:} positive: the identity row; negative: the label-form \code{verify_orthonormality} fails exactly on the window generators with $\delta_i\ne0$, where the $\fq^0$ term of the label-form pairing is $\mu^{-\delta(a)}$. \emph{Evidence:} emergent. \emph{Where:} \code{battery/checks_extra_examples.py}, \code{src/cone/cone_kalgebra.py}. \emph{Note:} On the flavour-in-coefficients cone tier the contract's label-level rho is the bare permutation (the audit); the element form is the honest automorphism, as in tests/test\_e7\_seeds.py::test\_class. \emph{Run (web, 2026-09-26):} 5 checks, 0 failing; labels: 12; pairs: 144; K: 4; seconds: 16.6. & \CERTIFIED\newline {\footnotesize implementation, fast\newline web}\newline {\footnotesize run web 2026-09-26: 5/5} \\ \addlinespace
\code{comp:e7/traces}\newline companion \S\ref{comp:e7} & The class serves $\Tr1$ as the printed product \eqref{eq:comp-e7-vacuum} and all $90$ seeds as the printed recipes \eqref{eq:comp-e7-recipes} over $(\fq^2;\fq^2)_\infty\theta(\mu)$, transported by the rule \eqref{eq:comp-e7-rule}; the product equals the Nahm sum on the BPS spectrum through $\fq^{20}$ and the flow \code{E7RGKAlgebra} through $\fq^{12}$; the seeds equal the $u(1)$ orthonormality bootstrap through $\fq^{20}$ and, without the rule, the bootstrap over $\rho^2$-orbits through $\fq^6$; the printed heads \eqref{eq:comp-e7-heads} hold. & The vacuum product, $\theta$, $\psi_a$, $\mathrm{th}_r$, $\Theta_{ab}$, the nine recipes and the rule transcribed from the companion's text with their own exact series arithmetic (no import of \file{src/cone/e7_seeds.py}); compared with \code{FiniteE7KAlgebra().trace}, with the fixtures \code{PINNED_20} and \code{NAHM_VACUUM_20} of \file{tests/test_e7_seeds.py}, with \code{elem_traces._vacuum_rps("e7", 16)}, \code{E7RGKAlgebra().trace} at $K=12$ and \code{u1_bootstrap.generate_u1("e7", 6, fold="rho2")}. \emph{Population:} Tr 1 through $\fq^{30}$ and all 90 seeds through $\fq^{22}$ against the class; the pinned witnesses through $\fq^{20}$; the live Nahm sum through $\fq^{16}$, the flow through $\fq^{12}$, the $\rho^2$-orbit bootstrap through $\fq^6$. \emph{Controls:} positive: the transcription's $\theta(\mu)$ against the Jacobi triple product through $\fq^{30}$; negative: the vacuum with $10$ replaced by $8$ or $14$ ($u=4,7$) disagrees from $\fq^6$ and $\fq^8$, one sign flipped in $N_{13}$ is caught by the pinned bootstrap, and the rule with $\mu^{+\delta_i}$ is caught. \emph{Evidence:} transcription and cross-presentation (the Nahm sum, the RG flow, the orthonormality bootstrap). \emph{Where:} \code{battery/checks_extra_examples.py}, \code{src/cone/e7_seeds.py}, \code{src/cone/elem_traces.py}, \code{src/cone/u1_bootstrap.py}, \code{src/rg/e7_rgkalgebra.py}, \code{tests/test_e7_seeds.py}, \code{src/cone/vacuum_nahm.py}. \emph{Note:} Standing: the vacuum is certified through q\^{}20 (the Nahm sum) and q\^{}12 (the flow; q\^{}14 once, 162 s, 2026-09-26); the seed recipes are measured, not derived, and their witness bootstraps rest on the recorded mu-support hypothesis (the q\^{}k coefficient of a seed has mu-degree at most 4k). Recorded beyond this row (finite\_type\_trace\_machinery.md section M): the recipes equal the bootstrap through q\^{}22 when it takes the printed product as its Tr 1 (q\^{}21-q\^{}22 never seen by the fit). The reading of the ten products Theta\_ab as Bershadsky-Polyakov module characters is a literature hypothesis confirmed only for the vacuum; nothing checked here depends on it. Measured once, not by this row (2026-09-26, 537 s cold): the per-seed BPS engine, which does not use the mu-support hypothesis, gives seeds 5 and 13 equal to the recipes through q\^{}4. \emph{Run (web, 2026-09-26):} 14 checks, 0 failing; seeds: 90; K: 22; fixtures: q\^{}20; Nahm live K: 16; flow K: 12; rho2 bootstrap K: 6; seconds: 143.4. & \MEASURED\newline {\footnotesize identity, fast\newline web}\newline {\footnotesize run web 2026-09-26: 14/14} \\ \addlinespace
\code{comp:e7/gauged}\newline companion \S\ref{comp:e7} & \code{U1E7ConeKAlgebra}: $202$ generators ($E^{\pm1}$ and $200$ of magnetic charge $-3$ to $3$) and $4160$ cones from its stored tables; its traces follow \eqref{eq:comp-e7-gauged} with the printed $\Tr1$ and seeds, a label of nonzero magnetic charge has trace $0$, and $\Tr1=1-\fq^2+\fq^6+\fq^{12}+\dots$, $\Tr E^{\pm1}=-\fq^3+\dots$. & The formula \eqref{eq:comp-e7-gauged} transcribed (the measure $(\fq^2;\fq^2)_\infty^2$ and the coefficient of $\mu^{-r}$), with the vacuum and seeds of \code{comp:e7/traces}; labels built from the class's dictionary \code{_E7_GENERATOR_PREIMAGES}; the E-tower also against the class's gauged Nahm sum \code{_vacuum_mu}. \emph{Population:} the E-tower $n=-3..3$ through $\fq^{14}$; the 90 neutral generator labels shifted by $E^{-1},1,E$ through $\fq^{10}$; the 130 generators of nonzero magnetic charge through $\fq^8$. \emph{Controls:} positive: the E-tower against the gauged Nahm sum through $\fq^{10}$; negative: the $[\mu^{+r}]$ reading and the formula without the measure are caught. \emph{Evidence:} transcription (the class traces by this formula, so on the neutral labels it is a regression guard) and cross-presentation (the gauged Nahm sum for the E-tower). \emph{Where:} \code{battery/checks_extra_examples.py}, \code{src/cone/u1e7_cone_kalgebra.py}, \code{tests/test_u1e7_cone_kalgebra.py}. \emph{Note:} The products of U1E7ConeKAlgebra are stored tables (u1e7\_cone\_tables.pkl, u1e7\_rho\_tables.pkl) built once from the flow U1A1E7RGKAlgebra, declared self-contained under the reading of the term the author approved. The label map phi is certified by products in comp:e7/implementation; the traces equal the flow's on 16 held-out labels (13 through q\^{}8, 3 through q\^{}6), recorded in finite\_type\_trace\_machinery.md section M and tests/test\_u1e7\_cone\_kalgebra.py. \emph{Run (web, 2026-09-26):} 7 checks, 0 failing; E-tower: 7; neutral labels: 270; magnetic generators: 130; seconds: 4.9. & \CERTIFIED\newline {\footnotesize identity, fast\newline web}\newline {\footnotesize run web 2026-09-26: 7/7} \\ \addlinespace
\end{longtable}
